% La Longue Marche — Volume 140-4
% AI Transcription (Gemini 3.1 Pro, medium thinking)
% 280 of 280 pages transcribed successfully
% 58 diagram pages re-transcribed with tikz-cd enhanced prompt
% Rebuilt: 2026-04-17

%% ===== Page 1 =====

UNIVERSITÉ MONTPELLIER

IMAG
INSTITUT MONTPELLIÉRAIN
ALEXANDER GROTHENDIECK

La "Longue Marche" à travers la théorie de Galois

[suite à La "Longue Marche" à travers la théorie de Galois..., pages 650 bis à 787, dont
table des matières] : notes manuscrites (s.d.).

Cote n° 140-4, 279 pages

Alexander GROTHENDIECK

s.d.

\newpage

%% ===== Page 2 =====

Pages 650-787 + table des matières

NB La page 650 figure deux fois (erreur de numérotation)

\newpage

%% ===== Page 3 =====

50. Enveloppes schématiques de groupes discrets : cas métabélien 650

1°) [unclear: Sections unipotentes]
2°) Essai de caractérisation des sections unipotentes 652
3°) Schéma en groupes associé à certaines extensions de groupes discrets (noyau ab.) 655
4°) Enveloppe schématique d'un groupe métabélien 661
5°) Caract. universelle de l'env. schématique 661

51. Exemple : la tour de Fermat. 664

1°) Le groupe $\underline{G} = \operatorname{Sl}(2, \mathbb{Z})/\{\pm 1\}$, et $\hat{\underline{G}}$ 665
2°) $\mathbb{Z}_{p,\sigma}$ : version provisoire 673
3°) Les groupes $\widehat{\mathbb{Q}_N^{\star}}$ et variantes
4°) La suite $1 \to N_{p,\sigma}^{0+} \to N_{p,\sigma} \to \widehat{\mathbb{Q}_{12}^{\star}} \to 1$, et relations avec $\operatorname{Sl}(2, \mathbb{Z}/12\mathbb{Z})$ 677
5°) $M_{p,\sigma}$, $\underline{M}_{p,\sigma}$ via $N_{p,\sigma}$ 682
6°) Suites $1 \to G_{\sigma} \to \mathbb{Z}_{p,\sigma} \to \widehat{\mathbb{Q}_{12}^{\star}} \to 1$, et $1 \to \hat{\underline{G}}^+ \to \underline{G}_{p,\sigma} \to \widehat{\mathbb{Q}_{12}^{\star}} \to 1$ 684
7°) $N_p, N_p^S$ etc 686
8°) $N_{\sigma}, N_{\sigma}^S$ etc 689
9°) Retour sur $\underline{M}_{p,\sigma}$, $\underline{M}_{p,\sigma}^S$ etc 690
10°) $\underline{M}_{p,\sigma}(0)$, $\underline{M}_{p,\sigma}(1/2)$, $\underline{M}_{p,\sigma}(-j)$ etc 691
11°) L'hom. $M \to U_{p,\sigma} = \underline{M}_{p,\sigma}(\hat{\mathbb{Z}})^{\Delta}$, et relations avec la tour des courbes de Fermat : position du problème. 695
12°) Scindages de l'extension $G_{p,\sigma} / \mu_G$ de $\hat{\mathbb{Z}}^{\star} \times \mathfrak{S}_3$ par $\hat{M}$. 700
13°) [crossed out: Symétries de la tour de Fermat], et relations de compatibilité conjecturales sur $\hat{M} \to G_{p,\sigma}$ 707
14°) Relations avec jacobiennes généralisées 710
15°) Interprétation de $\widehat{\Gamma_Q} \to N_{\sigma}^S \simeq \hat{\Gamma}^* \hat{\mathbb{Z}}$, et digression sur extensions kummeriennes 725
16°) Réflexions sur le th. de Fermat 733

\newpage

%% ===== Page 4 =====

52. [crossed-out: Enveloppes schématiques des [unclear] et [unclear] des groupes] Approximations de Lie des systèmes de Galois-Teichmüller.
1º) Notion de l. d. Lie $\Pi$-groupe : cas discret 737
2º) Cas profini 744
3º) [crossed-out: Approximations de l. a. de Lie des systèmes de Galois-Teichmüller] Approximations de Lie des systèmes de Galois-Teichmüller 748
4º) Approximations en Lie induites par une app... en Lie $\Pi$-autom. 758

53. Le groupe de Teichmüller spécial $S\hat{\Gamma}_{1,1} \simeq \widehat{GL}(2, \mathbb{Z})^{+\sim}$, et le groupe $S\hat{N}_{1,1}$ 750
1º) Action différentiable du gr. de Teichm. spécial id. (carnet provisoire)
2º) Les groupes $\operatorname{Gl}(2, \mathbb{R})^{+\sim}$, $\operatorname{Sl}(2, \mathbb{R})^\sim$ 761
3º) Le groupe $\widehat{GL}(2, \mathbb{Z})^{+\sim}$ 763
4º) L'isomorphisme $\widehat{GL}(2, \mathbb{Z})^{+\sim} \simeq S\hat{\Gamma}_{1,1}$ 765
5º) Digression sur les revêtements des groupoïdes d'espaces, et relations entre group. top. et groupoïdes de fractions [unclear] de morphismes 768

\newpage

%% ===== Page 5 =====

650 bis

50. Enveloppes schématiques de groupes discrets.
(Cas ind-abélien)

1. Soit $G$ un schéma en groupes loc. de prés. finie
[MARGIN: Sections unipotentes]
sur un schéma de base $S$ (le plus simple pour nous sera $S = \operatorname{Spec}(\mathbb{Z})$). Une section $g$ de $G$ sur $S$ est dite unipotente si sur toute fibre géométrique, soit $G_{\bar{s}}$, l'élément $g_{\bar{s}} \in G_{\bar{s}}(\bar{k})$ ($\bar{k} = k(\bar{s})$) est unipotent, i.e. s'il existe un hom. de groupes

(1) $$ \mathbb{G}_{a k} \xrightarrow{\varphi_{\bar{s}}} G_k \qquad k = k(\bar{s}) $$

tel que

(2) $$ g_{\bar{s}} = \varphi_{\bar{s}}(1) $$

On dit que $g$ est quasi-unipotente si pour tt $\bar{s}$, il existe $n \in \mathbb{N}^*$ tel que $g_{\bar{s}}^n$ soit unipotente. Si le car. de $\bar{s}$ est $p > 0$, cela signifie [unclear: aussi] que $g_{\bar{s}}$ est d'ordre fini. On voit que si $S$ est quasi-compact, alors $g$ est quasi-unipotente ssi $\exists n \in \mathbb{N}^*$ tel que $g^n$ soit unipotente.

[unclear: crossed out text]

Pour que $g$ soit unipotente, il suffit qu'il existe un hom.

(3) $$ \varphi : \mathbb{G}_a \longrightarrow G $$

$$ g = \varphi(1) $$

\newpage

%% ===== Page 6 =====

tel qu'on ait
(4) $$ \varphi(1) = g $$
Je me demande si c'est aussi nécessaire,
dans des "bons" cas. Prenons déjà le
cas où $S = Spec \, K$. En caractéristique $0$,
on a un résultat d'unicité pour $\varphi$,
qui permet de [unclear] que le théorème de
descente [unclear] est bien [unclear] :
si après extension $S' \to S$ fid. plate (ou par. exp. [unclear])
il existe $\varphi'$ ayant la propriété voulue,
alors $\varphi$ existe (sur $S$ lui-m. =). Donc sur
un corps de car. 0 c'est OK. En car. p>0
la difficulté provient du fait qu'il y a
des hom. de groupes $\mathbb{G}_a \to \mathbb{G}_a$ (tels $x \mapsto x^p - x$)
qui s'annulent [unclear] non triviale, et qui peuvent
s'annuler sur 1. On s'en tirera [unclear]
obliger de [unclear] à [unclear] : l'[unclear] de $\varphi$ qui
soit mono. Tout sous-groupe de $G$
isomorphe à $\mathbb{G}_a$ a-t-il au moins [unclear]
contenant une section pointée non triviale ?
Est-ce au moins vrai sur un corps de
base ?
Je n'ai pas envie de m'éterniser dans
ces questions (certes intéressantes !), et préfère

\newpage

%% ===== Page 7 =====

m'y limiter : un cas plus proche de celui
- auquel j'avais affaire, $S = Spec \mathbb{Z}$, ou supposons
$S$ un sch. intègre, [crossed out] de pt gén. $\eta$ de
car. $0$. Cela s'applique déjà, d'après le
résultat sur un corps de car. $0$, qu'il
existe un ouvert $U$ non vide
$$ (5) \qquad U \subset S \qquad \text{et} \qquad \varphi_U : \hat{G}|U \to G|U $$
[MARGIN: ouvert non vide]
tel que l'on ait (4) sur $U$ - et (pour
$U$ choisi) ce $\varphi_U$ est unique. Par l'unicité,
on voit donc qu'il y a un plus grand $U$
sur lequel on peut trouver un
tel $\varphi_U$, et la question ouverte : a-t-on
nécessairement $U = S$.

Si on avait $U \neq S$, quitte à localiser
en un pt maximal de $S \setminus U$, on peut se
ramener à examiner le cas où $S$ est local,
de point fermé $s$ - et à prouver (par
absurde que dans ce cas, on peut prolonger $\varphi_U$ sur $S$
(la section $\varphi_U$ : on a contradiction ...). Dans
le cas qui nous intéresse, $S$ est régulier de
dim $1$, donc un localisé est un
trait. Je vais me borner à ce cas, crucial

\newpage

%% ===== Page 8 =====

sans doute pour élucider ce cas général,
et qui suffira pour mon propos.

Soit $y$ le pt générique de $S$, $s$ un pt fermé.
Considérons le sous-groupe $L_y = \varphi_y(G_{log}) \subset G_y$,
soit $L$ son adhérence schématique dans $G$,
qui est un sous-schéma en groupes réduit
de $G$. Soit $L^\circ = L \setminus (\text{réunion des composantes de } L_s \text{ autres que la composante neutre})$
sa composante neutre. On a [unclear: génériquement]
$L^\circ \simeq G_{log}$. En fait, on s'en
fout de savoir si ($L^\circ_s = (L_s)^\circ$) est réduit!
Mais pour que $\varphi : G_{log} \to G$ ait un
prolongement $\varphi_s$ [unclear: continu], il n'est pas nécessaire
sans doute que l'on ait $L^\circ \simeq G_s$.
car il est possible que $\varphi_s$ ne soit pas
mono, i.e. que $\varphi_s$ ne soit pas mono-
il pourrait même être l'hom. trivial (si $G_s = 1$)!
Mais même si $G_s \neq 1$, il n'est pas clair que
$\varphi_s$ soit mono.

Prenons l'exemple où
(6) $$G = \underline{Aut}_{\mathcal{O}_S}(M) \simeq \operatorname{Gl}(n)_S$$
$M$ un $\mathcal{O}_S$-module libre de rang $n$ sur $S$,
où on aura (en supposant $y$ de car 0)
(7) $$\varphi_y(t) = \sum \frac{t^i}{i!} X^i$$
$$X \in End(M_K) \quad K = k(y)$$

\newpage

%% ===== Page 9 =====

552

est un endomorphisme nilp. de $M_K$.
[unclear: défini par un élé] de $End_A(M)\otimes_A K$
(8) $$X \in End_A(M)\otimes_A K, \quad X^n = 0$$
($A$ l'anneau des fonct. sur $S$, qui est une courbe)
On a par hypothèse
$$g_y = \varphi_y(1) = \sum \frac{X^i}{i!} \in \operatorname{Aut}_A(M)$$
i.e.
(9) $$\sum \frac{X^i}{i!}, \sum (-1)^i \frac{X^i}{i!} \in End_A(M)$$
et la question si $\varphi_y$ provient d'un
hom. de groupes $\mathbb{G}_a \to G$ revient
à celle de savoir si
(10) $$\frac{X^i}{i!} \in End_A(M) \quad \forall i \le n-1$$
(pour $i \ge n$ c'est trivial)

Pour $n=2$ c'est trivial, Pour $n=3$ trivial
(dans (9) il y a un seul terme avec un dénominateur).
[unclear: $\varphi(\lambda) = 1 + X\lambda + \frac{X^2}{2}\lambda^2 \quad X \in M_q(K)$]
la [unclear: première] question est de savoir si

$$\varphi_y(\lambda) = 1 + X \lambda + \frac{X^2}{2} \lambda^2 + \frac{X^3}{6} \lambda^3 \quad X \in M_q(K)$$
$$X^4 = 0$$
par hyp. $\varphi_y(1)$ et $\varphi_y(-1)$ dans $M_q(A)$

i.e.
(11) $$\frac{X^2}{2} \pm \frac{X^3}{6} \in M_q(A)$$
se demande s'il en résulte
(12) $$\frac{X^2}{2}, \frac{X^3}{6} \in M_q(A) \quad ?$$

Si $2$ inv. dans $A$ c'est clair, on en tire
de (11) que $X^2 \in M_q(A)$ donc $\frac{X^2}{2} \in M_q(A)$,
donc aussi $X^3/6 \in M_q(A)$, OK. Dans [unclear: l'affirmative]

\newpage

%% ===== Page 10 =====

se présente ici en car. résiduelle 2. Supposons
pour simplifier que 2 est [unclear: régulier dans] $A$, et un $X \notin M_4(A)$
$$X = \frac{1}{2^\nu} X_0 \quad \dots \quad X_0 \in M_4(A)$$
$$X_0 \notin 2 M_4(A)$$
Alors la donnée de $X$ équivaut à celle de
$X_0 = (a_{ij})_{1 \le i,j \le 4}$, les $a_{ij} \in A$ pas tous
divisibles par 2, avec

(11 bis) $$X_0^4 = 0, \quad \frac{1}{2} \frac{X_0^2}{2^{2\nu}} \pm \frac{1}{6} \frac{X_0^3}{2^{3\nu}} \in M_4(A)$$
et la question est s'il en résulte que l'on ait

(12 bis) $$\frac{1}{2} \frac{X_0^2}{2^{2\nu}} , \frac{1}{6} \frac{X_0^3}{2^{3\nu}} \in M_4(A) \quad ?$$

Bien sûr, la réponse est trivialement oui
si $X_0^3 = 0$, [unclear: s'écrivant d'ailleurs le cas où]
$X_0^3 \neq 0$, i.e. $g_y$ est "unipotent principal".

[MARGIN: Essai de caractérisation des sections quasi-unipotentes]

Pour en avoir le coeur net, il faudrait
un contrexemple dans le contexte précédent
(avec une matrice nilpotente d'ordre 4) -
ou un théorème ! Mais pour mon
propos, il semble qu'il ne soit pas
utile d'avoir une solution de ce pb.
Par contre, je vais essayer le
théorème (?) Soit $G$ un schéma en groupes
loc. de présentation finie sur un schéma
$S$ quasi-compact, et soit $g$ une section de
$G$. Pour que $g$ soit quasi-unipotente, il

\newpage

%% ===== Page 11 =====

faut et il suffit qu'il existe un entier
$n > 0$, et un hom. de schémas en groupes
$\varphi : \mathbb{G}_a \to G$, tel qu'on ait
(13) $\varphi(1) = g^n$

La suf(fisance) est un truisme, il s'agit d'examiner
la nécessité. Pour simplifier, on va encore
supposer $S$ noethérien et intègre - une
fois $S$ supposé noethérien, je pense que la
réduction au cas intègre doit se faire par
des dévissages faciles, que je n'ai pas envie
de faire, faute de besoin essentiel ! Je vais
donc examiner le cas $S$ noeth. intègre, pt
générique $y$.

a) car $y$ est $> 0$. Alors l'hypothèse
sur $g_y$ est que $g_y$ est d'ordre fini - donc
$\exists n > 0$ tel que $g_y^n = 1$, donc $g^n = 1$. On
prendra dans (13) cet $n$, et $\varphi$ l'hom. trivial.

b) Cas $y$ de car. nulle. Soit $\varphi$
(14) $\varphi_{(y)} : \mathbb{G}_{a y} \to G_y$
défini par la condition
$g_y = \varphi_{(y)}(1)$

Soit $U$, un "ouvert de définition" dans $S$.

Soit
$\varphi_{n, y}(\lambda) = \varphi_{(y)}(n\lambda) : \mathbb{G}_{a y} \to G_y$
et soit $U_n$ son ouvert de définition. On a
(15) $\varphi_{n,y}(1) = g^n$ sur $U_n$

\newpage

%% ===== Page 12 =====

et
(16) $$n | n' \implies U_n \subset U_{n'}$$

Il faut prouver que la réunion des $U_n$
est $S$. On se ramène au cas $S$ local,
et où $\varphi_1$ est déjà défini dans
(17) $$U_1 = S \setminus \{s\}$$ , $s$ le pt fermé de $S$
et cela revient au

Lemme (?) Soit $S$ noeth. intègre, $G$ schéma en groupes
sép. et loc. de prés. finie sur $S$, $U$ ouvert non vide de $S$,
$\varphi_U : \mathbb{G}_m \to G_U$ un hom. de groupes, on
suppose que $\varphi_U(1) = g_U$ se prolonge en une
section $g$ de $G$ sur $S$. A noter l'unicité (suppo-
-sée, d'après des résultats connus généraux de
SGA 3) Alors il existe un entier $n > 0$ tel que
$\varphi_U^n : \lambda \mapsto \varphi_U(\lambda^n)$ se prolonge sur $S$ en un hom.
$$\varphi : \mathbb{G}_m \to G$$ .

Je ne vais pas essayer de le démontrer dans ce
contexte général des schémas en groupes généraux,
et vais me limiter, à titre de test décisif
(pour emporter conviction que ça marche)
au cas où $G = \operatorname{Gl}(n)_S$, comme tantôt,
et celui où $S$ de plus est régulier de dim. 1.
Donc on est ramené à la situation d'une jauge
$A$, et de $X \in M_n(K)$ ($K$ corps des fractions)
satisfaisant

\newpage

%% ===== Page 13 =====

(18) $$ X^n=0 \ , \ \sum \frac{X^i}{i!} \in M_n(A), \ \sum \frac{i X^i}{i!} \in M_n(A) $$
et on veut montrer qu'il existe un entier $d>0$
tel qu'on ait
(19) $$ \frac{d^i X^i}{i!} \in M_n(A) \quad (1 \leq i \leq n-1) $$

[MARGIN: le cas de car. nulle se résout par l'affirmation à l'aide de la fonction logarithmique :
posant $X+\frac{X^2}{2} \cdots (\text{i.e. } e^X-1) = Y$,
$e^X = 1+Y$, donc
$X = \log(1+Y) = Y - \frac{Y^2}{2} + \frac{Y^3}{3} \cdots$
Montre une [unclear] de base nulle...]

Si $A$ est de car $0$ i.e. s'il est une $\mathbb{Q}$-algèbre, ceci implique d'ailleurs :
$$ X \in M_n(A) $$

Il faudrait que dans le cas de $S$ de car. $0$, c'est $\varphi_U$ lui-même qui puisse s'exprimer sur $S$ tout entier. Le premier cas douteux est :
(20) $$ X \in M_5(A), \ X^5 = 0 $$
$$ \frac{X^2}{2} \pm \frac{X^3}{3!} + \frac{X^4}{4!} \in M_5(A) $$
$$ \stackrel{?}{\Rightarrow} X \in M_5(A) $$

Par contre, dans le cas envisagé au n° précédent ($n=4$, car. résiduelle $2$ - le seul qui fera des difficultés) il est clair qu'il suffit de prendre pour $d$ une puissance de $2$ assez grande.
Cette rem. gen. s'applique [unclear] en d'égales caractéristiques. On déduit même un [unclear]

ceci
Corollaire au lemme dubitatif : si $S$ est local, de car. résiduelle $p>0$, et si $\varphi$ est définie sur un ouvert qui contienne $S - V(p)$, (i.e. sur le complémentaire de $V(p)$ et des car. $p>0$ fait inutile dire ça... que $S$ soit local) $\exists q = p^m$ tel que $\varphi^q$ se prolonge.

\newpage

%% ===== Page 14 =====

Quand $G$ est affine, ceci devrait être trivial. On se ramène à

Proposition. Soit $S$ un préschéma [unclear: fin. qc?], $n$ un entier $\ge 1$, $S_0 = V(n) \subset S$ le sous-schéma fermé défini par $n.1_S$ (donc [unclear: crossed out text]), $U = S \setminus S_0$, $G$ un schéma en groupes affine sur $S$, et

(21) $$ \varphi_U : G_{a|U} \to G|_U $$

un hom. de groupes. On suppose $U$ schématiquement dense dans $S$. Alors il existe une puiss. $m = n^r$ de $n$, et un hom. de schémas en groupes

$$ \psi_m : G_a \to G , $$

tels qu'on ait

(22) $$ \psi_m|_U = (\varphi_U^m : \lambda \mapsto \varphi_U(m\lambda) = \varphi_U(\lambda)^m) $$

Considérons en effet l'endo.

$$ u_m : G_a \to G_a \qquad \lambda \mapsto m\lambda $$

Alors la formule (22) signifie qu'on a commutativité

\begin{tikzcd}
G_{a|U} \ar[r, "(u_m)_U"] \ar[dr, "\psi_m|_U"'] & G_{a|U} \ar[d, "\varphi_U"] \\
& G_U
\end{tikzcd}

et l'existence de $\psi_m$ pour un $m$ grand, signifie que pour un $m$ grand le hom. composé

$$ \varphi_U \circ u_m : G_{a|U} \to G_U $$

se prolonge en un hom de groupes $G_a \to G$.

\newpage

%% ===== Page 15 =====

685

Pour raisons de densité schématiques et [unclear: usuelles]
il suffit de prouver que $\varphi_U$ se prolonge en un
hom de schémas en groupes.

OPS $S$ affine, $S = Spec A$, $A$ un anneau, et soit
$B$ celui de $G$. Donc l'anneau de $U$ sera
$A_u = A[1/u]$, celui de $G_U$ sera
$B \otimes_A A_u = B_u$, par hyp on a un hom.
(correspondant à $\varphi_U$)
$$ (23) \quad B_u \xrightarrow{\varphi^*} A_u[T] $$
tandis que l'hom correspondant à
$\psi_{mU}$ sera [unclear: la composition des] morphismes

$$ (24) \quad A_u[T] \xrightarrow{\psi_m^*} A_u[T] $$
$$ T \mapsto mT \quad \text{donc} \quad T^i \mapsto m^i T^i $$

Soit $x = a_0 + a_1 T + \dots + a_N T^N \in A_u[T]$ où les $a_i \in A_u$
on a
$$ \psi_m^*(x) = \sum a_i m^i T^i = \sum (a_i m^i) T^i $$
donc si $m = u^\nu$ est assez grand, on a
$m^i a_i \in A$ i.e. $m^i a_i T^i \in A[T] \quad \forall i \ge 1$
d'où
$$ \psi_m^*(x) \equiv a_0 \pmod{A[T]} $$
ce qui montre
$$ (25) \quad \psi_m^*(x) \in A[T] \text{ pour } m \text{ assez grand} \iff a_0 \in A $$

\newpage

%% ===== Page 16 =====

Appliquant ceci aux $1$-gén. d'idéaux de
$B$ qui engendrent $B$ comme $A$-algèbre,...

Corollaire $S$ [crossed out], $S_0$, $U$ étant comme dans
la proposition, et $G$ étant un schéma affine
de t. f. sur $S$ (ici pas question de
structure de groupe), et $q_U: E^1_U \to G_U$ un
hom. de schémas, [crossed out: pour qu'il se factorise],
Pour qu'il existe un [unclear: morphisme] $u$
en un tel que [unclear: prolongeant] un morphisme
donné par $(24)$ il faut et il s. que
$q_U$ se prolonge en une section $g$ de $G$ sur $S$.

Remarque Les arguments précédents prouvent
en tous cas le "théorème" de densité p. 662!
dans le cas où $S$ est noeth. de dim. $1$
(le cas intéressant pour nous étant celui où
le pt générique de $S$ est de car. $0$).

[MARGIN:
Schémas en gr-ps
associés à des
groupes discrets]
Soit maintenant $M$ un $\mathbb{Z}$-module
libre de type fini, on va construire un
pro-schéma en groupes sur $S$, associé
au syst. projectif des sous-$\mathbb{Z}$-modules
d'indice fini de $M$. Pour faire la
construction, je vais développer [unclear]

\newpage

%% ===== Page 17 =====

656

[MARGIN: Y a une complication sur l'existence d'une construction de [unclear] $M_{S_0}$ posera $\underline{G} = \underline{G}_{S_0} \times M_{S_0}$]

Lemme préliminaire, relatif à une extension

(26) $$1 \to M \to G \to H \to 1$$

suite exacte de groupes discrets, avec $M$ commutatif — le cas qui nous intéressera pour la suite étant celui où $M$ est un $\mathbb{Z}$-module libre de t. f. En tous cas $M$ est un $\mathbb{Z}$-module, et à ce titre définit un Groupe sur $\operatorname{Spec}\ \mathbb{Z}$, savoir 
$$A \longmapsto M \otimes_{\mathbb{Z}} A$$
($A$, $\mathbb{Z}$-algèbre commutative) — qu'on désigne par $\underline{M}$. En fait, $\underline{M}$ est représentable, i.e. un schéma en groupes sur $\mathbb{Z}$ si $M$ est libre de type fini. J'ai envie de définir aussi des schémas en Groupes sur $S_0$ (qui seront des schémas en groupes si $M$ l'est) $\underline{G}, \underline{H}$, et une suite exacte

(27) $$1 \to \underline{M} \to \underline{G} \to \underline{H} \to 1$$

de Groupes sur $S_0$, et un isomorphisme de suites exactes

(28)
\begin{tikzcd}
1 \ar[r] & M \ar[r] \ar[d, equal] & G \ar[r] \ar[d, "\simeq"'] & H \ar[r] \ar[d, "\simeq"'] & 1 \\
1 \ar[r] & \underline{M}(\mathbb{Z}) \ar[r] & \underline{G}(\mathbb{Z}) \ar[r] & \underline{H}(\mathbb{Z}) \ar[r] & 1
\end{tikzcd}

On prendra

(29) $$\underline{H} = H_{S_0} \text{ Groupe constant de val. } H$$

et les flèches verticales extrêmes de (28)

\newpage

%% ===== Page 18 =====

seraient des isom. canoniques.

On peut songer à définir $\underline{\mathfrak{S}}$ comme
foncteur (cov.) $A/\mathbb{Z}-alg$ ou (contravariant) sur $S/S_0$,

par
$$(30) \qquad \underline{\mathfrak{S}}(A) = \mathfrak{S} \times_M M(A)$$

$\underline{\mathfrak{S}}(A)$ est donc défini comme l'ext. de
$H$ par $M(A)$ déduite de l'ext. $\mathfrak{S}$
de $H$ par $M$ par l'hom.
$$(31) \qquad M \longrightarrow M \otimes_{\mathbb{Z}} A ,$$

qui a un sens. comme on a une
str. de $H-\mathbb{Z}$-module, on fait opérer $H$
sur $M \otimes_{\mathbb{Z}} A$ via son opération sur $M$.

La petite difficulté est que le foncteur-
groupes ainsi défini par (30) n'est pas
lui-même représentable - il ne donne l'
expression correcte de $\underline{\mathfrak{S}}(A)$ que si
$\operatorname{Spec} A$ est connexe. il correspond -
à un faisceau de
$$S \mapsto \underline{\mathfrak{S}}(A) = \underline{\mathfrak{S}}(S) \quad (S = \operatorname{Spec} A)$$
faisc. [unclear: en plat] [unclear: $\underline{\mathfrak{S}}$] au des. de $\underline{H}_S$ pour que la compos.

$$(32) \qquad S \longrightarrow \underline{H}_S$$
(qui correspond à une fonction loc. constante
sur $S$ à valeurs dans $H$) soit constante.

\newpage

%% ===== Page 19 =====

L'expression correcte s'obtient en partant
du préfaisceau "incorrect"
$$[\underline{\mathfrak{S}}] : A \longmapsto \mathfrak{S}_{M}(M \otimes_{\mathbb{Z}} A)$$
ou mieux, du contrafoncteur
$$S \longmapsto \mathfrak{S}_{M}(\Gamma(S, M \otimes_{\mathbb{Z}} \mathcal{O}_S)) \quad ,$$
et en prenant le faisceau Zariskien assoc-
-ié, soit

(33) $$ \underline{\mathfrak{S}}(S) = \lim_{\substack{\longrightarrow \\ \mathcal{U} = (U_i)_{i \in I} \\ \text{(où } (U_i)_{i \in I} \text{ une partition} \\ \text{de S en sous-schémas} \\ \text{ouverts } S \simeq \coprod U_i)}} \prod_{i \in I} [\underline{\mathfrak{S}}](U_i) $$

Voici une autre façon de présenter la
construction de $\underline{\mathfrak{S}}$, en utilisant la
construction générale de §47 -- dans le
contexte faisceautique. On a un système

(34) $M$, $\underline{\mathfrak{S}}_{S_0} \triangleright M_{S_0}$, $Op.$ de $\underline{\mathfrak{S}}_{S_0}$ sur $M$
Gr. sur $S_0$  Gr. pro(const) sur $S_0$, $M_{S_0}$ distingué  hom $\varphi : M_{S_0} \to M$

(où l'op. de $\underline{\mathfrak{S}}_{S_0}$ sur $M$, qui [unclear: en fait]
un op. de $\underline{\mathfrak{S}}$ sur $M$, est celle déduite
de l'op. de $\underline{\mathfrak{S}}$ sur $M$, par fonctorialité)

ces données satisfont aux conditions de
compatibilité a) b) de loc. cit. Par la
construction de loc. cit., on construit

\newpage

%% ===== Page 20 =====

à l'aide de ces données un système

(35) $\underline{\mathfrak{S}} \supset \underline{M} \quad , \text{ et } \Phi : \underline{\mathfrak{S}}_{S_0} \to \underline{\mathfrak{S}}$
Groupe sur $S_0$,
son $\underline{M}$ sous-
sous-groupe distingué

tel que (34) se reconstitue à partir de
(35) en prenant $\underline{M}$, $\underline{\mathfrak{S}}_{S_0}$ dans (35),
$M_{S_0}$ l'image inverse de $\underline{M} \subset \underline{\mathfrak{S}}$ par $\Phi$,
$\varphi$ induit par $\Phi$, et $\Theta$ l'opération déduite
de $\Phi$. De plus, $\Phi$ induit un iso.
$$ \underbrace{\underline{\mathfrak{S}}_{S_0} / M_{S_0}}_{H_{S_0}} \xrightarrow{\sim} \underline{\mathfrak{S}} / \underline{M} $$

- et on a là une description "axiomatique"
d'un système (35) en termes de (34).
La suite exacte de Groupes sur $S_0$

(36) $$ 1 \to \underline{M} \to \underline{\mathfrak{S}} \to \underline{\mathfrak{S}} / \underline{M} \to 1 $$
$$ \simeq $$
$$ H_{S_0} $$

[MARGIN: L'hyp. p. être de fait inutile, si dans les constructions (plus générales), on se [unclear: borne] au cas les $M$-torseurs sur $S_0$ définis par des images inverses [unclear: évidentes] sont triviaux]

Il se scinde sur base affine $S$
($S$ tel que $H^1(S, M \otimes_{\mathbb{Z}} \mathcal{O}_S) = 0$)

$$ 1 \to M(S) \to \underline{\mathfrak{S}}(A) \to H_{S_0}(S) \to 1 $$
[unclear: $\simeq M \otimes_{\mathbb{Z}} A$] \quad [unclear: si $S = Spec A$] \quad [unclear: $\simeq H \times S$]

\newpage

%% ===== Page 21 =====

En fait, on peut [unclear: construire] un diagr.
commut. de suites exactes

(38)
\begin{tikzcd}
1 \ar[r] & \underline{M}_{S_0} \ar[r] \ar[d] & \underline{\mathfrak{S}}_{S_0} \ar[r] \ar[d] & \underline{H}_{S_0} \ar[r] \ar[d] & 1 \\
1 \ar[r] & \underline{M} \ar[r] & \underline{\mathfrak{S}} \ar[r] & \underline{H}_{S_0} \ar[r] & 1
\end{tikzcd}

qui par passage aux sections donne
pour $S$ quelconque affine (disons) connexe
un diagr. comm.

[MARGIN: NB En [unclear: déduisant] de (39)
$\underline{\mathfrak{S}}(A)$ est isomorphe
au produit semi-direct
de $H$ par
$M \otimes_{\mathbb{Z}} A$
[unclear]
[unclear]]

(39)
\begin{tikzcd}
1 \ar[r] & M \ar[r] \ar[d] & \mathfrak{S} \ar[r] \ar[d] & H \ar[r] \ar[d, equal] & 1 \\
1 \ar[r] & M \otimes_\mathbb{Z} A \ar[r] & \underline{\mathfrak{S}}(A) \ar[r] & H \ar[r] & 1
\end{tikzcd}

pour $S=S_0$ on retrouve l'iso (28)

- . - . - .

Cette construction du $S_0$-schéma en
groupes $\underline{\mathfrak{S}}$, à partir d'une
suite exacte (27), est [unclear: manifestement]
fonctorielle par r. à cette dernière : un
hom. de suites exactes

(40)
\begin{tikzcd}
1 \ar[r] & M \ar[r] \ar[d] & \mathfrak{S} \ar[r] \ar[d] & H \ar[r] \ar[d] & 1 \\
1 \ar[r] & M' \ar[r] & \mathfrak{S}' \ar[r] & H' \ar[r] & 1
\end{tikzcd}

donne naissance à un diagr. comm.
de suites exactes

\newpage

%% ===== Page 22 =====

(41)
\begin{tikzcd}[row sep=1.5em, column sep=2em]
& 1 \ar[r] & M_{S_0} \ar[r] \ar[dd] \ar[dl] & \underline{G}_{S_0} \ar[r] \ar[dd] \ar[dl] & H_{S_0} \ar[r] \ar[dd] \ar[dl] & 1 \\
1 \ar[r] & M'_{S_0} \ar[r] \ar[dd] & \underline{G}'_{S_0} \ar[r] \ar[dd] & H'_{S_0} \ar[r] \ar[dd] & 1 & \\
& 1 \ar[r] & \underline{M} \ar[r] \ar[dl] & \underline{\underline{G}} \ar[r] \ar[dl] & H_{S_0} \ar[r] \ar[dl] & 1 \\
1 \ar[r] & \underline{M}' \ar[r] & \underline{\underline{G}}' \ar[r] & H' \ar[r] & 1 &
\end{tikzcd}

[MARGIN: Exemple 4. réduction discrète des groupes métabéliens]

Reprenons donc un $\mathbb{Z}$-module libre
de type fini $M$, et pour $H$ un
groupe d'ordre fini, considérons
la suite exacte

(42) $$1 \longrightarrow M_i \longrightarrow M \longrightarrow M/M_i \longrightarrow 1$$

Plus gén., si on part d'une suite
exacte (27), pour un sous-groupe $M_i$
d'indice fini de $M$, stable par
$\underline{\underline{G}}$ i.e. stable par $H$, on a une
suite exacte

(43) $$1 \longrightarrow M_i \longrightarrow \underline{\underline{G}} \longrightarrow \underline{\underline{G}}/M_i \longrightarrow 1$$

qui montre que $H_i = \underline{\underline{G}}/M_i$
est l'extension de

(44) $$1 \longrightarrow M/M_i \longrightarrow H_i \longrightarrow H \longrightarrow 1$$

déduite de (27) par
$M \longrightarrow M/M_i$.

\newpage

%% ===== Page 23 =====

[MARGIN: prendre les $n M$ !]
On note que si $H$ est un groupe
de type fini, la famille
des $M_i \subset M$ d'indice fini et
stables sous $H$ est cofinale dans l'ens.
de tous les $M_i$ d'indice fini de $M$.

Considérons dans ce cas la suite exacte de
schémas en groupes associés

(45) $$1 \to \underline{M}_i \to \underline{G}_i \to \underline{H}_i \to 1$$

Pour $i$ variable, elles forment un syst.
projectif d'extensions, d'où une
extension de pro-groupes sur $S_0 = \operatorname{Spec} \mathbb{Z}$.

On considère le pro-objet $(\underline{G}_i)$
comme pro-objet de la catégorie des
schémas en groupes (lisses etc ...)
sur $S_0$. On note que les morphismes
de transition

$$\underline{G}_j \to \underline{G}_i$$

sont affines - ce qui [unclear: équivaut] à la [unclear: surjectivité] pour l'hom

$$E(M_i, M_j) \to E(M_i, M_i)$$
[unclear: ext. de $(M_i/M_j)_{S_0}$ par $M_j$]

\newpage

%% ===== Page 24 =====

Or le premier est affine, comme somme
finie de schémas affines, le deuxième
est affine - donc le morphisme
est affine.

Ceci étant, la limite projective des
$\underline{\tilde{G}}_i$ dans la catégorie des schémas existe
(plus précisément, le foncteur limite
projective des $\underline{\tilde{G}}_i$ est représentable) -
il n'est peut être pas absurde de passer
à la limite au-dessus, et de travailler
avec un schéma en groupes - pro de
t.f.

Remarque. Le pro-groupe $(\underline{\tilde{G}}_i)$ sur $S_0 =$
$Spec \mathbb{C}$ ne dépend que du pro-
groupe discret $\tilde{G}_i$, muni d'une classe
[crossed out: immédiatement obtenue]
de pro-groupes [unclear] comme [unclear]
lité - l'hypothèse à faire étant que
pour un pro-groupe, il y a [unclear]
[unclear] dans $Z$ un modèle de type fini
(car le pro-groupe assez petit, il sera libre
de t.f.). Si $M_0$ est d'indice
fini dans $\tilde{G}_{i_0}$, si les $H_i$ sont finis,

\newpage

%% ===== Page 25 =====

alors ce proobjet ne dépend que de $G$ lui
-même - à la condition qu'il existe
un sous-groupe d'indice fini de $G$ qui
soit un groupe commutatif, et de t.f. en
tant que $\mathbb{Z}$-module. (On dit lors
que $G$ est "métabélien")

Regarder le cas typique
Ainsi, pour tout groupe discret
métabélien, on construit un proobjet
(syst. proj. de schémas en groupes affines),
que je vais noter

(46) $$Env(G) = (\underline{G}_i)$$

(l'ens d'indices étant formé des sous-groupes
distingués de $G$ qui sont comm. et sans
torsion). Ce proobjet dépend fonctoriellem.
de $G$. Il a des propriétés d'exacti-
tudes intéressantes et $\pm$ évidentes, p.ex.

il est multiplicatif

(47) $$Env(G \times G') \xrightarrow{\sim} Env(G) \times_{S_0} Env(G')$$

Si on a une suite exacte

(48) $$1 \to G'' \to G \to G' \to 1$$

de groupes métabéliens, on a une suite
exacte pour les enveloppes

(49) $$1 \to Env(G'') \to Env(G) \to Env(G') \to 1$$

\newpage

%% ===== Page 26 =====

en un sens qu'il convient de préciser :
l'occasion - pour les schémas en groupes
associés (obtenus en passant à la limite) on
aura bel et bien un lemme 'fid plat' dont
le noyau est de base ...

Le cas typique, auquel le cas général
se ramène de façon presque immédiate, est
celui où $G = M = \mathbb{Z}$. On a pour
tout entier $n \ge 1$ la suite exacte

(50) $$1 \to \mathbb{Z} \xrightarrow{n} \mathbb{Z} \to \mathbb{Z}/n\mathbb{Z} \to 1$$

Je vais expliciter la construction
de l'extension correspondante de
$\mathbb{Z}/n\mathbb{Z}$ par $\mathbb{Z} = G_a$, soit $\underline{G}_n$.
C'est un schéma affine, dont l'anneau
des coordonnées est isomorphe - de façon pas
bien canonique - au produit de
$n$ copies de celui de $\mathbb{Z} = G_a$, savoir
$\mathbb{Z}[T]$ - on trouve un tel isomorphisme
en choisissant une section ensembliste de
$\mathbb{Z}$ sur $\mathbb{Z}/n\mathbb{Z}$, par exemple par l'in-
clusion $\mathbb{Z} \cap [0, n-1]$. Les lois de

\newpage

%% ===== Page 27 =====

51

transition entre les $\mathbb{G}_n$ n'est rien moins que
plat bien sûr. Prenons : la limite projective,
- donne une extension de
$$(\hat{\mathbb{Z}})_{S_0} \stackrel{df}{=} \lim_{\leftarrow\atop n} (\mathbb{Z}/n\mathbb{Z})_{S_0}$$
- schéma en groupes profini sur $S_0$ -
par la limite projective des schémas
$\mathbb{G}_n$, s'envoyant dans lui-même par les hom.
vus plus haut (cf 20) : c'est le spectre
de l'anneau limite inductive des
$$\mathbb{Z}[T_n] \subset \mathbb{Q}[T]$$
$$T_n = \frac{1}{n} T$$
qui n'est - ici, chose étrange - que la sous $\mathbb{Z}$-algèbre
de $\mathbb{Q}[T]$ formée des polynômes
dont le terme constant est un entier.

On a une description analogue de
la limite projective de l'enveloppe
schématique des $\mathbb{Z}^n$ - c'est une
extension de $(\hat{\mathbb{Z}}^n)_{S_0}$ (attention : l'abus
de notation) par une composante neutre,
qui est le spectre du sous-anneau
de $\mathbb{Q}[T_1, ..., T_n]$ formé des polynômes dont
le terme constant est un entier...

\newpage

%% ===== Page 28 =====

C'est donc un schéma en groupes sur $S_0 = Spec \mathbb{Z}$
dont toutes les fibres sont isomorphes au groupe
additif, sauf la fibre générique, qui est
$\mathbb{G}_{a \mathbb{Q}}$.

Dans le cas d'un groupe abélien
de type fini quelconque $\mathfrak{G}$, désignons par $\hat{\mathfrak{G}}$ sa
complétion profinie, par $\hat{\mathfrak{G}}_{S_0}$ le schéma
en groupes profini sur $S_0$
$$ (50) \quad \hat{\mathfrak{G}}_{S_0} \overset{df}{=} \lim_{\leftarrow i} (H_i)_{S_0} $$
on a aussi une suite exacte canonique
$$ (51) \quad 1 \to \underline{Env}^0(\mathfrak{G}) \to \lim_{\leftarrow} \underline{Env}(\mathfrak{G}) \to (\hat{\mathfrak{G}})_{S_0} \to 1 $$
où $\underline{Env}^0(\mathfrak{G})$ désigne la comp. neutre
de $\lim_{\leftarrow} \underline{Env}(\mathfrak{G})$, i.e.
$$ (52) \quad \underline{Env}^0(\mathfrak{G}) \overset{df}{=} \lim_{\leftarrow i} \check{M}_i \simeq \lim_{\leftarrow n} n \cdot M_0 $$
où les $M_i$ sont les sous-groupes d'indice
fini dans $\mathfrak{G}$ [unclear], et (cf.
plus haut) pour un $\mathbb{Z}$-module $M$,
$$ (53) \quad \check{M} = \mathbb{V}(M) = \operatorname{Spec}(Sym_{\mathbb{Z}}(\check{M})) $$
désigne le fibré vectoriel [unclear: cov. sur $M$] qu'il définit.

\newpage

%% ===== Page 29 =====

où $M_0$ est un des $M_i$, de rang fini
- sym. (cofinal dans l'ens. de tous les
$M_i$) on fixe un $M_0$. Le choix d'une
base de $M_0$ sur $\mathbb{Z}$ donne un
iso

(54) $$Env^\circ(\hat{\mathbb{C}}) \simeq Env^\circ(M_0) \simeq (Env^\circ(\mathbb{Z}))^n$$
$$\simeq (\mathbb{E}')^n$$

[MARGIN: (où $(\mathbb{E}', \cdot)$ est le gr. multip.)]

(55) $$\mathbb{E}' \simeq Env^\circ(\mathbb{Z}) \simeq \operatorname{Spec}(\mathbb{Z}[T, T^{-1}])$$
$$\simeq \varprojlim (\mathbb{E}', n)$$
est le schéma en gr. pro-affine qu'on a
envisagé plus haut, qui représente le foncteur-groupe

(56) $$A \longmapsto \mathbb{E}'(A) = \{ (x_n) \in A^{\mathbb{N}^*} / x_n = x_{mn}^m \}$$
$$\forall m,n \in \mathbb{N}^*$$

[Si $A$ est tel qu'il existe $c \in \mathbb{N}^*$ avec
$c = 0$ dans $A$ - ce qui signifie
que les car. résiduelles de $A$ divisent $c$ ,
et même s'il existe dans $A$ un nil-
idéal dès que les car. résiduelles de $A$ sont
toutes $>0$ - alors $\mathbb{E}'(A) = \{1\}$ -
Si par contre $A$ est une $\mathbb{Q}$-alg., on a $\mathbb{E}'(A) \simeq A$ ]

\newpage

%% ===== Page 30 =====

## 5. Caractérisation universelle de l'enveloppe pro-unipotente (resp. quasi-unipotente) d'un groupe abstrait.

Soit $G$ un schéma en groupes loc. de prés. finie sur une base $S$.
On dit qu'un hom. de schémas en groupes
$$H \to G$$
est unipotent (resp. quasi-unipotent) si l'image de $H$ (i.e. l'image de $H_{red}$) est un sous-groupe unipotent (resp. quasi-unipotent) de $G$ - ce qui se vérifie fibre par fibre. Un hom. d'un groupe discret
$$\Gamma \to \Gamma(G/S)$$
est dit unipotent (resp. quasi-unipotent) si l'hom. de schémas en groupes qui lui correspond
$$\Gamma_S \to G$$
est unipotent (resp. quasi-unipotent).

\newpage

%% ===== Page 31 =====

Notons que d'après les propriétés des
$\varprojlim$ de syst. projectifs filtrants (à morphis-
mes de transition affines) entre des
schémas, si $\mathfrak{G}$ est nildélien,
d'où un système projectif de
schémas en groupes affines sur $S_0$, donc
sur $S$, noté $Env_S(\mathfrak{G})$, et une
limite projective
(57) $$LEnv_S(\mathfrak{G}) = LEnv(\mathfrak{G}) \times_{S_0} S$$
$$\simeq \varprojlim Env_S(\mathfrak{G}) \quad ,$$
on a, dès que $G$ est loc de
prés. finie
(58) $$\operatorname{Hom}_{progr/S}(Env_S(\mathfrak{G}), G) \xrightarrow{\sim} \operatorname{Hom}(LEnv_S(\mathfrak{G}), G)$$
$$\parallel def$$
$$\varinjlim \operatorname{Hom}(\mathfrak{G}_{i,S}, G)$$

Notons d'autre part qu'on a un
hom canonique
(59) $$\mathfrak{G}_S \to LEnv_S(\mathfrak{G})$$
déduit par ch. de base de l'hom
sur la base universelle $S_0$
(59 bis) $$\mathfrak{G}_{S_0} \to LEnv_{S_0}(\mathfrak{G})$$
provenant du système d'hom can.

\newpage

%% ===== Page 32 =====

$$ \underline{\mathfrak{G}}_{S_0} \longrightarrow \underline{\mathfrak{G}}_i $$

Donc, via (59), on trouve

(60) $$ \operatorname{Hom}_{S-gr}(L Env_S(\mathfrak{G}), G) \longrightarrow \operatorname{Hom}_{S-gr}(\mathfrak{G}_S, G) $$
$$ \simeq $$
$$ \operatorname{Hom}_{gr}(\mathfrak{G}, G(S)) $$

Comparant (58) et (60), on trouve

(61) $$ \operatorname{Hom}_{pro-S-gr}(Env_S(\mathfrak{G}), G) \simeq \operatorname{Hom}_{S-gr}(L Env_S(\mathfrak{G}), G) $$
$$ \longrightarrow \operatorname{Hom}_{S-gr}(\mathfrak{G}_S, G) \simeq \operatorname{Hom}_{gr}(\mathfrak{G}, G(S)) $$

Ceci posé, on a le résultat suivant

Théorème. Sous les conditions précédentes ($\mathfrak{G}$ groupe discret unithérien, $G$ schéma en groupes loc. de prés. finie sur schéma de base $S$), l'application canonique (61) est injective, et a comme image l' [unclear] des hom. $u : \mathfrak{G} \to G(S)$ [unclear] impliquant l'existence, pour $S$ [unclear] et $G$ affine sur $S_0$.

\newpage

%% ===== Page 33 =====

5 Exemple : La "tour" de Fermat.

1. Les groupes $\tilde{\Pi}_0, \tilde{\mathcal{C}}$.

[MARGIN: $\Pi_0$ engendré par les $\lambda_i = -l_i$ 
$l_0 = \begin{pmatrix} 1 & -2 \\ 0 & 1 \end{pmatrix}$ 
$l_1 = \begin{pmatrix} 1 & 0 \\ 2 & 1 \end{pmatrix}$]

Je m'intéresse au quotient de $\tilde{\mathcal{C}} = \widehat{Sl(2,\mathbb{Z})}$
obtenu en prenant le quotient par
$[\Pi_0, \Pi_0]$ — ou encore une extension
$\tilde{\mathcal{C}}$ de $\tilde{\mathcal{C}}/\Pi_0$ par le groupe
(1) $$M \overset{dfn}{=} \Pi_{0 ab} \simeq \left( \hat{\mathbb{Z}}_{(0)} + \hat{\mathbb{Z}}_{(1)} + \hat{\mathbb{Z}}_{(\infty)} \right) / diagonale$$

Ainsi on a une suite exacte de groupes
pro-finis
(2) $$1 \longrightarrow M \longrightarrow \tilde{\mathcal{C}} \longrightarrow \tilde{\mathcal{C}}/\Pi_0 \longrightarrow 1$$
[Sous $M$] : $\simeq \Pi_{0 ab}$, un $\tilde{\mathcal{C}}$-mod.
[Sous $\tilde{\mathcal{C}}/\Pi_0$] : $\simeq \operatorname{Sl}(2, \mathbb{Z}/4\mathbb{Z}) / \langle l_0, l_1 \rangle$

Par générateurs et relations on a.
(3) $$\tilde{\mathcal{C}} = \left\{ \tau, \rho, \sigma \ \Bigg| \ \begin{array}{l} \rho^6 = \sigma^4 = \tau^2 = 1 \\ \rho^3 = \sigma^2, \ ^\tau(\rho) = \rho^{-1}, \ ^\tau(\sigma) = \sigma^{-1} \\ \langle [\lambda_0, \lambda_1] = 1 \rangle, \text{ où} \\ \quad \begin{cases} \lambda_0 = -l_0 = - \sigma \rho \sigma \rho = \sigma (\rho) \sigma \dots \\ \lambda_1 = -l_1 = - \rho \sigma \rho \sigma = \rho \sigma (\rho) \dots \end{cases} \end{array} \right.$$

d'où la relation de commutateurs $[\lambda_0, \lambda_1] = 1$ s'écrit
aussi
(4) $$\sigma \rho (\sigma \rho \sigma) = \rho \sigma (\rho \sigma \rho)$$
(5) $$- \sigma \rho \sigma^{-1} \rho \sigma = \rho \sigma^{-1} \sigma \rho$$ [unclear]

où le sens des signes - peu significatifs
dans le formalisme des gén. et rel. en supprimant
dans la formule ces signes, et en remplaçant les
inverses un [unclear] par un signe [unclear] de
la formule par des $\sigma^{-1}$. Par exemple, la formule

\newpage

%% ===== Page 34 =====

s'écrit [unclear]

[MARGIN: NB $\lambda_{\infty} \overset{dfn}{=} -l_{\infty}$]
(6) $\lambda_{\infty} = \rho \sigma \rho^{-1} \sigma \rho = \rho (\sigma) \sigma \rho = - \rho (\sigma \rho)^{-1} \rho = \rho \sigma \rho^{-1}(\sigma)$

dans la forme

[MARGIN: NB $\sigma = \sigma_{\infty}$]
(7) $[\sigma, \lambda_{\infty}] = 1$

[MARGIN: NB On trouve
$[l_0, l_1] = [\lambda_0, \lambda_1] = (\sigma(\rho) \rho^{-1})^3$
$[l_{\infty}, l_0] = [\sigma_{\infty}, \lambda_{\infty}] = (\sigma (\rho^{-1}) \rho)^3 = int(\sigma \rho^{-1}) \cdot [l_0, l_1]$]

On va définir une extension de groupes
$Z$, à un isomorphisme (dépendant d'un choix) près
$\tilde{\mathfrak{S}}^\wedge$. On pose

[MARGIN: Il n'y a pas de moyen d'avoir $M$ - il [unclear] tjrs le]
(8) $M = \tilde{\mathfrak{S}}^\circ = E^{\{0, 1, \infty\}} / diagonale$

On va faire opérer $\widehat{Gl}(2, \mathbb{Z})$

(9) $\widetilde{Gl}(2, \mathbb{Z}) \simeq \widehat{Gl}(2, \mathbb{Z}) / [\Pi_0, \Pi_0]$

sur $M$, décrivant ainsi

(s) $\tilde{\mathfrak{S}}_3 \times \mu_2 \simeq \tilde{\mathfrak{S}} / M \simeq \tilde{\mathfrak{S}}_{0,3} / \Pi_0$

Désignons par $\lambda_0^\circ, \lambda_1^\circ, \lambda_{\infty}^\circ$ les images des
$\lambda_0, \lambda_1, \lambda_{\infty} = -l_i$ dans $\tilde{\mathfrak{S}}$, identifiés
à des éléments de $M = \tilde{\mathfrak{S}}^\circ = \tilde{\mathfrak{S}}^\circ(\hat{\mathbb{Z}})$

en fait...

(6) $\lambda_0^\circ, \lambda_1^\circ, \lambda_{\infty}^\circ \in M(\mathbb{Z}) \subset M(\hat{\mathbb{Z}}) = M \otimes \hat{\mathbb{Z}}$
$\lambda_0^\circ + \lambda_1^\circ + \lambda_{\infty}^\circ = 0$ (On a noté additivement la loi de groupe de $M = \tilde{\mathfrak{S}}^\circ$)

[MARGIN: Alors $\tilde{\mathfrak{S}}_3$ opère sur $M = \tilde{\mathfrak{S}}^\circ$]
en permutant $\lambda_0^\circ, \lambda_1^\circ, \lambda_{\infty}^\circ$, d'où

(7) $\sigma(\lambda_0^\circ) = \lambda_1^\circ, \sigma(\lambda_1^\circ) = \lambda_{\infty}^\circ, \sigma(\lambda_{\infty}^\circ) = \lambda_0^\circ$

[MARGIN: ceci se voit d'abord pour les $l_i$ et on passe au quotient pour les $\lambda_i^\circ$]
$\tau(\lambda_i^\circ) = - \lambda_{-i}^\circ$ par définition
où $\tau = \tau_{\infty} = \tau'$

\newpage

%% ===== Page 35 =====

685

L'extension $\underline{\mathbb{S}}$ annoncée est bien sûr
celle du $§ 50$ c), développée entre-temps (depuis
les premières lignes du présent $\S$, constatant
justement que j'étais gêné aux entou-
rnures pour expliciter la construction) :

(9) $$1 \to \underline{M} \to \underline{\mathbb{S}} \to H_{S_0} \to 1$$
[under $\underline{M}$ is $||$ and $\check{V}(M)$]
où $S_0 = Spec \mathbb{Z}$.

Pour tt schéma $S$, on posera

(10) $$\underline{\mathbb{S}}^\Delta(S) \subset \underline{\mathbb{S}}(S)$$
[under $\underline{\mathbb{S}}^\Delta(S)$ is $||$ defn]
Image inverse de la "diagonale" de
$\Gamma(H_S/S)$ (ie de l'image de $H \to \Gamma(H_S/S)$)
par l'hom. canonique $\underline{\mathbb{S}}(S) \to H_S(S)$,
de sorte que par la construction des
loc. cit. on aura un iso. canonique

(11) $$\underline{\mathbb{S}}^\Delta(S) \simeq [unclear: crossed out] \Gamma(S, \underline{G}_m \otimes_{\mathbb{Z}} M)$$
$\simeq M \otimes_{\mathbb{Z}} A^*$
[MARGIN: ds $S$ affine, d'anneau $A$]

et une suite exacte can.
(12) $$1 \to \Gamma(S, \underline{G}_m \otimes_{\mathbb{Z}} M) \to \underline{\mathbb{S}}^\Delta(S) \to H \to 1$$
[under $\Gamma(\dots)$ is $A^* \otimes_{\mathbb{Z}} M$]

Prenant $A = \hat{\mathbb{Z}}$, on trouve en parti-
culier (et ceci comme un fait
général sur les schémas en groupes
associés à des extensions de groupes
discrets à noyau abélien)
(13) $$\underline{\mathbb{S}}^\wedge \simeq \underline{\mathbb{S}}^\Delta(\hat{\mathbb{Z}})$$

\newpage

%% ===== Page 36 =====

2. Le groupe $Z_{\rho, \sigma}$ : version provisoire.

Je vais maintenant expliciter pour $\hat{\underline{S}}$ et $\tilde{\underline{S}} \to \hat{\underline{S}}$ les constructions générales du § 49 VI, n° / pour $\underline{S}$.

(14)
$$ \underline{S}^\Delta \overset{df}{=} \underline{M} = \hat{\underline{S}}^\circ \quad (\text{composante neutre de } \hat{\underline{S}})$$

Dans un cas comme celui-ci, le $\underline{S}^\Delta$ est justement le sous-groupe fermé du $(\rho, \sigma)$-groupe $\underline{S}_{\rho, \sigma}$ ($= \hat{\underline{S}}$ ici) engendré par $\lambda_{\dots}, \lambda_{\dots}$, on se dispense dans les notations - là où il n'y a pas ambiguïté - d'inclure la notation $\underline{S}$. Par là il y a lieu de distinguer p. ex. entre $C_{\rho, \sigma}$ et $C_{\rho, \sigma}^\Delta$. Nous ferons cependant les constructions dans le cadre restreint ci-dessus, plutôt que propre.

[MARGIN: [unclear: Je ferai le fin des $\dots$ dans $\dots$ cadre restreint, les $\dots$ possibles sans exiger qu'on se limite au $\dots$ groupes.]]

Le pas essentiel est de déterminer le groupe $Z_{\rho, \sigma}$

$$ \underline{Z}_{\rho, \sigma} = \underline{Aut}(\underline{S}^+) $$
$\overset{df}{=}$

[MARGIN: ici, "$\exists \alpha$" "$\exists \beta$" signifie bien sûr l'existence locale de $\alpha,\beta$...]

(15)
$$ \{ u \in \underline{Aut}(\underline{S}^+) \left| \begin{array}{l} a) \; \exists \alpha \in \hat{\underline{S}}_{\hat{\mathbb{Z}}} = \{1, z\}.\underline{M}, \quad u(\rho) = \alpha(\rho) \\ b) \; \exists \beta \in \dots \quad \quad u(\sigma) = \beta(\sigma) \\ c) \; u(\tilde{\underline{L}}_0^\circ) = \tilde{\underline{L}}_0^\circ \\ \quad (\iff u(\underline{L}_0^\circ) = \underline{L}_0^\circ) \\ [ d) \; (\text{automatique}) \quad u(\underline{S}^\Delta) \subset \underline{S}^\Delta ] \end{array} \right\} $$

Notons que $u$ est déterminé par les valeurs

(16)
$$ \begin{cases} u(\rho) = \xi \rho \\ u(\sigma) = \eta \sigma \end{cases} \quad \xi, \eta \in \Gamma \underline{\Pi} $$
où $\underline{\Pi} = \underline{M} \dots$
($\dots$ sous-gr. d'indice 2 dans $\underline{\Pi}$)

satisfaisant les conditions suivantes ;

\newpage

%% ===== Page 37 =====

(17) 
$$ \begin{cases} 
a) & \xi, \rho \text{ (localement) conjugués : } \rho \text{ via } \underline{S} \hookrightarrow \underline{M} \\ 
b) & \eta, \sigma \text{ -- -- -- -- : } \sigma \text{ -- -- -- } \\ 
c) & \text{équation des lacets} \\ 
& \xi = \eta h_1^\nu \quad \text{où } \nu \in \Gamma \underline{M} \text{ convenable} \\ 
& \quad \quad \quad \quad \text{(tel que } \mu = 2\nu+1 \in \Gamma_M \text{ )} 
\end{cases} $$

[crossed out: [unclear] il faudra]

Remarque. Il n'est plus possible ici de désigner
par $-1$ l'élément non trivial du centre $w$ de
$\widehat{\underline{S}}^\wedge$, et de désigner par $-x$ l'op. $x \in \widehat{\underline{S}}^\wedge$ le
produit de $x$ par celui-ci, à cause de la
confusion avec la notation $-x$ (pour l'inverse $x \mapsto -x$ dans $\underline{M}$)
du groupe abélien $M$ pour l'inverse de $x$ --
mais d'ailleurs $w x \neq -x$, puisque $w x \notin M$, $-x \in M$...

On pose bien sûr

(18) $\varepsilon_3(u) = \varepsilon_{\underline{S}}(\alpha)$ , $\varepsilon_2(u) = \varepsilon_{\underline{S}}(\beta)$

(bien défini, car $\rho$ et $\rho^{-1}$ ne sont conjugués
sur aucune fibre, idem pour $\sigma, \sigma^{-1}$ ),
comment les expliciter en termes de
$\xi, \eta$ ? On aura

(19)
$$ \begin{cases} 
\xi = \alpha \rho(\alpha)^{-1} \begin{cases} 
\in \Gamma M \text{ si } \alpha \in \Gamma M = \Gamma \underline{S} \text{ i.e. } \varepsilon_3(u) = 1 \\ 
= \underbrace{\alpha \tau}_{\in \Gamma M} \rho(\alpha \tau)^{-1} = \alpha_0 (- \lambda_1^{-1}) \rho(\alpha_0)^{-1} \in \Gamma M \text{ si } \varepsilon_3(u) = -1 \\ 
\quad \quad \quad \quad (\lambda_1^{-1} \in \underline{M}) 
\end{cases} \\ 
\eta = \beta \sigma(\beta)^{-1} \begin{cases} 
\in \Gamma M \text{ si } \beta \in \Gamma M = \Gamma \underline{S} \text{ i.e. } \varepsilon_2(u) = 1 \\ 
= (\underbrace{\beta_0 \tau}_{\in \Gamma M}) \sigma(\beta_0 \tau)^{-1} = w \beta_0 \sigma(\beta_0)^{-1} \in w M \text{ si } \varepsilon_2(u) = -1 
\end{cases} 
\end{cases} $$

donc dans tous les cas

(19 bis) $$ \begin{cases} 
\xi \in \Gamma M \\ 
\eta \in \Gamma(w^{\frac{1-\varepsilon_2(u)}{2}} M) 
\end{cases} $$

\newpage

%% ===== Page 38 =====

675

d'où on tire, par l'équation des lacets (17)

(20) $l_1^\nu = (\omega)_{1}^\nu \in \Gamma(\underline{\omega}^{\nu_2(u)} \underline{M})$ i.e. $\omega^\nu \in \Gamma(\underline{\omega}^{\nu_2(u)} \underline{M})$

On peut l'écrire sous la forme

(21) $\underline{\nu} = 2 \underline{\nu}_1' + \underline{\nu}_2$

où

(22) $\begin{cases} \underline{\nu}_1' : Z_{\hat{\rho}, \sigma} \to \mathbb{G}_a \\ \underline{\nu}_2, \underline{\nu}_3 : Z_{\hat{\rho}, \sigma} \rightrightarrows (\mathbb{Z}/2\mathbb{Z})_{S_0} \end{cases}$

sont des morphismes de schémas bien définis, en même temps que les homomorphismes de groupes

$\mu : Z_{\hat{\rho}, \sigma} \to \mathbb{G}_m, \quad \underline{\varepsilon}_2, \underline{\varepsilon}_3 : Z_{\hat{\rho}, \sigma} \to \{\pm 1\}_{S_0}$

par les formules (21) et

(23) $\begin{cases} \text{[unclear crossed out: } \mu = 2 \underline{\nu}_1' \dots (= 4 \underline{\nu}_1' + 2 \underline{\nu}_2 + 1) \text{]} \\ \underline{\varepsilon}_2 = (-1)^{\underline{\nu}_2}, \quad \underline{\varepsilon}_3 = (-1)^{\underline{\nu}_3} \end{cases}$

On aura donc, en vertu de (21)

(24) $l_1^\nu = l_1^{2\underline{\nu}_1'} l_1^{\underline{\nu}_2} = (\omega_1^2)^{\underline{\nu}_1'} l_1^{\underline{\nu}_2} = \lambda_1^{\underline{\nu}_1'} l_1^{\underline{\nu}_2}$

[unclear crossed out: Comme $l_1 = \lambda_1 \omega$, on s'aperçoit que les éléments $l_1 = a_1 \omega$ où $a_1 \in \underline{M} \dots$]

posant $l_1 = a_1 \omega$, on peut écrire (24) sous la forme

(26) $l_1^\nu = \underbrace{(\lambda_1^{\underline{\nu}_1'} a_1^{\underline{\nu}_2})}_{\in \Gamma \underline{M}} \omega^{\underline{\nu}_2}$

Pour expliciter les $a_i$, il faut les développer en combinaisons linéaires de $\lambda_0, \lambda_1, \lambda_\infty$ (ou de deux d'entre eux, si on veut \dots).

\newpage

%% ===== Page 39 =====

687

(24) $$l_1^\nu = \lambda_1^\nu w^{\nu_2}$$

Posant (cf (19))
(25) $$\begin{cases} \beta = \beta_0 \tau^{\nu_2} \quad (\beta \in \Gamma \underline{M}) \\ \eta_0 = \beta_0 \sigma(\beta_0)^{-1} \end{cases}$$

on aura donc
$$\eta = \beta \sigma(\beta)^{-1} = \beta_0 \sigma(\beta_0)^{-1} w^{\nu_2} = \eta_0 w^{\nu_2}$$
et l'équation des lacets prend la forme d'une équation dans $\underline{M}$
(26) $$\xi = \eta_0 \lambda_1^\nu$$ , i.e. compte tenu de (21)
(26 bis) $\xi, \eta_0$ sont deux sections de $\underline{M}$ qui se déterminent mutuellement (pour $\nu$ donné).

On a ainsi exprimé de façon commode l'équation des lacets, i.e. la condition c) de (17). Il reste à exprimer a) et b), qui se formulent :
(27) $$\begin{cases} a) \; \xi \text{ locale de la forme } \alpha_0 \rho(\alpha_0)^{-1} \; (\nu_3=0) \text{ ou } \alpha_0 \lambda_1 \rho(\alpha_0)^{-1} \; (\nu_3=1) \\ \text{avec } \alpha_0 \in \Gamma \underline{M} \\ b) \; \eta_0 \text{ locale de la forme } \beta_0 \sigma(\beta_0)^{-1} \end{cases}$$

(Constatons que la formulation est indépendante de $\nu_1, \nu_2$ [crossed out], mais dépend [unclear: précisément] de $\nu_3$ ). Pour expliciter ces conditions, on va écrire $\underline{M}$ additivement, les relations s'écrivent

[MARGIN: On a de même posé $\alpha = \alpha_0 \tau^{\nu_3}$]

(28) $$\begin{cases} \xi = \alpha_0 - \rho(\alpha_0) \quad (\nu_3=0) \quad \text{ou} \quad \xi = \alpha_0 - \rho(\alpha_0) - \lambda_1 \quad (\nu_3=1) \\ \eta_0 = \beta_0 - \sigma(\beta_0) \end{cases}$$

\newpage

%% ===== Page 40 =====

ce qui s'exprime aussi par
$$
(29) \quad \begin{cases} \zeta \in \operatorname{Im}(Id_{\underline{M}} - \rho_{\underline{M}}) \quad (\nu_3=0) \quad \text{ou} \quad (\zeta - \lambda_1 \in \operatorname{Im}(Id_{\underline{M}} - \rho_{\underline{M}})) \\ \eta_0 \in \operatorname{Im}(Id_{\underline{M}} - \sigma_{\underline{M}}) \end{cases}
$$

C'est le moment de passer à des calculs matriciels explicites, en identifiant $\underline{M}$ à la base $\lambda_0, \lambda_1$. Dans cette base, les opérations de $\rho_{\underline{M}}$, $\sigma_{\underline{M}}$, $\tau_{\underline{M}}$ sont donnés par

$$
(30) \quad \begin{cases} \rho_{\underline{M}} = \begin{pmatrix} 0 & -1 \\ -1 & -1 \end{pmatrix} & id_{\underline{M}} - \rho_{\underline{M}} = \begin{pmatrix} 1 & 1 \\ 1 & 2 \end{pmatrix} \\ \sigma_{\underline{M}} = \begin{pmatrix} 0 & 1 \\ 1 & 0 \end{pmatrix} & id_{\underline{M}} - \sigma_{\underline{M}} = \begin{pmatrix} 1 & -1 \\ -1 & 1 \end{pmatrix} \\ \tau_{\underline{M}} = \begin{pmatrix} -1 & 0 \\ 0 & -1 \end{pmatrix} \end{cases}
$$

Bien on a
$$det(id_{\underline{M}} - \rho_{\underline{M}}) = 3, \quad det(id_{\underline{M}} - \sigma_{\underline{M}}) = 0$$
ce qui montre que les endomorphismes $id_{\underline{M}} - \rho_{\underline{M}}$ et $id_{\underline{M}} - \sigma_{\underline{M}}$ de $\underline{M}$ ne sont [unclear: l'un ni l'autre] des automorphismes
[MARGIN: i.e. des isomorphismes à des [unclear: isogénies] près, de [unclear: degré] 3 pour le premier, [unclear: nul] pour le second.]
On aura donc plus intérêt à prendre comme "paramètres" non $\zeta$ et $\eta_0$, mais bien carrément $\alpha_0$ et $\beta_0$ (en plus des "paramètres" discrets $\nu_2, \nu_3$ et $\nu'$), et d'éliminer donc $\zeta, \eta_0$ (si possible...) avec les systèmes

\newpage

%% ===== Page 41 =====

(31) $\widetilde{Z}_{1,\sigma} = \{ (\lambda_1, \nu_2, \nu_3, \alpha_0, \beta_0) \mid \lambda_1 \in \Gamma \underline{G}_a, \nu_2, \nu_3 \in \Gamma \mathcal{O}_{S_0}, \alpha_0, \beta_0 \in \Gamma \underline{M}$
$\qquad \alpha_0 - \rho(\alpha_0) = \nu_3 \lambda_1 = \beta_0 - \sigma(\beta_0) + \nu_2 \lambda_1 \}$

[MARGIN: NB L'élimination de $\nu_1$ qui est redondant]

L'équation des lacets, en $\alpha_0, \beta_0$, prend alors la forme

(32) $(id - \rho_{\underline{M}})\alpha_0 = (id - \sigma_{\underline{M}})\beta_0 = (\nu_2 + \nu_3)\lambda_1$

Remarque. Le schéma $\widetilde{Z}_{1,\sigma}$ lui-même, tel qu'il
était défini au §49, n'est pas lisse
il s'en faut - se cassant en car. 2 et 3.
Par contre le schéma $\widetilde{Z}_{1,\sigma}$, plus fin,
introduit dans le cours des
[unclear: calculs] précédents (à la suite de $\xi_0, \eta$)
semble bien lisse. Je suis amené à
modifier un peu la présentation des
constructions du §49, en substituant
au groupe $\widetilde{Z}_{1,\sigma}$ (lui-même para-
métré par $\xi_0, \eta$), et en le remplaçant par
un groupe qui s'approche d'avantage
des autom. de $L_0$ (ou de $L_0^\#$) de $\mathcal{M}_{1,\sigma}$
(ou $\mathcal{M}_{1,\sigma}^\sim$). On fera une présentation dans
ce sens dans la suite (§ suiv. ?), et
supprime ici déjà l'esquisse de cette construction,

[MARGIN: Cela tient au fait que l'hom. $\underline{M} \times \underline{M} \to \underline{M}$, $(\alpha_0, \beta_0) \mapsto (id - \rho_{\underline{M}})\alpha_0 - (id - \sigma_{\underline{M}})\beta_0$ n'est bien s. pas un épimorphisme. Mais c'en est un si $3 \neq 0$ !]

\newpage

%% ===== Page 42 =====

D'où le \underline{schéma en groupes} $\widetilde{Z}_{g,\nu}$ muni d'une structure de groupe [unclear: additive], et opère sur $\underline{S}^+$ via

$$ (33) \left\{ \begin{array}{l} u(\rho) = (\alpha_0 - \rho(\alpha_0) + \nu_3 \lambda_1)\rho \\ u(\sigma) = (\beta_0 - \sigma(\beta_0) ) \omega^{\nu_2} \sigma \end{array} \right. $$

Mais il faut quand même vérifier que les formules (33) déterminent bel et bien un automorphisme $u$ de $\underline{S}^+$. La vérification essentielle à faire me semble que la relation (7)

(34) $\quad [\sigma, \lambda_\infty]$

est respectée par $u$.

Je vais expliciter l'action de $u$ sur $\underline{M}$, en explicitant ses valeurs sur

$$ \left\{ \begin{array}{l} \lambda_0 = \omega l_0 = \omega \sigma \rho \sigma \rho \\ \lambda_1 = \omega l_1 = \omega \rho \sigma \rho \sigma \end{array} \right. $$

en termes de $\xi_0, \eta_0$, où

$$ (35) \quad \left\{ \begin{array}{l} \xi_0 = \alpha_0 - \rho(\alpha_0) + \nu_3 \lambda_1 \\ \eta_0 = \beta_0 - \sigma(\beta_0) \end{array} \right. \quad \in \Gamma \underline{M} $$

Soit :

(36) $\quad u \rho = \xi_0 \rho$, $\quad u \sigma = \omega^{\nu_2} \eta_0 \sigma$.

[unclear: On trouve donc] pour la valeur des lacets

$u(\varepsilon_0) = \xi_0^{(\mu)}$ $\quad$ $(\mu = 2\nu_1+1 = 4\nu' + 2\nu_2 + 1)$

d'où

\newpage

%% ===== Page 43 =====

$$u(l_0) = u(\varepsilon_0^2) = l_0^\mu$$
et
$$u(\lambda_0) = u(\omega l_0) = \underbrace{u(\omega)}_{= \omega} u(l_0) = \omega l_0^\mu = \omega l_0 \frac{(l_0^\mu)}{l_0} = \lambda_0 \lambda_0^{\mu-1} = \lambda_0^\mu$$
comme il convient. (Pour me rassurer, j'ai vérifié
que la deuxième formule (33) donne bien
$$u(\omega) = (y_0 \sigma)(y_0 \sigma) = (y_0 \sigma(y_0)) \omega = \omega \quad \text{i.e.}$$
(37) $$u(\omega) = \omega$$
On aura
$$u(\lambda_1) = u(\rho(\lambda_0)) = [unclear: (Int \rho u)](\lambda_0) = [unclear: (Int \rho)](u \lambda_0) = [unclear: (Int \rho)](\lambda_0^\mu) = \lambda_1^\mu$$
donc finalement, comme prévisible

[MARGIN: $\mu = 2\nu+1 = 4\nu'+2\nu_2+1$ (38)]
(38) $$u(\lambda) = \mu.\lambda \quad \text{pour } \lambda \in \Gamma \underline{M}$$

On aura bien que $$u(\sigma) = (\beta_0 - \sigma(\beta_0)) \omega^{\nu_2} \sigma$$
commute avec $$u(\lambda_0) = \mu.\lambda_0$$

Tous ces calculs sont un peu vaseux,
logiquement, car à la fois [unclear: pour] prouver
l'existence d'un [unclear: prolongement] dans les
calculs précédents je l'admets déjà (implicite-
ment), dans un groupe "plus gros", savoir le
schéma en groupes [unclear: (caractéristiques)] libre engendré par $\rho, \sigma$...
Une façon de les rendre rigoureux serait
d'utiliser la description intrinsèque de $\underline{\tilde{G}}^+$ :
(39) $$\underline{\tilde{G}}^+ = \underline{\tilde{G}}^+_{S_0} *_{\underline{M}_{S_0}} \underline{M}_{S_0}$$
donc définir un hom. de $\underline{\tilde{G}}^+$ dans un
$S_0$-groupe $g$ (ici $g = \underline{\tilde{G}}^+$) revient à

\newpage

%% ===== Page 44 =====

définir des homomorphismes
(40) $$M_{S_0} \xrightarrow{u_M} g^\natural, \quad \underline{\mathfrak{S}}^+_{S_0} \xrightarrow{u_{\underline{\mathfrak{S}}}} g^\natural \quad (\text{i.e. } \underline{\mathfrak{S}} \to g^{c.s.})$$
satisfaisant les deux conditions suivantes :

a) Commutativité du diagramme

(41)
\begin{tikzcd}
M_{S_0} \ar[r] \ar[d] & \underline{\mathfrak{S}}^+_{S_0} \ar[d, "u_{\underline{\mathfrak{S}}}"] \\
M_{S_0} \ar[r, "u_M"'] & g
\end{tikzcd}

et b)

(42) $u_M$ est compatible avec les actions
de $\underline{\mathfrak{S}}^+_{S_0}$ sur $M_{S_0}$ (via son action sur $M_{S_0}$)
et sur $g$ (via $u_{\underline{\mathfrak{S}}}$).

Dans le cas naturel, on aura $g = \underline{\mathfrak{S}}^+$

(43) $u_M(\lambda) = i(\mu . \lambda)$

où $i : M \longrightarrow \underline{\mathfrak{S}}^+$
est l'inclusion évidente, et

(44) $u_{\underline{\mathfrak{S}}} : \underline{\mathfrak{S}}^+_{S_0} \longrightarrow g \quad \text{défini par (33)}$

Il faudra d'abord vérifier que (33)
définit bien un hom. $\underline{\mathfrak{S}}^+_{S_0} \longrightarrow g$,
ce qui revient à $= \underline{\mathfrak{S}} \to g^{c.s.}$.

\newpage

%% ===== Page 45 =====

Pour [unclear: utiliser] la présentation [crossed out text] de $\underline{\mathcal{S}}^+$
[unclear crossed out equations]

(44) $$\underline{\mathcal{S}}^+ = \left\{ \rho, \sigma \mid \rho^6 = \sigma^4 = 1, \rho^3 = \sigma^2, [\sigma, \lambda_0] = 1 \text{ où } \lambda_0 = \rho \sigma \rho^{-1} \sigma \rho \right\},$$

en posant
(45) $\rho' = \tilde{u}(\rho), \sigma' = \tilde{u}(\sigma) \quad \text{(cf (33))}$
[unclear crossed out text]
il faut vérifier les relations
(46) $\rho'^6 = \sigma'^4 = 1, [\sigma', \lambda'_0] = 1, \rho'^3 = \sigma'^2 \quad \text{avec} \quad \lambda'_0 = \rho' \sigma' \rho'^{-1} \sigma' \rho'$

On commence par vérifier que
(47) $$\rho'^3 = \sigma'^2 = \omega$$

ce qui implique [crossed out text] $\exists \underline{\mathcal{S}}^+ \xrightarrow{u_{\underline{\mathcal{S}}}} \mathcal{G}_{es}$
(48) $\rho'^6 = \sigma'^4 = 1$, donc $u_{\underline{\mathcal{S}}} : \rho \mapsto \rho', \sigma \mapsto \sigma'$

Ensuite, prenant note que
$$\lambda'_0 = u_{\underline{\mathcal{S}}}(\lambda_0),$$
il suffira en effet de vérifier
(49) $$\lambda'_0 = \mu \lambda_0$$
Il sera évident que $\lambda'_0$ commute à :
$$\sigma' = u(\sigma) = y_0 \sigma \quad \text{[unclear]}$$
puisque $\sigma$ et $y_0 \in M$ y commutent. D'où
pour construire $u_{\underline{\mathcal{S}}}$, les vérifications
essentielles sont (47) et (49) - d'ailleurs
une fois (47) prouvée, les calculs précédents
aboutissant à (38) se justifient (au moins
là où ils ont un sens, i.e. pour $u_{\underline{\mathcal{S}}}$ et $\lambda \in M$),

\newpage

%% ===== Page 46 =====

donc (49) sera bien vérifiée.

On a 
$$\rho'^3 = (\xi \rho(\xi) \rho^2(\xi)) \rho^3 = (\xi \rho(\xi) \rho^2(\xi)) \omega$$
$$\sigma'^2 = (\eta_0 \sigma(\eta_0)) \sigma^2 = (\eta_0 \sigma(\eta_0)) \omega$$
et les formules (47) se réduisent à :

(50) $$ \xi \rho(\xi) \rho^2(\xi) = 1 , \quad \eta_0 \sigma(\eta_0) = 1 $$

soit, additivement

$$ \xi + \rho(\xi) + \rho^2(\xi) = 0, \quad \eta_0 + \sigma(\eta_0) = 0 $$

dont la vérification est triviale, compte tenu que

$$ \rho_M^3 = \sigma_M^2 = id_M , \quad (1 + \rho_M + \rho_M^2) = \text{[unclear]} 0_M . $$

On a donc bien construit un hom.

$$ u_{\mathfrak{S}_0} : \widehat{\mathfrak{S}}^+ \to \widehat{\mathfrak{S}}^+(S_0) , $$

et il reste à vérifier (41) et (42).

Pour (41), nous l'avons déjà noté au niveau de l'existence d'un $u_{\mathfrak{S}_0}$ - cela revient à prouver :

( $$ u_{\mathfrak{S}_0}(\lambda_0) = \text{[unclear]} \lambda_0 , \quad u_{\mathfrak{S}_0}(\lambda_1) = \mu \lambda_1 $$

je vais quand même refaire la vérification, pour être sûr.

$$ \varepsilon_0' = u_{\mathfrak{S}_0}(\varepsilon_0) = u_{\mathfrak{S}_0}(\sigma) u_{\mathfrak{S}_0}(\rho) = \eta_0 \sigma \xi \rho = \eta_0 {}^{\sigma}\xi \sigma \rho $$
$$ = \eta_0 \sigma(\xi) \varepsilon_0 $$

\newpage

%% ===== Page 47 =====

Je dis que
$y'_0 = \varepsilon_0^{(M-1)}$ i.e. $y_0 \sigma(y) \omega^{V_2} = \underbrace{\varepsilon_0^{(\mu-1)}}_{l_0^{\nu'}} \implies \mu-1=2\nu_1$, soit
$y_0 \sigma(y) \omega^{V_2} = l_0^{\nu'}$ ; $\nu = 2\nu'$ i.e.
$y_0 \sigma(y) = l_0^{2\nu'} (\omega l_0 \lambda_0^{V_2}) = \underbrace{\lambda_0^{2\nu'+V_2}}_{\lambda_0^{2\nu'} \lambda_0^{-V_2}}$ i.e.
$y_0 \sigma(y) = \lambda_0^{\nu'}$ soit, en appliquant $\sigma$ et notant
que $\sigma(y_0) = y_0^{-1}$, $\sigma^2(y)=y$, $\sigma(\lambda_0)=\lambda_1$,
$y_0^{-1} y = \lambda_1(y)$ soit enfin $y = y_0 \lambda_1(y)$ $\sim$
$y = y_0 + \nu \lambda_1$ (écrit additivement)

ce qui est vrai par hypothèse (cf (31) ).
Reste à [unclear: s'assurer] (pour achever [unclear: l'étude] de $u_{\mathcal{E}_0}$ pour
les [unclear: éléments annulés] (pour [unclear: abréger on l'écrira simplement])
$u(\lambda_1) = u\rho(\lambda_0) = (\rho \mu)(\lambda_0) = \rho(u(\lambda_0)) = \rho(\lambda_0 \mu) = \lambda_1 \mu$
OK.

Reste à vérifier (42), ce qui [unclear: se réduit] à :
commutativité de $u_M = \mu\ id_M$ à
$(\rho_M, \rho'_M)$, $(\sigma_M, \sigma'_M)$ - ce qui signifie que
$\rho_M$ et $\rho'_M$ d'une part, $\sigma_M$ et $\sigma'_M$ d'autre part,
agissent [unclear: également], ce qui est trivialement évident.

Ainsi, le groupe $\tilde{\mathcal{E}}_{\rho, \sigma}$ est bel et
bien construit - du moins, en tant que
relevant un [unclear: sous-groupe] $Z$, muni des [unclear: morphismes]
de structure
(51)
\begin{tikzcd}
\tilde{Z}_{\rho, \sigma} \ar[r, shift left, "f_0"] \ar[r, shift right, "f_{S_0}"'] & \mathcal{E}_{0, S_0} \ar[r] & \mathcal{E}_{S_0}
\end{tikzcd}
d'où $\mu = 4\nu' + 2 + V_2 : \tilde{Z}_{\rho, \sigma} \to \mathcal{E}_{S_0}$
$\varepsilon_2, \varepsilon_5 : \tilde{Z}_{\rho, \sigma} \rightrightarrows \mu_{S_0} = \mu\ id_{S_0}$

\newpage

%% ===== Page 48 =====

et un morphisme de faisceaux sur $S_0$

(52) $$\tilde{Z}_{\rho,\sigma} \longrightarrow \underline{Aut}_{gr} (\underline{\mathbb{G}}^+)$$
$$u \longmapsto u$$
donnant lieu à un morphisme

(53) $$\tilde{Z}_{\rho,\sigma} \longrightarrow \underline{Aut}_{gr} (\underline{\mathbb{G}}^+) \times \mathbb{G}_m \times \mu_{S_0} \times \mu_{S_0}$$
$$u \longmapsto u, \mu(u), \varepsilon_2(u), \varepsilon_3(u)$$

- des calculs déjà faits tant de fois
dans des contextes divers donnent assurance
que ce morphisme (qui se ramène dans
cas. # 3,3) est un morphisme de groupes,
pour une structure de groupe convenable
sur $\tilde{Z}_{\rho,\sigma}$, uniquement déterminée par elle
(ex. un # 1,3, et coulant par voie
de descente galoisienne). On aura

[MARGIN: un peu conjecturé car $\tilde{Z}_{\rho,\sigma} \to \mathbb{G}_m \times \mu_{S_0} \times \mu_{S_0}$ n'est pas un morphisme de gr. de $\underline{M} \times \underline{M}$ par [unclear] (d'où la [unclear] $\alpha_0 + (1-\sigma_M)\beta_0$ [unclear] vient [unclear] brouillon [unclear] voir p. (55))]

une suite exacte
$$1 \longrightarrow \tilde{Z}_{\rho,\sigma}^{+0} \longrightarrow \tilde{Z}_{\rho,\sigma} \longrightarrow \mathbb{G}_m \times \mu_{S_0} \times \mu_{S_0} \longrightarrow 1$$
explicite
$$\tilde{Z}_{\rho,\sigma}^{+0} = \text{sous-schéma en groupes}$$
$$= \underline{M} \times \underline{M} \times \mathbb{G}_a$$
$$\{ (\alpha_0, \beta_0, \gamma') \mid \begin{array}{l} \alpha_0 \dots = 0 \\ (1+\rho_M)\alpha_0 = (1-\sigma_M)\beta_0 + 2\gamma'\lambda_1 \end{array} \}$$

On aura donc

(55) $$1 \longrightarrow \tilde{Z}_{\rho,F}^{+0} \longrightarrow \tilde{Z}_{\rho,\sigma}^{+0} \longrightarrow \mathbb{G}_a \longrightarrow 1$$

\newpage

%% ===== Page 49 =====

072

$$ (57) \quad \widetilde{\mathbb{Z}}_{\rho,\sigma}^{+00} \overset{d\acute{e}f}{=} \left\{ \underline{u} \in \Gamma \widetilde{\mathbb{Z}}_{\rho,\sigma} \mid \nu'(u) = \nu_0(u) = \nu_1(u) = 0 \right\} $$
$$ \simeq $$
$$ \left\{ (\alpha_0, \beta_0) \in \Gamma(\underline{M} \times \underline{M}) \mid (1 - \rho_M)\alpha_0 = (1 - \sigma_M)\beta_0 \right\} $$

[MARGIN: en car 3 le [unclear] de $2 \lambda_1' = 3!$ relation et a de [unclear]]

Le schéma en groupes $\widetilde{\mathbb{Z}}_{\rho,\sigma}^{+00}$ est lisse
en car. $\neq 3$
et [unclear: sans] relation $2$ - il s'en déduit de façon canonique
un isomorphisme [unclear: $\underline{M}$-module],
$$ (58) \quad \widetilde{\mathbb{Z}}_{\rho,\sigma}^{+00} \simeq \mathbb{V}(\underline{T}_{\rho}) $$
où $\underline{P}$ est le $\mathbb{Z}$-module lisse de rang $2$
défini par la suite exacte
$$ (59) \quad 0 \longrightarrow \underline{M} \longrightarrow \underline{M} \times \underline{M} \longrightarrow \underline{M} \longrightarrow 0 $$
$$ (\alpha_0, \beta_0) \longmapsto (1-\rho_M)\alpha_0 - (1-\sigma_M)\beta_0 $$

L'hom. $\underline{M} \times \underline{M} \longrightarrow \underline{M}$ s'explicite, en
posant
$$ (60) \quad \begin{cases} \alpha_0 = a_0 \lambda_0 + a_1 \lambda_1 & d'o\grave{u} \quad \alpha_0 - \rho \alpha_0 = (a_0 + a_1)\lambda_0 + (2a_1 - a_0)\lambda_1 \\ \beta_0 = b_0 \lambda_0 + b_1 \lambda_1 & d'o\grave{u} \quad \beta_0 - \sigma \beta_0 = (b_0 - b_1)\lambda_0 + (b_1 - b_0)\lambda_1 \end{cases} $$

comme
$$ (61) \quad (\alpha_0, \beta_0) \longmapsto (\underbrace{a_0 + a_1 + (b_0 - b_1)}_{c_0})\lambda_0 + (2a_1 - a_0 - (b_1 - b_0))\lambda_1 $$
C'est un hom. surjectif [unclear],
les éléments de l'image sont les
$$ (62) \quad \{ c_0 \lambda_0 + c_1 \lambda_1 \mid \underbrace{c_0 + c_1}_{= 3a_1} \equiv 0 \quad (3) \} $$

Remarque. Si on prend un $\underline{u}$ défini sur $\hat{\mathbb{Z}}$
provenant d'un élément de $M[0]$, alors

\newpage

%% ===== Page 50 =====

[crossed out: on aura fait]

$\mu \in \hat{\mathbb{Z}}^*$, $\mu \equiv (-1)^{\nu_3} \equiv 1-2\nu_3 \pmod{3}$
et la condition sur $(\mu, \alpha_0, \beta_0)$ pour être
dans l'image de $\tilde{\tilde{Z}}_{\rho,\sigma} \to \mathbb{G}_m \times p_{S_0} \times p_{S_0}$,
qui s'écrit en général (vu relation 162), §2)
$$2\nu' - \nu_2 + \nu_3 \equiv 0 \ (3)$$
devient, en multipliant par 2
$$4\nu' - 2\nu_2 + 2\nu_3 \equiv 0 \ (3)$$
i.e. $\mu-1$
$$\mu \equiv 1-2\nu_3 \ (3)$$
et elle est automatiquement satisfaite.

C'est pourquoi j'ai envie de simplifier
la définition de $\tilde{\tilde{Z}}_{\rho,\sigma}$, en définissant

[crossed out block: 
$$ \tilde{\tilde{Z}}_{\rho,\sigma} = \left\{ (\nu',\nu'',\nu_2,\nu_3,\alpha_0,\beta_0) \in \hat{\mathbb{Z}} \times \hat{\mathbb{Z}} \times \{0,1\}_{S_0} \times \{0,1\}_{S_0} \times M \times M \left| \begin{array}{l} \mu = 4\nu' - 2\nu_2 + 1 = 6\nu'' - 2\nu_3 + 1, \quad 2\nu' + \nu_3 - \nu_2 = 3\nu'' \\ (1-\rho_M)(\alpha_0) - (1-\sigma_M)(\beta_0) = 3\nu'' \lambda_1 \end{array} \right. \right\} $$
]

[MARGIN: 
en expr $\nu', \nu''$ par :
$\mu \equiv 1-2\nu_2 \ (4)$
$\mu \equiv 1-2\nu_3 \ (6)$
$\mu = 1-2\nu_2 + 4\nu'$
$\mu = 1-2\nu_3 + 6\nu''$
$-2\nu_2 + 2\nu' = -\nu_3 + 3\nu''$ i.e.
$2\nu' + \nu_3 - \nu_2 = 3\nu''$
]

ici
$2c = 6\nu''$, et on en fait
$c = 3\nu''$, on aura

(63) $$ \tilde{\tilde{Z}}_{\rho,\sigma} \simeq \mathbb{G}_m \times \hat{\mathbb{Z}} \times \{0,1\}_{S_0} \times \{0,1\}_{S_0} \times M \times M $$
$$ \left\{ (\mu,\nu'',\nu_2,\nu_3,\alpha_0,\beta_0) \left| \begin{array}{l} a) \text{ [crossed out] } \\ b) \mu \in \hat{\mathbb{Z}}^* \\ c) \text{ équation dans } M \\ (1-\rho_M)(\alpha_0) - (1-\sigma_M)(\beta_0) = 3\nu'' \lambda_1 \end{array} \right. \right\} $$

\newpage

%% ===== Page 51 =====

3. Le groupe [unclear: de Galois] $\mathbb{O}_N^*$.

On va introduire une construction auxiliaire,
en définissant un schéma en groupes qui
correspond à la limite des quantités
$\mu \in \mathbb{G}_m, \nu', \nu'' \in \hat{\mathbb{Z}}, \nu_2, \nu_3 \in \hat{\mathbb{Z}}|_{S_0}$, satis-
faisant

$$(63) \quad \mu = 3\nu' - 2\nu_2 + 1 = 6\nu'' - 2\nu_3 + 1 \quad , \quad 2\nu' - \nu_2 = 3\nu'' - \nu_3$$

De façon générale, soit $N \in \mathbb{N}^*$ un entier fixé
(le cas qui nous intéresse pour la suite est
$N=12$), et considérons le schéma sur $\mathbb{Z}$ associé
à l'extension de groupes discrets

$$(64) \quad 1 \to \mathbb{Z} \xrightarrow{N} \mathbb{Z} \to \mathbb{Z}/N\mathbb{Z} \to 1$$

qui s'inscrit dans une suite exacte

$$(65) \quad 1 \to E^1 \to \mathbb{O}_N \to (\mathbb{Z}/N\mathbb{Z})_{S_0} \to 1$$

Rappelons la [unclear: dite construction] simple

$$(66) \quad \mathbb{O}_N = (\mathbb{Z})_{S_0} \times [unclear: N\mathbb{Z}]$$

Mais je dis qu'il y a sur $\mathbb{O}_N$ une
structure d'anneau unique, dont l'
addition est la structure de groupe déjà
envisagée, et dont la multiplication est
obtenue ainsi : Une section de $\mathbb{O}_N$ est
localement de la forme $n + N[\lambda]$, avec
la règle de calcul

$$(67) \quad \begin{matrix} a) & (n + N[\lambda]) + (n' + N[\lambda']) = (n+n') + N[\lambda+\lambda'] \\ b) & n + N[\lambda] = n' + N[\lambda'] \iff n+N\lambda = n'+N\lambda' \text{ dans } G_0 \end{matrix}$$

\newpage

%% ===== Page 52 =====

On multiplie alors les sections par

(68) $$ (n+N[\lambda])(n'+N[\lambda']) = nn'+N[n\lambda'+n'\lambda+\lambda\lambda'] $$

on vérifie que c'est compatible avec
relations (67 b), que c'est commutatif ,
associatif, que $1+N[0]$ est
unité... Autre façon de voir la
structure d'anneau. Soit plus
généralement $A$ une $\mathbb{Z}$-algèbre
unitaire associative (disons) munie d'un idéal
bilatère $J$ qui soit $\mathbb{Z}$-module libre de t.f., considérons
l'idéal $NA$ et l'extension d'

(69) $$ 0 \to J \to A \to A/J \to 0 $$

d'où un schéma en groupes comm.
noté $A_{[N]}$, qui s'obtient comme
quotient de $A_{S_0} \times J$ par le

semi-direct)

(70) $$ 0 \to J \to A_{S_0} \times J \to A_N \to 0 $$
$$ \lambda \mapsto (\lambda, - \lambda) $$

et on trouve que $A_{S_0} \times J$ a une structure
d'Anneau produit semi-direct, et que

\newpage

%% ===== Page 53 =====

l'image de $\underline{J}$ [MARGIN: 67a] est un idéal ...
On trouve bien une structure d'anneau
sur $\underline{A}_{[N]}$, et une suite exacte de

(71) $$0 \to \underline{J} \to \underline{A}_{[N]} \to \underline{B}_{S_0} \to 0$$

où $\underline{A}_{[N]} \to \underline{B}_{S_0}$ est un hom. de
schémas en anneaux. Passant aux
groupes multiplicatifs, on trouve
une suite exacte

(72) $$1 \to (1+\underline{J})^* \to \underline{A}_{[N]}^* \to (\underline{B}_{S_0})^* \to 1$$
$$||$$
$$(\underline{B}^*)_{S_0}$$

où pour $\underline{A}_{[N]}^*$, $(\underline{B}_{S_0})^*$ et $\underline{B}^*$, le
$*$ désigne le foncteur [unclear: sous]-groupe multiplicatif
et où $(1+\underline{J})^*$ a le sens évident

(73) $$(1+\underline{J})^* \subset 1+\underline{J}$$
$$||$$
$$\{ (1+\lambda) \mid 1+\lambda \text{ inv. dans } \underline{A} \}$$
$$\simeq$$
$$\{ (\lambda,\lambda') \in \underline{J} \times \underline{J} \mid \lambda+\lambda'+\lambda\lambda'=0 \}$$
dans $\underline{J}$

Ainsi on aura, dans le cas particulier
$A = \mathbb{Z}$, $J = N\mathbb{Z}$
(74) $$1 \to (1+N\mathbb{Z})^* \to \mathbb{O}_N^* \to (\mathbb{Z}/N\mathbb{Z})_{S_0}^* \to 1$$
$\simeq \{1\}$

\newpage

%% ===== Page 54 =====

$(1+N\mathbb{Z})^*$ est défini comme

(75) $(1+N\mathbb{Z})^* \stackrel{déf}{=} \{ \lambda \in B_0 \mid 1+N\lambda \text{ inversible} \}$

[texte barré: Or on a un hom. canonique]

Les $\mathcal{O}_N$ pour $N$ variable forment
un syst. projectif d'Anneaux,
les $\mathcal{O}_N^*$ forment un syst. projectif
de schémas en groupes commutatifs.

Je ne vais pas regarder de plus
près la structure de ces systèmes
projectifs, ni examiner la limite
projective - question qui se pose
plus généralement pour un
$A$ comme ci-dessus. On va
expliciter les homomorphismes canoniques

(76) $\mathcal{O}_N^* \to \mathcal{O}_1^* = \mathcal{O}^* = G_m$

[texte barré: dont le noyau est ... qui s'identifie au foncteur ... le noyau est formé ... du système des]

(77) $\{ 1+N\lambda \mid 1+N\lambda=1 \text{ i.e. } N\lambda=0 \}$

qui est donc un sous schéma clos de $G_a$, isomorphe à ${}_N G_a$,

\newpage

%% ===== Page 55 =====

qui est un schéma en groupes [unclear: fini plat]
(si $N \neq 1$) qui est nul sur $\operatorname{Spec}(\mathbb{Z}[1/N])$,
et dont la restriction à la fibre en
un $p | N$ est $\simeq \mathbb{G}_a$.

Je me rends compte maintenant
que c'est une maladresse de sortir des
groupes $\mathbb{O}_N^\ast$, qui m'a fait un tas
d'ennuis, pour [unclear: s'y reporter] : $\mathbb{G}_m$.

Je voudrais en particulier examiner

(78) $$1 \longrightarrow (1+12\mathbb{Z})^\ast \longrightarrow \mathbb{O}_{12}^\ast \longrightarrow (\mathbb{Z}/12\mathbb{Z})^\ast_{S_0} \longrightarrow 1 .$$

Notons que
$$(\mathbb{Z}/12\mathbb{Z})^\ast \simeq (\mathbb{Z}/4\mathbb{Z})^\ast \times (\mathbb{Z}/3\mathbb{Z})^\ast = \mu \times \mu .$$

Donc (78) prend la forme
[crossed out: de trouver une section semblant à l'évaluation]

(79) $$1 \longrightarrow (1+12\mathbb{Z})^\ast \longrightarrow \mathbb{O}_{12}^\ast \xrightarrow{\varepsilon_2 \times \varepsilon_3} \mu \times \mu \longrightarrow 1$$

Les quantités $\varepsilon_2, \varepsilon_3$ sont des applications [crossed out: définies] régulières
sur $\mathbb{O}_{12}^\ast$

(80) $$\mathbb{O}_{12}^\ast \xrightarrow{(v_2, v_3)} \{0,1\}_{S_0} \times \{0,1\}_{S_0}$$

Je ne sais [unclear: au juste]

(81) $$\mathbb{O}_{12}^\ast \xrightarrow{(\varepsilon_2, \varepsilon_3)} \mu \times \mu \simeq \{0,1\}_{S_0} \times \{0,1\}_{S_0}$$
déduit de $1 \mapsto 0$, $-1 \mapsto 1$

\newpage

%% ===== Page 56 =====

que
et on [unclear: munit] celles ci de [unclear: la] loi de groupe pour l'add-
-ition mod 2 sur $\{0,1\}_{S_0}$ . On a ainsi
$$ \varepsilon_2 = (-1)^{\nu_2} = 1 - 2\nu_2 \ , \ \varepsilon_3 = (-1)^{\nu_3} = 1 - 2\nu_3 $$
Les "quatre choix de signes" de $(\mathbb{Z}/12\mathbb{Z})^*$
se décrivent dans le tableau

(82)
$$
\begin{array}{c|c|c}
\varepsilon_2 \backslash \varepsilon_3 & \begin{matrix} +1 \\ \text{i.e. } \nu_3=0 \end{matrix} & \begin{matrix} -1 \\ \text{i.e. } \nu_3=1 \end{matrix} \\
\hline
\begin{matrix} +1 \\ \text{i.e. } \nu_2=0 \end{matrix} & 1 & +5 \\
\hline
\begin{matrix} -1 \\ \text{i.e. } \nu_2=1 \end{matrix} & -5 & -1
\end{matrix}
$$

donc on trouve une section ensembliste
de (79)
en prenant les sections de $\mathcal{O}_{12}^*$ de la forme

(83) $$ \mu_0 = (-5)^{\nu_2} (+5)^{\nu_3} \begin{cases} \equiv (-1)^{\nu_3} & (3) \\ \equiv (-1)^{\nu_2} & (4) \end{cases} $$

Donc une section générale de $\mathcal{O}_{12}^*$ s'écrit

(84) $$ \mu = \underbrace{(-5)^{\nu_2} (+5)^{\nu_3}}_{\mu_0} + 12 [\text{unclear}] $$
$$ \left\{ \begin{array}{l} (\nu_2, \nu_3) \in \Gamma(\{0,1\}_{S_0} \times \{0,1\}_{S_0}) \\ \lambda_0 \in \Gamma(\mathbb{G}_a) \end{matrix} \right. $$

[crossed out: Notons d'autre part que 4]

$$ \mu = (4 + 1 \dots $$
$$ 5^{\nu_3} = 1 + 4\nu_3 $$

[MARGIN: NB $\nu_2 = \nu_2^2$ et]

$$ (-5)^{\nu_2} = [(-1)^{\nu_2} 5^{\nu_2}] = (1 - 2\nu_2)(1 + 4\nu_2) $$
$$ = 1 + 2\nu_2 - 8\nu_2^2 = 1 - 6\nu_2 $$

donc

(85) $$ \mu_0 = (1 + 4\nu_3)(1 - 6\nu_2) \text{ [crossed out terms]} $$
[crossed out: $= 1 + 2\nu_2 + (-8\nu_2^2 + 4\nu_3 (1 + 2\nu_2 - 8\nu_2^2))$]
$$ = 1 + 4\nu_3 - 6\nu_2 - 24\nu_2\nu_3 $$

[crossed out: On peut prendre $\mu_0' = 1 + 4\nu_3 - 6\nu_2$]

\newpage

%% ===== Page 57 =====

Si on [unclear] ( (84) et (85) )
$$ \mu = \mu_0 + 12 [\nu_0] = 1 + 4 \nu_3 - 6 \nu_2 + 12 (\nu_0 - 2 \nu_2 \nu_3) $$

on constate bien s'exprimer par [unclear: des] congruences
mod 4)
(86) $$ \mu = 1 - 2 \nu_2 + 4 \underbrace{[ \nu_3 - \nu_2 - 3 \nu_2 \nu_3 + 3 \nu_0 ]}_{\nu'} $$
et par congruences mod 6
(87) $$ \mu = 1 - 2 \nu_3 + 6 \underbrace{[ \nu_3 - \nu_2 - 2 \nu_2 \nu_3 + 2 \nu_0 ]}_{\nu''} $$

où $\nu'$ et $\nu''$ soulignés dans [unclear: la formule] désignent la section [MARGIN: (correspondante)]
de $\mathbb{O}_{12}^*$, le [unclear] s'exprimant la [unclear] section
de $\mathbb{G}_m = \mathbb{O}_1^*$. On aura bien
$$ 2 \nu' - \nu_2 = 2 \nu_3 - 3 \nu_2 - 6 \nu_2 \nu_3 + 6 \nu_0 $$
$$ 3 \nu'' - \nu_3 = 3 \nu_3 - 3 \nu_2 - 6 \nu_2 \nu_3 + 6 \nu_0 $$

En résumé [MARGIN: le schema en groupes lisse]

Proposition Introduisons sur $\mathbb{O}_{12}^*$ les $\mathbb{G}_a =$
fonctions $\nu_2, \nu_3$ (à valeurs dans $\{0,1\}_{S_0} \subset \mathbb{F}_2$)
$\nu', \nu''$ (à valeurs dans $\mathbb{G}_a$) et $\mu$ (à valeurs
dans $\mathbb{G}_m = \mathbb{O}_1^*$) définies par les conditions
$\mu =$ im. pr. can. de $\mu$ de $\mathbb{O}_1^*$ par
l'hom. de groupes $\mathbb{O}_N^* \to \mathbb{O}_1^*$

(88) $\begin{cases} (-1)^{\nu_2} = \text{im. de ce } \mu \text{ dans } (\mathbb{Z}/4\mathbb{Z})^*_{S_0} \\ (-1)^{\nu_3} = \text{ " " " } \mu \text{ " } (\mathbb{Z}/6\mathbb{Z})^*_{S_0} \end{cases}$

d'où
(89) $\begin{cases} \mu = (-5)^{\nu_2} (+5)^{\nu_3} + 12 [\nu_0] \\ = (1+4\nu_3)(1-6\nu_2) + 12 \nu_0 \\ = 1 + 4\nu_3 - 6\nu_2 + 12 [ \nu_0 - 2\nu_2\nu_3 ] \end{cases}$

pour une section $\nu_0$ bien déterminée de $\mathbb{G}_a = \mathbb{O}_1$

\newpage

%% ===== Page 58 =====

Posons de plus

(90)
$$ \begin{cases} \nu' = \nu_3 - \nu_2 - 3 \nu_2 \nu_3 + 3 \nu_0 \\ \nu'' = \nu_3 - \nu_2 - 2 \nu_2 \nu_3 + 2 \nu_0 \end{cases} $$

On a alors les relations

[MARGIN: où $\varepsilon_p = (-1)^{\nu_p}$
d'où $\varepsilon_p = 1 - 2\nu_p$
$(p \in \{2,3\})$]

(91)
$$ \begin{cases} \mu = 1 - 2 \nu_2 + 4 \nu' & (= \varepsilon_2 + 4 \nu') \\ \mu = 1 - 2 \nu_3 + 6 \nu'' & (= \varepsilon_3 + 6 \nu'') \\ -\nu_2 + 2 \nu' = -\nu_3 + 3 \nu'' & \text{i.e. } 2\nu' - \nu_2 + \nu_3 = 3 \nu'' \end{cases} $$

NB La condition $\mu = 1$ équivaut à la condition

(92) $$ \nu_2 = \nu_3 = \nu_0 = 0 $$

et implique

(93) $$ \nu = \nu' = \nu'' = 0 $$

\newpage

%% ===== Page 59 =====

67)

4. Le groupe $\widetilde{Z}_{\bar{\rho}, \sigma}$ [crossed out text] et la suite exacte 
$$ 1 \to N_{\bar{\rho}, \sigma} \to \widetilde{N}_{\bar{\rho}, \sigma} \to \mathbb{O}_{12}^* \to 1 . $$
[unclear: Rappelons] qu'il se définit comme plongé dans
$$ \mathbb{O}_{12}^* \times M \times M \quad : $$

$$ (94) \quad \widetilde{Z}_{\bar{\rho}, \sigma} \subset \mathbb{O}_{12}^* \times M \times M $$
$$ \{ (\mu, \alpha_0, \beta_0) \mid (1 - \rho_M)(\alpha_0) - (1 - \sigma_M)(\beta_0) = 3 \nu''(\mu)\lambda_1 \} $$

Notons qu'on a
$$ (95) \quad (1 - \rho_M)(\lambda_1 - \lambda_0) = 3 \lambda_1 $$
donc la relation du cocycle dans (94)
prend aussi la forme
$$ (96) \quad (1 - \rho_M)(\alpha_0 - \nu''(\lambda_1 - \lambda_0)) - (1 - \sigma_M)(\beta_0) = 0 $$

Je vais admettre [unclear: qu'on a] une
structure de groupe naturelle sur $\widetilde{Z}_{\bar{\rho}, \sigma}$ ,
pour laquelle la projection $\widetilde{Z}_{\bar{\rho}, \sigma} \to \mathbb{O}_{12}^*$
est un hom. de groupes. On trouve
maintenant une suite exacte

$$ (97) \quad 1 \to \widetilde{Z}_{\bar{\rho}, \sigma}^{+0} \to \widetilde{Z}_{\bar{\rho}, \sigma} \to \mathbb{O}_{12}^* \to 1 $$
où
$$ (98) \quad \widetilde{Z}_{\bar{\rho}, \sigma}^{+0} = \operatorname{Ker}(\widetilde{Z}_{\bar{\rho}, \sigma} \to \mathbb{O}_{12}^*) $$
$$ \simeq $$
$$ \{ (\alpha_0, \beta_0) \mid (1 - \rho_M)(\alpha_0) - (1 - \sigma_M)(\beta_0) = 0 \} $$

Ce schéma en groupes est lisse [unclear: en car.]
$\neq 3$, mais pas en car. 3, où sa [unclear: dim. pure]

\newpage

%% ===== Page 60 =====

de la dim. relative $2$ : la dim relative $3$.
Ce fait d'une fibre exceptionelle de dim.
plus grande ne peut, il me semble, être
éliminé par une entourloupette, [unclear] en
introduisant des définitions ad-hoc pour
un "meilleur" groupe $\tilde{Z}_{\hat{\rho}, \sigma}$ : il y a
"plus" d'automorphismes de $\widehat{\mathcal{C}}_k$ (pour $k$
un corps alg. clos donné), ayant les propriétés
envisagées, si $k$ est de car. $3$, que dans les
autres cas de caractéristique.

Cependant il me semble qu'on peut
construire un sous-schéma en groupes de
$\tilde{Z}_{\hat{\rho}, \sigma}$, que j'ai envie de noter
$N_{\hat{\rho}, \sigma}$, qui est lisse sur $S = \operatorname{Spec}(\mathbb{Z})$ et
coïncide avec $\tilde{Z}_{\hat{\rho}, \sigma}$ au dessus de $S_0$,
$S_0 = Spec \mathbb{Z}[1/6] = S \setminus \{2, 3\}$. Pour le découvrir
je note que les équations de bonté, (96)
la première non-tautologique
dans le sous-module de $M$ correspondant
à la congruence
$$c_0 + c_1 \equiv 0 \ (3) \quad (cf \ (62) \text{ [unclear]}$$
et si on introduit le sous-groupe [unclear]

\newpage

%% ===== Page 61 =====

) correspondant $M_0$ de $M$

(99)
$$ M_0 \subset M $$
$$ \| defn $$
$$ \{ c_0 \lambda_0 + c_1 \lambda_1 \mid c_0 + c_1 \equiv 0 (3) \} , $$

d'où un fibré vectoriel $\underline{M}_0$ sur $S_{0,1}$ —
et en fait un homomorphisme de schémas

(100)
$$ \widetilde{Z}_{p,\sigma} \longrightarrow \underline{M}_0 $$
$$ (\mu, \alpha_0, \beta_0) \longmapsto (1-\rho_M)(\alpha_0 + \gamma''(\lambda_1-\lambda_0)) + (1-\sigma_M)(\beta_0) $$

qu'on peut expliciter ainsi : on a un
hom. de $Z$-modules

(101)
$$ M \times M \times Z \longrightarrow M_0 $$
$$ (\alpha_0, \beta_0, \gamma'') \longmapsto (1-\rho_M)(\alpha_0 + \gamma''(\lambda_1-\lambda_0)) + (1-\sigma_M)(\beta_0) $$

d'où un hom. pour les schémas en
modules correspondants

(102)
$$ \underline{M} \times \underline{M} \times \underset{\substack{\| \\ \underline{G}_a}}{\underline{Z}} \longrightarrow \underline{M}_0 $$

qu'on compose avec

(103)
$$ \widetilde{Z}_{p,\sigma} \longrightarrow \underline{M} \times \underline{M} \times \underline{Z} $$
$$ (\mu, \alpha_0, \beta_0) \longmapsto (\alpha_0, \beta_0, \gamma''(\mu)) $$

On peut alors considérer l'équation des
cocycles (96) comme équation dans
$\underline{M}_0$ — elle devient ainsi plus stricte,

\newpage

%% ===== Page 62 =====

d'où un sous-schéma fermé de $\tilde{Z}_{\rho, \sigma}$,

(104) $$N_{\rho, \sigma} = $$
$$\{ (\mu, \alpha_0, \beta_0) \in \tilde{Z}_{\rho, \sigma} \mid \text{l'équation des cocycles (96) valable dans } \underline{M}_0 \}$$

[MARGIN: NB
Le groupe $N_{\rho, \sigma}$ devrait pouvoir se noter aussi $M_{\rho, \sigma}[\Gamma^{(0)}]$, sauf erreur ou plutôt ce serait $(M_{\rho, \sigma}/M_{\rho})[\Gamma^{(0)}]$ d'après les notations du §9.]

Considérons alors le morphisme
$$N_{\rho, \sigma} \longrightarrow \underline{O}_{12}^*$$
il est clair que ce morphisme est un épi, la fibre d'un point $s = (\nu_1, \nu_2, \nu_0)$ étant formée des $(\alpha_0, \beta_0) \in \Gamma \underline{M} \times \underline{M}$ satisfaisant (96) dans $\underline{M}_0$, [unclear: c'est à dire des] $(\tilde{\alpha}_0, \beta_0)$ satisfaisant l'équation linéaire dans $\underline{M}_0$

(105) $$(1 - \rho_M)(\tilde{\alpha}_0) - (1 - \sigma_M)(\beta_0) = 0$$

Or l'hom. de Modules
$$M \times M \longrightarrow M_0$$
(106) $$(\tilde{\alpha}_0, \beta_0) \longmapsto (1 - \rho_{\underline{M}})(\tilde{\alpha}_0) - (1 - \sigma_{\underline{M}})(\beta_0)$$
est un épi. [unclear]

$$\underline{M} \times \underline{M} \longrightarrow \underline{M}_0$$
correspondant qui s'en déduit est un épi.
Donc $N_{\rho, \sigma}$ est lisse sur $S_0$
et c'est un sous-schéma fermé de $\tilde{Z}_{\rho, \sigma}$

\newpage

%% ===== Page 63 =====

qui coïncide sur lui-dessus de $S_0 \setminus \{3\}$, il
s'ensuit que c'est l'adhérence schématique

(107) $$N_{\rho,\sigma} = \widetilde{Z}_{\rho,\sigma} | (Spec \mathbb{Z} \setminus \{3\})$$
(adh. schématique)

ce qui implique que c'est un
sous-schéma en groupes.

C'est peut-être le moment de regarder
d'un peu plus près le sous-$\mathbb{Z}$-module
$M_0$ de $M$,

(108) $$M_0 = \mathbb{Z}(\lambda_1 - \lambda_0) + \mathbb{Z} \rho_M(\lambda_1 - \lambda_0)$$
$$= \mathbb{Z} \underbrace{(\lambda_1 - \lambda_0)}_{\delta} + \mathbb{Z} \rho_M \underbrace{(\lambda_1 - \lambda_0)}_{\delta}$$
$$= \mathbb{Z} \delta \oplus \mathbb{Z} (1-\rho_M)(\delta)$$
$$= \mathbb{Z} \delta + \mathbb{Z} \rho_M \delta$$

Notons que $M_0$ est stable par $\rho$, de façon
générale par construction, pour tout $x \in M$, on a

(109) $x \equiv \rho_M(x) \pmod{M_0}$ \hspace{1cm} (i.e. $(1-\rho_M)(x) \in M_0$)
et $x \equiv \sigma_M(x) \pmod{M_0}$ \hspace{1cm} ($(1-\sigma_M)(x) \in M_0$)

Il est bien sûr stable aussi par $\tau_M = -id_M$.

Donc $M_0$ est invariant par $\widetilde{\mathfrak{S}}$,
au titre que sous-groupe de $M$ ($\widetilde{\mathfrak{S}}$ y opère [unclear]).
C'est un sous-groupe distingué dans $\widetilde{\mathfrak{S}} \ltimes M$,
groupe quotient [unclear] de $\mathfrak{S}_0 = \operatorname{Aut}(\mathbb{Z}, \mathbb{Z})$,

\newpage

%% ===== Page 64 =====

contenant $[\Pi_0, \Pi_0]$, d'indice $3$ dans $\Pi_0$.

Je présume qu'il contient le sous-groupe
de congruences relatif à $12$ - i.e. distingué,
qu'il correspond à un sous-groupe de
$S\ell(2, \mathbb{Z}/12\mathbb{Z})$. Il est d'indice
$3 \cdot 12 = 36$ dans $S\ell(2, \mathbb{Z})$, [unclear: rayé] $G\ell(2, \mathbb{Z})$,
tandis que $$S\ell(2, \mathbb{Z}/12\mathbb{Z}) \simeq \underset{48}{S\ell(2, \mathbb{Z}/4\mathbb{Z})} \times \underset{24}{S\ell(2, \mathbb{Z}/3\mathbb{Z})}$$
il s'agit donc d'un sous-groupe (c'est même un sous-groupe distingué)
$32$ de $S\ell(2, \mathbb{Z}/12\mathbb{Z})$, engendré par les
éléments images de
$$ (\sigma(\rho)\rho^{-1})^3 = [\lambda_0, \lambda_1], \text{ ou } [\sigma_0, \lambda_\infty] $$
cf (7)
$\delta = \lambda_1 \lambda_0^{-1}$
$\rho(\delta)$

ou encore le sous-groupe distingué engendré par
le seul élément dans $S\ell(2, \mathbb{Z}/12\mathbb{Z})$

$$ (110) \quad \left\{ \begin{array}{l} \delta = \lambda_1 \lambda_0^{-1} = l_1 l_0^{-1} = \left(\begin{array}{cc} 1 & 0 \\ 2 & 1 \end{array}\right) \left(\begin{array}{cc} 1 & +2 \\ 0 & 1 \end{array}\right) \\ = \left(\begin{array}{cc} 1 & +2 \\ 2 & 5 \end{array}\right) \end{array} \right. $$

La question est de si ce sous-groupe
distingué de $S\ell(2, \mathbb{Z}/12\mathbb{Z})$ est bel et
bien d'ordre $32$, donc d'indice $36$ - un
ordre qui pourrait être plus grand, donc l'indice
plus petit. Déjà quel est l'ordre de $\bar{\delta}$ ?

\newpage

%% ===== Page 65 =====

mod 12 ? On a

$$ \delta^2 = \begin{pmatrix} 1 & 2 \\ 2 & 5 \end{pmatrix} \begin{pmatrix} 1 & 2 \\ 2 & 5 \end{pmatrix} = \begin{pmatrix} 5 & 12 \\ 12 & 29 \end{pmatrix} \equiv \begin{pmatrix} 5 & 0 \\ 0 & 5 \end{pmatrix} \quad (12) $$

d'où $\delta^4 \equiv 1 \quad (12)$ ,

c'est déjà un bon point ! On a d'ailleurs
$$ (112) \qquad \begin{cases} \delta \equiv 1 \quad (2) \\ \delta^2 \equiv 1 \quad (4) \end{cases} \quad \delta \neq 1 \quad (4), \delta \neq 1 \quad (3) $$
$$ \delta^2 \equiv -1 \quad (3) $$

Considérons
$$ (113) \qquad \operatorname{Sl}(2, \mathbb{Z}/12\mathbb{Z}) \simeq \underbrace{\operatorname{Sl}(2, \mathbb{Z}/4\mathbb{Z})}_{48} \times \underbrace{\operatorname{Sl}(2, \mathbb{Z}/3\mathbb{Z})}_{24} \xrightarrow{u \times v} \underbrace{\mathfrak{S}_4^+/\Pi_0}_{12} \times \underbrace{\mathfrak{S}_3^+}_{3} $$

où

$u : \operatorname{Sl}(2, \mathbb{Z}/4\mathbb{Z}) \longrightarrow \mathfrak{S}_4^+/\Pi_0$

est l'hom. canonique surjectif, dont le noyau
est d'ordre 4 - les trois éléments
non nuls du noyau sont justes
$\lambda_0, \lambda_1, \lambda_2$ permutés entre eux par $\rho$, et
$\lambda_1 \lambda_0^{-1} = \lambda_1 \lambda_0 = \lambda_2$ etc... ici les $\lambda_i$ sont
notés pour qu'ils soient permutés par $\rho$ - et

$v : \operatorname{Sl}(2, \mathbb{Z}/3\mathbb{Z}) \longrightarrow \mathfrak{S}_3^+$

est l'hom. composé
$$ \operatorname{Sl}(2, \mathbb{Z}/3\mathbb{Z}) \longrightarrow \mathfrak{S}_4 \xrightarrow{\text{épi}} \mathfrak{S}_3^+ $$
sur $\mathbb{P}^1(\mathbb{Z}/3\mathbb{Z})$

dont le noyau est d'ordre 8. Le noyau
de $u \times v$ est d'ordre $4 \cdot 8 = 32$, et il contient
évidemment $\delta$. Reste à m'assurer que c'est bien
le sous-groupe distingué engendré par $\delta$ (dont l'ordre
a priori était [unclear: un sous-multiple] de 32), et

\newpage

%% ===== Page 66 =====

ce sous-groupe image de $\Pi_0 \subset \operatorname{Sl}(2,\hat{\mathbb{Z}})$, c'est-à-dire dans $\Pi_0$,
que l'image inverse d'un sous-groupe de
l'ét. précédente de $\operatorname{Sl}(2, \mathbb{Z}/12\mathbb{Z})$.

Résumons les résultats obtenus :

Proposition. Considérons dans $\hat{\mathfrak{S}}_0^+ = \operatorname{Sl}(2, \hat{\mathbb{Z}})$
la tour de sous-groupes invariants (et dans $\hat{\mathfrak{S}}_0$)
(114)
\begin{tikzcd}
\hat{\mathfrak{S}}_0^+ \ar[r, hookleftarrow, "6"] & \Pi = \Pi[2] \ar[r, hookleftarrow, "2"] \ar[d, equal, "\text{def}"'] & \Pi_0 \ar[r, hookleftarrow, "3"] \ar[d, phantom, "\text{sous-groupe N}"'] & \Pi_{00} \ar[r, hookleftarrow, "2!"] \ar[d, phantom, "\text{sous-groupe inv.}"'] & (\Pi_0, \Pi_0) \\
& \{u \in \hat{\mathfrak{S}}_0^+ \mid u \equiv 1 (2)\} & \text{engendré par } \rho, \dots & \text{engendré par les} & \\
& & \text{[unclear]} & \lambda_1 \lambda_0^{-1} = \rho_1 \rho_0^{-1} &
\end{tikzcd}

Considérons d'autre part le groupe quotient
(115) $$\operatorname{Sl}(2, \mathbb{Z}/12\mathbb{Z}) \simeq \hat{\mathfrak{S}}_0^+ / \Pi[12]$$
$$ \simeq \operatorname{Sl}(2, \mathbb{Z}/4\mathbb{Z}) \times \operatorname{Sl}(2, \mathbb{Z}/3\mathbb{Z}) $$

et la tour de sous-groupes suivants

[MARGIN: NB $(\hat{\mathfrak{S}}_0^+)^{ab} \simeq \mathbb{Z}/12\mathbb{Z}$ (cf 42 bis)
l'identif... en $H/\Pi_{00}^\bullet$ donne un homom (116)
d'ordre 3 $-$ surj. d'un ? en un
$\operatorname{Sl}(2, \mathbb{Z}/4\mathbb{Z}) / \operatorname{Sl}(2, \mathbb{F}_2)' \times \operatorname{Sl}(2, \mathbb{Z}/3\mathbb{Z}) / Syl_{2,3} \simeq \hat{\mathfrak{S}}_0^+ / \Pi_{00}$]

\begin{tikzcd}
H = \ar[d, phantom, "\cup" description, "6" left] & \operatorname{Sl}(2, \mathbb{Z}/4\mathbb{Z}) \times \operatorname{Sl}(2, \mathbb{Z}/3\mathbb{Z}) & \text{ordre } 48 \times 24 \\
\Pi^\bullet = \ar[d, phantom, "\cup" description, "2" left] & \operatorname{Sl}(2, \mathbb{F}_2) \times \operatorname{Sl}(2, \mathbb{Z}/3\mathbb{Z}) & \text{ordre } 8 \times 24 \\
\Pi_0^\bullet = \ar[d, phantom, "\cup" description, "3" left] & \operatorname{Sl}(2, \mathbb{F}_2)' \times \operatorname{Sl}(2, \mathbb{Z}/3\mathbb{Z}) & \text{ordre } 4 \times 24 \\
\Pi_{00}^\bullet = & \operatorname{Sl}(2, \mathbb{F}_2)' \times Syl_{2,3} \ar[d, "\simeq"'] & \text{ordre } 4 \times 8 \\
& \text{[unclear]} \simeq D_4 &
\end{tikzcd}

(117)
a) $\operatorname{Sl}(2, \mathbb{F}_2) \simeq \operatorname{Ker}(\operatorname{Sl}(2, \mathbb{Z}/4\mathbb{Z}) \to \operatorname{Sl}(2, \mathbb{Z}/2\mathbb{Z}))$
b) $\operatorname{Sl}(2, \mathbb{F}_2)' \subset \operatorname{Sl}(2, \mathbb{F}_2)$, sous-groupe distingué $\simeq \mathbb{F}_2 \times \mathbb{F}_2$ dont les trois él. $\neq 0$ sont $\lambda_0, \lambda_1, \lambda_\infty$ (images dans $\operatorname{Sl}(2, \mathbb{Z}/4\mathbb{Z})$) - c'est un supplémentaire dans la [unclear] du sous-groupe multiplicatif $\mu_2 \simeq \mathbb{F}_2$ des unit. [unclear] de $\operatorname{Sl}(2, \mathbb{Z}/4\mathbb{Z})$.
c) $Syl_2 \subset \operatorname{Sl}(2, \mathbb{Z}/3\mathbb{Z})$, sous-groupe de Sylow, d'ordre 8 (qui est distingué)

\newpage

%% ===== Page 67 =====

681

Ceci posé, la tour de sous-groupes (114) se
plonge isom- dans la tour (116).

Par acquis de conscience, je vais reprendre
l'esquisse de démonstration précédente. On
note que l'hom. $\bar{G}_0^+ \to H$ applique
les groupes figurant dans (114) dans
ceux de (116) - pour $\pi$ et $\pi_0$, on sait
d'ailleurs depuis un moment que ce
sont les images isom. de $\pi', \pi_0'$ !
(pour ceci, il suffirait d'examiner
$\operatorname{Sl}(2, \mathbb{Z}/4\mathbb{Z})$, au lieu de $\operatorname{Sl}(2, \mathbb{Z}/12\mathbb{Z})$...).
Ainsi l'image isom de $\pi_{00}$ dans $\bar{G}_0^+$
contient $\pi_{00}^\circ$. Comme les deux sont d'
indice 36, ces deux groupes sont égaux,
cqfd. (le même argument pourrait s'appliquer
pour $\pi_0, \pi_0'$...)

[MARGIN: NB / Mais $\bar{G}_0^+ \to \operatorname{Sl}(2, \mathbb{Z}/12\mathbb{Z})$ / se factorise par [unclear] / car s'il n'en était / pas ainsi, le / quotient de $\pi_0 / [\pi_0, \pi_0]$ / qui serait [unclear] / aurait deux [unclear] / l'un [unclear] de l'autre]

Corollaire L'homomorphisme canonique
$$ \bar{G}_0^+ = \operatorname{Gl}(2, \mathbb{Z}) \longrightarrow \operatorname{Gl}(2, \mathbb{Z}/12\mathbb{Z}) / (\pi_{00}^\circ \simeq Fl(2, \mathbb{F}_2)' \times Syl_3) $$
se factorise par $\bar{G} = \bar{G}_0^+ / [\pi_0, \pi_0]$
étudié dans le présent § 5, et l'image
inverse de $\pi_{00}^\circ$ n'est autre que le sous-
groupe $M_0$ de (99) (108).

\newpage

%% ===== Page 68 =====

Les considérations précédentes amènent
naturell^t la considération des revête-
[MARGIN: c'est ce revête^t $M_{1,1}$] ments d'ordre 3 de $U_{0,3 \mathbb{Q}}$ correspon-
dant au sous-groupe $\Pi_{00}$ de $\Pi_0 \simeq \Pi_{1}(U_{0,3 \overline{\mathbb{Q}}})^{\wedge}$.
On vérifie tout de suite que c'est
une courbe de genre 1 - si l'on
- [unclear] de travailler sur
$[\Pi_{00}, \Pi_{00}]$ au lieu de $[\Pi_0, \Pi_0]$, on
trouve un groupe extension de
$\Pi_0 / [\Pi_{00}, \Pi_{00}] \simeq M$, par un groupe de rang 2
[MARGIN: c'est une courbe elliptique, d'invariants 0 (courbe de Fermat dirait-on)] que j'appelle "elliptique", et sur
lequel $M$ va opérer de façon
certainement non triviale - par la suite
il me faudrait examiner cette action,
et la mettre en relation avec l'action
qui se trouve, subordonnée à
l'hom. $\operatorname{Sl}(2, \mathbb{Z})^{\wedge} \longrightarrow \operatorname{Sl}(2, \mathbb{Z})$.

\newpage

%% ===== Page 69 =====

5) Les groupes $\underline{M}_{\rho,\sigma}$, $\underline{M}_{\rho,\sigma}^+$ (via $\underline{N}_{\rho,\sigma}$).

À tout $l \in \hat{\Pi}_0$ correspond un autom
intérieur $int(l)$, qui est l'identité sur
$\underline{M}_0$. On aura, tautologiquement

(118) $$ \begin{cases} int(l)(\rho) = \alpha(\rho) \\ int(l)(\sigma) = \beta(\sigma) \end{cases} \quad \alpha = \beta = l $$

Mais [unclear] (rel. de lacets C relative : $f_1=1$ d. $\nu=0$)
$$ (1-\rho_M)(\underset{l}{\alpha}) = (1-\sigma_M)(\underset{l}{\beta}) $$
i.e.
$$ \rho_M(l) = \sigma_M(l) \quad ? $$

Posant 
(119) $$ l = c_0 \lambda_0 + c_1 \lambda_1 $$
la relation [unclear] s'écrit
$$ c_0 \lambda_1 + c_1 \underset{-\lambda_0-\lambda_1}{\lambda_\infty} = c_0 \lambda_1 + c_1 \lambda_0 $$
i.e.
$$ - c_1 \lambda_0 + (c_0 - c_1) \lambda_1 = c_1 \lambda_0 + c_0 \lambda_1 $$
i.e. $c_1=0$ i.e. $l = c_0 \lambda_0$. Donc c'est
sans remords qu'on se restreint à
considérer des $l \in \widehat{L}_0$ ; correspondant

(120) $$ l = c_0 \lambda_0 \in \widehat{L}_0 = Im (\hat{\mathbb{Z}} \xrightarrow{u} \underline{M}) $$
$$ c_0 \mapsto c_0 \lambda_0 $$

On a un demi diagramme
(121) 
\begin{tikzcd}
\mathbb{G}_a \simeq \widehat{L}_0 \ar[r] \ar[d, hook] & \underline{N}_{\rho,\sigma} \\
\underline{G}^+
\end{tikzcd}

\newpage

%% ===== Page 70 =====

l'image de
(122) $$L_0^{\underline{M}_0} \to \underline{N}_{\rho,\sigma}$$
est invariante - l'hom (122) étant en
fait compatible avec l'action de $\underline{N}_{\rho,\sigma}$
- via $p_r$ sur $L_0^{\underline{M}_0} \simeq \tilde{G}_0$ (i.e. induits
par les opérations de $\underline{N}_{\rho,\sigma}$ sur $\tilde{G}^+ \supset \underline{M}_a \to L_0^{\underline{M}}$
et opérations de [unclear] sur $\tilde{G}^+$ déduites
de (121) sont identiques. Donc d'après
le topos général on trouve un
schéma en groupes $\underline{\underline{M}}_{\rho,\sigma}$ et un
diagramme de suites exactes df $\underline{\Gamma}_{\rho,\sigma} = \underline{N}_{\rho,\sigma}/L_0^{\underline{M}}$

[MARGIN: avec les notati. du $a$ $§ 99$ $\underline{VI}$, le gr. noté $\underline{\underline{M}}_{\rho,\sigma}$ serait sans doute noté $\underline{M}_{\rho,\sigma}/H_p$ ; ici $\underline{M}_{\rho,\sigma}$ serait $\underline{M}_{\rho,\sigma}$ (loc. cit. 8°)]

(123)
\begin{tikzcd}
1 \ar[r] & L_0^{\underline{M}} \ar[r] \ar[d] & \underline{N}_{\rho,\sigma} \ar[r] \ar[d] & \underline{\Gamma}_{\rho,\sigma} \ar[r] \ar[d, "\simeq"'] & 1 \\
1 \ar[r] & \tilde{G}^{+\Delta} \ar[r] & \underline{\underline{M}}_{\rho,\sigma} \ar[r] & \underline{\Gamma}_{\rho,\sigma} \ar[r] & 1
\end{tikzcd}

Remplaçant $\tilde{G}^+$ par $\underline{M}$ , on trouve un $\underline{M}_{\rho,\sigma}$ , et
(124)
\begin{tikzcd}
1 \ar[r] & L_0^{\underline{M}} \ar[r] \ar[d] & \underline{N}_{\rho,\sigma} \ar[r] \ar[d] & \underline{\Gamma}_{\rho,\sigma} \ar[r] \ar[d, "\simeq"'] & 1 \\
1 \ar[r] & \underline{M} \ar[r] & \underline{M}_{\rho,\sigma} \ar[r] & \underline{\Gamma}_{\rho,\sigma} \ar[r] & 1
\end{tikzcd}

D'ailleurs on a dans (124)
dans (123) (comparer $§ 99$ $\underline{VI}$ 5° ...), le
(125) $$M_{\rho,\sigma} \subset \underline{M}_{\rho,\sigma} \dots$$

\newpage

%% ===== Page 71 =====

69bis
Il faut préciser la structure de $\underline{\Gamma}_{\rho,\sigma}$,
en notant que $\underline{L}_0^M \subset \underline{N}_{\rho,\sigma}^{o+}$, donc

(126) $$1 \to \underline{\Gamma}_{\rho,\sigma}^{o+} \to \underline{\Gamma}_{\rho,\sigma} \to \mathfrak{S}_2 \to 1$$

(127) $$1 \to \underline{L}_0^M \to \underline{N}_{\rho,\sigma}^{o+} \to \underline{\Gamma}_{\rho,\sigma}^{o+} \to 1$$

Considérons le noyau de l'hom. surj. de $\mathbb{Z}$-modules

(128) $$\underline{M} \times \underline{M} \to \underline{M}$$
$$(\alpha_0, \beta_0) \mapsto (1-\rho_M)(\alpha_0) - (1-\sigma_M)(\beta_0)$$

noyau formé des $(\alpha_0, \beta_0)$ tels que
$c_0 = a_0 + a_1 + (b_1 - b_0) = 0$
$c_1 = 2a_1 - a_0 - [unclear: 4b_0 - 3b_1] = 0 \quad | \text{ cf } (61)$

Soit $\underline{N}^{o+}$ ce noyau, de sorte qu'on a

(129) $$\underline{N}_{\rho,\sigma}^{o+} \simeq \underline{N}^{o+}$$

Considérons l'hom.

(130) $$\mathbb{Z} \to \underline{N}^{o+}$$
$$n \mapsto (n\lambda_0, n\lambda_0)$$

l'image est un [unclear: sous-groupe pur], car formé des
systèmes $(a_0, a_1, b_0, b_1)$ tels que
$a_1 = b_1 = 0, a_0 = b_0,$

$$0 \to \mathbb{Z} \to \underline{N}^{o+} \to \mathbb{Z}^3$$
$$(a_0, a_1, b_0, b_1) \mapsto (a_1, b_1, a_0 - b_0)$$

donc $\underline{N}^{o+} / \mathbb{Z} \simeq \mathbb{Z}^3$, c. à. d. [unclear: Libre, de Rg 4]

\newpage

%% ===== Page 72 =====

donc une suite exacte de $\mathbb{Z}$-modules

(131)
$$1 \longrightarrow \mathbb{Z} \longrightarrow N^{0+} \longrightarrow \Gamma^{0+} \longrightarrow 0$$
[MARGIN: $n \longmapsto (n\lambda_0,n\lambda_0)$]
[MARGIN: $\simeq N^{0+}/\mathbb{Z}$]

avec $\Gamma^{0+}$ $\mathbb{Z}$-module libre de rg 1.

On en conclut

(132)
$$\Gamma_{p,\sigma}^{0+} \simeq \underline{\Gamma}^{0+} \quad (\simeq \mathbb{G}_a)$$

Ainsi, tous les schémas en groupes
qu'on vient de construire sont lisses
sur $Spec \mathbb{Z}$.

Peut-on définir une section
naturelle de l'extension
[crossed out: $N_{p,\sigma}$ de $\Gamma_{p,\sigma}$ par $L_0^M$, donc des]
tions $\underline{M}_{p,\sigma}$ de $\underline{\Gamma}_{p,\sigma}$ par $\underline{M}$ ?
[MARGIN: lisse de dim relative 2]

Compte tenu de l'équation en
[unclear: les a,b] qui décrit $N_{p,\sigma}$, il semble
qu'une telle section
doive choisir un scindage des
[crossed out: suites exactes] $N^{0+} \longrightarrow M \times M \longrightarrow M_0$
et de

(133)
$$0 \longrightarrow \mathbb{Z} \longrightarrow N^{0+} \longrightarrow \Gamma^{0+} \longrightarrow 0$$

(extensions de $\mathbb{Z}$-modules). Mais il
faut vérifier que ces sections "ensemblistes"
des schémas considérés permettent de définir
[crossed out: des sections] multiplicatives...

\newpage

%% ===== Page 73 =====

5°) $Z_{\rho,\sigma}$ et $G_{\rho,\sigma} = \mathbb{O}_{12}^* G_{12}^+$ + $G_{\rho,\sigma}^+ \simeq \mathbb{O}_{12}^* \cdot \underline{M} = \underline{M}_{\rho,\sigma} \simeq \underline{M}$.
Une fois intr. duit $\mathbb{O}_{12}^*$ et $M_0 \subset M$, il me semble que le groupe $Z_{\rho,\sigma}$ ten lieu et place de $\tilde{Z}_{\rho,\sigma}$, dès l'équa - tion tirant compte de $\alpha, \beta$ et via ses données de $\xi, \eta$ redevient sympathique :

(134)
$$ \underline{Z}_{\rho,\sigma} = \operatorname{Aut}(\underline{\mathcal{S}}^+) \times \mathbb{O}_{12}^* $$
$$ \{ (u, \mu) \mid $$
a) $u(\rho) = \xi \rho$ (loc.) (conjugué à $\rho$, par un $\alpha \in \widehat{G_{12}^+}$ tel que $\chi(\alpha) = \varepsilon_3(\mu)$)
b) $u(\sigma) = \eta \sigma$ (loc.) (conjugué à $\sigma$, par un $\beta \in \widehat{G_{12}^+}$ tel que $\chi(\beta) = \varepsilon_2(\mu)$)
c) équation des bouts $\xi_0 = \eta_0 \nu''(\mu)$ soit $\xi_0 = \eta_0 + 3 \nu'' \lambda_1$ dans $M_0$

où $\nu'' = \nu''(\mu) \in \dots$ (cf pur. p. 685, et ci-dess.!)

(135)
$$ \begin{cases} \xi_0 = [- \nu_3 \lambda_1] \in \Gamma M_0 & \nu_3 = \nu_3(\mu) \\ \eta_0 = [\nu_2 \lambda_2] \in \Gamma M_0 & \nu_2 = \nu_2(\mu) \end{cases} $$

On doit supposer de plus que

(136)
$$ \begin{cases} \xi_0 \text{ locale dans l'im. de } (1 - \rho_M)(M) \\ \eta_0 \text{ locale dans l'im. de } (1 - \sigma_M)(M) \end{cases} $$

et on aura

(137)
$$ \underline{Z}_{\rho,\sigma} \simeq \tilde{\Sigma}_{\rho,\sigma} = \left\{ (\mu, \xi_0, \eta_0) \in \mathbb{O}_{12}^* \times M_0 \times M_0 \;\middle|\; \begin{matrix} 1) \; \xi_0 = \eta_0 + (\nu'' \lambda_1) \text{ dans } M_0 \\ \quad \text{où } \nu'' = \nu''(\mu) \\ 2) \text{ on a (136)} \end{matrix} \right\} $$

$G$

$$ (1 - \rho_M) : M \to M_0 $$
$$ (1 - \sigma_M) : M \to \mathbb{Z} = C_a \to M_0 $$
$$ (\lambda_0, \lambda_0+d_1, \lambda_1) \mapsto (b_1 - b_0) \mapsto (\lambda_1 - \lambda_0) = \lambda \delta $$

\newpage

%% ===== Page 74 =====

donc (136) [crossed out: est] est une condition sur $\eta_0$ seule, savoir

$$ (136 \text{ bis}) \qquad \eta_0 \in \operatorname{Im} ( G_a \to M_0) = \operatorname{Ker} (M_0 \to C_a) $$
$$ \lambda \mapsto \lambda \delta \qquad c_0 \lambda_0 + c_1 \lambda_1 \mapsto c_0+c_1 $$

[crossed out: Ceci dit, on trouve donc comme]
[crossed out: Notons qu'on a]
[crossed out: (138) $Z_{\rho,\sigma}^{eff} \simeq$]

$$ (138) \quad \begin{cases} M_0 = \mathbb{Z} \delta \oplus \mathbb{Z}.(3\lambda_1) & d'où \\ (\lambda_1-\lambda_0) \\ \underline{M}_0 = [unclear: crossed out] \mathbb{Z} \bar{\delta} \times \mathbb{Z}.(3\lambda_1) \end{cases} $$

[crossed out: et par suite on trouve]

$$ (139) \quad Z_{\rho,\sigma}^{eff} \simeq \hat{\mathbb{O}}_{12}^* \times \mathbb{Z} \bar{\delta} \quad (\in \bar{G}_a) $$
$$ (\mu, \xi_0, \eta_0) \mapsto (\mu, \bar{\eta}_0) $$

car $\xi_0$ se récupère sur $\mu, \bar{\eta}_0$ par l'équation des lacets

$$ \xi_0 = [crossed out: (1-\sigma_M)^{-1} (\eta_0 + \nu''(3\lambda_1))] $$
[crossed out: $= \nu'' \delta + \dots$ ]

NB Si on se donne par contre $\xi_0$, on en tire par l'équation des lacets et (138) la valeur de $\eta_0$, et celle de $\nu''$ - donc si on connaît en plus $\nu_2$ et $\nu_3$, on en déduit la valeur de $2\nu_0$ ou (sur un [unclear: schéma de base] où 2 est section régulière de $\mathcal{O}_S$) $\nu_0$, i.e. finalement $\mu$. Donc sur un

\newpage

%% ===== Page 75 =====

bon sur 2 congruences, de sorte que

(140) $$ \Sigma_{\rho,\sigma}^{1t}(S) \hookrightarrow (\underline{M_0} \times \{0,1\} \times \{0,1\})(S) $$
$$ (\underline{\mu}; \xi_0, \eta_0) \longmapsto (\xi_0, \varepsilon_2(\underline{\mu}), \varepsilon_3(\underline{\mu})) $$

l'image étant formée des systèmes
$(\xi_0, \nu_2, \nu_3)$ tels que, posant
$$ \xi_0 = \nu_0 \delta + \nu''(3\lambda_1) \quad (\nu_0, \nu'' \in \Gamma \mathbb{G}_a(S)) $$
grâce à (138), on ait

(141) $$ \begin{cases} \nu'' + \nu_2 - \nu_3 [- 2 \nu_2 \nu_3] \text{ [unclear: divisible par 2]} \\ (-5)^{\nu_2} (+5)^{\nu_3} + 6 (\nu'' + \nu_2 - \nu_3 + 2 \nu_2 \nu_3) \text{ inv.} \end{cases} $$
$$ 1 - 2 \nu_2 - 2 \nu_3 + 6 \nu'' $$

En tout état de cause on note que
$\Sigma_{\rho,\sigma}^{1t}$ au dessus de $Z_{\rho,\sigma}$ est lisse de
dim relative $2$

(142) $$ 1 \to Z_{\rho,\sigma}^{+0} \to Z_{\rho,\sigma} \to \mathbb{O}_{12}^\ast \to 1 $$
$$ || $$
$$ \mathbb{G}_a \simeq \underline{\mathbb{Z}} \delta $$

(L'isom. $\mathbb{G}_a \simeq Z_{\rho,\sigma}^{+0} \simeq \Sigma_{\rho,\sigma}^{1t +0}$ étant donné par
$$ \lambda \longmapsto (\underline{\mu}=1, \xi_0 = \lambda \delta, \eta_0 = \lambda \delta) $$

On posera
(143) $$ L \overset{\text{def}}{=} \{ \lambda \cdot \lambda_0 \mid \lambda \in \Gamma(\mathbb{G}_a) \} \subset M $$
et on définit
(144) $$ L \longrightarrow Z_{\rho,\sigma}^{+0} \quad (\overset{\sim}{\longrightarrow} \Sigma_{\rho,\sigma}^{1t +0}) $$
$$ l \longmapsto (\underline{\mu}(l), (1-\sigma_M)(l), (1-\rho_M)(l)) $$
[MARGIN: $"\lambda \lambda_0"$] $$ = (\underline{\mu}(l), \lambda \delta, - \lambda \delta) $$

\newpage

%% ===== Page 76 =====

d'où

(145) $$L^\natural \overset{\sim}{\longrightarrow} Z_{\rho,\sigma}^{+0}$$

ce qui donne

(146) $$\Gamma_{\rho,\sigma} \overset{dfn}{=} Z_{\rho,\sigma}/L^\natural \simeq \mathbb{O}_{12}^* \quad .$$

Si on construit maintenant $G_{\rho,\sigma}$
comme

(147) $$G_{\rho,\sigma} = (Z_{\rho,\sigma} \cdot \underline{\underline{C}}^+) / L^\natural$$

on a le diagramme

(148)
\begin{tikzcd}
& & & \mathbb{O}_{12}^* \ar[d, equal] \\
1 \ar[r] & L^\natural \ar[r] \ar[d, hook] & Z_{\rho,\sigma} \ar[r] \ar[d, hook] & \Gamma_{\rho,\sigma} \ar[r] \ar[d, equal] & 1 \\
1 \ar[r] & \underline{\underline{C}}^+ \ar[r] & G_{\rho,\sigma} \ar[r] & \Gamma_{\rho,\sigma} \ar[r] & 1
\end{tikzcd}

D'ailleurs on a un scindage canonique

de $Z_{\rho,\sigma}$ sur $\Gamma_{\rho,\sigma} = \mathbb{O}_{12}^*$

(149)
\begin{tikzcd}
\mathbb{O}_{12}^* \ar[r] & Z_{\rho,\sigma} \ar[r, "\sim"] & \Sigma_{\rho,\sigma}^{1/\natural} \\
\mu \ar[r, mapsto] & \tilde{\mu} \ar[r, mapsto] & (\mu, \xi_0 = \nu''(3\lambda_1), \eta_0 = 0)
\end{tikzcd}

d'où

(150) $$ \begin{cases} Z_{\rho,\sigma} \simeq \mathbb{O}_{12}^* \cdot L^\natural \\ G_{\rho,\sigma} \simeq \mathbb{O}_{12}^* \cdot \underline{\underline{C}}^+ \end{cases} $$

L'opération de $\mathbb{O}_{12}^*$ sur $L^\natural$ se fait
via $\mathbb{O}_{12}^* \to \operatorname{Aut}(M)$ homothéties, celle
de $\mathbb{O}_{12}^*$ sur $\underline{\underline{C}}^+$ se fait par

(151) $$ \begin{cases} u_\mu(c) = c & (= \delta_1 c \text{ en écriture add. pour } C^+) \\ u_\mu(x) = \mu x & \text{si } x \in M \end{cases} $$

[MARGIN:
NB On pose
$G_{\rho,\sigma} = \mathbb{O}_{12}^* \cdot \underline{\underline{C}}^+$
$\cup$
$Z_{\rho,\sigma} = \mathbb{O}_{12}^* \cdot L^\natural$
]

\newpage

%% ===== Page 77 =====

79) Les groupes $\mathcal{N}_{\rho}$, $\mathcal{N}_{\rho, \sigma}$ etc.

Étudions le sous-groupe

(152) $\mathcal{N}_{\rho} \subset G_{\rho, \sigma}$
$$= \{ u \in G_{\rho, \sigma} \mid u(\rho) = \rho^{\varepsilon_3(u)} \}$$

[MARGIN: Notation provisoire, cf plus bas p. 697]

Et l'hom. composé
$$\mathcal{N}_{\rho} \hookrightarrow G_{\rho, \sigma} \xrightarrow{\underline{\chi}} \mathbb{Q}_{12}^*$$
est épi, [unclear: en voici p-ê le pour nous rassurer].

Soit [unclear: $\mu$] section locale de $\mathbb{Q}_{12}^*$, [unclear: correspondante à] $\dots$ On aura

[crossed out text]

$$u_\mu(\rho) = \lambda_1^{\mu} \rho = (\lambda_1^{3\nu''} \lambda_1^{-\nu_3}) \rho = int({\delta''}^{\nu''} \tau_1^{\nu_3})(\rho)$$

d'où

(153) $$int(({\delta''}^{\nu''} \tau_1^{\nu_3})^{-1} u_\mu)(\rho) = \rho$$

[MARGIN: NB on note $\underline{\chi}$ l'hom can (154) $G_{\rho, \sigma} \to \mathbb{Q}_{12}^*$]

i.e.
$$({\delta''}^{\nu''} \tau_1^{\nu_3})^{-1} u_\mu = \tau_1^{-\nu_3} {\delta''}^{-\nu''} u_\mu \in \mathcal{N}_\rho$$
$$\underline{\chi}(u) = \underline{\chi}(({\delta''}^{\nu''} \tau_1^{\nu_3})^{-1}) \cdot \underline{\chi}(u_\mu)$$
$$\underline{\chi}(\tau_1^{\nu_3}) = (-1)^{\nu_3} = \varepsilon_3$$

Donc posant
$$v = (\tau_1^{\nu_3} {\delta''}^{\nu''})^{-1} u_\mu = (\tau_1^{\nu_3})^{-1} {\delta''}^{-\nu''} u_\mu$$
on a

(154) $$\underline{\chi}(v) = \mu \quad \text{, avec } \quad v = \sigma^{-\nu_3} {\delta''}^{-\nu''} u_\mu \in \mathcal{N}_\rho$$
$$(\delta = \lambda_1 \lambda_0^{-1}, \nu_3 = \nu_3(\mu), \nu'' = \nu''(\mu))$$

\newpage

%% ===== Page 78 =====

On aura donc une suite exacte

(155) 
$$1 \longrightarrow Z_\rho \longrightarrow N_\rho \longrightarrow \widehat{\mathbb{O}}_{12}^* \longrightarrow 1$$
[MARGIN: $\| \Gamma_{\rho,\sigma}$ sous $\widehat{\mathbb{O}}_{12}^*$]

et

$$Z_\rho = \underline{Cent}_{\underline{\mathbb{S}}^+} (\rho)$$

On a, si

(156)
$$H_{S_0}^+ \overset{df}{=} \underline{\mathbb{S}}^+ / \underline{M}_{S_0}^\circ \simeq (H^+)_{S_0}$$
$$H^+ = \underline{\mathbb{S}}^+ / M$$

un hom. canonique

(157) $$Z_\rho \longrightarrow \underline{Cent}_{H_{S_0}^+} (\rho) = \underline{L}_\rho \quad (\simeq (\mathbb{Z}/12\mathbb{Z})_{S_0})$$

qui est un épi (puisqu'il y a une section, provenant de l'inclusion $\underline{L}_\rho \subset Z_\rho$) et dont le noyau est

(158) $$Z_\rho^\circ = Z_\rho \cap \underline{M} = \{ x \in \underline{M} \mid \rho x = x \text{ i.e. } ((1-\rho_M)x = 0) \}$$

Ce noyau est nul au dessus de $Spec \mathbb{Z}[1/3]$ mais a dim 1 en car 3. Par suite on sent bien que la "bonne" définition de $Z_\rho^\circ$ doit être donnée par l'équation

(159) $$(1 - \rho_M) x = 0 \quad \text{dans } \underline{M}_0$$

\newpage

%% ===== Page 79 =====

ce qui revient en fait à remplacer
$\underline{\mathcal{N}}_{\rho}$ défini dans (153), par l'extension
schématique de sa restriction au dessus
de $S_0 \setminus \{3\}$. On trouve donc un
[unclear: sous-groupe] [unclear: univalent], que désormais
je note $\underline{\mathcal{N}}_{\rho}$, inséré dans la suite exacte

(159) $$1 \to \underline{L}_{\rho} \to \underline{\mathcal{N}}_{\rho} \xrightarrow{\chi_{\rho}} \underline{\hat{\mathbb{O}}}_{12}^* \to 1$$

Posons
(160)
$$ \begin{cases} \underline{G}_{\rho,\sigma}^{\natural} = \underline{\hat{\mathbb{O}}}_{12}^* \underset{\gamma_{\rho,\sigma}}{\cdot} \underline{\underline{\mathfrak{S}}}_{4}^{\natural} \subset \hat{\mathbb{O}}_{12}^* \cdot \underline{\mathfrak{S}}_{4}^{+} = G_{\rho,\sigma} \\ \underline{\mathcal{N}}_{\rho}^{\natural} = \underline{\mathcal{N}}_{\rho} \cap \underline{G}_{\rho,\sigma}^{\natural} \end{cases} $$

on trouve alors
(161) $$\underline{\mathcal{N}}_{\rho}^{\natural} \cap \underline{L}_{\rho} = \underline{M} \cap \underline{L}_{\rho} = \{1\}$$
donc
(161) $$\underline{\mathcal{N}}_{\rho}^{\natural} \hookrightarrow \underline{\hat{\mathbb{O}}}_{12}^*$$

Je dis que l'image est
[MARGIN: donc un $\underline{\mathcal{N}}_{\rho}^{\natural} \subset \underline{\mathcal{Z}}_{\rho}$, ils ont même centre $\underline{\mathcal{Z}}_{\rho}^{\circ}$]
(162) $$\gamma_{\rho,\sigma}^{\natural} = \underline{\hat{\mathbb{O}}}_{12}^{*\natural} \overset{df}{=} \{ x \in \Gamma \underline{\hat{\mathbb{O}}}_{12}^* \mid \varepsilon_3(x) = 1 \text{ i.e. } v_3(x) = 0 \}$$

Que l'image contienne $\underline{\hat{\mathbb{O}}}_{12}^{*\natural}$ résulte
de (154), montrant qu'elle [unclear: la contient].

Soit $$u = f^{-1} u_{\rho} f$$ une section locale [unclear: de] $\underline{\mathcal{N}}_{\rho}^{\natural}$, avec $$f \in \underline{G}_{\rho,\sigma}^{\natural} = \underline{\hat{\mathbb{O}}}_{12}^* \cdot \Gamma \underline{M} \cdot \underline{\underline{\mathfrak{S}}}_{4}^{\natural}$$, [unclear: écrite]
qui est dans $\underline{\mathcal{N}}_{\rho}$ d'où

\newpage

%% ===== Page 80 =====

$$int(f^{-1})u_\mu(f) = \rho^{\varepsilon_3} \quad i.e. \quad \lambda_1^\nu \rho = int(f)\rho^{\varepsilon_3}$$

en réduisant mod $\underline{M}$ on trouve
$$ \rho = \rho^{\varepsilon_3} \quad \text{dans } H$$
ce qui implique $\varepsilon_3=1$, OK. Donc - - -

(163) $$ \underline{N}_{\rho}^\natural \xrightarrow{\sim} \underline{\gamma}_{\rho,\sigma}' = \underline{D}_{12}^{*\prime} $$

- ainsi $\underline{N}_\rho^\natural$ est une section partielle
au dessus de $\underline{D}_{12}^{*\prime}$, de l'extension
$\underline{N}_{\rho,\sigma}$ de $\underline{D}_{12}^*$ par $\underline{L}_\rho$.
Dans cette section partielle on

(164) $$ \underline{N}_{\rho,\sigma}^\natural = \{1,\tau'\} $$

centrée par $\tau'$, qui est une section
de $\underline{N}_{\rho,\sigma}$. Est-ce que $\tau'$ normalise
$\underline{N}_\rho^\natural$ ? - - comme il normalise $\underline{N}_\rho$, il suffi-
rait qu'elle normalise $\underline{G}_{\rho,\sigma}^\natural \underline{M}$.
Comme $\tau' = \tau \sigma$, il normalise $\underline{G}_{\rho,\sigma}^\natural$
il vient - - - c à dire
que $\sigma$ normalise $\underline{G}_{\rho,\sigma}^\natural = \underline{D}_{12}^* \underline{G}_\rho^\natural$.
Comme $\sigma$ normalise $\underline{G}_\rho^\natural = \underline{G}^\natural$ qui est
distingué, il suffit de regarder l'action
de $\sigma$ sur $\underline{D}_{12}^* \subset \underline{G}_{\rho,\sigma}$, défini par
les $u_\mu$ de (151). La question est si
les $\sigma u_\mu \sigma^{-1} u_\mu^{-1} \in \Gamma \underline{G}_\rho^\natural$, or

\newpage

%% ===== Page 81 =====

$$ \sigma u_{\hat{\rho}} \tau^{-1} u_{\hat{\rho}}^{-1} = \sigma \overbrace{u_{\hat{\rho}}(\tau)^{-1}} = \sigma (\omega^{\nu_{2}} \sigma)^{-1} = \omega^{\nu_{2}} $$
donc

Proposition. Pour une section $x$ de $\underline{G}_{\rho,\sigma}^\natural$,
$\sigma(x)$ est une section de $\underline{G}_{\rho,\sigma}^\natural \cdot \{1, \omega_2\}$ et son
image dans $\{1, \omega_2\}$ est égale : $\varepsilon_2(x)$ —
donc elle est dans $\underline{G}_{\rho,\sigma}^\natural$ ssi $\varepsilon_2(x) = 1$ $\leadsto$
section de
(164) $$ \underline{Q}_{12}^{* \prime\prime} \overset{dfn}{=} \{ H \in \underline{\Gamma} \underline{Q}_{12}^* \mid \varepsilon_2(H) = 1 \} $$

Corollaire 1. $\underline{G}_{\rho,\sigma}^\natural$ n'est pas normalisé par $\sigma$ ni $\tau'$,
mais
(165) $$ \underline{G}_{\rho,\sigma}^{\natural \prime\prime} \overset{dfn}{=} \{ x \in \underline{\Gamma} \underline{G}_{\rho,\sigma}^\natural \mid \varepsilon_2(x) = 1 \} $$
est égal : $\underline{G}_{\rho,\sigma}^{\natural \prime\prime} = \underline{G}_{\rho,\sigma}^\natural \cap \sigma(\underline{G}_{\rho,\sigma}^\natural) = \underline{G}_{\rho,\sigma}^\natural \cap \tau'(\underline{G}_{\rho,\sigma}^\natural)$
et est normalisé par $\sigma$ et $\tau'$.

[MARGIN: NB On a $N_{\hat{\rho}}^{\prime\prime \natural} = (N_{\hat{\rho}}^\natural)^\circ$ (un peu montré) A-t-on $Z_{\hat{\rho}}^{\natural \prime\prime} = (Z_{\hat{\rho}}^\natural)^\circ$ ? Ce serait plus "naturel".
Corollaire 3
comparaison le [unclear] alors (168) (38) p. 622']

Corollaire 2. $\underline{Z}_{\hat{\rho}}^\natural$ n'est pas normalisé par
$\sigma$ ni $\tau'$, mais
(166) $$ \underline{Z}_{\hat{\rho}}^{\natural \prime\prime} \overset{dfn}{=} \underline{Z}_{\hat{\rho}}^\natural \cap \underline{G}_{\rho,\sigma}^{\natural \prime\prime} \subset \underline{Z}_{\hat{\rho}}^\natural $$
est normalisé par $\sigma$ et $\tau'$ — induisant des
(167) $$ \underline{Z}_{\hat{\rho}}^{\natural \prime\prime} \overset{\sim}{\to} \underline{Q}_{12}^{* 0} \quad (= \underline{Q}_{12}^{* \prime} \cap \underline{Q}_{12}^{* \prime\prime}). $$

Soit
$$ N_{\hat{\rho}}^{\prime\prime \natural} = \{1, \tau'\}_{\hat{\rho}} \cdot \underline{Z}_{\hat{\rho}}^{\natural \prime\prime} \subset N_{\hat{\rho}}^\natural , $$
$$ ( \{1, \tau'\}_{\hat{\rho}} \cdot (\underline{Z}_{\hat{\rho}}^\natural)^\circ ) $$
(169) $$ N_{\hat{\rho}}^{\prime\prime \natural} \overset{\sim}{\to} \underline{Q}_{12}^{* \prime\prime\prime} = \{ H \in \underline{\Gamma} Q_{12}^* \mid (\varepsilon_2 \varepsilon_3)(H) = 1 \} $$

\newpage

%% ===== Page 82 =====

Donc $N_{\rho}^{\prime\prime \natural}$ est une section partielle de $N_{\rho}$
sur $\underline{O}_{12}^{*\prime\prime}$. Enfin, posons

(170) $$Z_{\rho}^{!} = Z_{\rho}^{0 \natural} \times \mu_{S_0} \quad (\mu = \{1, -1\})$$

alors $Z_{\rho}^{!}$ est centralisé par $\tau'$, et

[unclear: introduisons]

(171) $$N_{\rho}^{?} \overset{def}{=} (\mu_{S_0}) \cdot Z_{\rho}^{!} \subset N_{\rho}$$

de sorte que $Z_{\rho}^{0 \natural}$ est d'indice 4 dans $N_{\rho}^{?}$. On a alors

Corollaire 4 On a une suite exacte

[MARGIN: NB Il-y-a lien logique (172) avec la suite $1 \to \mu_{\rho} \to N_{\rho} \to \underline{O}_{12} \to 1$ au moyen de l'extension de $\underline{O}_{12}$ qu'on se propose de [unclear: dégager]]

(172) $$1 \longrightarrow \mu_{S_0} \longrightarrow N_{\rho}^{?} \longrightarrow \underline{O}_{12}^{*} \longrightarrow 1$$
$$ \Vert $$
$$ \mu_{\rho} $$

et une bi-section de l'extension
$N_{\rho}^{?}$ de $\underline{O}_{12}^{*}$ par $\mu_{\rho}$.

Il y aurait lieu de développer
le diagramme des sous-groupes
remarquables de $N_{\rho}$, calqué sur
le diag. (38) de p. 622.

\newpage

%% ===== Page 83 =====

689

8) Les groupes $N_{\sigma}, N_{\sigma}^S$ ...

Posons
(173) $$ \begin{cases} N_{\sigma} \subset \Gamma_{\rho,\sigma} \\ =^{dfn} \{ u \in \Gamma_{\rho,\sigma} \mid u(\sigma) = \varepsilon_2(u)\sigma \quad (= \omega^{\nu_2(u)} \sigma) \} \end{cases} $$

(174) $$ N_{\sigma}^S =^{dfn} N_{\sigma} \cap G_{\rho,\sigma}^S $$

L'hom. composé
\begin{tikzcd}
N_{\sigma}^S \ar[r, hook] \ar[rr, bend right, "\chi_{\sigma}"'] & S_{\rho,\sigma} \ar[r, "\chi"] & \mathbb{O}_{12}^*
\end{tikzcd}
est épi. En effet, les $u_f$ sont dans $N_{\sigma}^S$ et donnent un scindage $\mathbb{O}_{12}^* \hookrightarrow N_{\sigma}^S$, d'où des décompositions en pro. $1/2$ directs
(175) $$ \begin{cases} N_{\sigma}^S \simeq \mathbb{O}_{12}^* \cdot S Z_{\sigma}^S \\ N_{\sigma} \simeq \mathbb{O}_{12}^* \cdot S Z_{\sigma} \end{cases} $$

où on a posé
(176) $$ \begin{cases} Z_{\sigma} = \operatorname{Ker}(N_{\sigma} \xrightarrow{\varepsilon_2 \circ p} \mu_{S_0}) = \operatorname{Cent}_{G_{\rho,\sigma}}(\sigma) \\ S Z_{\sigma} = Z_{\sigma} \cap S^+ = \operatorname{Cent}_{S^+_{\rho,\sigma}}(\sigma) \\ Z_{\sigma}^S = Z_{\sigma} \cap G_{\rho,\sigma}^S \\ S Z_{\sigma}^S = S Z_{\sigma} \cap G_{\rho,\sigma}^S \\ \quad = S^{\circ} = M \end{cases} $$

D'ailleurs
\begin{tikzcd}
S Z_{\sigma} \ar[r, "p_{S_0}^+"] & \operatorname{Cent}_{H_{S_0}^+}(\sigma) \ar[d, "\wr"'] \\
& L_{\sigma} \simeq (\mathbb{Z}/q\mathbb{Z})_{S_0} \ar[ul, dashed]
\end{tikzcd}
qui a une section, donc

(177) $$ S Z_{\sigma} \simeq L_{\sigma} \times S Z_{\sigma}^S $$

\newpage

%% ===== Page 84 =====

Enfin

(178)
$$S Z_{\sigma}^{\mathcal{G}} = \underline{M}$$
$$\parallel$$
$$\{ x \in \Gamma \underline{M} \mid (1 - \sigma_{\underline{M}})(x) = 0 \} = \{ [unclear crossed out] \}$$
$$= \{ a (\lambda \sigma \lambda) \mid a \in \Gamma G_{\cdot} \} \simeq G_{\cdot}$$

d'où, compte tenu de 177 :

(179)
$$S Z_{\sigma} \simeq L_{\sigma} \times \underbrace{G_{\cdot}}_{S Z_{\sigma}^{\mathcal{G}}}$$
$$(S Z_{\sigma}^{\mathcal{G}} = (S Z_{\sigma})^{\circ} \text{ (comp. neutre)} \simeq G_{\cdot})$$

On a donc

[unclear crossed out]
$$1 \to \underset{\parallel S Z_{\sigma}^{\mathcal{G}}}{G_{\cdot}} \to \mathcal{N}_{\sigma}^{\mathcal{G}} \underset{\gamma_{p,\sigma}}{\to} \mathbb{O}_{12}^* \to 1$$

plus précisément les isom. (175) (177) (179) donnent

(180)
$$\mathcal{N}_{\sigma}^{\mathcal{G}} \simeq \mathbb{O}_{12}^* \cdot G_{\cdot}$$
$$\mathcal{N}_{\sigma} \simeq \mathbb{O}_{12}^* \cdot (L_{\sigma} \times G_{\cdot})$$

Dans ces formules $\mathbb{O}_{12}^*$ opère sur $G_{\cdot}$ via $\chi : \mathbb{O}_{12}^* \to \mathbb{G}_m$ , et sur $L_{\sigma}$ via $\varepsilon_2 : \mathbb{O}_{12}^* \to (\pm 1)_{S_0}$ :

(181)
[unclear crossed out: $\mu(\sigma) = \sigma^{\varepsilon_2(\mu)}$]
$$\begin{cases} \mu(\lambda) = \mu . \lambda \\ \mu(\sigma) = \sigma^{\varepsilon_2} \end{cases} \quad - \; \mu = \chi(\mu) , \varepsilon_2 = \varepsilon_2(\mu)$$

\newpage

%% ===== Page 85 =====

9°) Retour sur $\underline{M}_{p,\sigma}$.

Nous inspirant des développements du § 49 (notamment VI, 5)), nous poserons

(182)
$$ \begin{cases}
S_0 \underline{M}_{p,\sigma} = Z_{p,\sigma} \times_{\gamma} N_{p}^? \times_{\gamma} N_{\sigma}^S \times_{\gamma} G_{p,\sigma} \\
S_0 \underline{M}_{p,\sigma}^S = Z_{p,\sigma} \times_{\gamma} N_{p}^? \times_{\gamma} N_{\sigma}^S \times_{\gamma} G_{p,\sigma}^S
\end{cases} $$

(183)
$$ \begin{cases}
\underline{M}_{p,\sigma} = N_{p}^? \times_{\gamma} N_{\sigma}^S \times_{\gamma} G_{p,\sigma} \\
\underline{M}_{p,\sigma}^S = N_{p}^? \times_{\gamma} N_{\sigma}^S \times_{\gamma} G_{p,\sigma}^S
\end{cases} $$

(184)
$$ \begin{cases}
S_0 \Gamma_{p,\sigma} = Z_{p,\sigma} \times_{\gamma} N_{p}^? \times_{\gamma} N_{\sigma}^S \\
\Gamma_{p,\sigma} = N_{p}^? \times_{\gamma} N_{\sigma}^S
\end{cases} $$

où
$$ \gamma = \gamma_{p,\sigma} = \hat{\mathbb{Z}}_{12}^* $$

En fait, les groupes $S_0 \underline{M}_{p,\sigma}$, $S_0 \underline{M}_{p,\sigma}^S$, $S_0 \Gamma_{p,\sigma}$ jouent plutôt un rôle technique auxiliaire, ce sont des extensions des groupes $\underline{M}_{p,\sigma}$, $\underline{M}_{p,\sigma}^S$ et $\Gamma_{p,\sigma}$ par $L_{\sigma}^{\underline{M}} = \operatorname{Ker}(Z_{p,\sigma} \xrightarrow{\gamma} \hat{\mathbb{Z}}_{12}^*)_{\lambda}$ (déduites de l'extension $S_0 \Gamma_{p,\sigma}$ de $\Gamma_{p,\sigma}$ par $L_{\sigma}^{\underline{M}} \simeq \mathbb{G}_a$, par image inverse via les épimorphismes $\underline{M}_{p,\sigma} \to \Gamma_{p,\sigma}$ et l'épi. [unclear: induit] $\underline{M}_{p,\sigma}^S \to \Gamma_{p,\sigma}$.
Nous attacherons notre attention surtout aux autres groupes (sans $S_0$), reliés par une suite exacte (comparer p. 83g (179)) :

\newpage

%% ===== Page 86 =====

(190)
\begin{tikzcd}
1 \ar[r] & \underline{C}^\natural \ar[r] \ar[d] & \underline{M}^\natural_{\rho,\sigma} \ar[r] \ar[d] & \Gamma_{\rho,\sigma} \ar[r] \ar[d, equal] & 1 \\
1 \ar[r] & \underline{C}^+ \ar[r] & \underline{M}_{\rho,\sigma} \ar[r] & \Gamma_{\rho,\sigma} \ar[r] & 1
\end{tikzcd}

et des suites exactes

(191)
\begin{tikzcd}
1 \ar[r] & S N^?_P \times S\mathbb{Z}^?_\rho \times \underline{C}^\natural \ar[r] \ar[d] & \underline{M}^\natural_{\rho,\sigma} \ar[r] \ar[d] & \mathbb{O}_{12}^* \ar[r] \ar[d, equal] & 1 \\
1 \ar[r] & S N^?_P \times S\mathbb{Z}^?_\rho \times \underline{C}^+ \ar[r] & \underline{M}_{\rho,\sigma} \ar[r] & \underset{\gamma_{\rho,\sigma}}{\mathbb{O}_{12}^*} \ar[r] & 1
\end{tikzcd}

(192)
$$1 \to S N^?_P \times S\mathbb{Z}^?_\rho \to \Gamma_{\rho,\sigma} \to \underset{\gamma_{\rho,\sigma}}{\mathbb{O}_{12}^*} \to 1$$

où

(193)
$$ \begin{cases} S N^?_P = \mu_P \simeq \{\pm 1\} & (172) \\ S\mathbb{Z}^?_\rho \simeq C_+ & (179) \end{cases} $$

Ces suites exactes se précisent, grâce
aux décompositions $1/2$ directes :

[MARGIN: crossed out lines]
$??? \simeq \mathbb{O}_{12}^* \cdot C_+$ (150, 151)
$??? \simeq \mathbb{O}_{12}^* \cdot C_+$ (180)
$N^?_P / \mu_P \simeq \mathbb{O}_{12}^*$ (172)

(194)
$$ \begin{cases} \mathbb{Z}_{\rho,\sigma} \simeq \mathbb{O}_{12}^* \cdot C_+ & (150, 151) \\ \underline{C}_{\rho,\sigma} \simeq \mathbb{O}_{12}^* \cdot \underline{C}^+ & (15) \\ \underline{C}^\natural_{\rho,\sigma} \simeq \mathbb{O}_{12}^* \cdot \underline{C}^\natural & (160) \\ N^\natural_S \simeq \mathbb{O}_{12}^* \cdot \underline{C}^\natural & (180) \\ \text{et enfin} \\ N^?_P / \mu_P \simeq \mathbb{O}_{12}^* & (172) \end{cases} $$

\newpage

%% ===== Page 87 =====

On trouve alors des isom. canoniques

(195) 
$$M_{\rho,\sigma}/\mu_p \simeq \mathbb{O}_{12}^* \cdot ( \mathbb{G}_a \times \underline{\mathfrak{S}}^+ )$$
$$M_{\rho,\sigma}^S/\mu_p \simeq \mathbb{O}_{12}^* \cdot ( \mathbb{G}_a \times \underline{M} )$$
$$\Gamma_{\rho,\sigma} \simeq \mathbb{O}_{12}^* \cdot \mathbb{G}_a$$

[MARGIN: $\uparrow$ action de $\mathbb{O}_{12}^*$ sur $\mathbb{G}_a \times \underline{\mathfrak{S}}^+$ d[onn]ée ... vue (181) et (151). Sur $\mathbb{G}_a \times \underline{M}$, $\mathbb{O}_{12}^*$ agit via [unclear] morphisme dans $\mathbb{G}_a \times \underline{\mathfrak{S}}^+$ ...]

( $M_{\rho,\sigma}$ et $M_{\rho,\sigma}^S$ , un [unclear] précédemment,
les $M_{\rho,\sigma}/\mu_p$ et $M_{\rho,\sigma}^S/\mu_p$ p[ar] r[apport à] $\mu_p$ )
se déduisent de l'extension

(196) $$1 \rightarrow \mu_p \rightarrow N_\rho^? \rightarrow \mathbb{O}_{12}^* \rightarrow 1$$

de $\mathbb{O}_{12}^*$ par $\mu_p$, par images inverses
des homomorphismes canoniques des groupes
(195) vers le facteur $\mathbb{O}_{12}^*$. Il
faudrait revenir sur l'extension
(196), qui est sûrement l'image
inverse d'un calcul de [unclear]
via $$\mathbb{O}_{12}^* \xrightarrow{(\varepsilon_2, \varepsilon_3)} (\mu_{4p})_{S_0}$$

Comparer § 49 V 8°). L'opération
L'opération de $\mathbb{O}_{12}^*$ sur $\mathbb{G}_a \times \underline{\mathfrak{S}}^+$
(id. sur $\mathbb{G}_a \times \underline{\mathfrak{S}} = \mathbb{G}_a \times \underline{M}$ ) se fait
via des opérations sur $\mathbb{G}_a$ et $\underline{\mathfrak{S}}^+$ que
j'ai explicitées dans [unclear] (181) et (151) !

\newpage

%% ===== Page 88 =====

10) Les sous-groupes $\underline{M}_{\rho, \sigma}[0]$, $\underline{M}_{\rho, \sigma}(0)$, $\underline{M}_{\rho, \sigma}(1/2)$, $\underline{M}_{\rho, \sigma}'(-J)$ $[\underline{M}_{\rho, \sigma}'''(-J) \supset \underline{M}_{\rho, \sigma}'(-J)]$.

On utilise le plongement

(195) $\underline{M}_{\rho, \sigma} = N_\rho^? \times_\gamma N_\sigma^? \times_\gamma \underline{G}_{\rho, \sigma} \subset \mathcal{N}_\rho^? \times \mathcal{N}_\sigma^? \times \underline{G}_{\rho, \sigma}$
$$ \Big\| $$
$$ \{(a, b, u) \in \mathcal{N}_\rho^? \times \mathcal{N}_\sigma^? \times \underline{G}_{\rho, \sigma} \mid \gamma(a) = \gamma(b) = \underline{\chi}(u)\} $$

On définit donc

(196) $\underline{M}_{\rho, \sigma}[0] \subset \underline{M}_{\rho, \sigma}$
[MARGIN: NB on aura $\underline{M}_{\rho, \sigma}[0] \subset \underline{M}_{\rho, \sigma}^S$]
$$ \| $$
$$ \{(a, b, u) \in \underline{M}_{\rho, \sigma} \mid u \in \underline{Z}_{\rho, \sigma}\} $$

donc
$$ \underline{M}_{\rho, \sigma}[0] \simeq N_\rho^? \times_\gamma N_\sigma^? \times_\gamma \underline{Z}_{\rho, \sigma} $$

(197) Soit $\underline{Z}_{\rho, \sigma}(0) \subset \underline{Z}_{\rho, \sigma}$
le facteur $1/2$ direct $\mathbb{Q}_{12}^*$ de (150) (151),
alors, on a $\underline{M}_{\rho, \sigma}(0) \simeq N_\rho^? \times_\gamma N_\sigma^? \times_\gamma \underline{Z}_{\rho, \sigma}(0)$

[MARGIN: NB $\underline{Z}_{\rho, \sigma}(0) \subset \underline{Z}_{\rho, \sigma} \subset \underline{G}_{\rho, \sigma}$ est en fait le fact. $1/2$ direct $\mathbb{Q}_{12}^*$ de $\underline{G}_{\rho, \sigma}$ ds (150), il proviendrait la notation $\underline{G}_{\rho, \sigma}(0)$]

(198)
\begin{tikzcd}
\underline{M}_{\rho, \sigma}(0) \ar[r, hook] \ar[d, "\downarrow \simeq"'] & \underline{M}_{\rho, \sigma}[0] \ar[d, "\downarrow \simeq"] \\
S \underline{\Gamma}_{\rho, \sigma}(0) \ar[r, hook] \ar[d] & S_0 \underline{\Gamma}_{\rho, \sigma} \\
\underline{\Gamma}_{\rho, \sigma}(0) &
\end{tikzcd}

(199) $\underline{M}_{\rho, \sigma}(0) = \{(a, b, u) \in \underline{M}_{\rho, \sigma} \mid u \in \underline{Z}_{\rho, \sigma}(0)\}$
[MARGIN: dans la descript. (195) de $\underline{M}_{\rho, \sigma}$ comme prod. $1/2$ direct.]
$$ = \underset{\cap \atop \Gamma_Q}{\underline{\mu}} \cdot ( \underset{\cap \atop \Gamma_{\underline{Q}^+}}{\underline{\lambda}} , 0 ) $$

\newpage

%% ===== Page 89 =====

[MARGIN: N.B. $H_\rho \subset \underline{M}_{\rho, \sigma}(0) \subset \underline{M}_{\rho, \sigma}[0]$ i.e.]

$$ (200) \quad \underline{M}_{\rho, \sigma}(0) / H_\rho \simeq \mathbb{Q}_{12}^* \cdot G_a \hookrightarrow \underline{M}_{\rho, \sigma} / H_\rho \simeq \mathbb{Q}_{12}^* (G_a \times \underline{S}^+) $$

(inclusion canonique de $G_a$ dans $G_a \times \underline{S}^+$)

D'autre part on a

$$ (201) \quad \underline{M}_{\rho, \sigma}[0] / H_\rho \simeq \mathbb{Q}_{12}^* \cdot (G_a \times L_0^M) \hookrightarrow \underline{M}_{\rho, \sigma} / H_\rho \simeq \mathbb{Q}_{12}^* (G_a \times \underline{S}^+) $$

On définit d'autre part

$$ (202) \quad \underline{M}_{\rho, \sigma}(1/2) = \{ (a, b, U) \in \underline{M}_{\rho, \sigma} \mid U = b \} $$

(donc $\quad \underline{M}_{\rho, \sigma}(1/2) \subset \underline{M}_{\rho, \sigma}^\natural$ )

D'où les

$$ (203) \quad \underline{M}_{\rho, \sigma}(1/2) \overset{\simeq}{\longrightarrow} \underline{\Gamma}_{\rho, \sigma} $$
$$ \underline{M}_{\rho, \sigma}^\natural \simeq \underline{M}_{\rho, \sigma}(1/2) \cdot M $$
$$ \underline{M}_{\rho, \sigma} \simeq \underline{M}_{\rho, \sigma}(1/2) \cdot \underline{S}^+ $$

Dans la représentation de $\underline{M}_{\rho, \sigma}/H_\rho$, $\underline{M}_{\rho, \sigma}^\natural/H_\rho$ en produits $1/2$ directs, on a la description

[MARGIN: $H_\rho \subset \underline{M}_{\rho, \sigma}(1/2)$]

(204)
$$ \underline{M}_{\rho, \sigma}(1/2) / H_\rho = \left\{ \underline{\mu} \cdot (\lambda, \lambda(\lambda_1 + \lambda_2)) \mid \lambda \in \Gamma_{G_a} \right\} $$
$$ \wr\downarrow \qquad \qquad \qquad \qquad \Gamma_{G_a} \quad \to \quad \Gamma_M $$
$$ \mathbb{Q}_{12}^* \cdot G_a \hookrightarrow \underline{M}_{\rho, \sigma}^\natural / H_\rho \simeq \mathbb{Q}_{12}^* (G_a \times M) $$

associé à $G_a \to G_a \times M$
$\lambda \mapsto \lambda, \lambda(\lambda_1+\lambda_2)$

\newpage

%% ===== Page 90 =====

Enfin on pose

(205) $$M_{p,\sigma}(-\bar{j}) \subset M_{p,\sigma}$$
$$\|$$
$$\{ (a,b,U) \in M_{p,\sigma} \mid U=a \}$$

ici on n'aura pas néc. $U \in \Gamma(G_{p,\sigma}^\natural)$ car $N_p^\natural \notin G_{p,\sigma}^\natural$, donc $M_{p,\sigma}(-\bar{j}) \notin \underline{M}_{p,\sigma}^\natural$, contrairement à ce qui se passe pour $M_{p,\sigma}(0), M_{p,\sigma}(\infty), M_{p,\sigma}(1/2)$. On aura encore

(206) $$M_{p,\sigma}^\natural(-\bar{j}) \simeq \underline{\Gamma}_{p,\sigma}$$
$$M_{p,\sigma} \simeq M_{p,\sigma}(-\bar{j}) . \underline{S}^+$$

Mais on a $H_p \not\subset M_{p,\sigma}(-\bar{j})$
plus précisément
$$H_p \cap M_{p,\sigma}(-\bar{j}) = \{1\}$$
donc la [unclear] $M_{p,\sigma}(-\bar{j})$ ne définit pas (comme $M_{p,\sigma}(1/2)$) une section de $M_{p,\sigma}/H_p$ au dessus de $\underline{\Gamma}_{p,\sigma}/H_p (\simeq \underline{N}_\sigma^\natural \simeq \underline{O}_{12}^* . G_a)$. On a cependant
$$M_{p,\sigma}(-\bar{j}) . H_p = M_{p,\sigma}(-\bar{j}) . H_G$$
et en posant
(207) $$\bar{M}_{p,\sigma}(-\bar{j}) = M_{p,\sigma}(-\bar{j}) . H_p / H_p = \text{Im } (M_{p,\sigma}(-\bar{j}) \to M_{p,\sigma}/H_p)$$
(isomorphique à $\bar{\Gamma}_{p,\sigma} = \Gamma_{p,\sigma}/H_p \simeq (M_{p,\sigma}/H_p) / \underline{S}^+$)

\newpage

%% ===== Page 91 =====

$\tilde{89}$

[MARGIN: 
NB. On a $\tilde{\Gamma}_{p,\sigma} \twoheadrightarrow N_p/\pm 1$ 
(épi de noyau $\simeq \hat{\mathbb{Z}}$ $(= G_m)$ i.e. $\operatorname{Ker}(N_G \to N_p / \pm 1) \simeq \hat{\mathbb{Z}} \cdot u$) 
et $\tilde{\Gamma}_{p,\sigma} \twoheadrightarrow N_G / H_p \to 1$

la désignation de $M'_{p,\sigma}(-\tilde{j})$ est un peu [unclear] car il s'agit plutôt de $M'_{p,\sigma}(\tilde{j})$ \dots
]

(208) [unclear] $\mathcal{H}_G \longrightarrow \tilde{M}_{p,\sigma}(-\tilde{j}) \longrightarrow \tilde{\Gamma}_{p,\sigma} \longrightarrow 1$

$\tilde{M}_{p,\sigma}(-1/2)$ est une bissection de $\tilde{M}_{p,\sigma} \overset{def}{=} M_{p,\sigma} / H_p$ au dessus de $\tilde{\Gamma}_{p,\sigma}$. Posant

(209) $$M'_{p,\sigma}(-\tilde{j}) \overset{def}{=} \tilde{M}_{p,\sigma}(-\tilde{j}) \cap \tilde{M}_{p,\sigma}(-1/2)$$

[MARGIN: $M'_{p,\sigma}(-\tilde{j})$ s'identifie à une section partielle de $\tilde{M}_{p,\sigma}$ au dessus de $\tilde{\Gamma}_{p,\sigma}$, d'image $\tilde{\Gamma}'_{p,\sigma}$ d'indice 2 dans $\tilde{\Gamma}_{p,\sigma}$.]

(209') $$M'_{p,\sigma}(-\tilde{j}) \overset{\sim}{\longrightarrow} \tilde{\Gamma}'_{p,\sigma} \overset{def}{=} \operatorname{Ker}(\tilde{\Gamma}_{p,\sigma} \overset{\varepsilon_2}{\longrightarrow} \mu_2)$$

On aura donc

(210) 
$$M'_{p,\sigma} \overset{def}{=} \operatorname{Ker}(M_{p,\sigma} \overset{\varepsilon_2}{\longrightarrow} \mu_2) \simeq M'_{p,\sigma}(-\tilde{j}) \cdot \underline{\mathfrak{S}}^+$$
$$\tilde{M}'_{p,\sigma} \overset{def}{=} \operatorname{Ker}(\tilde{M}_{p,\sigma} \overset{\varepsilon_2}{\longrightarrow} \mu_2) \simeq M'_{p,\sigma}(-\tilde{j}) \cdot \underline{M}$$

Remarque. Je m'arrête là sur les développements de sous-groupes [unclear] et de $\tilde{M}_{p,\sigma}(-\tilde{j})$ - j'y reviendrai encore une fois, sans doute, dans le contexte des schémas en groupes $M_{p,\sigma}$ génériques. Il devient clair qu'il faudra clarifier,

\newpage

%% ===== Page 92 =====

pour la suite, entre les notations possibles
respectives pour les groupes notés ici
$\underline{M}_{p,\sigma}$ et $\underline{M}_{p,\sigma} / H_p \simeq M_{p,\sigma}$ (et idem avec $\sim$)
ce sont ces derniers, il est vrai, qui me
semblent les plus intéressants, car c'est
[MARGIN: plutôt le gr. profini [unclear] dans le recouv. restreint...]
eux qui recouvrent le groupe universel
$\mathcal{M} = \mathcal{N}_{1,1}^\sim$ (celui-ci a $\Gamma_{\mathbb{Q}}^\sim$ pour $\widehat{Sl(2,\mathbb{Z})}^\sim$).
Il semblerait que ces groupes méritent
la notation la plus simple, en l'occurrence
$M_{p,\sigma}$, et les groupes notés précédemment
$\underline{M}_{p,\sigma}$ prendraient des $\sim$, ou un
signe analogue. (Le $\sim$ est déjà utilisé
dans un sens assz différent dans $\Gamma_{\mathbb{Q}}^\sim$,
$\mathcal{M}_{0,3}^\sim$ etc - d'autre part les
exposants $', '', ''', \circ, +, !, \sharp, \flat$ déjà utilisés
dans des sens différents. Le double emploi
du $\sim$ risque d'ailleurs d'introduire
des conflits de notation, tôt ou tard,
il est peut-être préférable de remplacer
le $\sim$ dans $\Gamma_{\mathbb{Q}}^\sim$ par des lettres telles que $\Gamma_{\mathbb{Q}}^{\text{ét}}$
(suggérant "étendu") ou $\Gamma_{\mathbb{Q}}^{\text{gén}}$ "généralisé"
[unclear: géom] ... Il faudra que je résolve ce di-

\newpage

%% ===== Page 93 =====

dilemmes notationnels, pour l'instant je garde les notations introduites par la bande dans la présente section, qui seront revisées avec soin dans un $\S$ spécial ultérieur.

Il faudrait encore décrire des sous-groupes $\underline{M}_{p, r}^\sim (-J)$, $\underline{M}_{p, r}^{\prime \sim} (-j)$ dans $\underline{M}_{p, r}^\sim$ décrit comme produit semi-direct, savoir

(211) $$ \underline{M}_{p, r}^\sim (-J) \subset \underline{M}_{p, r}^\sim \simeq \mathbb{Q}_\ell^\times (G_a \times \widehat{\mathfrak{S}}_{p,r}^+) $$
$$ = \{ (\lambda, x) \cdot u_{\underline{\mu}} \in N_p^? \} $$

où $u_{\underline{\mu}}$ est défini dans (151), en particulier

$$ u_{\underline{\mu}}(\underline{\rho}) = \lambda_1^\nu \underline{\rho} \quad \text{[unclear]} \quad \nu = \nu(\underline{\mu}) , $$

[unclear crossed out line ending with (154)]

[MARGIN: $\therefore \delta = \lambda_1 \lambda_0^{-1}$]

$$ u = \sigma_3 \delta^{-\nu} u_{\underline{\mu}} \in \Gamma \underline{N}_p , \quad \chi(u) = 1 $$

et posant $\sigma = \omega \tau' \tau$, le premier membre s'écrit aussi

(212) $$ u = \omega^{1/3} ( \omega^{1/3} \cdot \tau^{1/3} \delta^{-\nu} u_{\underline{\mu}} ) $$
$$ \in \Gamma \underline{S}_{p,r}^\sim , \text{ car les trois facteurs le sont} $$
$$ \in \Gamma \underline{N}_p \cap \underline{G}_{p,r}^\sim = \Gamma \underline{Z}_p^S $$

\newpage

%% ===== Page 94 =====

et on voit que $v \in N_{\hat{p}}^?$, d'où la for-
mule [unclear] explicite de la
fibre des éléments de $N_{\hat{p}}^?$ au-dessus
de $\underline{\nu} = (\nu_1, \nu_3, \nu_0) \in \Gamma Q_{\hat{1}}^\times$
$$ (212) \quad (N_{\hat{p}}^?)_{\underline{\nu}} = \{ \omega^i \sigma^{\nu_3} \delta^{-\nu''} u_{\underline{\mu}} \mid i \in \Gamma \{0,1\}_{S_0} \} $$

où $\nu_3 = \nu_3(\underline{\nu})$, $\nu'' = \nu''(\underline{\nu})$, $\delta = \lambda_1 \lambda_0^{-1}$. Don-
c on trouve

[MARGIN: $\delta = \lambda_1 \lambda_0^{-1}$]
$$ (213) \quad M_{\hat{p}, \sigma}^\sim (-j) = \left\{ \underset{\begin{subarray}{c} \in \\ \mathbb{G}_a \end{subarray}}{\underline{\lambda}}, \underset{\begin{subarray}{c} \in \\ \hat{\mathbb{S}}^+ \end{subarray}}{x}) \cdot \underset{\begin{subarray}{c} \in \\ Q_{\hat{1}}^+ \end{subarray}}{\underline{\mu}} \mid \begin{array}{l} x = \omega^i \sigma^{\nu_3} \delta^{-\nu''} u_{\underline{\mu}} \text{ , où } \\ i \in \Gamma \{0,1\}_{S_0} \end{array} \right\} $$

On aura donc

$$ (214) \quad M_{\hat{p}, \sigma}^{\prime \sim} (-j) = M_{\hat{p}, \sigma}^\sim (-j) \cap M_{\hat{p}, \sigma}^\natural $$
$$ = \left\{ (\underline{\lambda}, \delta^{-\nu''}) \underline{\mu} \mid \begin{array}{l} \underline{\mu} \in \Gamma Q_{\hat{1}}^{\pm 1} \text{ (avec } \varepsilon_3(\underline{\mu})=0 \text{)} \\ \underline{\lambda} \in \Gamma \mathbb{G}_a \\ \nu'' = \nu''(\underline{\mu}) \end{array} \right\} $$

\newpage

%% ===== Page 95 =====

11) L'homomorphisme $M \to M_{\rho,\sigma}(\hat{\mathbb{Z}})^{\sim}$, relatif à la tour des courbes de Fermat.

Notons qu'on a pour tout entier $N \in \mathbb{N}^*$

[MARGIN: Le $\Delta$ indique qu'on prend (215) sections de [unclear] qui sont [unclear]]

$$ \mathbb{O}_N^{\Delta}(\hat{\mathbb{Z}}) \simeq \mathbb{Z} \quad \text{d'où} \quad \mathbb{O}_\Gamma^{\Delta}(\hat{\mathbb{Z}})^{\wedge} \simeq \hat{\mathbb{Z}}^* $$

On a donc, compte tenu des décompositions en produits semi-directs (150) et des isomorphismes

(216)
$$ Z_{\rho,\sigma}^\sim \overset{df}{=} Z_{\rho,\sigma}^\Delta(\hat{\mathbb{Z}})^\wedge \simeq \hat{\mathbb{Z}}^* . \hat{L}_0 \quad \text{où} \quad \hat{L}_0 \simeq \hat{\mathbb{Z}} . \lambda_0 \simeq \hat{\mathbb{Z}} $$
$$ G_{\rho,\sigma}^\sim \overset{df}{=} G_{\rho,\sigma}^\Delta(\hat{\mathbb{Z}})^\wedge \simeq \hat{\mathbb{Z}}^* . \underline{G}^{+\wedge} $$
$$ C_{\rho,\sigma}^\sim \overset{df}{=} C_{\rho,\sigma}^\Delta(\hat{\mathbb{Z}})^\wedge \simeq \hat{\mathbb{Z}}^* . M^\wedge $$
$$ N_{\rho}^\sim \overset{df}{=} N_\rho^\sim(\hat{\mathbb{Z}})^\wedge \simeq \mathbb{O}_\Gamma^\Delta(\hat{\mathbb{Z}})^\wedge \simeq \hat{\mathbb{Z}}^* $$
$$ N_{\sigma}^\sim \overset{df}{=} N_{\sigma}^\sim(\hat{\mathbb{Z}})^\wedge \simeq \hat{\mathbb{Z}}^* . \hat{\mathbb{Z}} $$
$$ M_{\rho,\sigma}^\sim \overset{df}{=} M_{\rho,\sigma}^\sim(\hat{\mathbb{Z}})^\wedge \simeq \hat{\mathbb{Z}}^* . (\hat{\mathbb{Z}} \times \underline{G}^{+\wedge}) $$
$$ M_{\bar{\rho},\bar{\sigma}}^\sim \overset{df}{=} M_{\bar{\rho},\bar{\sigma}}^\sim(\hat{\mathbb{Z}})^\wedge \simeq \hat{\mathbb{Z}}^* . (\hat{\mathbb{Z}} \times \hat{M}) $$

où les opérations de $\hat{\mathbb{Z}}^*$ sur les $\hat{\mathbb{Z}}$ facteurs à droite sont explicitées dans les sections précédentes. On notera que, modulo des ajustages de notations et de définitions, les groupes $Z_{\rho,\sigma}^\sim$ etc. précédents sont ceux qui serais construits directement, en termes de $(\hat{\underline{G}}^\wedge, \underline{G}^{+\wedge}, M^\wedge)$ (non indexé de $\rho,\sigma,\tau$), dans §49 VI.

On a des hom. canoniques

(217)
\begin{tikzcd}
M = [unclear] \ar[r] & M_{\rho,\sigma}^\sim \\
M' \ar[r] \ar[u, hook] & M_{\bar{\rho},\bar{\sigma}}^\sim \ar[u, hook]
\end{tikzcd}

qui comme d'habitude se décompose en

\newpage

%% ===== Page 96 =====

trois, compatibles en un hom.
$$M \longrightarrow \mathbb{G}_{\rho, \sigma}^* \overset{dfn}{=} \gamma_{\rho, \sigma}^*(\hat{\mathbb{Z}})^\Delta = \hat{\mathbb{Z}}^*$$
qui n'est autre que le caractère cyclotomique $\chi$. Le $2^e$ de ces trois,
$$M \longrightarrow N_\rho^\sim \simeq \hat{\mathbb{Z}}^*$$
se réduit ici à [unclear: l'ident.] du caractère, et ne pose donc aucune question intéressante. On a d'autre part

(218)
$$M \longrightarrow G_{\rho, \sigma} = \hat{\mathbb{Z}}^* . \mathfrak{S}^{+ \wedge}$$
$$\downarrow \qquad \downarrow$$
$$M^! \longrightarrow G_{\rho, \sigma}^\sim = \hat{\mathbb{Z}}^* . \mathbb{M}^\wedge$$

qu'il faudra examiner de plus près - comme (même sur $M^!$) le caractère cyclotomique $\chi$ est surjectif, on sait en tous cas que ces deux homs sont surjectifs. Enfin, on a un hom $M \longrightarrow N_\sigma^\sim$, qui se factorise en

(219)
$$\Gamma_{\mathbb{Q}}^{et} \longrightarrow N_\sigma^\sim \overset{dfn}{=} \hat{\mathbb{Z}}^* . \hat{\mathbb{Z}}$$
[MARGIN: (1/2 direct)]

qu'il faudra examiner également - notamment voir s'il est surjectif.
Bien entendu, je n'ai de chance de le faire que par "voie arithmétique", en passant par l'hom. naturel sur $\Gamma_{\mathbb{Q}}$ lui-même $\longrightarrow \mathfrak{S}_{[unclear]}$.

\newpage

%% ===== Page 97 =====

J'ai envie de mettre ces homomorphismes en
relation avec la "tour" des courbes algébriques
définies sur $\mathbb{Q}$, indexées par $n \in \mathbb{N}^\ast$

(220) $X_n : X_0^n + X_1^n + X_\infty^n = 0$ \qquad \qquad (courbe de
$\subset \mathbb{P}^2_{\mathbb{Q}}$ \qquad \qquad \qquad \qquad Fermat d'indice $n$)

Soit $U_n$ l'ouvert des pts où $X_0 X_1 X_\infty \neq 0$

Sur $X_n$ le groupe fini $\mu_n^{\{0,1,\infty\}} / \mu_n$ (diag.)
opère ($\mu_n \subset \mathbb{G}_m$), et on a aussi
que $\mathfrak{S}_3 = \mathfrak{S}_{\{0,1,\infty\}}$, de sorte que le
produit semi-direct

[MARGIN: Notation $F_n$ pas très heureuse - il est possible qu'il y aura un besoin ultérieur de la notation $F_n$ pour les [unclear] de $U_n, U_\infty$ avec un in-dice $n$? (comparer p. 719)]

(221) $F_n = \mathfrak{S}_3 \cdot \mu_n^{\{0,1,\infty\}} / \mu_n = F_n^0$

opère sur $X_n$.
On a

(222) $X_n / F_n^0 \simeq X_1 \simeq U_{0,3}$ ,

isomorphismes compatibles aux opérations de
$\mathfrak{S}_3$. Plus généralement, si
$n = n'd$, on a

(223)
\begin{tikzcd}
X_n \ar[r] & X_{n'} \\
x_i \ar[r, |->] & x_i^d
\end{tikzcd}

et de cette façon $X_n$ devient un revête-
ment principal (= torseur) sur $X_{n'}$, de
groupe $\mu_d^{\{0,1,\infty\}} / \mu_d$, de façon plus précise,
[unclear], (223) est compatible avec
l'épimorphisme

(224)
\begin{tikzcd}
F_n \ar[r, "\text{épi}"] \ar[d, equal] & F_{n'} \ar[d, equal] & g \ar[r, |->] & g \\
\mathfrak{S}_3 \cdot \mu_n^{\{0,1,\infty\}}/\mu_n \ar[r] & \mathfrak{S}_3 \cdot \mu_{n'}^{\{0,1,\infty\}}/\mu_{n'} & (x_0, x_1, x_\infty) \ar[r, |->] & (x_0^d, x_1^d, x_\infty^d)
\end{tikzcd}

\newpage

%% ===== Page 98 =====

et $\mu_n^{\{0,1,\infty\}}/\mu_n$, a [unclear: fortiori] $\mu_d^{\{0,1,\infty\}}/\mu_d$, opère libre-
-ment sur $X_n$, et [unclear: le morphisme de passage au]
quotient [unclear: induit] une suite exacte cour-
-te

(224)
$$1 \longrightarrow \mu_d^{\{0,1,\infty\}} / \mu_d \longrightarrow F_n \longrightarrow F_{n'} \longrightarrow 1$$
$$ \mathfrak{S}_3 \ltimes \mu_n^{\{0,1,\infty\}}/\mu_n \quad \mathfrak{S}_3 \ltimes \mu_{n'}^{\{0,1,\infty\}}/\mu_{n'} $$
$$ (x_0, x_1, x_\infty) \longmapsto (x_0, x_1, x_\infty) $$
$$ (x_0, x_1, x_\infty) \longmapsto (x_0^d, x_1^d, x_\infty^d) $$
$$ g \longmapsto g $$
pour $x_0, x_1, x_\infty \in \mu_n$
$g \in \mathfrak{S}_3$

et l'homomorphisme (223) est compa-
-tible avec (224), et induit un iso. (compa-
tible aux actions de $F_{n'}$)

(225) $$ X_{n'} \simeq X_n \Big/ \Big( \mu_d^{\{0,1,\infty\}}/\mu_d \Big) , $$

De plus, l'action de [unclear: le sous-groupe]
$$ F_n^\circ = \mu_n^{\{0,1,\infty\}}/\mu_n $$
sur $U_n$ est libre, et a fortiori celle
de $\mu_d^{\{0,1,\infty\}}/\mu_d$, de sorte que $U_n$ est
bien un torseur sur $U_{n'}$ de groupe
$$ F_d^\circ = \mu_d^{\{0,1,\infty\}}/\mu_d $$

Passant à la limite, on trouve

(226)
$$ X_\infty = \lim_{\longleftarrow n \in \mathbb{N}^\times} X_n , \quad U_\infty = \lim_{\longleftarrow} U_n $$
$$ F_\infty = \lim_{\longleftarrow} F_n = \left( \mathfrak{S}_3 \ltimes \hat{\mathbb{Z}}(1)^{\{0,1,\infty\}} / \hat{\mathbb{Z}}(1) \right) $$

\newpage

%% ===== Page 99 =====

avec opération de $F_\infty$ sur $U_\infty$ [unclear]

$$(227) \quad F_\infty^\circ = \lim_{\longleftarrow} F_n^\circ \simeq \widehat{T.C.}( \dots )$$

qui est librement agissante sur $U_\infty$, le quotient étant

$X_\infty \simeq U_{0,3}$.

Pour la theorie de Galois de $(U_{0,3}, \underline{S}_3)$,
le groupe profini $F_\infty$ sur $\mathbb{Q}$ [unclear crossed out text]

[MARGIN: N.B. L'action de $M_{0,3}$ sur $M \simeq \hat{\mathbb{Z}}^*$ (qui [unclear] l'action par automorphismes sur $F_\infty^\circ$) se factorise par l'homomorphisme $M_{0,3} \to \underline{S}_3$ et par le [unclear] de $\underline{S}_3 \to \operatorname{Aut}(M)$]

[crossed out] $(227) \quad M_{0,3} = \Pi_1(U_{0,3 \mathbb{Q}}, \underline{S}_3 ; \vec{-j}) \to$

le groupe [unclear]

$$(228) \quad F_\infty^\circ = F_\infty^\circ(\bar{\mathbb{Q}}) \simeq \hat{\mathbb{Z}}^{\{0,1,\infty\}} \rtimes \dots$$

sur lequel $\Gamma_\mathbb{Q}$ opère [unclear]

[unclear paragraph with heavy corrections] ... fait opérer $\underline{S}_3$ sur l'image ... l'action de $\mathfrak{S}_3$ sur $U_{0,3 \mathbb{Q}}$ [unclear] ... $M_{0,3} = \Pi_1(U_{0,3 \mathbb{Q}}, \underline{S}_3 ; \vec{-j})$

$$(229) \quad F_\infty \to [unclear]$$

$$(230) \quad U_\infty(\vec{-j}) \text{ un torseur sous } F_\infty^\circ, \text{ sur}$$

lequel $M_{0,3} = \Pi_1(U_{0,3 \mathbb{Q}}, \underline{S}_3 ; \vec{-j})$ opère
de façon compatible avec ses opérations
sur $F_\infty^\circ$.

Le dévissage de ce $(F_\infty^\circ, M_{0,3})$-torseur
[unclear] : [unclear] un revêtement
principal $M_\infty$ de $U_{0,3 \mathbb{Q}}$, de groupe $F_\infty^\circ$,

\newpage

%% ===== Page 100 =====

avec les opérations de $\mathfrak{S}_3$ dessus compatibles
avec le [unclear] de $\mathfrak{S}_3$ sur $F_\infty^\circ$.

D'autre part, la catégorie des torseurs
[unclear: sur des groupes abéliens] de la
forme $M \otimes_Z T$, où $T$ est un $\hat{Z}$-module
inversible, est équivalente à celle des
torseurs sous le groupe affine [unclear: restreint]
$$ \hat{Z}^* \cdot M^\wedge $$
La donnée d'une classe d'isomorphie
de tels torseurs, avec opération de
$M_{0,3}$ dessus compatible avec l'action
sur $M^\wedge$ [unclear] Considérons d'autre part
la suite exacte
(231) $$ 1 \to M \to G_{S_0/\mu_G} \to \mathcal{O}_{12}^* \rtimes (\mathfrak{S}_3)_{S_0} \to 1 $$
elle donne, par passage à
$( \ )(\hat{Z})^\Delta$ , une suite exacte
(232) $$ 1 \to M^\wedge \to G_{S_0/\mu_G}^\wedge \to \hat{Z}^* \rtimes \mathfrak{S}_3 \to 1 $$
qui redonne l'opération de [unclear]
de $\hat{Z}^* \rtimes \mathfrak{S}_3$ sur $M^\wedge$, Elle contient
la suite exacte
(233) $$ 1 \to M^\wedge \to [unclear] \to \mathfrak{S}_3 \to 1 $$
qui est un quotient de la suite exacte
dont [unclear]

\newpage

%% ===== Page 101 =====

698

$$ (234) \qquad 1 \to \Pi_1(U_{0,3 \overline{\mathbb{Q}}}, \overline{y}) \to \Pi_1(U_{0,3 \mathbb{Q}}, \mathfrak{S}_3; \overline{y}) \to \mathfrak{S}_3 \to 1 $$

par le sous-groupe invariant des commutateurs
du groupe $\Pi_1(U_{0,3 \overline{\mathbb{Q}}}, \overline{y})$. Cette structure
[MARGIN: si, elle est triviale, cf. plus bas...]
d'extension (233) de $\mathfrak{S}_3$ par $M^\wedge$ n'est pas
triviale, sauf erreur, bien que l'extension
induite de $\{1, \rho, \rho^2\} \simeq \mathbb{Z}/3\mathbb{Z}$ par $M^\wedge$, - ou
$\{1, \sigma\}$ par $\hat{M}$, le soit. A fortiori (232) n'est
pas triviale (bien qu'elle ait des scindages
canoniques au-dessus de $\hat{\mathbb{Z}}$, au-dessus de
$\{1, \rho, \rho^2\}$ et au-dessus de $\{1, \sigma\}$). Donc
(sous réserve de vérifier ces points tantot...)
le groupe $C_{\rho,\sigma}/\mu_a$ ne peut s'identifier
à un sous-groupe du groupe 
$$ \boxed{ Aff_{\hat{\mathbb{Z}}}(\hat{M}) \simeq \operatorname{Aut}_{\hat{\mathbb{Z}}}(\hat{M}). \hat{M} } $$
comprenant le sous-groupe $\hat{\mathbb{Z}}^\times \times \mathfrak{S}_3$ de $\operatorname{Aut}_{\hat{\mathbb{Z}}}(\hat{M})$.

Considérons le sous-groupe
$$ (235) \qquad Aff \simeq (\hat{\mathbb{Z}}^\times \times \mathfrak{S}_3) . \hat{M} $$
On a un diagr.

\begin{tikzcd}
(237) & 1 \ar[r] & \hat{M} \ar[r] & Aff \ar[r] & \hat{\mathbb{Z}}^\times \times \mathfrak{S}_3 \ar[r] & 1 \\
& & & & \mathcal{M}_{0,3} \ar[u]
\end{tikzcd}

: ligne exacte. Le donné de la structure
(230) équivaut à isoler (à un conjugué près, à l' [unclear]

\newpage

%% ===== Page 102 =====

donnée d'une classe de $\widehat{M}$-conjugaison d'un
relèvement de $M_{0,3} \to \widehat{\mathfrak{S}}_3 \times \hat{\mathbb{Z}}^*$ en un

[MARGIN: NB Cett hom. induit sur $\Pi_{0,3}$ l'hom [unclear] $1 \to (\Pi_{0,3}) \xrightarrow{\sim} \widehat{M}$ propriété inaltérée]

(238) $$M_{0,3} \longrightarrow Aff$$ .

Le choix d'un tel relèvement dépend du
choix d'un point de $U_0(\overline{\mathbb{Q}})$ au dessus du
point $\vec{1}$ de $U_{0,3}(\overline{\mathbb{Q}})$ - le chgmt de ce
point origine se fait par translation par un
$x \in \widehat{M}$, et le relèvement (238) est alors
changé par autom. int. par cet élément $x$.
Dans le [unclear] de F[unclear] fournit une
telle classe de conjugaison de relèvements (238)
dont la connaissance équivaut en principe
à celle de le [unclear] de F[unclear] à [unclear: isom.]
près. On voudrait mettre, si possible,
en [unclear] un [unclear] avec
les hom. (237), et plus particulièrement
sans doute un (238), qui fournit un
hom de $M_{0,3}$, non dans $Aff$, mais
dans le groupe analogue $G_{P^1 \setminus H}$
de (232), qui relève l'hom. $M_{0,3} \to \hat{\mathbb{Z}}^* \times \widehat{\mathfrak{S}}_3$

\newpage

%% ===== Page 103 =====

699

(239)
\begin{tikzcd}
1 \ar[r] & M^\wedge \ar[r] & G_{p\Gamma}/H_G \ar[r] & \hat{\mathbb{Z}}^* \times \mathfrak{S}_3 \ar[r] & 1 \\
& & \hat{\mathbb{Z}}^* \times G^{+ \wedge} / H_G \ar[u, "\simeq"'] & & \\
& & & \mathcal{M}_{0,3} = \pi_1(\mathcal{U}_{0,3 \overline{\mathbb{Q}}, \mathfrak{S}_3} ; -\vec{1}) \ar[uu] \ar[ul, dashed]
\end{tikzcd}

Sur $G_{0,3}^{+ \wedge} \simeq \pi_1(\mathcal{U}_{0,3 \overline{\mathbb{Q}}, \mathfrak{S}_3}, -\vec{1}) \subset \mathcal{M}_{0,3}$ , cet hom. induit l'hom. canonique du passage au quotient de $G_{0,3}^{+ \wedge}$ obtenu en rendant $\Pi_{0,3}^\wedge$ abélien , soit $G_{p\Gamma}/H_G \simeq G^{+ \wedge}/H_G$ , extension de $\mathfrak{S}_3$ par $M^\wedge$.

Reprenons l'hom. $\mathcal{M}_{0,3} \to \text{Aff}$ (238) , il s'insère dans un diag. commutatif

(240)
\begin{tikzcd}
1 \ar[r] & \Pi_{0,3}^\wedge \ar[r] \ar[d] & \mathcal{M}_{0,3} \ar[r] \ar[d] & \Gamma_{\mathbb{Q}} \times \mathfrak{S}_3 \ar[r] \ar[d, "\chi \times \text{id}"] & 1 \\
1 \ar[r] & \hat{M} \ar[r] & \text{Aff} \ar[r] & \hat{\mathbb{Z}}^* \times \mathfrak{S}_3 \ar[r] & 1
\end{tikzcd}

où $\Gamma_{\mathbb{Q}} \times \mathfrak{S}_3 \to \hat{\mathbb{Z}}^* \times \mathfrak{S}_3$ est l'hom. de (237) , et
où $\Pi_{0,3}^\wedge \to \hat{M}$ est l'hom. canonique

(241) $$ \Pi_{0,3}^\wedge \to (\Pi_{0,3}^\wedge)_{ab} \simeq ((\Pi_{0,3})_{ab})^\wedge \simeq \hat{M} \ . $$

Si on désigne par $\mathcal{M}_1$ le quotient de $\mathcal{M}_{0,3}$ par $[\Pi_{0,3}^\wedge, \Pi_{0,3}^\wedge]$ , on trouve de même

\newpage

%% ===== Page 104 =====

un hom. canonique de suites exactes

(242)
\begin{tikzcd}
1 \ar[r] & M^\wedge \ar[r] \ar[d, equal] & M_1 \ar[r] \ar[d] & N_{0,3} \ar[r] \ar[d] & 1 \\
1 \ar[r] & M^\wedge \ar[r] & Aff \ar[r] & \hat{\mathbb{Z}}^\times \times \mathfrak{S}_3 \ar[r] & 1
\end{tikzcd}

(dépends du choix d'un pt de $\mathcal{U}_{0,3}(\overline{\mathbb{Q}})$ au-dessus de $-\vec{j}$)
qui montre que l'extension $M_1$ de $\hat{\Gamma}_Q$ par $M^\wedge$
est [unclear: isome] : l'image inv. de l'extension
triviale $Aff$ de $\hat{\mathbb{Z}}^\times \times \mathfrak{S}_3$ par $M^\wedge$, donc est
triviale. Plus précisément encore, considérons
le quotient $M_2$ de $M_1$ par le noyau de $M_1 \to Aff$,
on trouve un diagramme

(242 bis)
\begin{tikzcd}
1 \ar[r] & M^\wedge \ar[r] \ar[d, "\simeq"'] & M_2 \ar[r] \ar[d, hook] & N_2 \ar[r] \ar[d, hook] & 1 \\
1 \ar[r] & M^\wedge \ar[r] & Aff \ar[r] & \hat{\mathbb{Z}}^\times \times \mathfrak{S}_3 \ar[r] & 1
\end{tikzcd}

(où $N_2 \overset{def}{=} M_2 / \text{Im } M^\wedge$). On voit donc
que l'extension de $N_2$ par $M^\wedge$ est triviale
- [unclear: avec un choix de] $M^\wedge$ - conjugaison -
de scindages de cette extension - les
scindages correspondant aux pts de
$\mathcal{U}_{0,3}(\overline{\mathbb{Q}})$ au-dessus du pt $-\vec{j} \in \mathcal{U}_{0,3}(\mathbb{Q})$

\newpage

%% ===== Page 105 =====

On aurait vraiment envie de montrer que
cet homomorphisme $M_{0,3} \to Aff$ est isomorphe
à l'hom. $M_{0,3} \to G_{\rho,0}/\mu_{\hat{\mathbb{Z}}}$ de (239) , moyennant
un isomorphisme canonique de
$Aff$ avec $G_{\rho,0}/\mu_{\hat{\mathbb{Z}}}$. Cela impliquerait que
cette dernière extension est scindée. Les
isomorphismes de cette extension avec $Aff$
correspondraient aux scindages de cette
extension.

12°) Scindages de l'extension $G_{\rho,0}/\mu_{\hat{\mathbb{Z}}}$ de $\hat{\mathbb{Z}}^* \times \mathfrak{S}_3$ par $M$.
Comme on a déjà un relèvement
$$ \hat{\mathbb{Z}}^* \hookrightarrow G_{\rho,0} $$
(150, 151), je vais regarder les scindages
de l'extension (231) qui prolongent celui-ci -
ce qui revient à, les
scindages de l'extension induite sur $\mathfrak{S}_3$ :

[MARGIN: partout : il faut se donner ce scindage qui $\mathfrak{S}_3 \hookrightarrow \underline{S}^+ / \mu_{\hat{\mathbb{Z}}}$ ($G_{\rho,0} \simeq \hat{\mathbb{Z}}^* \dots$)]
[MARGIN: On regarde les scindages [unclear] correspondent bien uniquement aux repères]

(241) $$ 1 \to M \to \underline{S}^+ / \mu_{\hat{\mathbb{Z}}} \to (\mathfrak{S}_3)_{S_0} \to 1 $$

\newpage

%% ===== Page 106 =====

(244) \quad (\rho', \sigma') = (\zeta \rho, \eta \sigma) \quad avec \quad \zeta, \eta \in \Gamma \underline{M}
satisfaisant
\rho'^3 = \sigma'^2 = 1, \quad \sigma'(\rho') = \rho'^{-1} .

Nous allons interpréter ces conditions en
termes de $\zeta, \eta$ . La condition $\rho'^3 = 1$ s'écrit
comme $\rho^3 = 1$ dans $(\mathfrak{S}_3)_{S_0}$
(245) $\quad \zeta \rho(\zeta) \rho^2(\zeta) = 1 \quad i.e. \text{ additivement } \quad (1+\rho_M+\rho_M^2)(\zeta) = 0$
cette condition est tjours satisfaite, car on a
(245) $\quad 1 + \rho_M + \rho_M^2 = 0$ .

La condition $\sigma'^2 = 1$ s'écrit
(246) $\quad \eta \sigma(\eta) = 1 \quad i.e. \text{ additivement } \quad (1+\sigma_M)(\eta) = 0$
et elle signifie que on a
(247) $\quad \eta = b \delta \quad \text{avec } b \in \Gamma \underline{G}_a$
$\delta = \lambda_1 - \lambda_0$

Nous poserons
(248) $\quad \zeta = a_0 \lambda_0 + a_1 \lambda_1 \quad \quad a_0, a_1 \in \Gamma \underline{G}_a$
et nous allons exprimer en termes
de $a_0, a_1, b$ la condition $\sigma'(\rho') = \rho'^{-1}$.

On a
$\sigma'(\rho') = (int \, \eta) \sigma(\zeta \rho) = (int \, \eta) [\sigma(\zeta) \underset{\rho^{-1}}{\sigma(\rho)}]$
la condition $\sigma'(\rho') = \rho'^{-1}$ s'écrit
$\sigma(\zeta) \rho^{-1} = \eta^{-1} \rho^{-1} \zeta^{-1} \eta$
soit, en multipliant par $\rho$, additivement :
$\rho(\sigma(\zeta)) - \eta = -\rho(\eta) - \zeta$

\newpage

%% ===== Page 107 =====

701

$$ \underset{\varepsilon_1 M}{\underbrace{\rho_M \sigma_M}} (\xi) + \lambda_1 = -\rho_M(\eta) - \xi + \eta \quad \text{ou encore} $$

(249) $$ \underset{(2a_1 - a_0)\lambda_1}{\underline{(1+(\varepsilon_1)_M)(\xi)}} + \lambda_1 = \underset{3b\lambda_1}{\underline{(1-\rho_M)(\eta)}} $$

ou encore, par un calcul immediat

(250) $$ 2a_1 - a_0 + 1 = 3b \quad \text{ou} \quad 2a_1 - a_0 - 3b = -1 $$
[MARGIN: (lire pour $\rho'$ donné)]

Ceci montre que pour $\xi$ [raturé: $\in \Gamma \underline{M}$] pour qu'il existe $\eta$ (lire: $\xi,\eta$) satisfaisant aux conditions (244) - i.e. de sorte que $(\rho' = \xi \rho, \sigma' = \eta \sigma)$ définisse un scindage de (243) - il f. et s. que l'on ait

(251) $$ 2a_1 - a_0 + 1 \equiv 0 \quad (3) $$

et alors $\eta$ est déterminé de façon unique par (250), au moins si 3 est section régulière de $\underline{O}_S$ sur la base $S$ envisagée (ce qui est le cas si $S = Spec \hat{\mathbb{Z}}$).

Ainsi les scindages de (243) forment un schéma $\underline{Sc}$ sur $S_0$, qui est canoniquement isomorphe par la

(252) $$ \underline{Sc} \simeq \{ (a_0, a_1, b) \in \Gamma E^3 \mid 2a_1 - a_0 - 3b = -1 \} $$
$$ \simeq \{ (\xi, \eta) \in \Gamma(\underline{M} \times \underline{M}) \mid (1+\varepsilon_M)(\xi) - (1-\rho_M)(\eta) = -\lambda_1 \} $$

ce schéma est de façon naturelle un...

\newpage

%% ===== Page 108 =====

[unclear: torseur sous le schéma , on en déduit]

[MARGIN: $K \simeq \mathbb{Z}^2$]
$$ (253) \qquad \underline{S}_c \simeq \underline{K} \quad \text{, où } K \text{ est le noyau de} $$
la suite exacte de
$\mathbb{Z}$-modules
$$ 0 \longrightarrow K \longrightarrow \mathbb{Z}^3 \longrightarrow \mathbb{Z} \longrightarrow 0 $$
$$ (\alpha_0, \alpha_1, \beta) \longmapsto 2\alpha_1 - \alpha_0 - 3\beta $$

On a d'autre part une opération
de $\underline{M}$ sur $\underline{S}_c$ par automorphismes intérieurs,
que je vais expliciter. À $\lambda \in \Gamma\underline{M}$
correspond l'opération
$$ (\xi, \eta) \longmapsto \lambda \xi (\rho_M \lambda)^{-1}, \lambda \eta \sigma_M(\lambda)^{-1} $$
soit, additivement

$$ (254) \qquad (\xi, \eta) \longmapsto ( \xi + (1 - \rho_M)(\lambda), \eta + (1 - \sigma_M)(\lambda) ) $$

Ainsi, l'opération en question provient de
la structure de $\underline{S}_c$ - torseur sur
$\underline{S}_c$ , et d'un homomorphisme de
groupes

$$ (255) \qquad \begin{cases} \underline{M} \longrightarrow \underline{S}_c \\ \lambda \longmapsto ( (1 - \rho_M)(\lambda), (1 - \sigma_M)(\lambda) ) \\ ((c_0+c_1)\alpha_0 + (2c_1-c_0)\alpha_1, (c_1-c_0)\beta) \quad \text{où } \lambda = c_0\lambda_0 + c_1\lambda_1 \end{cases} $$

qui en fait est associé à l'hom. de $\mathbb{Z}$-modules

$$ (256) \qquad \begin{cases} \underline{M} \longrightarrow \underline{K} \subset \underline{\mathbb{Z}}^3 \\ c_0\lambda_0 + c_1\lambda_1 \longmapsto (\alpha_0 = c_0+c_1, \alpha_1 = 2c_1-c_0, \beta = c_1-c_0) \end{cases} $$

\newpage

%% ===== Page 109 =====

702

On voit tout de suite que cet hom. est
universellement injectif (car la composante
$M \to \mathbb{Z}^3$ l'est) donc [unclear: passe] iso pour les
questions de torseurs. On a ainsi prouvé :

Proposition L'extension $\underline{\widehat{C}}^+ / H_{\underline{C}}$ (243) de
$(\underline{G}_3)_{S_0}$ par $\underline{M}$ est scindée sur $S_0$
— ou sous une autre forme — les scindages
sur $S_0$ d'un $\underline{Z}$ - [unclear: torseur] correspondent 1-1 (aux
scindages de $\underline{\widehat{C}}^+ / H_{\underline{C}}$ comme
extension de $\underline{G}_3$ par $\underline{M}$. Sur une
base quelconque $S$, deux scindages
sont conjugués par une unique
section de $\underline{M}$ sur $S$, donc le schéma
des scindages de (243) est de
façon naturelle un torseur sous $\underline{M}$.

Notons qu'il est immédiat, sur une
base quelconque $S$, que les scindages
de (243) correspondent biunivoquement
aux scindages de l'extension de
groupes

(244) $$1 \to \underline{M}(S) \to \underline{\widehat{C}}^+ / H_{\underline{C}}(S) \to \underline{G}_3 \to 1$$ ,

qui s'identifie ainsi à l'extension des

\newpage

%% ===== Page 110 =====

$\mathfrak{S}_3$ par $M(S) = \Gamma(S, \underline{\mathbb{Z}}) \otimes_{\mathbb{Z}} M$, déduite
de l'extension universelle $\widehat{\mathfrak{S}}^+ / H_{\widehat{\mathfrak{S}}}$ de
$\mathfrak{S}_3$ par $M$, via l'hom.
$$M \to M(S) = M \otimes_{\mathbb{Z}} \Gamma(S, \underline{\mathbb{Z}})$$

Appliquons ceci au cas $S = \operatorname{Spec}(\mathbb{Z})$,
— on trouve le

Corollaire L'extension $\widehat{\mathfrak{S}}^{+\wedge} / H_{\widehat{\mathfrak{S}}}$ de
$\mathfrak{S}_3$ par $M^\wedge$ est scindée, et ses
scindages forment un torseur sous
$M^\wedge$ (pour autom. intérieurs).

Étudions maintenant les
scindages de $\mathfrak{S}_{\rho, \sigma} / H_{\mathfrak{S}}$ [unclear: tout au moins]
Il s'agit de déterminer, pour
un scindage donné

(245) $\quad (\rho' = \xi \rho, \sigma' = \eta \sigma) \quad \begin{cases} \xi = a_0 \lambda_0 + a_1 \lambda_1 \\ \eta = b \delta = b(\lambda_1 - \lambda_0) \end{cases}$
de $\widehat{\mathfrak{S}}^+ / H_{\widehat{\mathfrak{S}}}$ au dessus de $(\mathfrak{S}_3)_{f_0}$
les relèvements
$$D_{12}^* \to \mathfrak{S}_{\rho, \sigma} / H_{\mathfrak{S}}$$
compatibles avec le scindage précédent, i.e.
qui correspondent à $\rho', \sigma'$. Pour ceci

\newpage

%% ===== Page 111 =====

70)

je vais déterminer le commutant de $(\rho',\sigma')$ dans $\underline{G}_{\rho,\sigma}/K_{\underline{G}}$, i.e. le commutant du sous-groupe section partielle
(258) $$H_{(\rho',\sigma')} \subset \underline{G}_{\rho,\sigma} / K_{\underline{G}}$$

L'image de ce commutant dans $(\underline{G}_{\underline{S}})_{\underline{S}_0} \simeq (\underline{G}_{\rho,\sigma} / \underline{M}) / O_{\hat{\mathbb{Z}}}^*$ est contenue dans le centre de $(\underline{G}_{\underline{S}})_{\underline{S}_0}$, donc est trivial, donc on conclut
(259) [unclear: crossed out text] $$H_{(\rho',\sigma')} \subset O_{\hat{\mathbb{Z}}}^* . \underline{M} = \underline{G}_{\rho,\sigma}^\natural$$
$$ \left\{ \lambda u_\mu \mid \lambda \in \underline{M}, \mu \in O_{\hat{\mathbb{Z}}}^* \right\} $$

qui est défini ds (150, 151). On a donc à déterminer le centralisateur de $H_{(\rho',\sigma')}$ en exprimant la condition pour un $\lambda u_\mu$ d'y être, i.e. de commuter à $\rho',\sigma'$.

On aura
$$(\lambda u_\mu)(\rho') = int(\lambda) u_\mu(\zeta \rho) = int(\lambda) u_\mu(\zeta) u_\mu(\rho) = int(\lambda)(\zeta^\mu \lambda_1^\nu \rho)$$

où $\nu = \nu(\mu)$, donc la relation $(\lambda u_\mu)(\rho') = \rho'$ s'écrit aussi
$$\lambda \zeta^\mu \lambda_1^\nu \rho(\lambda^{-1}) \rho = \zeta \rho \quad \text{i.e.} \quad \lambda \zeta^\mu \lambda_1^\nu \rho(\lambda^{-1}) = \zeta ,$$
ou en notation additive
$$\lambda + \mu \zeta + \nu \lambda_1 - \rho(\lambda) = \zeta \quad \text{avec } \mu-1=2\nu$$

\newpage

%% ===== Page 112 =====

(260) $(1-\rho)(\lambda) + 2\nu\zeta + \nu\lambda_1 = 0$

La relation $(\lambda u_{\mu})(\sigma') = \sigma'$ s'explicite, via
$(\lambda u_{\mu})(\sigma') = int(\lambda) u_{\mu}(\eta\sigma) = int(\lambda)(u_{\mu}(\eta) u_{\mu}(\sigma)) = \lambda \eta^\mu \sigma \lambda^{-1}$
$= \lambda \eta^\mu \sigma(\lambda^{-1}) \sigma$
[MARGIN: $(\lambda u_{\mu})(\sigma') = \sigma'$]
[MARGIN: $\eta^\mu \omega^2 \sigma = \sigma$ car $\omega=1$ dans $G_{\rho,\sigma}/H_G$]

donc la relation s'écrit $\lambda \eta^\mu \sigma(\lambda^{-1}) = \eta$

ou encore, additivement
$\lambda + \mu\eta - \sigma(\lambda) = \eta$

(261) $(1-\sigma)\lambda + 2\nu\eta = 0$

En termes de $a_0, a_1, b$ de (257), ces
équations s'explicitent comme des conditions
sur
$$ \lambda = x\lambda_0 + y\lambda_1 $$

par les équations
(262)
$$ \begin{cases} -x+y = - 2\nu b & (exprime (261)) \\ x+y = - 2\nu a_0 & \} (exprime (260)) \\ -x+2y = -\nu(2a_1+1) \end{cases} $$

où les quantités $a_0, a_1, b$ figurent
d'ailleurs [unclear: liées par des facteurs]
multiplicatifs $\nu = \nu(\mu)$ satisfaisant (250)
(263) $-a_0 + 2a_1 - 3b + 1 = 0$

La première et troisième équation (262)
se résolvent en $x, y$
(264)
$$ \begin{cases} y = \nu(-2a_1+b-1) \\ x = y + 2\nu b = \nu(-2a_1+3b-1) \end{cases} $$

et la deuxième équation (262) devient, via

\newpage

%% ===== Page 113 =====

704

y portant ces valeurs de $x, y$, une relation
liant $a_0, a_1, b, \nu$ :
$$ \nu(-4a_1+6b-2) = -2\nu a_0 \quad \text{i.e.} $$
$$ 2\nu(a_0-2a_1+3b-1)=0 $$
qui est une conséquence de (263)! Ainsi
on trouve : [unclear: que pour un scindage]

Proposition. Pour tout scindage de l'ext.
$\widehat{\mathbb{S}}^+/\mu_2$ de $(\widehat{\mathfrak{S}}_3)$ par $M$ ([unclear: C sen bin]
[MARGIN: ds $G_{\rho,\sigma}/\mu_G$]
quelconque), le commutant (dans ledit
groupe section est un sous-groupe
section de l'extension $G_{\rho,\sigma}$ de $\widehat{\Phi}_{12}^*$
par $M$. [Le relèvement
(264) $$ u_{\rho',\sigma'} : \widehat{\Phi}_{12}^* \to G_{\rho,\sigma} $$
qui est associé à $\rho',\sigma'$ donnés par (257)
$$ \rho' = \xi \rho, \sigma' = \eta \sigma \qquad \begin{cases} \xi = a_0 \lambda_0 + a_1 \lambda_1, \eta = b \delta = b(\lambda_1 - \lambda_0) \\ a_0, a_1, b \in \hat{\mathbb{Z}} \end{cases} $$
est donné par
(265) $$ u_{\rho',\sigma'}(t) = \lambda u_{\rho,\sigma} $$
où $u_{\rho,\sigma}$ est défini en (150, 151), et
où $\lambda = x \lambda_0 + y \lambda_1$, avec $x, y \in \hat{\mathbb{Z}}$ donnés
par 264, et où $\nu = \nu(t)$ ).

Corollaire. Le morphisme canonique
au dessus de
(266) $$ Scindages (G_{\rho,\sigma}/\mu_G \text{ sur } \widehat{\Phi}_{12}^* \times \widehat{\mathfrak{S}}_3) $$
$$ \longrightarrow Scindages(\widehat{\mathbb{S}}^+ \text{ au-dessus de } \widehat{\mathfrak{S}}_3) $$

\newpage

%% ===== Page 114 =====

est un isomorphisme de faisceaux
dans $(Sch)_{et}$ , et par suite (prop. p. 702)
le premier est un $\underline{M}$ torseur via
conjugaison.

Cet énoncé s'explicite en un énoncé
correspondant sur les extensions ensemblistes
comme dans le cor. à la prop. précédente.
On trouve en particulier, sur $S = \operatorname{Spec}(\mathbb{Z})$

Corollaire Les scindages de l'extension
$G_{\rho,\sigma}/k_G$ de $\hat{\mathbb{Z}}^\times \times \hat{\mathfrak{S}}_3$ par $\hat{M}$ forment
(par restriction) un torseur sous $\hat{M}$, isomorphe au torseur
des scindages de l'ext. $\hat{\mathfrak{S}}^+ / k_G$ de
$\hat{\mathfrak{S}}_3$ par $\hat{M}$.

Variante de l'étude précédente pour
l'extension $G_{\rho,\sigma}$ de $\mathbb{Q}_{12}^* \times H_{S_0}^+$ par $\underline{M}$
(sans division par $k_G = \{1, \omega_{S_0}\}$ ). Conduit
On réduit la question des scind. à ceux
de la sous-extension $\hat{\mathfrak{S}}^+$ de $H_{S_0}^+$ par
$\underline{M}$ - ils correspondent aux couples

\newpage

%% ===== Page 115 =====

705

(267) $$ \begin{cases} \rho' \xi \rho' \sigma' = \eta \sigma' \quad (\xi, \eta \in \Gamma \underline{M}) \quad \text{tels qu'on ait} \\ \rho'^3 = \sigma'^2 = \omega, \quad \sigma'(\rho') = \rho'^{-1} \end{cases} $$

(On notera qu'un scindage doit transformer $\omega$, qui opère trivialement sur $\underline{M}$ et est central dans $\underline{H}^+$, [unclear: dans une section centrale ou l'autre] on vérifie qu'on a :
(268) $$ \begin{cases} \underline{M} \cap Cent(\underline{\widehat{G}}^+) = \underline{M}^{\widehat{G}_3} = \{1\} \quad \text{, d'-} \\ Cent(\underline{\widehat{G}}^+) \simeq \mathcal{K}_{\underline{G}} = \{1, \omega\} \end{cases} $$
— ce qui montre qu'un scindage doit transformer $\omega$ en $\omega$, d'où finalement la condition (267) sur $\rho', \sigma'$ correspondante à un scindage). On a alors
$$ \rho'^3 = (\xi \rho)^3 = \xi \rho(\xi) \rho^2(\xi) \rho^3 = \xi \rho(\xi) \rho^2(\xi) \omega $$
$$ \sigma'^2 = (\eta \sigma)^2 = \eta \sigma(\eta) \sigma^2 = \eta \sigma(\eta) \omega $$
donc les deux conditions $\rho'^3 = \omega$ et $\sigma'^2 = \omega$ équivalent respectivement à :
$$ \begin{cases} \xi \rho(\xi) \rho^2(\xi) = 1 & \text{i.e.} \quad (1 + \rho_M + \rho_M^2)(\xi) = 0 \\ \eta \sigma(\eta) = 1 & \text{i.e.} \quad (1 + \sigma_M)(\eta) = 0 \end{cases} $$

i.e. cela se réduit aux équations (245) (246), et l'on a vu que la première condition sur $\xi$ est automatiquement satisfaite, celle sur $\eta$ s'écrit : (247) $\eta = b \delta = b (1-\lambda_0)$, $\lambda \in \Gamma \widehat{G}_3$.
La condition $\sigma'(\rho') = \rho'^{-1}$ s'explicite d'autre part comme à la p. 700', et redonne

\newpage

%% ===== Page 116 =====

sur l'équation (249), qui s'explicite par (250).

Cela montre qu'il existe un scindage
[unclear: de un $S_0$-type $Z$] - ce qui revient à un
scindage [unclear: de la suite] donnant $\underline{\mathbb{G}}^+$ ext. de
[MARGIN: $\simeq \underline{\mathbb{G}}_0 / \underline{M}_0$] $\underline{H}^+ = \underline{\mathbb{G}}^+ / \underline{M}$ (où $\underline{M} \simeq (\mathbb{F}_p)_S$) - et que
les scindages en question forment un
torseur sur $\underline{K}$, où $\underline{K}$ est défini dans
(253), et en fait [unclear: $\underline{K}$ est isom à]
$\underline{M} \simeq \underline{K}$ (254), cet $\underline{M} \simeq \underline{K}$ correspond
à l'opération de $\underline{M}$ sur le schéma des
scindages, par conjugaison. Ainsi on trouve

Proposition ([unclear: Sur une base] $S$ quelconque).
L'extension $\underline{\mathbb{G}}^+$ de $\underline{H}^+$ par $\underline{M}$
est scindée - deux scindages sont
conjugués par une unique section
de $\underline{M}$ - donc l'application de
"passage au quotient par $\mu_2$" définit une
bijection de l'ens. des scindages
de $\underline{\mathbb{G}}^+$ sur $\underline{H}^+$, sur l'ens des
scindages de $\underline{\mathbb{G}}^+ / \mu_2$ sur $\underline{H}^+ / \mu_{H^+} \simeq \underline{\mathfrak{S}}_3$.

\newpage

%% ===== Page 117 =====

705

Etudions maintenant le centralisateur
dans $\underline{G}_{\rho,\sigma}$ d'un sous-groupe section
au dessus de $H^+$, disons noté $H(\rho',\sigma')$.
On doit avoir maintenant

$$ (269) \quad Cent_{\underline{G}_{\rho,\sigma}} (H(\rho',\sigma')) \subset \Phi_{12}^* \cdot (M \times \mu_2) $$
$$ \text{avec } \mu_2 = \{1, \omega\}_{S_0} $$
$$ = \left\{ \lambda \omega^i \mu \Big| \begin{array}{l} \lambda \in \Gamma M \\ \mu \in \Gamma \Phi_{12}^* \\ i \in \mathbb{Z}/2\mathbb{Z} \end{array} \right\}_{S_0} $$

D'ailleurs $\lambda \omega^i \mu$ centralise $H(\rho',\sigma')$ ssi $\lambda \mu$ centralise. La commutation à $\rho'$
s'explicite encore par la condition (260),
le calcul p. 703 s'appliquant sans variation. Pour
la condition de commutation à $\sigma'$, le
calcul de p. 703 devient en modifiant les
- [unclear: vu que plus ici] $\omega^2 = 1$, la condition
s'écrit
$\omega^i \lambda \eta^{H_0}(\lambda^{-1}) \sigma = \eta \sigma$ i.e. $\omega^i \lambda \eta^{H_0}(\lambda^{-1}) = \eta$
et équivaut à la conjonction de la
condition

$$ (270) \quad v_2(\mu) = 0 \quad i.e. \quad \mu \in \Gamma \Phi_{12}^{* \prime \prime} $$

et de la relation (261). En présence
de $v_2(\mu) = 0$, les calculs de p. 703, 704
s'appliquent, [unclear: fournissant le théorème] :

Corollaire Pour tout sous-groupe section de $\underline{G}^+$
au dessus de $H^+$, le centralisateur dans
$\underline{G}_{\rho,\sigma}$ du sous-groupe section $H(\rho',\sigma')$

\newpage

%% ===== Page 118 =====

On a alors une suite exacte

(271) $$1 \to K \underline{C}^+ \to T \to \mathbb{Q}_{12}^{* \prime\prime} \to 1$$

et posant $T^{\natural} = T \cap G_{\rho, \sigma}^{\natural}$, on a

(272) $$T^{\natural} \simeq \mathbb{Q}_{12}^{* \prime\prime},$$

i.e. $T^{\natural}$ est une section partielle de
$G_{\rho, \sigma}^{\natural}$ au dessus de $\mathbb{Q}_{12}^{* \prime\prime}$.

Corollaire L'extension $G_{\rho, \sigma}$ de $\mathbb{Q}_{12}^* \times \underline{H}^+$
par $\underline{M}$ n'est scindée sur aucune base $S$
non vide. L'extension $G_{\rho, \sigma}^{\prime\prime}$ de $\mathbb{Q}_{12}^{* \prime\prime} \times \underline{H}^+$
par $\underline{M}$ est scindée part. sur tte base $S$,
et deux scindages sont conjugués par une
unique section de $\underline{M}$. L'application restriction
fournit une bijection de l'ens des
scindages de $G_{\rho, \sigma}^{\prime\prime}$ sur $\mathbb{Q}_{12}^{* \prime\prime} \times \underline{H}^+$, avec
l'ens des scindages de $\underline{C}^+$ sur $\underline{H}^+$.

Corollaire L'ext. $G_{\rho, \sigma}$ de $\hat{\mathbb{Z}}^* \times H^+$ par
$M^{\wedge}$ n'est pas scindée. Celle de $G_{\rho, \sigma}^{\prime\prime}$ en
$\hat{\mathbb{Z}}^{* \prime\prime} \times H^+$ par $M^{\wedge}$ est scindée, et deux
scindages sont conjugués par un unique
élément de $\hat{M}$. L'application restriction
définit une bijection de l'ens des
scindages de $G_{\rho, \sigma}^{\prime\prime}$ sur $\hat{\mathbb{Z}}^{* \prime\prime} \times H^+$ avec l'ens des
scindages de $\underline{C}^{\wedge}$ sur $H^+$.

\newpage

%% ===== Page 119 =====

707

Explicitons un scindage de l'extension
$\underline{G}^+$ de $H^+$ par $M$ - correspondant à un
scindage de $\underline{\tilde{G}}^+$ de $H^+$ par $\underline{\hat{M}}$ - , via
une solution particulière de l'équation
(273)
$$a_0 - 2a_1 + 3b - 1 = 0$$
en prenant
$$a_1 = b = 0 \quad \text{d'où} \quad a_0 = 1$$

i.e.
(274)
$$ \begin{cases} \xi = \lambda_0, \eta = 0 \text{ (écriture additive)} \\ \rho' = \lambda_0 \rho, \quad \sigma' = \sigma \end{cases} $$
Ces formules fournissent un scindage de $\underline{G}^+ / \mu_2$
en dessus de $\underline{S}_3 \simeq H^+ / \mu_{H^+}$. Explicitons
les scindages correspondants de $\mathbb{Q}_{12}^*$
dans $G_{\rho,\sigma} / \mu_a$ resp. de $\mathbb{Q}_{12}^{*\ \prime\prime}$ dans $G_{\rho,\sigma}$,

(275)
$$ \begin{cases} u_{\rho,\sigma} : \mathbb{Q}_{12}^* \to G_{\rho,\sigma} / \mu_a \quad \text{resp. } \mathbb{Q}_{12}^{*\ \prime\prime} \to G_{\rho,\sigma} \\ \mu \longmapsto \lambda_0^\nu u_\mu \quad \quad \quad (\nu = \nu(\mu)) \end{cases} $$

où $u_\mu$ défini dans (151). (NB (275) est contenu
dans prop. p. 704 (265) et (264), qui donne
$$u_{\rho,\sigma} = \lambda u_\mu \quad \lambda = x \lambda_0 + y \lambda_1, \quad x = y = -\nu$$
d'où $\lambda = -\nu(\lambda_0 + \lambda_1) = \nu \lambda_0$.

On aura donc, posant
$$u'_\mu = u_{\rho,\sigma}(\mu)$$

(276)
$$ \begin{cases} u'_\mu(\rho) = \rho', \ u'_\mu(\sigma) = \sigma' \\ u'_\mu(\lambda) = \mu \lambda \ (\lambda \in \Gamma \underline{M}) \ (\text{où } \mu = \mu(\mu)) \end{cases} \Bigg\} \text{(ce qui d'ailleurs suffit à caractériser } u'_\mu) $$

\newpage

%% ===== Page 120 =====

et on en déduit aussi
$$ (277) \begin{cases} u'_{\underline{\mu}}(\rho) = \lambda_0^{-2\nu} \rho \\ u'_{\underline{\mu}}(\sigma) = \omega^{\nu_2} \sigma \end{cases} \quad (= \sigma \text{ dans } \underline{G}_{g,r}/\mu_G \dots). $$

\newpage

%% ===== Page 121 =====

70bis

13°) Retour sur la tour de Fermat.

Commençons par dégager une suite topologique générale, qui était un peu embrouillée entre les lignes dans 110).

Soit $\Gamma$ (ici $\Gamma_{\mathbb{Q}} \times \mathfrak{S}_3 = \mathcal{N}_{0,3}$) opérant sur un topos $\mathcal{X}$ (ici $\underline{M}_{0,3} \simeq M_{0,4 \overline{\mathbb{Q}}}$) conn. loc. connexe, muni d'un "pt" (géométrique) $x_0$ (ici : le point $\vec{1}$ à valeurs dans $\mathbb{Q}$), soient

[MARGIN: en fait on discute [unclear: localement] ($\mathcal{X}$ loc. simplemt [unclear: connexe] ou profini, il faudra y revenir...)]

$$ (278) \begin{cases} \Pi = \Pi_1(\mathcal{X}, x_0), \mathcal{M} = \Pi_1(\mathcal{X}, \Gamma; x_0) \\ M = \Pi_{ab} = \Pi/[\Pi,\Pi] \quad, \tilde{\mathcal{M}} = \mathcal{M}/[\Pi,\Pi] \end{cases} $$

d'où des suites exactes et hom. de suites exactes

(279)
\begin{tikzcd}
1 \ar[r] & \Pi \ar[r] \ar[d, "epi"'] & \mathcal{M} \ar[r] \ar[d, "epi"'] & \Gamma \ar[r] \ar[d, "\sim"'] & 1 \\
1 \ar[r] & M \ar[r] & \tilde{\mathcal{M}} \ar[r] & \Gamma \ar[r] & 1
\end{tikzcd}

(Au lieu de prendre $M = \Pi/[\Pi,\Pi] = \Pi_{ab}$ on aurait pu prendre un quotient de $\Pi_{ab}$ invariant par action de $\Gamma$, mais peu importe...). Considérons "le rev. universel abélien" $\tilde{\mathcal{X}}_{ab}$ de $\mathcal{X}$ en $x_0$, i.e. le quotient par $[\Pi,\Pi]$ du rev. universel de $\mathcal{X}$ en $x_0$ - c'est donc un rev. principal de $\mathcal{X}$, de groupe $M$. On

\newpage

%% ===== Page 122 =====

Je propose d'examiner la façon de relever l'action de $\Gamma$ sur $\mathcal{H}$ en une action sur $\mathcal{H}'_0$, compatible avec l'action de $\Gamma$ sur le groupe d'opérateurs $M$
[unclear crossed out text]

Plus généralement, si $\mathcal{H}'$ est un $M$-torseur sur $\mathcal{H}$ isomorphe à $\mathcal{H}'_0$, se déduit de $\mathcal{H}'_0$ par twist : à l'aide d'un $M$-torseur $E$ ($=$ fibre de $\mathcal{H}'$ en $x_0$) étudier les relèvements analogues de l'action de $\Gamma$ sur $\mathcal{H}$ en une action de $\Gamma$ sur $\mathcal{H}'$, compatible avec l'action de $\Gamma$ sur $\mathcal{H}$.

Notons que, si $\tilde{\mathcal{H}}$ désigne le rev. universel de $\mathcal{H}$ en $x_0$, $\mathcal{M}$ opère sur $\tilde{\mathcal{H}}$ de façon compatible avec sa projection canonique sur $\mathcal{H}$, en prolongeant l'action de $\Pi$ sur $\tilde{\mathcal{H}}$ - comme $[\Pi,\Pi] \subset \Pi$ est distingué dans $\mathcal{M}$, cette action passe au quotient en une action canonique de $\tilde{\mathcal{M}} = \mathcal{M}/[\Pi,\Pi]$ sur $\mathcal{H}'_0 = \tilde{\mathcal{H}}/[\Pi,\Pi]$, qui satisfait aux

\newpage

%% ===== Page 123 =====

[MARGIN: 70^g]

[crossed out: deux] conditions analogues. Soit $\Theta_0$ cette opération-ci
de $\overline{\mathcal{G}}$, et $\Theta$ une opération quelconque satisfai-
sant aux dites conditions.

On veut donc faire de $\mathcal{X}'$, qui est un
rev. de $\mathcal{X}$, donné par sa fibre $E$ en $x_0$
et l'action de $\pi = \pi_1(\mathcal{X}, x_0)$ sur $E$ (par
translations, grâce à une structure de
$M$ torseur sur $E$) un revêtement
de $(\mathcal{X}, \overline{\mathcal{G}})$ - i.e. prolonger l'action
donnée de $\pi$ sur la fibre $E$, en une
action de $\mathcal{M}$ sur $E$ - qui d'ailleurs
se factorisera nécessairement par $\overline{\mathcal{M}}$.
Je dis que cette action se fera néc^t à
travers le groupe $\mathcal{A}ut$ du torseur
$E$ - de façon précise, $\Gamma$ opère sur
$M$, soit $\gamma_M$ l'opération définie par $\gamma \in \Gamma$,
et soit $g \in \mathcal{M}$ au dessus de $\gamma \in \Gamma$,
- on aura alors, pour $x \in E$, $\lambda \in M$

(280) $$g(\underset{\in M}{\lambda} . \underset{\in E}{x}) = \gamma_M(\lambda) . g x$$

ce qui [unclear: est] trivial par le calcul
$$g(\lambda . x) = (g \lambda) x = (g \lambda g^{-1} g) x = (\underbrace{g \lambda g^{-1}}_{\gamma_M(\lambda)}) \cdot (g x)$$

\newpage

%% ===== Page 124 =====

Donc si on introduit le groupe $Aff(E)$
on aura un homomorphisme de
suites exactes

(281)
\begin{tikzcd}
1 \ar[r] & M \ar[r] \ar[d, "id_M"] & \tilde{\Gamma} \ar[r] \ar[d, "\theta_{\tilde{\Gamma}}"] & \Gamma \ar[r] \ar[d, "\Theta_\Gamma"] & 1 \\
1 \ar[r] & M \ar[r] & Aff(E) \ar[r] & \operatorname{Aut}(M) \ar[r] & 1
\end{tikzcd}

où les flèches verticales extrêmes $id_M$
et $\Theta_\Gamma$ sont connues, et où la flèche médiane
est à déterminer. On peut dire
que la catégorie des $\tilde{X}'$ au dessus
d'un [unclear: l'é.p.] de $\Gamma$ dessus satisfaisant
aux autres [unclear: axiomes] ) équivaut
à celle des couples $(E, \theta_{\tilde{\Gamma}})$, où $E$
est un $M$-torseur, et $\theta_{\tilde{\Gamma}}$ un
hom. de groupes $\tilde{\Gamma} \to Aff(E)$
compatible avec $id_M$ et $\Theta_\Gamma$ i.e. qui
rend commutatif (281). On
peut dire aussi une fois $E$ choisi,
(et ce choix est unique à isom. [unclear: unique] près -)
qu'on a défini une extension
$Aff(E)$ de $\operatorname{Aut}(M)$ par $M$, d'où
[unclear: tirée par] $\Gamma \to \operatorname{Aut}(M)$ une extension

\newpage

%% ===== Page 125 =====

[MARGIN: les scindages de $Aff(E)$ déterminés par les pts de $E$ diffèrent de l'origine d'un scindage par l'aut. intérieur...]

$\widetilde{71^\circ}$ qui est d'ailleurs scindable.

par $M$, (Ceci dit, les relèvements
de $\Gamma$ au-dessus [unclear] correspondent aux
isomorphismes de cette extension
de $\Gamma$ par $M$, avec l'extension $\widetilde{G}$ elle-
même. Donc elle existe ssi l'extension de
$\widetilde{G}$ de $\Gamma$ par $M$ est scindée, i.e. si sa classe
(282) $c^2(\widetilde{G}) \in H^2(\Gamma, M)$ est nulle.
S'il existe un tel relèvement, i.e. un tel
isomorphisme, ceux-ci forment un torseur
sous le groupe des automorphismes de
l'extension (triviale) $\widetilde{G}$ de $\Gamma$ par $M$
qui s'identifie au groupe
$$Z^1(\Gamma, M)$$
des homomorphismes croisés de $\Gamma$ dans $M$.
Les classes de $M$-conjugaison de
tels relèvements (= isomorphismes)
correspondent donc aux éléments de
$$H^1(\Gamma, M)$$

Une autre façon de le dire, en prenant
$E = \text{torseur trivial } \simeq M$, d'où $Aff(E)$ scindée
est que les relèvements cherchés corres-
pondent aux scindages de l'extension
$\widetilde{G}$ de $\Gamma$ par $M$, et les classes de $M$-conju-
gaison de relèvements d'actions de $\Gamma$ sur $\widetilde{X}$ !

\newpage

%% ===== Page 126 =====

correspondent aux classes de $M$-conjugaison de scindages de ladite extension $\tilde{G}$ de $\Gamma$ par $M$.

Ceci bien posé, revenons à la situation de $(U_{0,3 \overline{Q}}, \Gamma = \Pi_Q \times \mathcal{S}_3)$, les $\Gamma$-revêtements de $\mathcal{X}$ correspondent aux $\mathcal{S}_3$-revêtements de $U_{0,3 Q}$, i.e. aux [unclear] de $(U_{0,3 Q}, \mathcal{S}_3)$. 

[MARGIN: NB le groupe qu'on avait noté $M$ devient ici $M_{0,3}$, $\hat{M}$ devient $\hat{M}_{0,3}$, $\Pi$ devient $\Pi_{0,3}$]

À l'aide des constructions générales du $§ 4$ $\underline{VI}$, [unclear: on déduit des quotients]
$$ \hat{\mathcal{C}}^+ / \mu_2 \simeq \hat{S}_{0,3} / [\hat{\Pi}_{0,3}, \hat{\Pi}_{0,3}] = \tilde{M} $$

[MARGIN: on prendra ici la notation $\tilde{M}$ !]

(avec les notations de la page précédente, l'action de $\mathcal{S}_3 \subset \Gamma (= \Pi_Q \times \mathcal{S}_3)$ sur $\tilde{M}$, on en déduit un groupe profini, $\tilde{G}$ qui était noté $\hat{G}_{g, \nu}$ dans les sections précédentes, qui y apparaissait comme un produit semi-direct
$$ \tilde{G}_{g, \nu} = \hat{\mathbb{Z}}^* \cdot \tilde{M} \quad \text{, ext. de } \hat{\mathbb{Z}}^* \times \mathcal{S}_3 \text{ par } \tilde{M} $$
et un hom. canonique
$$ \tilde{M} = M_{0,3} / [\Pi_{0,3}, \Pi_{0,3}] \longrightarrow \tilde{G}_{g, \nu} $$
s'insérant dans un [unclear] de suites exactes

\newpage

%% ===== Page 127 =====

711

(28bis)
\begin{tikzcd}
1 \ar[r] & \hat{M} \ar[r] \ar[d] & \tilde{\mathcal{M}} \ar[r] \ar[d] & \Gamma \ar[r] \ar[d, "N_{0,3}=\Gamma_{\mathbb{Q}} \times \mathfrak{S}_3" description, "\hat{\mathbb{Z}}^\times \times id_{\mathfrak{S}_3}"] & 1 \\
1 \ar[r] & \hat{M} \ar[r] & G_{\tilde{\rho}} \ar[r] & \hat{\mathbb{Z}}^\times \times \mathfrak{S}_3 \ar[r] & 1
\end{tikzcd}

D'ailleurs comme l'on attend - on a trouvé que
l'extension $G_{\tilde{\rho}}$ de $\hat{\mathbb{Z}}^\times \times \mathfrak{S}_3$ par $\hat{M}$ était
scindée, d'où il résulte que
l'ext. $\tilde{\mathcal{M}}$ de $\Gamma$ par $\hat{M}$ l'est aussi
(ce qui au demeurant est aussi conséquence de
l'existence de la tour de Fermat, en
y appliquant les réflexions générales
des pages précédentes). Bien mieux,
on a vu qu'il y a une unique classe
de $\hat{M}$-conj. de scindages de l'extension
$G_{\tilde{\rho}}$, ce qui induit une classe de
$\hat{M}$-conjugaison canonique (mais non
unique!) de scindages de $\tilde{\mathcal{M}} \to$ 
de $\Gamma$. Cela définit donc une
$\hat{M}$-conjugaison canonique de relèvements
sur $\tilde{X} = \widetilde{U_{0,3 \overline{\mathbb{Q}}}}$
de l'action de $\Gamma$ (à $\simeq$ près, ou ce
qui revient au m., un foncteur sur
$(U_{0,3 \mathbb{Q}}, \mathfrak{S}_3)$, de fibre un groupe profini

[MARGIN: NB $H^1(\Gamma, \hat{M})$ est un groupe fort gros! Mais voir cf th. p. 716]

\newpage

%% ===== Page 128 =====

commutatif au dessus de $(\mathcal{U}_{0,3\mathbb{Q}}, \mathfrak{S}_3)$
qui au dessus de $\mathcal{U}_{0,3\overline{\mathbb{Q}}}$ soit isomorphe
au rev. $\hat{\mathscr{X}}^!$ envisagé plus haut.

Revenons pour un instant à
ces généralités de tantôt. Notons que
la donnée d'un diagramme (281)
définit un quotient de $\tilde{M}$ à
savoir le coimage de $\theta_{\tilde{M}}$
donnant lieu à une suite
exacte

$$ (283) \quad 1 \longrightarrow M \longrightarrow \tilde{M} \longrightarrow \underset{\stackrel{def}{=} \tilde{M}/Im M}{\tilde{\Gamma}} \longrightarrow 1 $$

où on a

$$ (284) \quad \tilde{\Gamma} \hookrightarrow \operatorname{Aut}(M) $$

$\tilde{\Gamma}$ s'identifie à l'image de $\Gamma$ dans $\operatorname{Aut}(M)$
par l'action de $\Gamma$ sur $M$ provenant de
la structure d'ext. $\tilde{M}$ de $\Gamma$ par $M$.
Ce quotient $\tilde{\Gamma}$ de $\tilde{M}$, donnant
naissance à une ext. (283), n'est
pas quelconque - à part le fait que
$M$ s'y envoie injectivement, et que $\tilde{\Gamma} = \tilde{M}/Im M$

\newpage

%% ===== Page 129 =====

opère [unclear: fidèlement] sur $M$ : en plus, l'
extension (283) est triviale, et le fait
qu'elle provienne d'un $\tilde{\theta}$ implique
une structure supplémentaire dans
cette extension, à savoir la donnée d'une
classe de $M$-conjugaison de scindages.
Plus précisément, on a une application
canonique

$$ (285) \qquad E \longrightarrow Scind(\tilde{M} \text{ comme ext. de } \tilde{\Gamma} \text{ par } M) $$

compatible avec les opérations de $M$
sur $E$ (par translation) et sur $Scind$
(par conjugaison sur les scindages).
La classe de $M$-conjugaison envisagée des scindages est l'image.
Pour que $x$ et $y \in E$ aient même image
dans $Scind$, il f. et s. que l'élément
$u = y-x \in M$ soit dans $M^\Gamma = M^{\tilde{\Gamma}}$ ,
ce qui montre que

$$ (286) \qquad \text{l'hom. } E \longrightarrow Scind \text{ est injectif} $$

ssi $M^\Gamma = \{0\}$ .

Dans le cas où $M^\Gamma = \{0\}$ , on voit
donc que $E$ s'identifie [unclear: exactement] à
une classe de $M$-conjugaison de
scindages de $\tilde{M}$.

\newpage

%% ===== Page 130 =====

Résumons ces considérations, dans une

Proposition Soient $\mathcal{X}$ un topos connexe
loc. connexe muni d'un pt géom.
$x_0$, ou de plus dans le contexte discret
($\mathcal{X}$ loc simple connexe) ou profini
(revêtements [unclear] profinis).
On a un groupe $\Gamma$ d'autom. de $\mathcal{X}$
($\Gamma$ profini dans le cas profini - il faudrait
expliciter un peu bien faire les hypothèses
raisonnables à faire pour tenir compte
de la topologie profinie de $\Gamma$...). Posons
$\pi = \pi_1(\mathcal{X}, x_0)$, $M = \pi_1(\mathcal{X}, \Gamma; x_0)$ d'où
(*) $1 \to \pi \to M \to \Gamma \to 1$
où on se donne un quotient $M'$ de $\pi$, sur
form. un sous-groupe distingué de $M$,
ce qui revient à un quotient de
$\pi_{ab} = \pi/[\pi,\pi]$ stable par l'action de $\Gamma$,
d'où [unclear] extension quotient de (*),

(287) $1 \to M \to \tilde{M} \to \Gamma \to 1$

où $M$ est commut. et $\Gamma$ opère sur $M$.
On s'intéresse à la catégorie des revêtements
$\mathcal{X}'$ de $(\mathcal{X}, \Gamma)$, satisfaisant la condition

\newpage

%% ===== Page 131 =====

suivante : après l'oubli des opérations de $\Gamma$ on trouve un revêtement [unclear: connexe & abélien] galoisien de $\mathfrak{X}$, correspondant à un groupe quotient $M$ de $\pi = \pi_1(\mathfrak{X}, x_0)$ - de sorte que $\mathfrak{X}'$ devient de façon naturelle un torseur sous $M_{\mathfrak{X}}$.

Le foncteur
$$ \mathfrak{X}' \longmapsto \mathfrak{X}'(x_0) \quad (\text{fibre en } x_0) $$
de la catégorie de ces $\mathfrak{X}'$ dans la catégorie des ensembles, s'en va en fait vers les ensembles $E$ munis des structures suivantes

a) $E$ est un $M$-torseur, et en tant que tel donne naissance à une extension
$$ 1 \to M \to \text{Aff}(E) \to \operatorname{Aut}(M) \to 1 $$
d'où une sous-extension induite

(288)
$$ 1 \to M \to \text{Aff}_\Gamma(E) \to \tilde{\Gamma} \to 1 $$

où $\tilde{\Gamma} = \operatorname{Im}(\Gamma \to \operatorname{Aut}(M))$

b) On a un hom d'extensions

(289)
\begin{tikzcd}
1 \ar[r] & M \ar[r] \ar[d, "id"'] & \tilde{M} \ar[r] \ar[d, "\theta_{\tilde{M}}"'] & \Gamma \ar[r] \ar[d, "can"] & 1 \\
1 \ar[r] & M \ar[r] & \text{Aff}_\Gamma(E) \ar[r] & \tilde{\Gamma} \ar[r] & 1
\end{tikzcd}

La catégorie des $\mathfrak{X}'$ (à un iso...) équivaut à celle des couples $(E, \theta_{\tilde{M}})$, où $E$ est un $M$-torseur, et $\theta_{\tilde{M}}$ un hom. d'extensions (289).

\newpage

%% ===== Page 132 =====

Les autom. d'un objet $\overline{X}'$ correspondent
aux "translations" par des éléments
de $M^\Gamma$, ou (en termes de $E, \theta_{\tilde{M}}$)
aux translations de $E$ par de tels
éléments de $M^\Gamma$, vus comme des
éléments de $Aff_\Gamma(E)$ (ils sont centraux,
et en conjuguant par ces éléments on
trouve le $\underline{\underline{\text{même}}}$ hom. $\theta_{\tilde{M}}$ ... ). Les classes
d'isomorphie répondent
biunivoquement aux classes de $M$-conjugaison
des scindages de l'extension (287) de
$\Gamma$ par $\tilde{M}$. De façon précise, à un
$(E, \theta_{\tilde{M}})$ est associé canoniquement un
homomorphisme de $M$-ensembles
(290) $$E \to Scind(\tilde{M} \text{ ext. de } \Gamma \text{ par } M)$$
dont l'image est la classe de $M$-conjugaison
décrivant la classe d'isomorphie envisagée.
L'hom. (290) se factorise par $M^\Gamma$ en
un mono
(291) $$M^\Gamma \backslash E \hookrightarrow Scind(\tilde{M} \text{ ext. de } \Gamma \text{ par } M)$$
et si $M^\Gamma = 0$, on voit que $E$ se reconstitue
canoniquement à partir de la classe de

\newpage

%% ===== Page 133 =====

714

conjugaison envisagée. L'opération [unclear: (action)]
(définie par $\Theta_{\tilde{M}}$)
de $\tilde{M}$ sur $E$ n'est autre que l'opération
(induite par la)
de conjugaison sur les scindages
(qui bien entendu ne fait pas sortir
d'une classe de $M$-conjugaison, - on
conjugue par des éléments quelcon-
ques de $\tilde{M}$).

Corollaire. Supposons que l'on ait
(292) $$H^0(\tilde{\Gamma}, M) = H^1(\tilde{\Gamma}, M) = 0$$
($= M^\Gamma$)
i.e. que deux scindages d'une cat.
triviale de $\Gamma$ par $M$ soient conjugués
par un élément unique de $M$ (car $M^\Gamma = 0$). Ce
implique que la catégorie des $X'$, ou
encore celle des couples $(E, \Theta_{\tilde{M}})$, est
rigide, i.e. n'a d'autom. que l'id. (un objet est défini à un 
isomorphisme près). Ceci posé, les classes
d'isomorphie en question (i.e. les classes
de conjugaison de scindages de
l'ext. $\tilde{M}$ de $\Gamma$ par $M$) correspondent
biunivoquement aux p.-torseurs quotients de $\tilde{M}$.

\newpage

%% ===== Page 134 =====

de $M$ satisfaisant aux conditions suivantes

a) $M \to \tilde{M}$ ([unclear: façon] $M \simeq \bar{M} \to \tilde{M}$ ) est injectif,
d'où extension
(293) $$1 \to M \to \tilde{M} \to \tilde{\Gamma} \to 1$$
[MARGIN: $\simeq \tilde{M}/M$ au dessus de $\tilde{\Gamma}$]
b) $\tilde{\Gamma} \simeq \tilde{M}/M$ opère fidèlement sur $M$,
i.e. $\tilde{\Gamma} \simeq \bar{\Gamma} = \operatorname{Im}(\tilde{\Gamma} \to \operatorname{Aut}(M))$.

c) L'extension $\tilde{M}$ de $\tilde{\Gamma}$ par $M$ est
triviale (i.e. sa classe dans $H^2(\tilde{\Gamma}, M)$ est
nulle).

Pour un tel [unclear: quotient] $\tilde{M}$ de $\tilde{\tilde{M}}$ on récupère
$E$ comme l'ensemble de tous les
scindages de l'extension $\tilde{M}$ de $\tilde{\Gamma}$ par $M$,
sur laquelle on fait opérer $\tilde{M}$
via $\tilde{M}$ par conjugaison - la restriction
de cette opération : $M \subset \tilde{M}$ fait de $E$
un $M$-torseur (par l'hyp. (292)).

Revenons à nouveau au cas particulier
qui nous intéresse. Au n° précédent (120)
nous avons vérifié les relations (292),
en [unclear: constatant] en même temps que cette
fameuse extension $G_{\mathbb{Q}}$ de $\hat{\Gamma} \simeq \hat{\mathbb{Z}}^* \times \mathfrak{S}_3$
par $M^{\wedge}$ (noté $M$ dans les considérations
générales)

\newpage

%% ===== Page 135 =====

715

était triviale, et que deux scindages
étaient conjugués par un élément de
$M$ et un seul (c'est cette deuxième
moitié d'existence et unicité de conju-
gaison qui est le contenu géom. de (292))
De plus, on a donc construit justement,
en [unclear: - II,], un hom surjectif
$$ M = M_{0,3} \longrightarrow \widehat{Cg}_{,\sigma} $$
faisant de ce $\widehat{Cg}_{,\sigma}$ un quotient de
$M$ et de $\tilde{M}$, qui a les propriétés
a) b) c) des corollaires précédents. Cela
définit donc, à isomorphisme unique
près, un objet de la catégorie des $\vec{X}'$,
i.e. un prorevêtement de
ou ce qui revient au même, un rev.
proétale de $(M_{0,3 \mathbb{Q}}, \mathfrak{S}_3) \simeq M_{1,1 \mathbb{Q}}$,
que, après oubli des op. de $\mathfrak{S}_3$ et
ext de la base de $\mathbb{Q}$ à $\overline{\mathbb{Q}}$, devienne
le rev. universel abélien de $M_{0,3 \overline{\mathbb{Q}}}$.

Mais d'autre part la tour des
revêtements de Fermat nous fournit
très exactement un tel objet - et la

\newpage

%% ===== Page 136 =====

je n'avais pas eu à ma disposition
cette tour de Fermat, je n'aurais
certes pas suspecté que [crossed out: l'extension] $\widetilde{C}^{+\wedge} = \widetilde{C}^{+\wedge}_{H}$ de
$H^+ \simeq \widehat{\mathfrak{S}}_3$ par $M^\wedge$ soit scindée !

La conjecture qui se présente tout de
suite, c'est que ces deux revêtements
pro-étales sont isomorphes - l'isomor-
phisme entre les deux est alors
unique.

Pour l'autre voie de détermination
des classes d'isomorphie [unclear: (c.a.d. $\tilde{M}$ sur $M_{0,3} / [\Gamma_{\mathbb{Q}}, \Gamma_{\mathbb{Q}}]$)]
de scindages de $\Gamma = \Gamma_{\mathbb{Q}} \times \mathfrak{S}_3$ par
$M^\wedge$, on a par la suite exacte associée
à la suite spectrale de Hochschild

$$ (294) \quad 0 \to H^1(\Gamma_{\mathbb{Q}}, H^0(\widehat{\mathfrak{S}}_3, \widehat{M})) \to H^1(\Gamma, \widehat{M}) \to H^0(\Gamma_{\mathbb{Q}}, H^1(\widehat{\mathfrak{S}}_3, \widehat{M})) $$

Or on a vu que
$$ (295) \quad H^0(\widehat{\mathfrak{S}}_3, \widehat{M}) = 0 \quad (\text{[unclear: car] } H^0(L_f, \widehat{M}) = 0) $$

[crossed out: Et] d'autre part
$$ H^1(\widehat{\mathfrak{S}}_3, \widehat{M}) = H^1(\mathfrak{S}_3, M) \otimes_{\mathbb{Z}} \hat{\mathbb{Z}} \quad , $$

[MARGIN: donc $H^1(\Gamma, \widehat{M}) \simeq H^0(\Gamma_{\mathbb{Q}}, H^1(\widehat{\mathfrak{S}}_3, \widehat{M}))$]

\newpage

%% ===== Page 137 =====

715

ce qui nous ramène à calculer
$$H^1(\hat{\mathfrak{S}}_3, M) \quad ?$$
Mais l'étude faite de l'extension $\hat{\mathfrak{S}}^+ / \mu_2$ de
$\hat{\mathfrak{S}}_3$ par $M$ - qui nous a [unclear: montré] que cette
extension était scindée, et que deux
scindages étaient conjugués par un
unique élément de $M$, nous donne
en fait
$$(296) \quad \begin{cases} H^2(\hat{\mathfrak{S}}_3, M) = H^1(\hat{\mathfrak{S}}_3, M) = 0 \\ H^0(\hat{\mathfrak{S}}_3, \hat{M}) = H^1(\hat{\mathfrak{S}}_3, \hat{M}) = 0 \end{cases}$$
d'où finalement, grâce à (295), la
relation triomphale
$$(297) \quad | \quad H^1(\Gamma_{\mathbb{Q}} \times \hat{\mathfrak{S}}_3, \hat{M}) = 0 \quad | \quad !$$

On trouve donc - à nouveau
contre toute attente, je l'avoue ! -
le
Théorème Il existe un prorevêtement [unclear: $= M_{0,3 \mathbb{Q}}$]
proétale [unclear] $U_\infty$ de $(U_{0,3 \mathbb{Q}}, \hat{\mathfrak{S}}_3)$ , ayant
la propriété suivante : L'image inverse
de $U_\infty$ sur $U_{0,3 \bar{\mathbb{Q}}}$ (oubliant l'op. de $\hat{\mathfrak{S}}_3$)
est "le" revêtement abélien universel
de $U_{0,3 \bar{\mathbb{Q}}}$. Pour deux tels revête-

136

\newpage

%% ===== Page 138 =====

ments proétales de $(U_{0,3 \mathbb{Q}}, \mathfrak{S}_3) = M'_{1,\mathbb{Q}} ,$
il existe un unique isomorphisme
de l'un sur l'autre.

Nous avons obtenu l'existence de
deux façons : par la tour de Fermat
et par la théorie générale du § 49 III,
appliquée au cas présent - ce qui
nous donne un $U_\infty$ et un $U'_\infty$ ,
chacun décrit à isomorphisme [unclear: unique] près !

Le th. d'unicité résulte du
petit calcul précédent - j'avais
peur des "gros"
$$ H^1(\Gamma_{\mathbb{Q}}, \hat{M}) \simeq M_\infty \otimes_{\hat{\mathbb{Z}}} \underset{\simeq \widehat{\mathbb{Q}^*}}{H^1(\Gamma_{\mathbb{Q}}, T_{\hat{\mathbb{Z}}}(\mathbb{G}_m))} , $$
mais finalement c'est la nullité de
$H^i(\mathfrak{S}_3, M)$ pour $i=0,1$ qui implique
(sans rien connaître sur $\Gamma_{\mathbb{Q}}$ !)
celle de $H^1(\Gamma_{\mathbb{Q}} \times \mathfrak{S}_3, \hat{M})$, qui assure
l'unicité qu'il nous fallait.

\newpage

%% ===== Page 139 =====

717

[unclear: Tout] Le théorème précédent élucide
complètement, il me semble, la
signification authentico-géométrique
(en termes de la tour de Teich.) de
l'homomorphisme ((12.18) p. 695')
$$M = M_{1,1} \longrightarrow G_{\rho, \sigma} \simeq \hat{\mathbb{Z}}^* \cdot \mathfrak{S}_3$$
obtenu dans la section 11º) par voie purement
algébrique — ou du moins celui de l'
homomorphisme un peu plus faible
$$M \longrightarrow G_{\rho, \sigma} / \mu_{\hat{\mathbb{Z}}} \simeq \hat{\mathbb{Z}}^* \cdot \mathfrak{S}_3 / \mu_{\hat{\mathbb{Z}}}$$
qui se factorise par un hom.
$$\tilde{M} = M / (\mu_M = \mu_{\hat{\mathbb{Z}}}) \longrightarrow \tilde{G}_{\rho, \sigma}$$
[MARGIN: extension de $\tilde{\Gamma} = \hat{\mathbb{Z}}^* \times \mathfrak{S}_3$ par $M^\wedge$]

noté $M$ dans les pages précédentes.

Le passage de $G_{\rho, \sigma}$ à $\tilde{G}_{\rho, \sigma}$, dans [unclear: le stade]
actuel, avait comme rôle
essentiel (a posteriori!) de fournir
une extension par $M^\wedge$ ([unclear: à une composante neutre près]) d'un
groupe opérant fidèlement sur $M^\wedge$, ce
qui n'était pas le cas de l'ext. $G_{\rho, \sigma}$
de $\hat{\mathbb{Z}}^* \times H^+$ par $M^\wedge$, puisque le dit
centre $\mu_{\hat{\mathbb{Z}}}$ y opérait trivialement.

\newpage

%% ===== Page 140 =====

Ceci, joint aux relations
$$H^\circ(\hat{\mathbb{Z}}^* \times \mathfrak{S}_3, \hat{M}) = H^1(\hat{\mathbb{Z}}^* \times \mathfrak{S}_3, \hat{M}) = 0 \quad,$$
permettait d'interpréter le groupe
$G_{\rho, \tilde{\sigma}}$ comme un groupe de transformations
sur l'âme $E$ de ce [unclear: pro-jeu singulier],
ce qui permettait d'interpréter
l'hom $\mathcal{M} \to G_{\rho, \tilde{\sigma}}$ (ou $\mathcal{M}^\sim \to G_{\rho, \tilde{\sigma}}$)
comme une spécification de $\mathcal{M}$
(ou $\tilde{\mathcal{M}} = \Pi_1(U_{0,3 \mathbb{Q}}, \underline{\mathfrak{S}}_3) = \Pi_1(M_{1,1 \mathbb{Q}}, \dots)$
comme, [unclear], sur $E$ [unclear: correspondant]
à un revêtement de $(U_{0,3 \mathbb{Q}}, \underline{\mathfrak{S}}_3)$
$\simeq M_{1,1 \mathbb{Q}}$ - et ce revêtement
n'est autre que le prorevêtement
de Fermat.

Il resterait à interpréter de façon
analogue l'hom $\mathcal{M} = \mathcal{N}_{1,1} \to G_{\rho, \sigma}$,
s'insérant dans l'hom. de suites exactes

(298)
\begin{tikzcd}
1 \ar[r] & \mu_M \ar[r] \ar[dd] & \mathcal{M} \ar[r] \ar[dd] & \mathcal{M}^\sim \ar[r] \ar[dd] & 1 \\
& (\pm 1, \omega_f) & (\mathcal{N}_{1,1}) & (\tilde{\mathcal{N}}_{1,1} = \mathcal{M}_{0,3}) & \\
1 \ar[r] & \mu_G \ar[r] & G_{\rho, \sigma} \ar[r] & G_{\rho, \tilde{\sigma}} \ar[r] & 1 \\
& \pm 1, \omega_f & & &
\end{tikzcd}

Pour [unclear] d'interpréter d'ailleurs
$\mathcal{M}$, en tant qu'extension centrale de $\mathcal{M}^\sim$

\newpage

%% ===== Page 141 =====

718

par $\mu$, comme image inverse de l'extension
$G_{\rho,\sigma}$ de $G_{\rho,\tilde{\sigma}}$ par $\mu$, par l'homomorphisme
$\tilde{M} \to G_{\rho,\tilde{\sigma}}$ précédent. Il y a une
analogie frappante avec ce qui s'était
passé, d'ailleurs, quand on était
parti (au § 48, et dans certains développe-
ments qui l'ont précédé) du groupe
quotient ("de type Lie") $\operatorname{Sl}(2,\hat{\mathbb{Z}})'$ de
$\widehat{\Pi}_{0,3}^+$, pour construire un hom.
$\mathcal{M}_{0,3} \to \operatorname{Gl}(2,\hat{\mathbb{Z}})/\{\pm 1\}$
d'où il nous a été possible de reconstituer
l'ext. $M = \mathcal{N}_{1,1}$ de $\mathcal{M}_{0,3}$ par $\mu$ juste-
ment comme l'image inverse de
l'extension $\operatorname{Gl}(2,\hat{\mathbb{Z}})$ de $\operatorname{Gl}(2,\hat{\mathbb{Z}})/\{\pm 1\}$
par $\mu = \{\pm 1\}$. Du pt de vue actuel
- où j'essaye d'expliciter autant que
possible des objets authentico-géomé-
triques en termes d'automorphismes
de groupes profinis de lacets $\widehat{\Pi}_{0,3}$ - le fait
remarquable, c'est que les deux

\newpage

%% ===== Page 142 =====

extensions de $\tilde{M}$ par $\mu$, sortant des
calculs en apparence assez différents,
soient canoniquement isomorphes.

Je présume qu'une explication générale
de ce fait n'offre aucune difficulté,
dans le contexte actuel - du fait que
dans les deux cas, l'extension de
$\operatorname{Sl}(2, \hat{\mathbb{Z}})/\pm 1$ resp de $\tilde{C}_{g,\sigma}$ provenait
du fait que j'étais parti d'un
groupe quotient ($\operatorname{Sl}(2, \hat{\mathbb{Z}})$ resp.
$\hat{G}_{0,3}^{+\wedge}$) non pas de $\hat{G}_{0,3}^{+\wedge} \simeq \operatorname{Sl}(2, \hat{\mathbb{Z}})^{\wedge}/\pm 1$
mais de $\operatorname{Sl}(2, \hat{\mathbb{Z}})^{\wedge} \simeq N_{1,1}$ lui-même -

Si on cherche malgré tout une
interprétation (genre "tour de Fermat")
de l'hom. $M \to C_{g,\sigma}$, il faudrait
interpréter $C_{g,\sigma}$ (comme l'est $\tilde{C}_{g,\sigma}$)
comme un groupe
d'automorphismes d'un "fibré type"
"pas trop gros" - la difficulté est
d'inventer quelque chose sur quoi on
opère de façon non triviale.

\newpage

%% ===== Page 143 =====

719

Interprétation $M_{0,3} = N_{1,1}'$ comme
$\Pi_1(M_{1,1}')$, où $M_{1,1}'$ est le schéma modulaire / $\mathbb{Q}$
des courbes elliptiques à symétrie près, on
voit que la tour de Fermat revient à
ceci : sur tout schéma de base $S$ [MARGIN: (de car. 0)]
et tte quasi-courbe elliptique $\tilde{X}$ sur $S$,
considérons le revêtement étale de degré
trois $S'$ de $S$ correspondant [MARGIN: (des pts d'indice]
2 de $\tilde{X}$), posons par déf.

(299) $$M_0(\tilde{X}/S) = p_*(\mathbb{Z}_{S'})/\mathbb{Z}_S$$

le $\mathbb{Z}$-module [unclear: localement constant] libre
de rang deux correspondant, et
considérons

(300) $$M_1(\tilde{X}/S) = M_0(\tilde{X}/S) \otimes_{\mathbb{Z}} \hat{\mathbb{Z}}(1)$$
[MARGIN: $\| \text{déf.}$ $T_0(G_m)$]

qui est un pro-groupe [unclear: profini] sur $S$,
un syst de coefficients loc. const.
(isom. sur un [unclear] à $\hat{\mathbb{Z}}^2$. [MARGIN: Bien] Ceci posé,
Ainsi, $M_1(\tilde{X}/S)$ dépend fonctoriellement
de $\tilde{X}/S$ pour les iso et isog. inverses
(on a une [unclear: section] cartésienne d'une
cat. fibrée de base ...). Ceci posé,

\newpage

%% ===== Page 144 =====

la donnée de la "tour de Fermat"
équivaut à la donnée, pour tt $\tilde{X}/S$
comme dessus, d'un torseur

[MARGIN: [unclear: poser une] [unclear: condition] [unclear: cohomologique] cf (310) p. 724]
(301) $F_{\infty}(\tilde{X}/S)$ torseur sur $S$ sous $M_1(\tilde{X}/S)$ ,

compatible aux iso. et chgts de base.
Le théorème précédent dit essentiel-
ment qu'un tel $F_{\infty}(\tilde{X}/S)$ existe,
et est unique à iso. unique près,
en exigeant de plus que l'hom.
naturel
$$\hat{\Pi}_{0\hat{3}} \longrightarrow M^{\wedge}$$
associé à une telle donnée $F_{\infty}$ soit
l'hom canonique...

Soit d'autre part $X$ une courbe
elliptique relative au dessus d'une
base $S$ de car. $0$, et soit $S' \subset X$ le
sous-schéma fermé, étale de degré $3$ sur $S$,
correspondant aux pts d'ordre $2$. Considérons
la jacobienne généralisée $J$ de $X$
déduite de cette multisection, soit
$J^0$ sa partie de degré $0$, on a alors

\newpage

%% ===== Page 145 =====

720

une suite exacte

(302) $$1 \to J_0 \to J^0 \to X \to 1$$
$$\simeq$$
$$\mathbb{G}_m \otimes_{\mathbb{Z}} M_\circ$$

où $M_\circ = M_\circ(X/S)$ est défini plus haut
(299). On en déduit une suite
exacte

(303) $$1 \to M_1 \to T_\infty(J^0) \to T_\infty(X) \to 1$$
$$=$$
$$M_\circ \otimes_{\mathbb{Z}} T_\infty(\mathbb{G}_m)$$

(304) $$T_\infty(G) = \lim_{\leftarrow n} {}_n G$$

est le système local des $\pi_1$ des schémas
[unclear: à fibres connexes] en groupes lisses comm. sur $S$. Donc
la donnée d'un torseur sous $M_1$ équi-
vaut à la donnée d'un torseur sous
$T_\infty(J^0)$, muni d'une section du
torseur associé de groupe $T_\infty(X)$.
Y a-t-il une description remarquable du
torseur de Fermat $F_\infty(X/S)$, en termes
de la géométrie de la courbe elliptique
relative $X/S$ ?

\newpage

%% ===== Page 146 =====

14°) Relations avec jacobiennes généralisées.

Considérons la "jacobienne généralisée" de
$U_{0,3 \mathbb{Z}} = \mathbb{P}^1_{\mathbb{Z}} - \{0,1,\infty\}_{\mathbb{Z}}$ , c'est donc un
schéma en groupes $J^*_{U_{0,3}} = J^*$, [unclear: augmenté]
vers $\mathbb{Z}$, s'insérant dans la suite
exacte

[MARGIN: NB $S_0 = Spec \mathbb{Z}$]

(305) $$1 \longrightarrow J^\circ \longrightarrow J^* \xrightarrow{deg} \mathbb{Z}_{S_0} \longrightarrow 1$$
(avec $J^\circ \simeq \mathbb{G}_{mS_0} \otimes_{\mathbb{Z}} M$)

et muni d'un homomorphisme
canonique

(306) $$\varphi : U_{0,3 \mathbb{Z}} \longrightarrow J^1$$
(où $J^1$ est le sous-schéma de $J^*$ image inverse de la section 1 de $\mathbb{Z}_{S_0}$)

Les opérations de $\mathfrak{S}_3$ sur $U_{0,3}$ définis-
sent par fonctorialité des opérations de
$\mathfrak{S}_3$ sur $J^*$, compatibles avec la structure
d'extension (305), et telles que
(306) commute aux opérations de $\mathfrak{S}_3$.
Notons que $J^1$ est un $J^\circ$-torseur, et

\newpage

%% ===== Page 147 =====

721

même un $(\widehat{\mathfrak{S}}_3, J^0)$-torseur (un $J^0$-torseur sur lequel $\widehat{\mathfrak{S}}_3$ opère de façon compatible avec son op. sur $J^0$), donc classifié par un élément de $H^1(S_0, \widehat{\mathfrak{S}}_3; J^0)$. On présume qu'on a

(307) $$H^0(S_0, \widehat{\mathfrak{S}}_3, J^0) = H^1(S_0, \widehat{\mathfrak{S}}_3, J^0) = 0$$

(où
(308) $$S_0 = \operatorname{Spec}\ \mathbb{Z}, \quad J^0 = \mathbb{G}_m \otimes_{\mathbb{Z}} M \simeq \mathbb{G}_m^{\{0,1,\infty\}} / diag$$
[unclear: attention on a isomorph. obtenu via $(x_0,x_1,x_\infty) \mapsto (x_1 x_\infty^{-1}, x_\infty x_0^{-1}, x_0 x_1^{-1})$ qui ...]
$$\simeq \operatorname{Ker}(\mathbb{G}_m^{\{0,1,\infty\}} \xrightarrow{(x_0, x_1, x_\infty) \mapsto x_0 x_1 x_\infty} \mathbb{G}_m)$$

On a en effet
[MARGIN: Calcul correct, car il n'y a en vue [unclear] de $\widehat{\mathfrak{S}}_3$-torseur si ce n'est qu'il suppose connu $H^1(S, \widehat{\mathfrak{S}}_3) = \{1\}$ (d'où $\widehat{\mathfrak{S}}_3 \simeq \mathfrak{S}_3$ (cf. 122)), sans calcul général d'ailleurs. Dans cette suite spec. $H^p(\widehat{\mathfrak{S}}_3, H^q(S, J^0)) \Rightarrow \dots$]
$$H^0(S, \widehat{\mathfrak{S}}_3; J^0) = \left\{ (x_0, x_1, x_\infty) \in \mathbb{Z}^{* \{0,1,\infty\}} \left| \begin{array}{l} x_0 x_1 x_\infty = 1 \\ (x_0, x_1, x_\infty) \text{ inv. par } \widehat{\mathfrak{S}}_3 \end{array} \right. \right\}$$
$$= \{ (x, x, x) \mid x \in \mathbb{Z}^{* \{0,1,\infty\}}, \quad x^3 = 1 \}$$
$$= 1$$
seule racine cubique de $1$ dans $\mathbb{Z}$ (= dans $\mathbb{Q}$) est $1$. Ce qui établit la première relation 307). Pour la seconde, on utilise la suite spectrale

146

\newpage

%% ===== Page 148 =====

$$H^n(S_0, \mathfrak{S}_3 ; J^0) \stackrel{E^{p,q}_2=}{\Longleftarrow} H^p(\mathfrak{S}_3, H^q(S_0, J^0))$$

qui donne la suite exacte

$$0 \to H^1(\mathfrak{S}_3, H^0(S_0, J^0)) \to H^1(S_0, \mathfrak{S}_3 ; J^0) \to H^0(\mathfrak{S}_3, H^1(S_0, J^0)) \to H^2(\mathfrak{S}_3, H^0(S_0, J^0))$$
$$ \parallel $$
$$ (H^0(\mathfrak{S}_3, M_{\mathbb{Z}} \otimes H^1(S_0, \mathbb{G}_m))) $$

Or on a

[MARGIN: NB on n'a pas besoin de la relation $Pic(S)=0$]

$$H^1(S_0, J^0) \simeq H^1(S_0, \mathbb{G}_m^2) \simeq Pic(S_0)^2 = 0$$

donc il reste

$$H^1(S_0, \mathfrak{S}_3 ; J^0) \simeq H^1(\mathfrak{S}_3, H^0(S_0, J^0))$$
$$ \simeq (\mathbb{Z}^*)^{\{0,1,\infty\}} / \mathbb{Z}^* $$
$$ \simeq \mathbb{Z}^* \otimes M_{\mathbb{Z}} $$

Mais en fait, les calculs faits
par ailleurs montrent qu'on a

(309) $$H^q(\mathfrak{S}_3, M_{\mathbb{Z}} \otimes_{\mathbb{Z}} I) = H^1(\mathfrak{S}_3, M_A \otimes_{\mathbb{Z}} I) = 0$$

quel que soit le $\mathbb{Z}$-module $I$ (p. ex.
on a $I = H^0(S, \mathbb{G}_m) = Pic(S)$ ;
$I = J^0(S)$, $S$ un schéma quelconque), d'où

finalement,

(310) $$H^1(\mathfrak{S}_3 ; J^0) \simeq H^1(S, \mathfrak{S}_3 ; J^0) = 0$$

pour tout schéma $S$. On trouve ainsi :

\newpage

%% ===== Page 149 =====

722

Proposition Soit $J^0 = \mathbb{G}_m^{\{0,1,\infty\}} / \underset{\text{diag.}}{\mathbb{G}_m} \simeq \mathbb{G}_m \otimes_\mathbb{Z} M$
Alors pour tout $(\mathfrak{S}_3, J^0)$-torseur sur
une base $S$ quelconque, il existe une
unique trivialisation de celui-ci, i.e.
une unique section, invariante pour les
opérations de $\mathfrak{S}_3$.

Corollaire 1 Sur toute base $S$, il existe
un unique morphisme

(311) $$ \varphi: U_{0,3 S} \longrightarrow J^0_S $$

qui soit compatible avec les opérations
de $\mathfrak{S}_3$, et qui (via l'hom.
$\hat{J}_S \to J^0_S$ que la factorisation, par la
propriété universelle des fo jacobiennes
généralisée) induise l'hom.
$\hat{J}_S \to \hat{J}_S$ identique.

[MARGIN: NB Avec les notations de 695, $U_{0,3S}$ (où $U_{0,3}$ pour un $S$ est défini par : $U \subset \mathbb{G}_m^3 / \text{diag.} \mathbb{G}_m$, (312) $\mathcal{U}_1 = \{(x_0,x_1,x_\infty) \in \mathbb{G}_m^3 | x_0x_1x_\infty=1\} / \text{mod diag } \mathbb{G}_m$]

Corollaire 3 Soit $X$ un schéma de

\newpage

%% ===== Page 150 =====

Brauer-Sévéri sur $S$, de dim relative
$1$ (i.e. une forme de $\mathbb{P}^1_S$), muni
d'un sous-schéma $S' \subset Y$, étale de
degré $3$ sur $S$, et soit $U = Y \setminus S'$,
$J^*_{U/S}$ la jacobienne généralisée, de
sorte qu'on a

(313) $$1 \longrightarrow J^0_{U/S} \longrightarrow J^*_{U/S} \longrightarrow \mathbb{Z}_S \longrightarrow 1$$
[underneath $J^0_{U/S}$ is written: $\simeq \mathbb{G}_m \otimes_{\mathbb{Z}} M_0(S'/S)$]

où $M_0(S'/S) \simeq p_*(\mathbb{Z}_{S'})/\mathbb{Z}_S$ [underneath: en $\mathbb{Z}$-modules, diag.] est le $S$-
schéma (loc. isom. à $\mathbb{Z}^2$, déduit de
$M_0$ en le tordant (via action de $\mathfrak{S}_3$)
par le $\mathfrak{S}_3$-torseur sur $S$
[unclear: fibré principal] associé : $S'$.
Il existe alors un unique morphis-
me

(314) $$\psi: U \longrightarrow J^*_{U/S}$$

commutant à l'action de $\operatorname{Aut}(S'/S)$
$\simeq \operatorname{Aut}(Y, S')$, et induisant via la
propriété universelle l'hom. [unclear: canonique]
$J^*_{U/S} \to J^*_{U/S}$ qui soit l'identité.

\newpage

%% ===== Page 151 =====

723

Explicitons l'isomorphisme

(315) $$U_{0,3} = \mathbb{P}^1 \setminus \{0,1,\infty\}_{S_0} \stackrel{\psi}{\longrightarrow} \mathbb{G}_m^{\{0,1,\infty\}} / \mathbb{G}_m$$
$$ \simeq \mathbb{G}_m \setminus \{1\}_{S_0} \simeq \{z \in \mathbb{G}_m \mid z-1 \text{ inv}\} \qquad \hookleftarrow U_1 = \{ (x_0, x_1, x_\infty) \in \mathbb{G}_m^3 / \mathbb{G}_m \mid x_0 + x_1 + x_\infty = 0 \}$$

[crossed out: illegible] isomorphisme de $U_{0,3}$ [crossed out: avec] sur $U_1$ qui
[crossed out: se prolonge en] :

(316) 
$$ \bar{\psi} : \mathbb{P}^1_{S_0} \longrightarrow \bar{U}_1 \subset \mathbb{P}^2_{S_0}$$
$$ \qquad\qquad\qquad\qquad\qquad \| $$
$$ \qquad\qquad\qquad\qquad\qquad \{ (x_0, x_1, x_\infty) \in \mathbb{P}^2 \mid x_0 + x_1 + x_\infty = 0 \} $$
$$ \begin{cases} \bar{\psi}(0) = (0, 1, -1) \\ \bar{\psi}(1) = (-1, 0, 1) \\ \bar{\psi}(\infty) = (1, -1, 0) \end{cases} $$

il suffit de prendre $\qquad\qquad\qquad\qquad\qquad\qquad = (\frac{-z}{z-1}, 1, \frac{1}{z-1})$
(317) $$ \psi(z) = (-z, z-1, 1) = (1, 1/z - 1, -1/z) $$

Revenons maintenant à la donnée d'une courbe elliptique relative $X$ sur $S$, modèle symétrique (correspondant à un morphisme $S \to \text{[unclear: Espace]} \overline{\mathcal{M}}_{1,1}$).

\newpage

%% ===== Page 152 =====

On suppose $2$ inversible sur $S$, de sorte que
la donnée de $X$ (correspondant, par le
type de Legendre $X \longmapsto (X/\pm, \dots)$)
: la donnée d'un $(Y,S')$ comme dessus,
avec en plus une section de $U=Y \setminus S'$.
En fait, la donnée de $(Y,S')$ revient
à la donnée d' un $\mathcal{S}_3$-torseur
$T$ sur $S$ - et $Y$ s'obtient à
partir de l' [unclear] comme le schéma
torseur $T \wedge_{\mathcal{S}_3} \mathbb{P}^1_S$
avec $U = Y \setminus S' \simeq (T \wedge_{\mathcal{S}_3} \mathbb{P}^1_S) \setminus \underbrace{T \wedge_{\mathcal{S}_3} \{0,1,\infty\}_S}_{S'}$
plonge canoniquement dans
$T \wedge_{\mathcal{S}_3} \mathcal{U}_{0,3|S} \simeq \mathcal{Y}_{1|S}$. Le plongement
$U \subset \mathcal{Y}_{1|S}$ peut s'interpréter 
comme un morphisme canonique

$$ (318) \quad X \longrightarrow \mathcal{Y}_{1|S} = \mathcal{Y}^0_{\times|S} $$

[MARGIN: notation de p. 720]
qui doit être ainsi : [unclear]

\newpage

%% ===== Page 153 =====

[MARGIN: 724]

ment, [unclear: barré] à expliciter en termes
de "fonctions elliptiques" standard.
Sans m'attarder sur une telle
explicitation, et revenant plutôt au plonge-
ment correspondant

(319) $$U \hookrightarrow \mathcal{J}_{U/S}$$

et à la section marquée
(320) $$e \in U(S) \subset H^0(S, \mathcal{J}_{U/S})$$

[unclear: texte barré]

(321) $$1 \to {}_n\mathcal{J}_{U/S} \longrightarrow \mathcal{J}_{U/S} \stackrel{n \cdot id_{\mathcal{J}_{U/S}}}{\longrightarrow} \mathcal{J}_{U/S} \to 1$$

(pour $n$ variable), i.e. la tour d'isogénies
universelle déduite par passage à $\lim_{\leftarrow}$ :

(322) $$1 \to T_\infty(\mathcal{J}_{U/S}) \longrightarrow \tilde{\mathcal{J}}_{U/S} \stackrel{p}{\longrightarrow} \mathcal{J}_{U/S} \to 1$$
$$ \simeq T_\infty(\mathbb{G}_{m,S}) \otimes_{\hat{\mathbb{Z}}} \mathcal{M}_0(S/S) \quad =^{def} \lim_{\leftarrow} (\mathcal{J}_{U/S}) $$

pour déduire, par image inverse de $e$,
un torseur sous $T_\infty(\mathcal{J}_{U/S}) \simeq T_\infty(\mathbb{G}_{m,S}) \otimes_{\hat{\mathbb{Z}}} \mathcal{M}_0(S/S)$.
Ce posé, on aura cf (301) p. 719!)

(323) $$F_\infty(\tilde{X}/S) \simeq p^{-1}(e) \text{, le } T_\infty(\mathcal{J}_{U/S})\text{-torseur}$$
déduit de la section $e$
de $\mathcal{J}_{U/S}$ sur $S$

\newpage

%% ===== Page 154 =====

<u>Remarque</u> Le choix d'isomorphisme du
$T_\infty(\mathcal{J}_{U/S})$ - torseur obtenu est un
élément de $H^1(S, T_\infty(\mathcal{J}_{U/S}))$, image
de l'élément $e$ de $H^0(S, \mathcal{J}_{U/S})$, par
le morphisme canonique

(324) $$H^0(S, \mathcal{J}^0) \longrightarrow H^1(S, T_\infty(\mathcal{J}^0))$$ ,

dont le noyau est formé des sous-
groupe de $H^0(S, \mathcal{J}^0)$ des éléments
"infiniment divisibles". Dans le
cas où $S$ [unclear] de type fini sur $\mathbb{Z}$
- p. ex. [unclear] de type fini sur $\mathbb{Q}$ -
on sait (et c'est [unclear] "élémen-
taire") que ce noyau se réduit
à zéro. Ce qui signifie donc que
la section $e$ (i.e. la "fonction
elliptique" universelle) est alors entière-
ment déterminée, quand on connaît

\newpage

%% ===== Page 155 =====

725

1°) la trivialisation étale $S'$ de $S$ associé,
correspondant aux pts d'ordre $2$ de
$\tilde{X}/S$ - d'où $T^0 \pi_{1 U/S} \simeq T^0(\tilde{X}, S')$ - et

2°) Le $T_\infty(\pi_{1 U/S})$ - torseur $F_\infty(\tilde{X}/S)$ ,

\newpage

%% ===== Page 156 =====

15°) Interprétation de $\Gamma_{\mathbb{Q}} \to \mathcal{N}_\sigma^\sim$.

Au n° 11° nous avons introduit l'hom
canonique
$$ \begin{matrix} \mathcal{M} \overset{\text{déf}}{=} \mathcal{N}_{1,1}^{\text{ét}} & \longrightarrow & \mathcal{M}_{p,\sigma}^\sim \\ \cup & & \cup \\ \mathcal{M}^\cdot & \longrightarrow & \mathcal{M}_{p,\sigma}^{\sim \, \natural} \end{matrix} \quad cf \text{ (217) p.693)} $$
que nous avons "décomposé" bien plus tard
en deux homomorphismes (en plus d'un
hom. trivial $\mathcal{M} \to \Gamma_{\mathbb{Q}}^{\text{ét}} \simeq \hat{\mathbb{Z}}^\times$ ) en deux
hom.
$$ (325) \quad \begin{matrix} \mathcal{M} & \longrightarrow & G_{p,\sigma} \simeq \hat{\mathbb{Z}}^\times \cdot \mathfrak{S}^\wedge \\ \cup & & \cup \\ \mathcal{M}^\cdot & \longrightarrow & G_{p,\sigma}^\natural \simeq \hat{\mathbb{Z}}^\times \cdot M^\wedge \end{matrix} \quad (cf \text{ (218)}) $$
et
$$ (326) \quad \Gamma_{\mathbb{Q}}^{\text{ét}} \longrightarrow \mathcal{N}_\sigma^\sim \simeq \hat{\mathbb{Z}}^\times \cdot \hat{\mathbb{Z}} \quad (cf \text{ (219)}). $$

Le premier hom. (325) a été caracté-
risé en termes [crossed out: algébrico-] arithmético-géométriques
(pour ce qui est de sa restriction à
$\mathcal{N}_{1,1} \subset \mathcal{N}_{1,1}^{\text{ét}}$ lui-même, et on déduisait
mod $\mu_a$ pour former un hom
$$ [crossed out: (327)] \quad \mathcal{M}_{0,3} \longrightarrow G_{p,\sigma}^\sim \overset{\text{def}}{=} G_{p,\sigma} / \mu_a $$
dans la section 13°), via la tour des

\newpage

%% ===== Page 157 =====

725

courbes de Fermat. Du point de vue actuel,
il y a moyen d'expliciter la géométrie
algébrique et top. de type fini sur $\mathbb{Q}$ en
termes d'actions de gr. profinis,
je pense au moins; inversement, que
nous sommes arrivés à "décrire" la
tour des courbes de Fermat, via
l'hom. (327) construit ici "à la
force du poignet".

Il resterait maintenant à
trouver une explication arithmético-
géométrique de l'hom (326) de
$\Gamma_{\mathbb{Q}}$ dans $\hat{\mathbb{Z}}^\times \cdot \hat{\mathbb{Z}}$, induit par (325).
Mais [unclear: en] utilisant le plonge[ment]

( $$ N_\sigma^{\dots} \hookrightarrow M_{p,\sigma}^{\dots} \simeq \hat{\mathbb{Z}}^\times \cdot M $$ )

l'hom (323) apparait comme la comp.
$$ \Gamma_{\mathbb{Q}}^{ét} \hookrightarrow_{K_{1/2}} M_{1}^{!} \xrightarrow[\text{i.e. "Fermat"}]{(327)} G_{p,\sigma}^{\dots} $$ )
de façon plus précise -- [unclear: via] un diagr.
[unclear: commutatif :]

\newpage

%% ===== Page 158 =====

(328)
\begin{tikzcd}
{[\Gamma_\mathbb{Q} \hookrightarrow]} \Gamma_\mathbb{Q}^{et} \ar[r] & \mathcal{N}_\sigma^\S \simeq \hat{\mathbb{Z}}^* \cdot \hat{\mathbb{Z}} \\
{[\underset{\Pi_1(U_{0,3\mathbb{Q}})}{M_{0,3}^!} \hookrightarrow] M_{0,3}^{! et}} \ar[u, hook, "K_{1/2}"] \ar[r, "\text{''Fermat''}"', "(32\bar{5})"] & G_{p,\sigma}^\S \simeq \hat{\mathbb{Z}}^* \cdot \hat{M} \ar[u, hook]
\end{tikzcd}

Cela signifie donc que l'on trouve
l'hom. induit par (326)

(329) $$\Gamma_{\mathbb{Q}} \longrightarrow \mathcal{N}_\sigma^\S \simeq \hat{\mathbb{Z}}^* \cdot \hat{\mathbb{Z}}$$

en regardant l'image de l'hom.

(330) $$\Gamma_{\mathbb{Q}} \longrightarrow \hat{\mathbb{Z}}^* \cdot \hat{M} = Aff(\hat{M}) \ ,$$

dans $\mathcal{N}_\sigma^\S \subset \hat{\mathbb{Z}}^* \cdot \hat{M}$ , déduit de la
tour des courbes de Fermat en
regardant la fibre de la dite tour
au point $1/2$ , ce qui donne un
torseur sous $T_a(\hat{\mathbb{Z}}_{(m)}) \otimes_{\hat{\mathbb{Z}}} \hat{M}$ , s'explici-
tant galoisiennement justement par
l'hom. (326). Notons que le $K_{1/2}$
de $U_{0,3\mathbb{Q}}$ corresp. par l'iso. $\varphi$ (317)
de $U_{0,3\mathbb{Q}}$ avec $U_{1,\mathbb{Q}}$, au point [unclear: $(\xi_{1/2}, 1)$]

\newpage

%% ===== Page 159 =====

727

[crossed out: de]

(331) 
$$ s_{1/2} = (\frac{1}{2}, \frac{1}{2}) \in U_1 \subset {\mathbb{G}_m}_{\mathbb{Q}}^3 / {\mathbb{G}_m}_{\mathbb{Q}} \subset \mathbb{P}_{\mathbb{Q}}^2 $$
$$ \simeq \{ (x_0, x_1, x_\infty) \in \Gamma( {\mathbb{G}_m}_{\mathbb{Q}}^3 / {\mathbb{G}_m}_{\mathbb{Q}} ) \mid x_0 + x_1 + x_\infty = 0 \} $$

La fibre géom. de $U_n$ en ledit point s'identifie
ainsi aux systèmes de solutions des
équations ci-dessous :

(332) 
$$ U_n(s_{1/2}) \simeq \{ (x_0, x_1, x_\infty) \in \Gamma( {\mathbb{G}_m}_{\bar{\mathbb{Q}}}^3 / {\mathbb{G}_m}_{\bar{\mathbb{Q}}} ) \mid x_0^n = \frac{1}{2}, x_1^n = \frac{1}{2}, x_\infty^n = -1 \} $$
$$ \simeq \{ (x_0, x_1, x_\infty) \in \Gamma( {\mathbb{G}_m}_{\bar{\mathbb{Q}}}^3 ) \mid x_0^n = \frac{1}{2}, x_1^n = \frac{1}{2}, x_\infty^n = -1 \} $$
$$ \simeq \{ (x_0, x_\infty) \in \Gamma( {\mathbb{G}_m}_{\bar{\mathbb{Q}}}^2 ) \mid x_0^n = \frac{1}{2}, x_\infty^n = -\frac{1}{2} \} $$
$$ \simeq T_n(\frac{1}{2}) \times T_n(-\frac{1}{2}) $$

où l'on a posé, pour tt $\zeta \in \mathbb{G}_m(S)$ ($S$ schéma, par ex $S = \operatorname{Spec}(\bar{\mathbb{Q}})$)

(333) $$ T_n(\zeta) = \text{[unclear: la fibre en } \zeta \text{ de } x \mapsto x^n] $$

Passant à la limite on trouve
un iso.

(334) $$ U_\infty(s_{1/2}) \simeq T_\infty(\frac{1}{2}) \times T_\infty(-\frac{1}{2}) , $$

où l'on a posé, comme de juste

(335) $$ T_\infty(\zeta) = \lim_{\longleftarrow} T_n(\zeta) $$
$$ (\text{un } T_\infty(\mathbb{G}_m)\text{-torseur}) $$

\newpage

%% ===== Page 160 =====

Ainsi, l'hom (329) doit pouvoir s'interpréter en termes de l'hom galoisien

(335) $$ \Gamma_{\mathbb{Q}} \longrightarrow \operatorname{Gal}( \tilde{\mathbb{Q}} [\sqrt[\infty]{-\frac{1}{2}}] / \mathbb{Q} ) $$

où $\tilde{\mathbb{Q}}$ désigne l'extension cyclotomique universelle de $\mathbb{Q}$

(336) $$ \tilde{\mathbb{Q}} \simeq \mathbb{Q} [\sqrt[\infty]{1}] $$ .

Pour apprécier la "taille" de l'image de $\Gamma_{\mathbb{Q}}$ dans $\hat{\mathbb{Z}}^* \ltimes \hat{\mathbb{Z}} \simeq Aff_{\hat{\mathbb{Z}}}$, il y a lieu de faire une digression de nature générale, sur les groupes de Galois sur $\mathbb{Q}$ des extensions de la forme

$$ \tilde{\mathbb{Q}} [ \sqrt[\infty]{\gamma_1}, \dots, \sqrt[\infty]{\gamma_N} ] $$ .

Pour [unclear: asseoir] sur une base [unclear: solide] ce [unclear: qu'on a à dire] (on se place ici d'emblée dans le cas où $S = Spec \, \mathbb{Q}$) donnons nous un groupe abélien $A$ (le cas qui nous intéresse ici sera de prendre $A = \mathbb{Z}$), et

\newpage

%% ===== Page 161 =====

728

un hom.

(337) $$A \xrightarrow{\xi} \mathbb{G}_m(S) = H^0(S, \mathbb{G}_{mS})$$

(dans le cas présent, il est donné par $1 \mapsto -1/2$)

ou ce qui revient au même, un hom.

(338) $$A_S \xrightarrow{\xi} \mathbb{G}_{mS}$$ .

Considérons la tour des isogénies de Kummer

(338') $$1 \to \mu_n \to \mathbb{G}_m \xrightarrow{\lambda \mapsto \lambda^n} \mathbb{G}_m \to 1$$

d'où, par image inverse via (338), une tour d'extensions de schémas en groupes étales sur $S$

(339) $$1 \to \mu_n \to E_n(\xi) \to A_S \to 1$$

et par passage à la $\varprojlim$

(340) $$1 \to T_\infty(\mathbb{G}_{mS}) \to E_\infty(\xi) \to A_S \to 1 .$$

Supposons $S$ connexe, et prenons un pt géom. $\bar{s}$ sur $S$, on aura donc un hom. canonique

(341) $$\Pi_1(S, \bar{s}) \to \operatorname{Aut}_{\text{gr. pro. finis}}(E_\infty(\xi)_{\bar{s}})$$

\newpage

%% ===== Page 162 =====

où la fibre $E_{\infty}(\xi)_{\bar{s}}$ est elle-même
munie d'une structure d'extension
abélienne
(342) $$1 \to T_{\infty}(\bar{k}^*) \to E_{\infty}(\xi)_{\bar{s}} \to A \to 1 \ .$$

L'hom (341) est compatible avec cette
structure d'extension, quand on fait
opérer $\Pi_1(S, \bar{s})$ trivialement sur $A$,
et sur $T_{\infty}(\bar{k}^*)$ ($\hat{\mathbb{Z}}$-module libre de
rg 1) par multiplication par le
caractère cyclotomique
(343) $$\Pi_1(S, \bar{s}) \to \Pi_1(\mathbb{Q}, \bar{s}) \xrightarrow{\chi} \hat{\mathbb{Z}}^* \ .$$

Soit $G$ le groupe
(344) $$G = \operatorname{Aut}(E_{\infty}(\xi)_{\bar{s}})$$
$$= \left\{ u \in Aut \ \left| \ \begin{array}{l} u \text{ respecte la structure d'ext. en} \\ \text{induisant } id_A \text{ et ds } T_{\infty}(\bar{k}^*) \text{ une} \\ \text{homothétie par un scalaire} \\ \mu \in \hat{\mathbb{Z}}^* \text{ bien déterminé} \end{array} \right. \right\}$$

un hom canonique
qui ne dépend que de $A$
$G \to \hat{\mathbb{Z}}^*$, s'inscrivant ds un
diagramme
(345) $$1 \to \operatorname{Hom}_{\hat{\mathbb{Z}}}(A, T_{\infty}(\bar{k}^*)) \to G \to \hat{\mathbb{Z}}^*$$

\newpage

%% ===== Page 163 =====

74

où l'image de $G \to \hat{\mathbb{Z}}^*$ est formée des $\mu \in \hat{\mathbb{Z}}^*$ tels que l'homothétie $\mu_T$ ne change pas la classe de l'ext. (342) de $A$ par $T$. L'hom (341) peut lors s'interpréter comme s'insérant dans un diagramme

(345)
\begin{tikzcd}
1 \ar[r] & \operatorname{Hom}_\mathbb{Z}(A, T_\infty(\bar{k}^*)) \ar[r] & G \ar[r] & \hat{\mathbb{Z}}^* \\
& & \Pi_1 \ar[u] \ar[ur, "\text{"caractère cyclotomique"}"']
\end{tikzcd}

[MARGIN: NB.] Si $S$ est sch. de type fini sur $\mathbb{Q}$, alors l'image de $\Pi_1 \to \hat{\mathbb{Z}}^*$ est un sous-groupe ouvert de $\hat{\mathbb{Z}}^*$, donc il en est a fortiori ainsi de l'image de $G \to \hat{\mathbb{Z}}^*$ - qui sera surjectif si le car. cyclotomique sur $\Pi_1$ est surjectif, i.e. si les seules racines de 1 sur $S$ sont $\pm 1$ (p. ex. si $S = Spec \mathbb{Q} \dots$).

Le diagramme (345) a
une propriété de fonctorialité évidente
pour $A$ variable (disons), i.e. pour un

\newpage

%% ===== Page 164 =====

hom. de groupes abéliens
(347) $A' \xrightarrow{f} A$ ,

l'extension $E_{\infty}(A')_S$ s'identifie
à l'image inverse de l'ext. $E_{\infty}(A)_S$
par (347), d'où un hom.
$G_A \to G_{A'}$ (injectif si $f$ est surjectif)
s'insérant dans un diagramme
commutatif

(349)
\begin{tikzcd}
1 \ar[r] & \operatorname{Hom}_{\mathbb{Z}}(A', T) \ar[r] & G_{A'} \ar[r] & \hat{\mathbb{Z}}^* \\
1 \ar[r] & \operatorname{Hom}_{\mathbb{Z}}(A, T) \ar[r] \ar[u] & G_A \ar[r] \ar[u] & \hat{\mathbb{Z}}^* \ar[u, "1"'] \\
& & \Pi_1 \ar[u, "\gamma_A"'] \ar[uu, bend left=45, "\gamma_{A'}"] \ar[ur, "\text{car. cyclot.}"']
\end{tikzcd}

où la flèche sur les $\operatorname{Hom}_{\mathbb{Z}}$ est trans-
posée de (347). Cela montre que
pour étudier la structure de l'image de $\gamma_A$ (345),
il est commode, par fonctorialité, de se

\newpage

%% ===== Page 165 =====

730

étudier le cas où $j$ est injectif
(350) $$j : A \hookrightarrow \mathbb{G}_m(S) \ .$$

On peut d'ailleurs "diviser" encore un peu le diagramme (345), en introduisant

(351) $$\Pi_1^+ = \Pi_1^+(S, \bar{s}) \overset{df}{=} \operatorname{Ker}\Big(\Pi_1(S, \bar{s}) \overset{cyclot.}{\longrightarrow} \hat{\mathbb{Z}}^*\Big)$$

et l'hom. de suites exactes

(352)
\begin{tikzcd}
1 \ar[r] & \Pi_1^+ \ar[r] \ar[d, "\gamma_j^+"'] & \Pi_1 \ar[r, "\chi_{\Pi}"] \ar[d, "\gamma_j"'] & \hat{\mathbb{Z}}^* \ar[d] \\
1 \ar[r] & \operatorname{Hom}_{\hat{\mathbb{Z}}}(A, T) \ar[r] & G_j \ar[r] & \hat{\mathbb{Z}}^*
\end{tikzcd}
.

On peut considérer que l'étude de la "taille" de $\operatorname{Im}(\gamma_j : \Pi_1 \to G_j)$ est essentiellement ramenée à celle de la taille de $Im \chi_{\Pi}$ et de celle de $Im \gamma_j^+$; de façon précise, si on pose

(353) $$Im \gamma_j = \tilde{\Pi}_1 \ , \ Im \gamma_j^+ = \tilde{\Pi}_1^+ \ , \ Im \chi_{\Pi} = \Gamma$$

on a une suite exacte

(354) $$1 \to \tilde{\Pi}_1^+ \to \tilde{\Pi}_1 \to \Gamma \to 1 \ .$$

Je voudrais essayer de donner des

\newpage

%% ===== Page 166 =====

indications sur la taille de $\tilde{\Pi}_1^+ \subseteq \operatorname{Hom}(A, \hat{T})$
[MARGIN: $\simeq \operatorname{Hom}(A, \hat{\mathbb{Z}})$]
dans le cas où on suppose

(355)
1°) $A$ $\mathbb{Z}$-module de type fini
2°) $S$ est de type fini sur $\mathbb{Q}$ et normal
3°) $\jmath : A \longrightarrow G_m(S)$ est injectif

L'hyp. 2°) implique déjà — n'est-ce pas —
que $\Gamma \subseteq \hat{\mathbb{Z}}^*$ est un sous-groupe ouvert.

L'hyp. que $S$ est normal implique d'
ailleurs que les groupes cherchés
$\tilde{\Pi}_1^+$ et $\Gamma$ ne changent pas si on remplace
$S$ par son pt générique (on suppose
pour simplifier que $\bar{s}$ est au-dessus du
dit pt générique), donc ops que
$S = Spec \ K$, où $K$ est un corps de
type fini sur $\mathbb{Q}$. Et on peut p. ê.
supposer, si on y tient, que $S$
est affine, et lisse sur $\mathbb{Q}$...

Notons en outre, posant

(356) $\check{A} = \operatorname{Hom}_{\mathbb{Z}}(A, \mathbb{Z})$

\newpage

%% ===== Page 167 =====

731

que $\check{A}$ est un $\mathbb{Z}$-module libre de t.f.
(de rang $r$ / i.e. $r = rg_{\mathbb{Z}} A$), et qu'on a
un isom canonique
(357) $$ \operatorname{Hom}_{\mathbb{Z}}(A, T) \simeq \check{A} \otimes_{\mathbb{Z}} T \quad (\simeq \hat{\mathbb{Z}}^n \text{ non canon.}) $$

J'aurais envie de prouver que l'image
de $\tilde{\Pi}_1^+$ ou $\Pi_1$ dans $\operatorname{Hom}_{\mathbb{Z}}(A, T) \simeq \check{A} \otimes_{\mathbb{Z}} T$
est ouverte. Je me demande si,
dans le cas $S_1$ un schéma de type fini
dominant $\mathbb{Z}$,
et $A = H^0(S_1, \mathbb{G}_m)$, pour
$S = S_{1\mathbb{Q}}$, si l'hom
$$\Pi_1 = \Pi_1(S, \bar{s}) \to \operatorname{Hom}(A, T) \simeq \check{A}(1)$$
n'est pas nécessairement surjectif ??
(Plus généralement, il en est peut être
ainsi si $\mathbb{G}_m(S) / Im A$ est de
torsion, ce à quoi on doit pouvoir
toujours se réduire, dans l'hyp. (355), quitte
à agrandir $A$ en $A' \subset H^0(S, \mathbb{G}_m)$ avec $A'/A$
fini ...).

\newpage

%% ===== Page 168 =====

Je vais revenir au cas $S = Spec \, \mathbb{Q}$, alors

$\mathbb{G}_m(\mathbb{Q}) = \mathbb{Q}^* \simeq \{\pm 1\} \times \mathbb{Z}^{(\mathcal{P})}$

où $\mathcal{P} \sim$ l'ens des nbs premiers,

$\pi_1 = \operatorname{Gal}(\overline{\mathbb{Q}}/\mathbb{Q}) = \Gamma_{\mathbb{Q}}$ .

La question est donc d'étudier les hom

(358) $\hat{\zeta}_p : \Gamma_{\mathbb{Q}}^+ \to \hat{\mathbb{Z}}(1)$

associés aux $p \in \mathcal{P}$, et de répondre notamment
aux questions

a) Les $\hat{\zeta}_p$ sont-ils surjectifs, ou du moins l'image est elle un sous-groupe ouvert ?

b) Les $\hat{\zeta}_p$ sont ils linéairement indépendants sur $\hat{\mathbb{Z}}$ ? Plus précisément, si $p_1, \dots, p_n$ sont des nbs premiers distincts, l'hom
$$ \Gamma_{\mathbb{Q}}^+ \to \hat{\mathbb{Z}}^n(1) $$
$$ (\hat{\zeta}_{p_1}, \dots, \hat{\zeta}_{p_n}) $$
est-il surjectif, ou du moins l'image de $\Gamma_{\mathbb{Q}}^+$ est-elle un sous-groupe ouvert ?

Je vais commencer par regarder un peu la question a) (le cas [unclear: que l'on a en vue dans la descente] de Fermat étant celui où $p=2$). Pour ceci, on décompose $\hat{\zeta}_p$ en ses composantes

(359) $\hat{\zeta}_{p, \ell} : \Gamma_{\mathbb{Q}}^+ \to \hat{\mathbb{Z}}_{\ell}(1) \qquad (\ell \in \mathcal{P})$

\newpage

%% ===== Page 169 =====

732

j'aimerais prouver (au moins !) que ce sont
des hom. surjectifs, ou du moins surjectifs
sauf pour un nb fini de $\ell$, et à image
ouverte pour tt $\ell$. Etudier $\xi_{p,\ell}$, c'est
se restreindre à la tour des isogénies
$$ (339) \qquad 1 \to \mu_n \to \mathbb{G}_m \xrightarrow{\lambda \mapsto \lambda^n} \mathbb{G}_m \to 1 $$
pour $n$ puissance de $\ell$
$$ (360) \qquad \text{la tour des} \quad n=\ell^\nu \quad , $$
et étudier l'image inverse de $p \in \mathbb{G}_m(\mathbb{Q})$
par ces isogénies. Il est naturel alors
de regarder le groupe de décomposition
$$ (361) \qquad D_\ell = \Pi_1(\mathbb{Q}_p, \bar{\mathbb{Q}}_p) \hookrightarrow \Gamma_\mathbb{Q} $$
et le caractère cyclotomique $\ell$-adique induit
$$ (362) \qquad D_\ell \xrightarrow{\chi_\ell} \mathbb{Z}_\ell^\ast \quad , $$
ou mieux, le sous-groupe d'inertie
$$ (363) \qquad I_\ell \subset D_\ell $$
et le caractère induit par $\chi_\ell$ sur $I_\ell$.
On [unclear] en effet
$$ (364) \qquad I_\ell^+ \overset{\text{déf}}{=} \operatorname{Ker}(\chi_\ell : I_\ell \to \mathbb{Z}_\ell^\ast) \subset \Gamma_\mathbb{Q}^+ $$
car les $\chi_{\ell'}$ ($\ell' \neq \ell$) sont triviaux sur $I_\ell$
(ils ne sont pas triviaux cependant sur $D_\ell$ !).

\newpage

%% ===== Page 170 =====

On va donc regarder la restriction

(365) $\overline{\xi}_{pre} | I_l^+ : I_l^+ \to \hat{\mathbb{Z}}_l(1)$

et on peut se demander déjà si cet
homomorphisme est surjectif, ou si
du moins son image est ouverte ?

\newpage

%% ===== Page 171 =====

733

16°) Réflexions sur le théorème de Fermat.

Le "théorème de Fermat" peut s'énoncer en
disant que pour $n \ge 3$, la courbe de Fermat
affine sur $\mathbb{Q}$
(366) $$U_n \overset{\text{def}}{=} \left\{ (x_0, x_1, x_\infty) \in \mathbb{P}^2 \mid \begin{matrix} x_0, x_1, x_\infty \text{ inv.} \\ x_0^n + x_1^n + x_\infty^n = 0 \end{matrix} \right\}$$
$$ \simeq \{ (x_0, x_1, x_\infty) \in \mathbb{G}_m^{\{0,1,\infty\}} / \mathbb{G}_m \mid x_0^n + x_1^n + x_\infty^n = 0 \}$$

n'a pas de [unclear: soluti] pt rat / $\mathbb{Q}$. Nous
avions introduit cette "tour" des courbes
de Fermat dans § 11°) (p. 696 ff), et dans
14°) on a vu la description de $U_n$ comme
image inverse de
(367) $$ U_1 \subset \overbrace{\mathbb{G}_m \otimes_\mathbb{Z} M}^{T} \simeq \mathbb{G}_{m \mathbb{Q}}^{\{0,1,\infty\}} / \mathbb{G}_{m \mathbb{Q}} $$
$$ \Vert $$
$$ \{ (x_0, x_1, x_\infty) \in \mathbb{G}_m^{\{0,1,\infty\}} / \mathbb{G}_m \mid x_0 + x_1 + x_\infty = 0 \} $$

par l'isogénie de Kummer d'indice $n$
$$ x \mapsto x^n $$
(368) $$ 1 \to \mu_n \otimes_{\mathbb{Z}} M \to T \longrightarrow T \to 1 \ . $$

Calculons le genre $g_n$ de $\hat{U}_n$, qui est un
revêtement de degré $n^2$ de $\hat{U}_1 \simeq \overline{M}_{0,3}$ ,
[crossed out: ramifié au-dessus de] ramifié au dessus des pts $0, 1, \infty$ , et
étale en dehors. La formule de Hurwitz

\newpage

%% ===== Page 172 =====

donne

$$ \underbrace{\underbrace{(2-2g_n)}_{EP \, de \, \hat{U}_n} - 3n}_{EP \, de \, U_n} = n^2 ( \underbrace{2 - 3}_{EP \, de \, U_1} ) $$

i.e. $$ 2(g_n-1) = n^2-3n $$ i.e.

(369) $$ g_n = 1 + \frac{1}{2} n(n-3) $$

Donc ... 

(370) $$ g_1 = g_2 = 0 \quad , \quad g_3 = 1 \quad , \quad g_n > 1 \text{ si } n \ge 4 $$

En tout état de cause, tous les $U_n$ sont
anabéliens, et aussi "anabéliens" les $\hat{U}_n$ pour $n \ge 4$.

Si le yoga développé dans la première
partie de ces notes est valable, l'existence
(ou la non existence) de points rat / $\mathbb{Q}$ d'une
courbe $U_n$ correspond à un scindage
de l'extension de groupes profinis

(371) $$ 1 \to \Pi_1(U_{n \bar{\mathbb{Q}}}) \to \Pi_1(U_n) \to \Gamma_{\mathbb{Q}} \to 1 $$
[MARGIN: $\Pi_1(U_1) \simeq \hat{M}_{0,3}$ ]

qui soient distincts des "scindages à l'infini"
(localisés en les $3n$ pts à l'infini de $U_n$ ... ), ou
encore, les scindages tels que le centralisateur de
$$ (\Gamma_{\mathbb{Q}}^+ = \operatorname{Ker}(\Gamma_{\mathbb{Q}} \xrightarrow{\chi} \hat{\mathbb{Z}}^*) ) $$
dans $\Pi_1(U_n)$ soit réduit à $1$. Donc - si le "th.

\newpage

%% ===== Page 173 =====

734

de Fermat" est bel et bien valable (pour $U_n$, disons),
il devrait se refléter par la non-existence d'un
tel scindage. D'ailleurs, si nous arrivions à
prouver l'inexistence de scindages de (371) autres
que ceux "provenant des pts à l'infini",
le th. de Fermat serait prouvé pour $U_n$, même
sans avoir établi le "yoga anabélien" (dont une
démonstration semble bien hors d'atteinte !). Or
il n'est nullement exclu qu'une analyse assez
poussée de $M_{0,3}$ et $N_{1,1}$ dans la direction
actuelle - impliquant aussi des "approxi-
mations" de plus en plus fines de
$\Pi_1(U_n)$ et de la structure d'extension (371),
permette finalement de prouver l'inexistence
de scindages "non triviaux" de (371), pour
tout $n \ge 3$. Chaque nouvel affinement dans
les "approximations de Lie" de $M_{0,3}$ devrait
amener un réexamen de la situation (371)
et de la question de ses scindages, dans l'optique
d'une démonstration éventuelle du "th. de Fermat".

\newpage

%% ===== Page 174 =====

Il s'impose de partir d'approximations de
Lie au moins aussi fines que celle du
présent $\S$, qui apparait comme l'approximation
"minima" dans laquelle on puisse exprimer
"algébriquement" la tour des courbes de
Fermat, ~~via~~ comme on l'a vu dans 13°),
à partir de $M_{0,3} \subset \hat{M}_{0,3}$.

Tenant compte des opérations de $\mathfrak{S}_3$
sur $U_n$, la suite ex.te (371) se prolonge
d'ailleurs dans une suite ex.te plus
riche

$$(372) \quad 1 \longrightarrow \pi_1(U_{n,\overline{\mathbb{Q}}}) \longrightarrow \underset{\overset{\wedge}{\pi_1(U_1,\mathfrak{S}_3) \simeq M_{0,3}}}{\pi_1(U_n,\mathfrak{S}_3)} \longrightarrow \Gamma_{\mathbb{Q}} \times \mathfrak{S}_3 \longrightarrow 1$$

et les calculs faits dans le présent $\S$ semblent
suggérer que cette structure d'extension plus
riche jouera un rôle important.

Du point de vue de l'approximation
de Lie de $M_{0,3}$ développée dans le présent
$\S$, il y a lieu de remplacer la structure
d'extension (372) par une structure quotient
moins fine, en divisant par un sous-groupe
invariant de $\pi_1(U_n, \mathfrak{S}_3)$ contenu dans

\newpage

%% ===== Page 175 =====

$73r$

$\Pi_1(\mathcal{U}_n, \overline{\mathbb{Q}})$, savoir l'intersection dans $M_{0,3}$ de ce
groupe avec $[\overline{\pi_0, \hat{j}}, \overline{\pi_0, \hat{j}}]$. On trouve un
diagramme commutatif

(373)
\begin{tikzcd}
1 \ar[r] & M^\wedge[n] \ar[r] \ar[d, equal] & \widetilde{\Pi_1(\mathcal{U}_n, \mathfrak{S}_3)} \ar[r] \ar[d, "p_n"'] & \Gamma_{\mathbb{Q}} \times \mathfrak{S}_3 \ar[r] \ar[d, "\chi \times id"] & 1 \\
1 \ar[r] & M^\wedge[n] \ar[r] \ar[d] & \widetilde{G}_{\rho, \sigma}[n] \ar[r] \ar[d] & \hat{\mathbb{Z}}^* \times \mathfrak{S}_3 \ar[r] \ar[d, equal] & 1 \\
1 \ar[r] & M^\wedge \ar[r] & \widetilde{G}_{\rho, \sigma} \ar[r] & \hat{\mathbb{Z}}^* \times \mathfrak{S}_3 \ar[r] & 1
\end{tikzcd}

où l'image de $M^\wedge[n] \to M^\wedge$ est $n M^\wedge$.

Voici une façon de construire "algébriquement"
ces diagrammes (373) (pour $n$ variable)
à partir du diagramme

(374)
\begin{tikzcd}
1 \ar[r] & M^\wedge \ar[r] \ar[d, equal] & \overset{\mathcal{E}_1}{\widetilde{\Pi_1(\mathcal{U}_1, \mathfrak{S}_3)}} \ar[r] \ar[d] & \Gamma_{\mathbb{Q}} \times \mathfrak{S}_3 \ar[r] \ar[d, "\chi \times id"] & 1 \\
& & (M_{0,3}/[\overline{\pi_0, \hat{j}}, \overline{\pi_0, \hat{j}}]) \ar[d] & & \\
1 \ar[r] & M^\wedge \ar[r] & \widetilde{G}_{\rho, \sigma} \ar[r] \ar[d, "\simeq"'] & \hat{\mathbb{Z}}^* \times \mathfrak{S}_3 \ar[r] & 1 \\
& & (\hat{\mathbb{Z}}^* \times \mathfrak{S}_3).M^\wedge & & 
\end{tikzcd}

On choisit un scindage de la première,
qui donne [unclear] un scindage de la seconde
(scindage [unclear: par] un [unclear: sous-groupe] fini [unclear]
[unclear: engendré par] un unique élément de $M^\wedge \dots$). [unclear: En prenant]
[unclear] on peut déduire $U_n$ en-dessous de [unclear]

[MARGIN: On a vu qu'en un sens un peu tautologique les scindages de l'ext. $\mathcal{E}_1$ [unclear: équivalent] à ceux de l'ext. [unclear] (dans 12°), i.e. à un [unclear: ensemble] de choix tels [unclear: canoniques]]

\newpage

%% ===== Page 176 =====

pt base $\vec{a}_j$ de $U_{0,3 \mathbb{Q}} \simeq \mathcal{U}_{1 \mathbb{Q}}$ - il définit
un pt. géom. de chacun des $\mathcal{U}_{n \mathbb{Q}}$ au
dessus de $\vec{a}_j$. En termes de ce scindage
de $\mathcal{E}_1$, permettant d'écrire

$$ (375) \quad \mathcal{E}_1 \simeq (\Gamma_{\mathbb{Q}} \times \mathfrak{S}_3) . \hat{M} $$
(pr. 1/2 direct)

les $\mathcal{E}_n \subset \mathcal{E}_1$ s'identifient aux
$\simeq \Pi_1(\mathcal{U}_n, \mathfrak{S}_n)/...$

produits semi-directs

$$ (376) \quad \mathcal{E}_n = (\Gamma_{\mathbb{Q}} \times \mathfrak{S}_3) . (n\hat{M}) $$
$\overset{\text{déf}}{=} \hat{M}[n]$

(Ces interprétations sont des resucées des
généralités de la page-fleuve de p. 712
ff, et mériteraient de figurer dans
la scholie-proposition en question...)

Un pt rationnel $x_n / \mathbb{Q}$ de $\mathcal{U}_n$ définit
un scindage de l'extension (de $\Gamma_{\mathbb{Q}}$ par $n\hat{M} = \hat{M}[n]$)

$$ (377) \quad \mathcal{E}_n^! \simeq \Gamma_{\mathbb{Q}} . (n\hat{M}) \quad (\text{1/2 direct}) $$

quotient de (371), dont la connaissance

\newpage

%% ===== Page 177 =====

736

modulo $\hat{M}[n]$-conjugaison implique celle
des scindages correspondants de $\mathcal{E}_1$, donc
implique la connaissance des pts correspon-
-dant $x_n$ de $U_n(\mathbb{Q})$, ce qui signifie
que $x_n$ est connu mod. opération de $V_n =$
$\mu_n(\mathbb{Q})^{\{0,1,\infty\}} / \mu_n(\mathbb{Q})$. Or on a

$$ (378) \quad V_n \overset{df}{=} \mu_n(\mathbb{Q})^{\{0,1,\infty\}} / \mu_n(\mathbb{Q}) = 1 $$

si $n$ impair

car $\mu_n(\mathbb{Q}) = 1$ si $n$ impair - donc dans
ce cas $x_n$ est connu quand on connaît
le scindage correspondant, i.e. modulo la
donnée de $\hat{M}[n]$-conjugaison.

Ceci nous amène à déterminer toutes les
classes de $\hat{M}[n]$-conjugaison possibles de
tels scindages, afin d'examiner (si faire se
peut) lesquels peuvent raisonnablement
correspondre à un $x_n \in U_n(\mathbb{Q})$. Ces classes
correspondent aux éléments de

$$ (379) \quad H^1(\Gamma_{\mathbb{Q}}, _n\hat{M}) \simeq H^1(\Gamma_{\mathbb{Q}}, \hat{M}) \simeq H^1(\Gamma_{\mathbb{Q}}, \hat{\mathbb{Z}}(1)) \otimes_{\hat{\mathbb{Z}}} M $$
[unclear: $\simeq G(\mathbb{Q})^\wedge$] \qquad $(\mathbb{Q}^*)^{\wedge}$

\newpage

%% ===== Page 178 =====

$$G \overset{dfn}{=} G_{m\mathbb{Q}}^{\{0,1,\infty\}} / G_{m\mathbb{Q}}$$
est le tore de dim 2 sur $\mathbb{Q}$ dans lequel
se trouvent plongés les $\mathcal{U}_n$. On voit que
tout se passe comme si on étudiait
simplement l'existence (!) et le repérage
cohomologique des pts $\mathbb{Q}$-rationnels de
$G$ lui-même - sans nous préoccuper
de l'équation $x_0^n + x_1^n + x_\infty^n = 1$ qui définit
$\mathcal{U}_n$. Cela montre : à l'évidence que
l'approximation sous laquelle nous travaillons
est beaucoup trop grossière - en fait,
pratiquement nous restons dans le
domaine d'une analyse typiquement
"abélienne". Il est clair qu'il est
nécessaire de passer à des approximations
plus fines.

\newpage

%% ===== Page 179 =====

737

## 52 Enveloppes schématiques des L-groupes.

1) Notion de L-groupe, de LT-groupe (cas discret).
La notion de L-groupe est la généralisation
naturelle de celle de "groupe de lacets",
qui s'impose si on veut pouvoir englober
des groupes fondamentaux mixtes de
surfaces orientables (de "caractère fini", i.e.
hom. à un $X \setminus S$, $X$ surface compacte,
$S$ partie finie) ou un groupe fini d'auto-
morphismes.

Nous avons vu que si $\Pi$ est un
groupe de lacets (discret), alors tout sous-
groupe d'indice fini $\Pi_0$ hérite une
structure de lacets, en prenant comme
"sous-groupes lacets" de $\Pi_0$ les intersections
avec $\Pi_0$ des ss-gr. lacets de $\Pi$ . [N.B. On
s'aperçoit qu'il n'y a que dans les cas
$(g,\nu)$ avec $g=0$, $\nu \in \{0,1\}$, où la structure de lacets
n'est pas déterminée par la seule donnée d'une
famille de classes de conjugaison de ss-groupes
d'iner. ($\simeq \mathbb{Z}$) - le seul cas qui ici ferait pb

\newpage

%% ===== Page 180 =====

est le cas (0,1) de $\Pi \simeq \mathbb{Z}$ - dans le cas (0,0) $\Rightarrow \Pi=\{1\}$
et $\Pi_0 = \Pi$ ... ]. Inversement, la structure
à lacets de $\Pi$ se récupère à partir de
celle de $\Pi_0$ (quand elle existe !), deux
sous-groupes lacets de $\Pi$ étant les centralisateurs
dans $\Pi$ des sous-groupes-lacets de $\Pi_0$. Le cas
exceptionnel à engendrement fini,
Soit $\Pi$ un groupe discret, et
soit $\mathcal{L}_0(\Pi)$ l'ens des structures de groupes
[unclear: abélien i.e.]
à lacets sur $\Pi$, de type $(g,\nu)$ [unclear]
Pour les sous-groupes $\Pi_i$ d'indice fini,
les $\mathcal{L}_0(\Pi_i)$ forment un système inductif
d'ensembles. Si $\Pi_i$ est distingué,
(1) $\mathcal{L}_0^{G_i}(\Pi_i) = $ ens des structures à lacets
sur $\Pi_i$ invariantes par
$G_i = \Pi/\Pi_i$.

S'il existe un $\Pi_0$ d'indice fini tel que
pour le syst. des $\Pi_i \subset \Pi_0$,
$\mathcal{L}_0(\Pi_i) \neq \emptyset$, alors $\Pi$ admet
une génération finie et il s'ensuit
que l'ens. des sous-groupes d'indice fini
sont cofinaux dans l'ens des sous-groupes
d'indice fini. On posera alors

\newpage

%% ===== Page 181 =====

738

$$L^\natural(\Pi) = \varinjlim_{\substack{\Pi_i \text{ distingué} \\ \text{d'indice fini} \\ \text{de } \Pi}} L_0^\natural(\Pi_i)$$

[MARGIN: N.B. Si $\Pi_i$ n'est pas distingué dans $\Pi$, on peut définir $L_0^\natural(\Pi_i)$ comme l'ens. des structures de $L_0(\Pi_i)$ invariantes par $\Pi$, i.e. telles que pour tt $g \in \Pi$, on ait [unclear] $\in L_0(g \Pi_i g^{-1})$ (qui est [unclear] d'indice fini...)]

On en déduit des morphismes canoniques ($\Pi_i$ dist. d'indice fini dans $\Pi$)

$$(3) \quad L_0^\natural(\Pi_i) \to L(\Pi)$$

sont injectifs. Si $\lambda \in L(\Pi)$, on dit que $(\Pi, \lambda)$ est un $L$-groupe (ou que $\lambda$ est une $L$-structure sur $\Pi$, et si $\lambda \in \operatorname{Im}(L_0^\natural(\Pi_i) \to L(\Pi))$, on dit que $\Pi_i$ est un sous-groupe de définition du $L$-groupe $\Pi$ - on le considère comme un groupe [unclear] pour l'unique $\lambda_i \in L_0^\natural(\Pi_i)$ donnant $\lambda$.

Soient $\Pi$ un groupe discret, $\Pi'$ un sous-groupe de $\Pi$ d'indice fini. Alors on a un homomorphisme canonique d'induction des $L$-structures

$$(4) \quad L(\Pi) \hookrightarrow L(\Pi')$$

qui est injectif, et dont l'image est formée de l'ens des $L$-structures sur $\Pi'$ qui sont

\newpage

%% ===== Page 182 =====

"invariantes sous $\pi$" dans un sens évident.
En particulier, si $\pi$ est une extension
d'un groupe fini par un groupe
discret $\pi_0$, alors les $\mathcal{L}$-structures sur
$\pi$ correspondent aux $\mathcal{L}$-structures
sur $\pi_0$ invariantes par $\pi / \pi_0 \to$ [unclear: $Out(\pi_0)$]
opérant sur l'ensemble des $\mathcal{L}$-structures de $\pi_0$
de façon évidente. [NB $\pi \mapsto \mathcal{L}(\pi)$ est
un foncteur en les groupes extérieurs
$\pi$, avec les isomorphismes extérieurs
comme flèches ...]

Exemple standard : soient $X$ une
surface 0-conn. "de présentation finie",
$G$ un groupe fini opérant sur $X$,

(5) $\pi = \pi_1(X, G ; -)$

où $-$ est un pt base de $X$, de
sorte qu'on a une suite exacte

(6) 
\begin{tikzcd}[row sep=small]
1 \ar[r] & \pi_0 \ar[r] & \pi \ar[r] & G \ar[r] & 1 \\
& \pi_1(X,-) \ar[u, equal]
\end{tikzcd}

On a donc sur $\pi_0$ une $\mathcal{L}$-structure : [unclear: locale]

\newpage

%% ===== Page 183 =====

739

invariante par les op. ext. de $G$, d'où
une $L$-structure sur $\Pi$. Il semble
que tout $L$-groupe soit isomorphe à
un groupe obtenu ainsi - mais c'
est là un résultat qui semble
assez profond (esst. équivalent à celui de l'existence
de pts fixes des sous-groupes finis des groupes
de Teichmüller...), et qui n'est
pas compris totalement. Il faudrait
dégager un "bon énoncé" général, reliant la
2-catégorie des $L$-groupes, ou celle des
surfaces $X$ de caract. finie, $d$-pointées
(munies d'une "signature" $d$ à supports
finis (à valeurs dans $\mathbb{N}^*$)), en associant à
une telle surface un $\Pi_1(X, d, a)$,
relatif à un pt base $a \in X - supp \; d$.
A l'exclusion de quelques cas simplistes ("triviaux") on
devrait trouver alors un $L$-groupe $\Pi$,
isomorphe à un $\Pi_1(X, G, a)$ avec $G$
opérant fidèlement, ce qui devrait s'exprimer

\newpage

%% ===== Page 184 =====

en termes de $\Pi$ par le fait que
le centralisateur de tt sous-groupe
d'indice fini de $\Pi$ est réduit à $\{e\}$.
Ceci devrait donner une description
complète, en termes topologiques, des
L-groupes ayant ctt propriété supplémen-
taire. Il y aurait ensuite encore
un travail de raccordage à faire,
en termes de certaines extensions de
tels L-groupes par des groupes finis,
pour trouver le cas général. De
plus, il faudrait faire le lien
avec les "topos L-galoisiens" associés,
en paraphrasant les développements
correspondants sur les topos galoisiens etc
associés aux groupes loc. cts.

\newpage

%% ===== Page 185 =====

740

Il se pourrait que le "bon" contexte géné-
ral pour définir et étudier des groupes
de Teichmüller généraux soit celui
des L-groupes, plutôt que celui des
groupes locaux.

Prenant p. ex. le groupe 
[unclear: rayé] $G^+_{0,3} \simeq S\ell(2, \mathbb{Z}) / \{ \pm 1 \}$
(4) $$ \simeq \Pi_1( U_{0,3}(\mathbb{C}), \mathfrak{S}_3 - J ),$$
celui-ci est muni d'une L-structure,
qui peut être définie en termes de
l'unique classe de conjugaison du
sous-groupe L-engendré par $\varepsilon_0 = \begin{pmatrix} 0 & -1 \\ 1 & 1 \end{pmatrix}$,
ou du L-engendré par $\varepsilon_1 = \begin{pmatrix} 1 & 0 \\ 1 & 1 \end{pmatrix}$.
Le L-groupe (discret) a la propriété remar-
quable que tout L-automorphisme
"direct" (induisant l'identité sur le module d'orientation)
de celui-ci est intérieur. Il se
pourrait que cette propriété, ajoutée à
celle que le centre est trivial (ou
mieux, peut-être, que le centralisateur

\newpage

%% ===== Page 186 =====

de tout sous-groupe d'indice fini est
réduit à $\{1\}$) caractérise $\widehat{\mathbb{G}}_{0,3}^+$ à
isomorphisme (pas unique) près.

NB Pour prouver que tout L-autom.
[unclear] de $\widehat{\mathbb{G}}_{0,3}^+$ est intérieur, le
point essentiel me semble de prouver
qu'il conserve le sous-groupe
$\Pi_{0,3}$, noyau de
$$1 \to \Pi_{0,3} \to \widehat{\mathbb{G}}_{0,3}^+ \to \mathfrak{S}_3 \to 1.$$
En fait, tout automorphisme de
la structure de groupe de $\widehat{\mathbb{G}}_{0,3}^+$
(pas nécessairement un L-autom.)
conserve le sous-groupe distingué
$\Pi_{0,3}$. Cela provient du

Lemme Tout sous-groupe distingué $\pi \subset \widehat{\mathbb{G}}_{0,3}^+$ tel
que $\widehat{\mathbb{G}}_{0,3}^+ / \pi \simeq \mathfrak{S}_3$, alors $\pi = \Pi_{0,3}$.
Plus précisément, tout hom. surjectif
$p: \widehat{\mathbb{G}}_{0,3}^+ \to \mathfrak{S}_3$ est conjugué à l'hom.
canonique (dont le noyau est $\Pi_{0,3}$).

Dém. On utilise la présentation
de $\widehat{\mathbb{G}}_{0,3}^+$
$(\rho^3 = \sigma^2 = 1)$

\newpage

%% ===== Page 187 =====

741

qui implique qu'un hom. de $\widehat{\mathcal{G}}_{0,3}^+$ dans
un autre groupe, dont l'image n'est pas
d'ordre $1, 2$ ni $3$ i.e. d'ordre $\ge 4$, transforme
$\rho$ et $\sigma$ resp. en des éléments d'ordre $3$
et d'ordre $2$. Donc si l'image est $\mathfrak{S}_3$,
$\rho$ est transformé en $\rho$ ou $\rho^{-1}$ (OPS, quitte
à conjuguer, que c'est $\rho$) et $\sigma$ est trans-
formé en un des $\sigma_0, \sigma_1$ ou $\sigma_2$ (OPS, quitte à
conjuguer par $\rho$ ou $\rho^{-1}$, que c'est $\sigma_0$), et
on gagne.

On peut appeler groupe de Teichmüller
d'un $L$-groupe $\Pi$, le groupe de ses
$L$-automs extérieurs - ou peut-être,
mieux, le groupe des $L$-automor-
phismes extérieurs du sous-groupe $\Pi^+$
(d'indice $1$ ou $2$) de $\Pi$, formé des
$g \in \Pi$ tels que l'autom. $int(g)$ soit
direct. La propriété envisagée de $\widehat{\mathcal{G}}_{0,3}^+$
est que son groupe de Teichmüller
"direct" est réduit à $1$. La propriété

\newpage

%% ===== Page 188 =====

correspondante du L-groupe $\overline{G}_{0,3} \simeq \operatorname{Gl}(2, \mathbb{Z})/\pm 1$
est que tout L-autom du sous-
groupe $\overline{G}_{0,3}^+$ est induit par un unique
élément de $\overline{G}_{0,3}$. Cette propriété est
peut être caractéristique des L-groupes
discrets $\overline{G}_{0,3}$.

J'ai envie de présenter sous une
forme légèrement différente la notion
de L-groupe $\Pi$ en termes de la
donnée d'un $\mathbb{Z}$-module $T$
libre de rg 1 (le $\mathbb{Z}$-module
d'orientations) et d'une [ensemble fini rayé] famille
$(k_i)_{i \in I}$ de germes d'hom. extérieurs

(5) $$k_i : T \longrightarrow \Pi \quad (i \in I)$$
(germes d'hom. ext.) ,

où la notion de "germe" est relative
à celle de sous-groupes d'indices finis
de $T$. [Pour un tel, dirons-nous rayé]
[L'axiome à mettre est rayé]

\newpage

%% ===== Page 189 =====

742
et tout sous-groupe [unclear: ouvert] $\Pi_0$ de $\Pi$ d'indice
fini, on trouve un ens de germes
d'hom. extérieurs
$$T \longrightarrow \Pi_0$$
induites par les hom. [unclear] effectifs
$nT \longrightarrow \Pi$ de la classe $k_i \natural$ — [unclear]
un nb fini de germes d'hom. ext.
$T \longrightarrow \Pi_0$ pour $i$ fixé, d'où un nb fini
$i$ variable dans $I$. Ceci dit, l'unique
axiome d'une $L$-structure, c'est
qu'on puisse trouver un sous-
groupe $\Pi_0$ d'indice fini, tel que
la "structure induite" provienne
d'une structure : [unclear: Facets] de $\Pi_0$.

On peut généraliser la notion de
$L$-structure, en s'inspirant de la
définition précédente (p. ex.), et
en y laissant tomber l'exigence que
$\Pi_0$ soit d'indice fini - mais en

\newpage

%% ===== Page 190 =====

prenant soin par contre que $\Pi_0$ soit
distingué. On exige que pour tout
$i \in I)$ les germes d'hom $k_i : \hat{T} \to \Pi$
de la classe $K_i^\natural$ se factorisent par

[MARGIN: NB il n'est pas clair a priori qu'un ens. de germes d'hom. ext. soit fini]

ce qui donne un ens. de
germes extérieurs de $T$ dans $\Pi_0$ -
et on exige que cette structure-là
sur $\Pi_0$ provienne (pour $\Pi_0$ convenable) d'une structure
à lacets sur $\Pi_0$, évidemment
uniquement déterminée une fois
$\Pi_0$ choisi. Quant à $\Pi_0$, il est
défini à commensurabilité près -
i.e. si $\Pi_0, \Pi_0'$ conviennent, alors
$\Pi_0 \cap \Pi_0'$ est d'indice fini dans $\Pi_0$ et
$\Pi_0'$, et convient également.

M'inspirant de la terminologie
"antérieure", j'ai envie (à défaut d'
une meilleure suggestion) d'appeler
les groupes $\Pi$ munis d'une telle

\newpage

%% ===== Page 191 =====

793

structure - liée à une classe de
sous-groupes discrets commensurables
entre eux - des L-T-groupes.
(Ici L est initiale de "locats", et T
de "Teichmüller") - ce sont des
sortes de mixtions de groupes de
Teichmüller et de L-groupes. Pour
illustrer ce point, partons d'une
extension

(6) $$1 \to \pi \to \tilde{G} \to \mathcal{T} \to 1$$

où $\pi$ est muni d'une L-structure
invariante par $\tilde{G}$, alors $\tilde{G}$ est
un L-T-groupe admettant (si $\pi$ est
un L-groupe : locats) $\pi$ comme
sous-groupe de définition. Partant
d'un groupe : locats $\pi$, on peut
prendre p.ex.

$$ (7) \begin{cases} \tilde{G} = \operatorname{Aut}_L(\pi) \qquad \text{d'où} \\ \mathcal{T} = Autext_L(\pi) \end{cases} $$

\newpage

%% ===== Page 192 =====

et ainsi le groupe total $\mathfrak{G}$,
extension du groupe de Teichmüller
$\tau$ par le groupe à lacets $\pi$, apparait
comme un $L-T$-groupe.

\newpage

%% ===== Page 193 =====

744

2°) Cas profini

L'extension de la notion de L-groupe $\Pi$
-- ou LT-groupe -- en termes de
genres d'hom extérieurs
(8) $\quad$ [unclear: $K_S : T \to \Pi$] ,
est essentiellement évidente. Voici
l'exemple le plus important pour
nous. 

Soit d'abord $k$ un corps alg.
clos et soit $\mathcal{M}$ une
multiplicité algébrique sur $k$, qui
soit de t.f. sur $k$, séparée
et lisse en dim relative 1. On
doit pouvoir prouver que
$\mathcal{M}$ puisse se déduire par une
courbe lisse et quasi-projective
$\mathcal{U}$ sur $k$, un groupe fini $G$
qui opère sur $\mathcal{U}$, comme (quotient)
(9) $\quad \mathcal{M} = (\mathcal{U}, G) =$ Topos annelé des
faisceaux [unclear: abéliens] sur
$\mathcal{U}$ où $G$ opère...

\newpage

%% ===== Page 194 =====

Plus intrinsèquement, il doit être
vrai qu'il existe un rev. étale fini
de $M$ (qu'on peut supposer galoisien,
de groupe $G$) qui soit un
[crossed out: schéma algébrique] - lequel sera néc.
quasi-projectif, lisse, de dim relative 1
sur $k$. Supposons en tous cas
qu'il en soit ainsi, choisissons
un pt base géom. $a$ de $M$, et
considérons

(10) $$\pi = \pi_1(M, a) \simeq \pi_1(U, G; a) \ .$$

Nous supposerons désormais $k$ de
car. zéro. Supposons d'abord $k$
alg clos. Utilisons la suite exacte

(11) $$1 \longrightarrow \underset{\pi_0}{\pi_1(U_{\tilde{a}})} \longrightarrow \underset{\pi}{\pi_1(M, G; a)} \longrightarrow G \longrightarrow 1$$

et la structure : [unclear: toute] naturelle
sur $\pi_0$, on trouve une $L$-structure
sur $\pi = \pi_1(M, a)$, qui [unclear: manifeste] sur

\newpage

%% ===== Page 195 =====

745

dépend pas du choix du rev. étale fini galoisien $U$ de $M$, et du relèvement choisi du pt géom. $a$ de $M$ en un pt. géom. de $U$. Nous munirons tacitement $\Pi = \Pi_1(M, a)$ de cette L-structure.

Si à nouveau $k$ est un corps quelconque (le cas le plus intéressant étant celui où $k$ est ext. de type fini de $\mathbb{Q}$), on aura une suite exacte

$$ (12) \quad 1 \longrightarrow \Pi_1(M_{\bar{k}}, a) \longrightarrow \Pi_1(M, a) \longrightarrow \Gamma_{\bar{k}/k} \longrightarrow 1 $$

et $\Pi_1(M_{\bar{k}}, a)$ apparait comme muni d'une L-structure invariante par $\Pi_1(M, a) = \Pi$, de sorte que $\Pi$ apparait comme un $\Gamma_{\bar{k}/k}$ - groupe.

Je suis particulièrement intéressé au cas "universel" du type $(g,\nu)$ — mais la base naturelle n'est pas alors le

\newpage

%% ===== Page 196 =====

spectre d'un corps $k$, mais la multiplicité modulaire $M_{g,\nu}$. [A vrai dire, pour la structure de L-groupe discret $\Pi_0$, jouant le rôle de groupe de référence, on doit pouvoir introduire une multiplicité modulaire convenable sur $\mathbb{Q}$, et même sur $\mathbb{Z}$. Le fait qui risque de faire des difficultés conceptuelles pour une telle construction, est qu'une multiplicité "modulaire" (sur $\mathbb{Q}$, ou sur $\mathbb{Z}$...) décrit en principe une catégorie fibrée — telle celle des schémas en courbes relatifs de type $(g,\nu)$ — alors que les L-multiplicités relatives correspondantes à un "type" $\Pi_0$-donné forment (en groupoïdes) une $2$-catégorie fibrée sur $(Sch)$, et pas seulement une catégorie fibrée...]

\newpage

%% ===== Page 197 =====

745

On aura donc une suite exacte

(13) $$1 \longrightarrow \Pi_{g,\nu}^\wedge \longrightarrow \underbrace{\Pi_1(\underline{U}_{g,\nu_\mathbb{Q}})}_{\underline{U}_{g,\nu}} \longrightarrow \underbrace{\Pi_1(\underline{M}_{g,\nu_\mathbb{Q}})}_{\underline{M}_{g,\nu}} \longrightarrow 1$$

où $\underline{U}_{g,\nu}$ est la "courbe de type $g,\nu$ universelle" (sur car. $0$). Ici c'est $\Pi_1(\underline{U}_{g,\nu_\mathbb{Q}})$ qui apparaît comme un TL-groupe profini.

Je voudrais dégager une propriété commune à ce cas "universel" (13), [unclear: et au cas d'un "groupe de Galois relatif" - crossed out] : un corps de base $k$ extension de type fini de $\mathbb{Q}$ (le premier cas se ramène d'ailleurs essentielle-ment au second, en prenant un "point générique" de $\underline{M}_{g,\nu}$ ...). Pour ceci, considérons un LT-groupe $\mathcal{M}$, et un [unclear: auto - crossed out] homomorphisme

$$u : \Gamma \longrightarrow Out \mathcal{M}$$

où $\Gamma = \operatorname{Gal}(\bar{k}/k)$. On peut

\newpage

%% ===== Page 198 =====

p. ex. choisir un sous-groupe [unclear: de type fini]
distingué $\Pi$ dans $M$, qui [unclear: s'insère] dans
une suite exacte

$$(14) \quad 1 \to \Pi \to M \to N \to 1 \quad ,$$

et un hom.
$$(15) \quad K : T \to \Pi \quad , \text{ d'im } K(T) = L$$

[unclear: relatif] : la structure : [unclear: l'action] sur $\Pi$.
Considérons le normalisateur de
$L$ dans $M$

$$(16) \quad \operatorname{Norm}_M(L) \overset{df}{=} N(L)$$

et considérons l'action de $N(L)/L \subset N$
sur $L$. (NB On voit aisément
que l'image de $N(L)/L$ dans $N$ est
un sous-groupe distingué de $N$ ...) Cette
action se fait par un "caractère
cyclotomique", induit par un
caractère

$$(17) \quad N \xrightarrow{\chi_N} \hat{\mathbb{Z}}^* .$$

L'hypothèse essentielle est que l'image

\newpage

%% ===== Page 199 =====

747

de $N$ par $\chi$ (donc celle de $N(L)/L$ par
$\chi$) soit un sous-groupe ouvert de $\hat{\mathbb{Z}}^*$.
Cette hypothèse a un caractère intrinsè-
que au LT-groupe $M$ lui-même (indé-
pendamment du choix de $\Pi$) -
car le caractère cyclotomique
(18) $$ \chi_M : M \to \hat{\mathbb{Z}}^* $$
a un caractère intrinsèque à $M$, et l'
hypothèse signifie que l'image de $M$
dans $\hat{\mathbb{Z}}^*$ est ouverte.

Revenant à $L$ et $N(L)$, on voit
que l'hypothèse faite implique ceci :

(19) Tout sous-groupe ouvert de $N(L)/L$
opère non trivialement sur $L$.

On sait, par des arguments
standards que ceci implique que
pour tout corps $K$ [crossed out text]
(notamment un corps $\ell$-adique $\mathbb{Q}_\ell$),

\newpage

%% ===== Page 200 =====

et tout hom. de $\mathcal{M}$ dans un groupe
$G(K)$, où $G$ est un schéma en
groupes de t.f. affine sur $K$; &

(20) $$u: \mathcal{M} \longrightarrow G(K),$$

la restriction de $u$ à $L$ est "quasi-uni-
potente", ie.

(21) $\exists$ sous-groupe ouvert $L'$ de $L$, tel que
$u(L') \subset G(K)$ soit formé d'éléments
unipotents,

ce qui revient à dire, si $l$ est un
générateur de $L \simeq \hat{\mathbb{Z}}$, qu'on ait

(22) $u(l)$ est un élément quasi-
unipotent de $G(K)$, i.e. $\exists n \in \mathbb{N}^*$,
tel que $l^n$ soit unipotent.

\newpage

%% ===== Page 201 =====

748

30) Enveloppe schématique d'un Teich-groupe profini : les hom. admissibles.

Soient $M$ un $T$-$L$-groupe profini, dont le caractère structurel (18)
$$M \xrightarrow{\chi_M} \hat{\mathbb{Z}}^*$$
ait une image ouverte, $\Pi$ un sous-groupe de définition (c'à d dense à quotients finis), et $\Pi_0 \subset \Pi$ un sous-groupe discret de $\Pi$, qui en soit une "discrétification" - c-à-d un groupe [unclear: abstrait] discret, dont $\Pi$ s'identifie à son complété profini...

Soit $G$ un schéma en groupes [MARGIN: affine de t.f.] sur $\mathbb{Q}$. Soit $\mathbb{A}_\mathbb{Q}$ l'anneau des adèles de $\mathbb{Q}$
(23)
$$\mathbb{A}_\mathbb{Q} = \hat{\mathbb{Z}} \otimes_\mathbb{Z} \mathbb{Q} \subset \prod_{\ell \in P} \mathbb{Q}_\ell$$
et considérons le groupe des points adéliques
(24)
$$G(\mathbb{A}_\mathbb{Q}) \subset \prod_{\ell \in P} G(\mathbb{Q}_\ell) \ .$$
Si $G$ provient de un schéma en groupes

\newpage

%% ===== Page 202 =====

affine de t.f. $G$ sur $\mathbb{Z}$, alors ...

(25) $$ G(A_{\mathbb{Q}}) \simeq \{ (x_l)_{l \in P} \in \prod G(\mathbb{Q}_l) \mid x_l \in G(\hat{\mathbb{Z}}_l) \text{ pour presque tt } l \} $$

On se donne de plus un hom

(26) $$ G_{\mathbb{Q}} \xrightarrow{\chi_G} \mathbb{G}_{m \mathbb{Q}} $$

On s'intéressera à des homomorphismes de groupes

(27) $$ M \longrightarrow G(A_{\mathbb{Q}}) $$

satisfaisant à des propriétés que nous allons énumérer.

1°) Continuité. Celà équivaut à dire que a) les hom. $M \to G(\mathbb{Q}_l)$ sont continus, et b) $\exists S$ partie finie de $P$, telle que pour $l \notin S$, on ait une factorisation (continue)

(28) $$ M \longrightarrow G(\hat{\mathbb{Z}}_l) \quad (l \notin S) $$

- condition qui en fait ne dépend pas du "modèle" sur $\mathbb{Z}$  $G$ de $G_{\mathbb{Q}}$ choisi.

[MARGIN: NB (28) est un hom. de groupes profinis.]

\newpage

%% ===== Page 203 =====

741

ReNB Sauf erreur, la propriété 1) signifie qu'
on peut trouver un schéma en groupes $G$
affine de t.f. sur $\mathbb{Z}$, qui définit $G_\mathbb{Q}$, et
tel que le (27) se factorise par $G(\hat{\mathbb{Z}})$.
(29) $$M \longrightarrow G(\hat{\mathbb{Z}})$$

2°) Compatibilité avec $\chi$ :

(30)
\begin{tikzcd}
M \ar[r] \ar[d, "\chi_M"'] & G(\mathbb{A}_\mathbb{Q}) \ar[d, "\chi_G"] \\
\hat{\mathbb{Z}} \ar[r, hook] & \mathbb{G}_m(\mathbb{A}_\mathbb{Q}) \simeq \mathbb{A}_\mathbb{Q}^* \subset \prod_{l \in P} \mathbb{Q}_l^*
\end{tikzcd}

est commutatif.

3°) Compatibilité avec la [unclear: désintégration] $\pi_0$

(31)
\begin{tikzcd}
\Pi'_0 \ar[r, hook] \ar[d, "\text{factorisation}"'] & \Pi_0 \ar[r, hook] & M \ar[d] \\
G(\mathbb{Q}) \ar[rr, hook] & & G(\mathbb{A}_\mathbb{Q})
\end{tikzcd}
,

pour un sous-groupe d'indice fini
convenable $\Pi'_0$ de $\Pi_0$.

NB Quitte à changer les notations, [unclear: on peut supposer] $\Pi'_0 = \Pi_0$.

4°) Caractère minimal (irrédundance) :
Si $G'$ est un sous-schéma en groupes de

\newpage

%% ===== Page 204 =====

$G_Q$, tel que $\mathcal{M} \to G(\mathbb{A}_Q)$ se
factorise par $G'(\mathbb{A}_Q)$, alors $G=G'$.

NB Si cette hypothèse n'est pas satisfaite,
alors parmi tous les $G'_Q \subset G_Q$ qui la
mettent en défaut, il y en a un
(unique)
plus petit, et en remplaçant $\chi_G$
par $\chi_{G'} = \chi_G | G'_Q ...$ les hyp.
1º) à 4º) sont satisfaites pour $G', \chi'$,
et $\mathcal{M} \to G'(\mathbb{A}_Q)$.

[MARGIN: Topologie sur les s. g. d'un groupe top. - le cas de ses s. g. discrets]
Remarque Soit $\mathfrak{S}_0 \subset \mathcal{M}$ le
sous-groupe de $\mathcal{M}$ formé des $g \in \mathcal{M}$ qui
sont tels que $g(\Pi_0)$ et $\Pi_0$ soient
commensurables (i.e. qui normalisent
le "germe de discrétification" $\Pi_0$ de $\Pi$).
Il est possible qu'il faille renfor-
cer la condition 3º, en exigeant que
l'on ait
$$\mathfrak{S}_0 \longrightarrow G(Q)$$ .
L'avenir le décidera !

\newpage

%% ===== Page 205 =====

250

Réflexion faite, je préfère poser les choses de
façon légèrement différente, en introduisant,
dans le LT-gr. profini $M$, un ~~[unclear]~~
distingué $\pi$ avec une L-structure (par
ex. une structure : toute pièce,
si par ex $\pi = \operatorname{Sl}(2, \mathbb{Z})$) définissant la
L T-structure de $M$, une discrétisation
$\pi_0$ de $\pi$ (dans un sens évident, pour les
L-structures) un sous-groupe discret
$\mathfrak{S}_0 \supset \pi_0$ normalisant $\pi_0$ (NB. dans mon petit
exemple [unclear] on a $\mathfrak{S}_0 = \pi_0$), tel que
[MARGIN: son adh.] $\mathfrak{S} = \widehat{\mathfrak{S}}_0$ dans $M$ soit distingué dans $M$.

On a donc un diag. de groupes

(32)
\begin{tikzcd}
& \tau \ar[d, no head] & \Gamma \ar[d, no head] \\
1 \ar[r, phantom, "\subset"] & \pi \ar[r, phantom, "\subset"] \ar[d, equal] & \mathfrak{S} \ar[r, phantom, "\subset"] \ar[d, equal] & M \\
& \widehat{\pi}_0 & \widehat{\mathfrak{S}}_0 \\
& \uparrow & \uparrow \\
1 \ar[r, phantom, "\subset"] & \pi_0 \ar[r, phantom, "\subset"] \ar[uu] & \mathfrak{S}_0 \ar[uu]
\end{tikzcd}

où la première ligne est formée de
groupes profinis, la deuxième de groupes
discrets. On notera que $\mathfrak{S}$ s'identifie : au
quotient du complété profini de $\mathfrak{S}_0$, mais

\newpage

%% ===== Page 206 =====

je n'exige pas pour le moment qu'il soit
isomorphe à ce complété lui-même. [MARGIN: Dans $= \tilde{\mathfrak{S}}_0 / \Pi_0$]
idem pour $\Pi$ et $\Pi_0$. Le groupe $\Pi_0$
joue le rôle d'un groupe de Teichmüller
"devant" discret, le groupe $\tau$ celui d'un complété
de gr. de Teichmüller, le groupe
$\Gamma = M / \tilde{\mathfrak{S}}$ joue le rôle d'un groupe
Galois profini. Je suis intéressé
surtout par deux situations extrêmes

a) $\tilde{\mathfrak{S}}_0 = \Pi_0$ d'où $\tilde{\mathfrak{S}} = \Pi$ - situation provenant
d'une [unclear: variété] lisse relative sur
un corps $k$ ext. de t. f. de $\mathbb{Q}$, ou
plus généralement d'un schéma [unclear: affine]
(sans multiplicité) [MARGIN: [unclear: K est int.] [crossed out]]
ent. de type fini sur $\mathbb{Q}$,
le gr. de fond. modéré [unclear]
relative à la dite multiplicité,

b) $\tilde{\mathfrak{S}}_0 / \Pi_0 \simeq$ groupe de Teichmüller,
ou d'indice fini dans un gr. de T.
[crossed out: de $\Pi_0$ dans ce cas est intérieur]
semblable à une situation "modulaire", où

\newpage

%% ===== Page 207 =====

751

$\Pi$ s'interprète comme le groupe fondamental
d'une multiplicité modulaire convenable,
$\tilde{C}$ comme sa partie "géométrique", $\Pi_0$ comme
sa partie "transcendante", et $\Gamma$ comme
$\Gamma_{\bar{Q}/Q}$ ...

On fait sur (32) l'hypothèse que l'
opération cyclotomique sur $M$ (opérant sur
$\Pi$ par autom. de L-structures) soit
triviale sur $\mathfrak{G}$, a fortiori sur $\Pi$ (donc

(33) $\Pi = \Pi^+$, $\mathfrak{G} = \mathfrak{G}^+$ )

i.e. qu'il se factorise
(34) $\Gamma \xrightarrow{\chi_\Gamma} \hat{\mathbb{Z}}^*$
dont l'image soit un sous-groupe ouvert.
Si $(C_{Q,a})$ est comme plus haut
(p. 748), on s'intéresse aux hom. (32)

$$M \to G(\mathcal{A}_Q)$$

satisfaisant les conditions 1) 2) 3°) 4°)
comme ci-dessus, mais où 3°) se formule
de façon un peu plus forte
(3° fort) (Compatibilité aux compactifications des
Teichmüller $\mathfrak{G}_0$) : On a une factorisation

\newpage

%% ===== Page 208 =====

tion $\mathbb{G}_0 \to G(\mathbb{Q})$, rendant commutatif le diagramme

(35)
\begin{tikzcd}
\mathbb{G}_0 \ar[r, hook] \ar[d] & \mathcal{M} \ar[d] \\
G(\mathbb{Q}) \ar[r, hook] & G(\mathbb{A}_{\mathbb{Q}})
\end{tikzcd}

Les structures de la forme (32) 
[MARGIN: Quitte par la suite à introduire des variantes sur ces structures]
méritent sans doute un nom, je vais les appeler - ne serait-ce que provisoirement - structures de Galois-Teichmüller.

Si $G$ est un groupe algébrique affine sur $\mathbb{Q}$, un hom $\mathcal{M} \to G(\mathbb{A}_{\mathbb{Q}})$ satisfaisant les conditions 1) 2) $3^{\circ \text{forte}}$) 4) est appelé un "hom de Lie" de la structure de Galois-Teichmüller.

Je vais désigner par la suite par (en abrégé $M$) $M_{\mathbb{Q}}$ le schéma en groupes sur $\mathbb{Q}$ affine.
Je désigne par $\mathbb{G}_{\mathbb{Q}}$ le sous-schéma en groupes de $M_{\mathbb{Q}}$ engendré par l'

\newpage

%% ===== Page 209 =====

752

image de $\mathfrak{S}_0 \to G(\mathbb{Q})$, par $\Pi_Q$
celui engendré par l'image de $\Pi \to G(\mathbb{Q})$

On a ainsi des sous-groupes

fermés

(36) $$ \Pi_Q \subset \mathfrak{S}_Q \subset M_Q $$

distingués dans $M_Q$, donnant lieu à

des schémas en groupes quotients

(37) $$ \mathbb{T}_Q \simeq \mathfrak{S}_Q / \Pi_Q \, , \, \Gamma_Q \simeq M_Q / \mathfrak{S}_Q $$
$$ N_Q \simeq M_Q / \Pi_Q $$

donnant lieu à quatre suites
exactes de schémas en groupes, dont la
dernière est

(38) $$ 1 \to \mathbb{T}_Q \to N_Q \to \Gamma_Q \to 1 \ . $$

Notons enfin le [unclear: machin] $\chi_c : G_Q \to \mathbb{G}_m$
est épi, et se factorise par $\Gamma_Q$

(39) $$ \Gamma_Q \xrightarrow{\chi_\Gamma} \mathbb{G}_{m Q} $$

\newpage

%% ===== Page 210 =====

Le système des groupes sur $Q$ (36, 137), muni

de l'hom de

$$ \underline{\pi} \subset \underline{\mathfrak{S}} \subset \underline{\mathcal{M}} $$

dans $\mathcal{G}$

(40) $$ \underline{\pi}_Q(A_Q) \subset \underline{\mathfrak{S}}_Q(A_Q) \subset \underline{\mathcal{M}}_Q(A_Q) $$

induit par $\underline{\mathcal{M}} \to G(A_Q)$ et satisfaisant

aux conditions 1) à 4) ci-dessus (i.e.

d'un "hom de Lie") s'appellera une

"approximation de Lie" du

système de Galois-Teichmüller envisagé.

Considérons un hom

(41) $$ T_\bullet \to \pi_\bullet $$

où la L-structure de $T_\bullet$
[unclear: induite] par composition avec $\pi_\bullet \to \underline{\Pi}(Q)$

[MARGIN: NB pour [unclear] l'étude des TE [unclear], [unclear] L-structure [unclear] composition [unclear] est un [unclear] meilleur choix [unclear] déduit par [unclear] la L-str. sur $T_\bullet$]

(42) $$ T_\bullet \to \underline{\Pi}(Q) . $$

((21) p. 747')

On voit que cet hom

est quasi-unipotent, et qu'il définit par

suite un hom de $Q$-schémas en groupes

\newpage

%% ===== Page 211 =====

(43) $\underline{k}_i : \underline{T}_0 \xrightarrow{753} \underline{\Pi}_{\mathbb{Q}}$
$(\simeq \mathbb{G}_{m\mathbb{Q}} \otimes_{\mathbb{Z}} T_0)$ (pas un [unclear] en groupes !)
qui définit, par passage aux pts $\mathbb{Q}$-rationnels, le genre d'hom

(42). On trouve ainsi un ens. fini
(indexé par l'ens des génér. de sous-groupes
locaux de $\pi_0$, - ou de $\Pi$ ce qui revient
au m.) d'hom extérieurs (43) de schémas
en groupes. J'ai envie d'appeler l'appro-
-ximation [unclear: crossed out word] de Lie dégénérée $L_i$
les hom précédents aux tores non
triviaux, de [unclear]

(44) $\underline{k}_i : \underline{T}_{\mathbb{Q}} \simeq \mathbb{G}_{m\mathbb{Q}} \otimes_{\mathbb{Z}} T_0 \hookrightarrow \underline{\Pi}_{\mathbb{Q}}.$

Supposons que $\underline{M}_{\mathbb{Q}}$ provienne
d'un schéma en groupes (affine de t.f.)
$\underline{M}$ sur $\mathbb{Z}$, tel que l'on ait

(45) $\underline{M} \longrightarrow \underline{M}(\hat{\mathbb{Z}}).$

(il me semble qu'on peut tjs trouver un
tel $\underline{M}$...). On désigne alors par

\newpage

%% ===== Page 212 =====

$\underline{\Pi}$, $\underline{\mathfrak{S}}$ les adhérences schématiques
dans $\underline{M}$ de $\underline{\Pi}_\mathbb{Q}$, $\underline{\mathfrak{S}}_\mathbb{Q}$, d'où un diagramme
de schémas en groupes sur $\mathbb{Z}$

(46) $$ \underbrace{\underline{\Pi} \subset \underline{\mathfrak{S}}}_{\underline{\Gamma}} \subset \underline{M} $$
[MARGIN: with another brace under $\underline{\mathfrak{S}} \subset \underline{M}$ labeled $\underline{N}$]

(sous réserve de représentabilité des
quotients $\underline{\tau}, \underline{\Gamma}, \underline{N}$ - mais je crois que
pour Raynaud il n'y a jamais de pb...).

[unclear: On a des factorisations] [crossed out]

(47)
\begin{tikzcd}
\underline{\Pi} \ar[r, hook] \ar[d] & \underline{\mathfrak{S}} \ar[r, hook] \ar[d] & \underline{M} \ar[d] \\
\underline{\Pi}(\hat{\mathbb{Z}}) \ar[r, hook] & \underline{\mathfrak{S}}(\hat{\mathbb{Z}}) \ar[r, hook] & \underline{M}(\hat{\mathbb{Z}})
\end{tikzcd}

On suppose bien sûr également que
l'hom $\chi_{\underline{M}} : \underline{M}_\mathbb{Q} \to \mathbb{G}_{m \mathbb{Q}}$ se
prolonge en

(48) $$ \chi_{\underline{M}} : \underline{M} \to \mathbb{G}_{m \mathbb{Z}} $$

lequel nécessairement se factorise par

(49) $$ \chi_{\underline{\Gamma}} : \underline{\Gamma} \to \mathbb{G}_{m \mathbb{Z}} . $$

\newpage

%% ===== Page 213 =====

754.

Il serait intéressant de dégager un th. d'existence
(voire d'unicité, moyennant des conditions plus fortes
...) d'un "modèle" lisse $\underline{M}$ sur $\mathbb{Z}$ de $\underline{M}_{\mathbb{Q}}$,
satisfaisant aux conditions précédentes
($\underline{M} \rightarrow \underline{M}(\mathbb{Z})$, et $\underline{M} \xrightarrow{\chi_{\underline{M}}} G_{m \mathbb{Z}}$).

Il est un peu plus vrai en général, je
présume, que les hom. extérieurs $k_i$ (43) se
prolongent tels quels en des hom.
[unclear: $\underline{T}_i \tilde{\times} \underline{T}_{\mathbb{Z}} \rightarrow \underline{\Pi}$]. Mais je vais considérer
le centralisateur de $k_i : T_0 \rightarrow \Pi_0$ dans $\Pi_0$,
soit $\tilde{L}_{i,0}$, qui est un groupe
métabélien : car ($T_0 \xrightarrow{k_i} \Pi_0$) y est un
sous-groupe distingué d'indice fini. Considérons
l'hom.

(50)
$$ \tilde{L}_{i,0} \longrightarrow \underline{\Pi}(\mathbb{Z}) , $$
$$ = Cent_{\Pi_0}(k_i : T_0 \rightarrow \Pi_0) $$

qui est ent. unipotent - il se factorise
donc par un hom. de schémas en groupes

(51)
$$ Env(\tilde{L}_{i,0}) \rightarrow \underline{\Pi} $$

(cf § 50, 4°, p. 658' ff), se factorisant

\newpage

%% ===== Page 214 =====

par un des [schémas en] groupes lisses quotients $\underline{Z}_{i,n}$ de
$Env(\tilde{L}_{i,0})$, d'une [unclear] de la [unclear] :

[MARGIN: $S_0 = Spec \mathbb{Z}$]
(52) $$1 \to \underline{uT_0} \to \underline{Z}_{i,n} \to (H_{i,n})_{S_0} \to 1$$
$$ \simeq (\tilde{L}_{i,0}/\underline{uT_0}) $$

Mais bien entendu, la restriction de l'hom.
obtenu
$$ \underline{Z}_{i,n} \to \underline{\Pi} $$
à
$$ \underline{Z}_{i,n}^0 \simeq \underline{uT_0} $$
n'est pas en général un mono - s'il
l'est pour un $n$ fixé, dans un $n'$ ($n|n'$, $n \neq n'$);
il ne l'est plus pour $n'$. Un cas parti-
culièrement favorable serait celui où
on peut trouver un $M$ pour lequel l'hom.
induit

[MARGIN: cet $u$ est unique]
(53) $$ \underline{uT_0} \to \underline{\Pi} $$

soit un mono. Ceci implique que l'adhérence
schématique dans $\underline{\Pi}$ du "sous-groupe"
additif" à 1 paramètre $\underline{\mathbb{G}}_a \simeq \underline{L}_{i,\mathbb{Q}} = \underline{k}_i(\underline{T}_0)$
est un sous-schéma en groupes lisse de $\underline{\Pi}$.

\newpage

%% ===== Page 215 =====

755

Cette condition doit être [unclear: d'ailleurs] suffisante - du
moins si on exige que cette add. schématique
ait une comp. neutre qui soit un isom. : $\mathbb{G}_{a S_0}$.
On trouve alors, un homomorphisme
unique (à conjugaison près)

(54) $$ \underline{k}'_! : \underline{T}_0 \longrightarrow \underline{\Pi} $$

tel que $\underline{k}'_{! \mathbb{Q}}$ soit de la forme $\underline{k}_{! \mathbb{Q}} \circ (r, id)$,
pour un $r \in \mathbb{Q}^{*+}$ convenable. Je ne suis
pas sûr que sous ces conditions, l'hom déduit de $\underline{k}'_!$ en
passant aux $\mathbb{Z}$-pts
$$ \underline{T}_0 \simeq \underline{T}_0(\mathbb{Z}) \longrightarrow \underline{\Pi}(\mathbb{Z}) $$
se factorise en un hom

(55) $$ \underline{T}_0 \xrightarrow{\underline{k}'_!} \Pi_0 \xrightarrow{can} \underline{\Pi}(\mathbb{Z}) , $$

où $T_0 \to \Pi_0$ [unclear] un "germe" $\underline{k}'_{!-}$.

Je vais néanmoins le poser comme hypo-
thèse suppl., lorsque les $\underline{k}'_!$ existent.

J'ai envie d'énoncer quelques autres
"bonnes hypothèses", qui s'efforcent de [unclear: relier]

\newpage

%% ===== Page 216 =====

en travaillant sur des "fondations" de
Lie "entières" i.e. sur $Spec \ \mathbb{Z}$.
Considérons tous les hom. $\underline{k}_i^!$ possibles
($i$ variable, mais aussi en conjuguant les
$k_i : \underline{T}_0 \to \Pi_0$ au départ par des éléments
quelconques de $\Pi_0$), et les hom.
correspondants sur les algèbres de Lie
(56) $$Lie(\underline{k}_i^!) : \underline{T}_0 \hookrightarrow \underline{Lie}(\underline{\Pi}) = Lie(\Pi)$$
qui sont par hyp. des monomorphismes
[MARGIN: cette hyp. semble raisonnable du fait que $\Pi_0 = \operatorname{Sl}(2)$, mais je n'y pige ... !]
universels, on voudrait que $Lie(\underline{\Pi})$ soit
engendré (comme $\mathbb{Z}$-module) par
les images, ou bien (ce qui revient au même = ) par
les images d'un générateur de $\underline{T}_0$. Je
reviens à dire que si on désigne de $Lie(\underline{\Pi})$
par $A_i$ l'image de $\underline{k}_i^!$, disons l'image d'un dit
générateur de $\underline{T}_0$, les éléments obtenus
dans $Lie(\underline{\Pi})$, et leurs conjugués
par les éléments de $\Pi_0$, engendrent le
$\mathbb{Z}$-module $Lie(\underline{\Pi})$. Ceci signifierait que

\newpage

%% ===== Page 217 =====

756

Le schéma en groupes $\underline{\Pi}^\circ$ sur $S_0 = Spec \mathbb{Z}$, vu en tant
que faisceau étale sur le "gros topos étale" de
$S_0$, est engendré par ses sous-faisceaux en
sous-groupes
$\underline{k}'_!(\underline{G}_\circ)$ (avec répétition possible d'une même
classe de conjugaison), et encore $\underline{\Pi}$ est
(Et il y a [unclear] pratiquement quant aux
les $\underline{k}'_!(\underline{G}_\circ)$ lors d'un changement de base) et
le groupe discret $(\Pi_0)_{S_0}$. (Et il se pose
la question naturellement d'expliciter les
relations... suffit-il d'écrire les relations de
$\Pi_0$, plus les commutativités de

(57)
\begin{tikzcd}
T_0 \ar[r, "k'_!"] \ar[d, "can"'] & \Pi_0 \ar[d, "can"] \\
\underline{T}_0 \ar[r, "\underline{k}'_!"'] & \underline{\Pi}
\end{tikzcd}
? )

Autres conditions concernant les schémas
en groupes

(58)
$$ \gamma_{\underline{\Pi}} = \underline{\Pi} / \underline{\Pi}^\circ, \gamma_{\underline{\mathfrak{S}}} = \underline{\mathfrak{S}} / \underline{\mathfrak{S}}^\circ, \gamma_{\underline{M}} = \underline{M} / \underline{M}^\circ $$

- on voudrait que ces schémas soient
en groupes constants, correspondant aux
schémas groupes finis discrets qui décrivent les
quotients correspondants des groupes sur $\mathbb{Q}$,
lesquels groupes discrets sont des quotients

\newpage

%% ===== Page 218 =====

(59) $\gamma_{\Pi}$ (resp. $\gamma_{\mathfrak{S}}$, resp. $\gamma_{\mathcal{M}}$) quotients de $\Pi, \mathfrak{S}, \mathcal{M}$
$\Big\{ \gamma_{\Pi} \simeq (\gamma_{\Pi})_{S_0}, \gamma_{\mathfrak{S}} \simeq (\gamma_{\mathfrak{S}})_{S_0}, \gamma_{\mathcal{M}} \simeq (\gamma_{\mathcal{M}})_{S_0}$
qui s'insèrent dans une suite d'homo.

(60) $$ \gamma_{\Pi} \to \gamma_{\mathfrak{S}} \to \gamma_{\mathcal{M}} $$
et des des suites exactes

(61)
$$ \gamma_{\Pi} \to \gamma_{\mathfrak{S}} \to \gamma_{\mathbb{Z}} \to 1 $$
$$ \gamma_{\Pi} \to \gamma_{\mathcal{M}} \to \gamma_{\Gamma} \to 1 $$
$$ \gamma_{\mathfrak{S}} \to \gamma_{\mathcal{M}} \to \gamma_{F} \to 1 $$
$$ \gamma_{\mathbb{Z}} \to \gamma_{\Gamma} \to \gamma_{F} \to 1 $$

Cette hypothèse sur $\Pi, \mathfrak{S}, \mathcal{M}$ permet de
définir des groupes $\underline{\Pi}^{\Delta}(\hat{\mathbb{Z}}), \underline{\mathfrak{S}}^{\Delta}(\hat{\mathbb{Z}}), \underline{\mathcal{M}}^{\Delta}(\hat{\mathbb{Z}})$
et on aura nécessairement les deux premières
factorisations de
(62)
$$ \Pi \to \underline{\Pi}^{\Delta}(\hat{\mathbb{Z}}) $$
$$ \mathfrak{S} \to \underline{\mathfrak{S}}^{\Delta}(\hat{\mathbb{Z}}) $$
$$ \mathcal{M} \to \underline{\mathcal{M}}^{\Delta}(\hat{\mathbb{Z}}) $$

(car $\Pi_0 \to \underline{\Pi}(\mathbb{Z}) \subset \underline{\Pi}^{\Delta}(\hat{\mathbb{Z}})$
$\mathfrak{S}_0 \to \underline{\mathfrak{S}}(\mathbb{Z}) \subset \underline{\mathfrak{S}}^{\Delta}(\hat{\mathbb{Z}})$ ) . Quant

à la troisième factorisation, elle devrait être
incluse dans les conditions de départ l'
un des [unclear: termes] de la donnée d'une structure de
Teichmüller-Galois -- savoir que
$\underline{A}_{\mathbb{Q}} / \underline{A}_{\mathbb{Q}}^\circ$ est constant, ce qui permet
[MARGIN: l'est surement]

\newpage

%% ===== Page 219 =====

757

alors de définir $\underline{M}^\Delta_{\mathbb{Q}}(A_{\mathbb{Q}})$, et on exige

que $M \to \underline{M}_{\mathbb{Q}}(A_{\mathbb{Q}})$ se factorise par

(63) $$M \to \underline{M}^\Delta_{\mathbb{Q}}(A_{\mathbb{Q}})$$ .

Si de plus on suppose que les images

des trois hom (62) sont des sous-groupes

ouverts, (condition qui renforce considé-

rablement la condition [unclear: divisée de 4)]

(p. 749), [unclear: pouvant gréer le conjuguer avec]

la surjectivité des composées

(64)
$$ \begin{cases}
\Pi \to \underline{\Pi}^\Delta(\hat{\mathbb{Z}}) \to \Gamma_{\Pi} \\
\mathbb{G} \to \underline{\mathbb{G}}^\Delta(\hat{\mathbb{Z}}) \to \Gamma_{\mathbb{G}} \\
M \to \underline{M}^\Delta(\hat{\mathbb{Z}}) \to \Gamma_M
\end{cases} $$
) ,

on sera particulièrement heureux. On

dira qu'on a une [unclear: M-approximation] de Lie

pour $\mathbb{Z}$, dite "entière") parfaite [unclear: ou définitive]

si toutes les conditions suppl envisagées

précédemment sont satisfaites...

\newpage

%% ===== Page 220 =====

Remarques Dans la présentation précédente, le
groupe $\Pi_0$ $(=\pi)$ lui-même semble passer au second
plan, au profit de $\mathcal{S}_0 \subset \mathcal{S}$ $(= \mathcal{S}_{\bar{\pi}})$ et de sa $LT$-structure.
Il me semble en effet que l'importance du
rôle de $\Pi_0$ s'efface derrière celle des $\bar{\pi}$-groupes
$\subset \Pi_0$ d'indice fini $\Pi'_0$ (s'il en existe) qui serait
engendré par les $K_i^!(\Gamma_0) \subset \Pi_0$ (et leurs
[unclear: ou gr.ps : $Out_I$ et les $K_i^!(\Gamma_0)$ est sous-g.]
conjugués, soit $\Pi''_0$, etc. [unclear: et de là à trouver]
les sous-groupes locaux exactement. Il
n'est pas trop clair, même si un tel $\Pi'_0$ existe,
qu'il soit néc.t unique (prenons le cas
p.ex. où $I = \emptyset$ !). Le cas qui m'avait
occupé jusqu'à présent, $\Pi_0 = \mathcal{S}_0 = \mathcal{S}^+_{0,3}$ --
$\widehat{\mathcal{S}}_{0,3}^+$, a été un peu "misleading", car
on a pris alors $\Pi_0 = \mathcal{S}_0$ [unclear: extra] justesse. Si on
s'était avisé un peu mieux de ce fait au
début de ce travail, où (prenant en outre $I=\emptyset$) $\Pi_0 = \Pi_{0,3}$ comme
le groupe de départ, on aurait vu
(comme cela a d'ailleurs été le cas effective-
ment) le rôle de $\Pi_{0,3}$ s'amenuiser au
profit des sous-groupes plus petits (type $\Pi_{\underline{g}}$
donnant naissance aux $\mathcal{S} \hookleftarrow \mathcal{S}_{\underline{g}}$ du $§ 49$ VI ...)

\newpage

%% ===== Page 221 =====

40) [unclear: rayé: approximation] $\mathcal{T}_{S,\mathbb{Q}}^{S'}(\Pi_0,\mathbb{G}_0,\underline{M})$ [unclear: rayé: hom. de Lie] induite par une approximation de $\mathbb{G}_0$.

Partons d'un système discret $(\mathbb{G}_0, \Pi_0)$ d'un LT-groupe $\mathbb{G}_0$, et d'un sous-groupe [MARGIN: on suppose $\Pi_0^*=\{1\}$, l'hom. de syst. de LT-groupes déduit de la structure LT $\Pi_0 \to \mathbb{G}_0$ (l'un des 2) soit l'inclusion...] invariant $\Pi_0$ de $\mathbb{G}_0$. Supposons donné en outre un groupe $\mathbb{G}_{\mathbb{Q}}$ sur $\mathbb{Q}$, et un hom. 
$$ (65) \quad \varphi : \mathbb{G}_0 \to \mathbb{G}(\mathbb{Q}) $$
ayant la propriété que pour tout hom
$$ (66) \quad k_i : T_0 \to \mathbb{G}_0 $$
qui soit partie de la LT-structure l'hom. composé
$$ T_0 \to \mathbb{G}_0 \to \mathbb{G}(\mathbb{Q}) $$
soit quasi-unipotent. On suppose de plus que $\mathbb{G}_{\mathbb{Q}}$ soit engendré par l'ens de ses pts $\mathbb{Q}$-rationnels définis par ces éléments de $\mathbb{G}_0$ via (65).

On considère $\mathbb{G}_{\mathbb{Q}}$ comme muni de l'homomorphisme (65) et de la famille des homomorphismes extérieurs
$$ (67) \quad K_i : T_0 \to \mathbb{G}_{\mathbb{Q}} $$

\newpage

%% ===== Page 222 =====

donnant lieu à des compatibilités

(68)
\begin{tikzcd}
T_0 \ar[r, "K_i"] \ar[d, "can"'] & \underline{\mathfrak{S}}_0 \ar[d, "(65)"] \\
T_{0\,Q}(Q) \ar[r, "\underline{K}_i(Q)"'] & \underline{\mathfrak{S}}(Q) \ .
\end{tikzcd}

Je considère alors le schéma en groupes [MARGIN: extérieurs (resp. à un...)] sur $Q$ des (automorphismes de $\underline{\mathfrak{S}}_Q$, formé des $\mu \in \underline{Aut}(\underline{\mathfrak{S}}_Q)$, tels que pour tt $i \in I$, il existe loc. une section $f_i$ de $\underline{\Pi}_Q$ (engendré par l'image de $\Pi_0$ dans $\underline{\mathfrak{S}}(Q)$), telle que l'on ait

(69) $$ \mu(\underline{K}_i(\lambda)) = int(f_i)^{-1}(\mu \lambda) \ . $$

Je désigne par $\underline{M}_Q^\varphi$ ce schéma en groupes sur $Q$, s'envoyant vers $\underline{G}_Q$ par $(u,\mu) \mapsto \mu$, par un hom. qui est trivial sur $\operatorname{Im}(\underline{\mathfrak{S}}_Q \to \underline{M}_Q^\varphi)$, d'où se factorise par un hom.

(70) $$ \underline{\Gamma}_Q^\varphi \overset{df}{=} \underline{M}_Q^\varphi / Im \underline{\mathfrak{S}}_Q \xrightarrow{\lambda^\varphi} \underline{G}_Q \ . $$

Si maintenant $(\underline{\mathfrak{S}}_0, \Pi_0)$ fait partie d'un système de [unclear]-Teichmüller

\newpage

%% ===== Page 223 =====

$(\Pi_0 \subset \mathbb{G}_0 \subset \mathcal{M})$, et si $\mathbb{G}_0 \xrightarrow{\varphi} \underline{\mathbb{G}}^\varphi_\mathbb{Q}$ provient d'un
hom de Lie de ce système, $\mathcal{M} \to M(\mathbb{A}_\mathbb{Q})$
par la méthode vue la section précédente,
alors on trouve un hom. naturel
via op. de $\underline{M}_\mathbb{Q}$ sur $\underline{\mathbb{G}}_\mathbb{Q}$, et $\lambda_M$
(71) $$ \underline{M}_\mathbb{Q} \xrightarrow{\varphi} \underline{M}^\varphi_\mathbb{Q} $$
rendant commutatif le diagramme

(72)
\begin{tikzcd}
\mathbb{G}_0 \ar[r] \ar[d] & \mathcal{M} \ar[d] & \underline{\mathbb{G}}_\mathbb{Q} \ar[r] \ar[dr, "via\ \text{can.}\ int."'] & \underline{M}_\mathbb{Q} \ar[d, "\varphi"] \\
\underline{\mathbb{G}}(\mathbb{Q}) \ar[r] \ar[d, "int."'] & M(\mathbb{A}_\mathbb{Q}) \ar[d] & & \underline{M}^\varphi_\mathbb{Q} \\
\underline{M}^\varphi_\mathbb{Q}(\mathbb{Q}) \ar[r, hook] & \underline{M}^\varphi_\mathbb{Q}(\mathbb{A}_\mathbb{Q}) & &
\end{tikzcd}

Si le centre de $\underline{\mathbb{G}}_\mathbb{Q}$ est trivial (ce qui est
une hyp. "crédible" surtout si le centre de
$\mathbb{G}_0$ l'est ...) lors $\underline{\mathbb{G}}_\mathbb{Q} \to \underline{M}^\varphi_\mathbb{Q}$ est
injectif, et $\underline{\mathbb{G}}_\mathbb{Q}$ s'interprète comme
"le $\underline{\mathbb{G}}_\mathbb{Q}$" "procréé" par un hom. de
Lie
$$ (\Pi \subset \mathbb{G} \subset \mathcal{M}) \longrightarrow (\Pi \subset \underline{\mathbb{G}}_\mathbb{Q} \subset \underline{M}^\varphi_\mathbb{Q}) , $$
qui a un certain caractère universel
p.r. à l'hom. donné (65) $\varphi : \mathbb{G}_0 \to \underline{\mathbb{G}}(\mathbb{Q})$.

\newpage

%% ===== Page 224 =====

Quand on ne suppose pas d'avance l'exis-
tence d'un tel hom. de Lie dans lequel
s'insère $\underline{\mathfrak{S}}_\mathbb{Q}$ , il n'est pas automatique
qu'on puisse trouver un hom.

(73) $$\varphi_M : \mathcal{M} \longrightarrow \underline{\mathcal{M}}_\mathbb{Q}^\varphi (A_\mathbb{Q})$$

qui rende compte des opérations de
$\mathcal{M}$ sur $\underline{\mathfrak{S}}_\mathbb{Q}$ , par la commutativité
d'un diagramme , pour tout $g \in \mathcal{M}$

(74)
\begin{tikzcd}
\underline{\mathfrak{S}} \ar[r, "int(g)|\underline{\mathfrak{S}}"] \ar[d, "\hat{\varphi}"'] & \underline{\mathfrak{S}} \ar[d, "\hat{\varphi}"] \\
\underline{\mathfrak{S}}(A_\mathbb{Q}) \ar[r, "int(\varphi_M(g))|\underline{\mathfrak{S}}(A_\mathbb{Q})"'] & \underline{\mathfrak{S}}(A_\mathbb{Q})
\end{tikzcd}

On s'intéresse seulement aux $\varphi$ (65)
pour lesquels (73) existe - au moins
quitte à se restreindre à un sous-groupe
ouvert de $\mathcal{M}$...

\newpage

%% ===== Page 225 =====

53

700

1°) Action différentiable du groupe de Teichmüller spécial.

Soit $X$ une surface connexe (orientable, compacte dans le cas qui nous intéresse) $I$ une partie discrète de $X$. Pour $i \in I$, soit $S_i$ l'espace (homéomorphe à [unclear]) des demi-droites [unclear: issues de l'origine] dans l'espace tangent, et soit $\tilde{S}_i$ un rev. universel de $S_i$.

Soit $g$ un autom. diff. de $(X,S)$,
il induit un autom. de
(1) $$S = S_I \overset{def}{=} \coprod_{i \in I} S_i$$
noté $u_g$. Considérons l'ens des relèvements $\tilde{u}_g$ de $u_g$ en un autom. de
(2) $$\tilde{S} = \coprod_{i \in I} \tilde{S}_i$$ .

Notons, si $T_0$ est le module d'orientation de $X$, que $T_0$ opère sur chaque $\tilde{S}_i$,

\newpage

%% ===== Page 226 =====

de sorte que $T_0$ opère en fait sur $\tilde{S}$ et

(3) $$ \tilde{S}/T_0 \simeq S . $$

Ceci montre que (si on est comme ci-
dessus) le groupe $T_0^I$ opère à g.
(p.ex) sur l'ens des relèvements $\tilde{u}$.
qui devient ainsi un torseur sous
$T_0^I$.

Soit $G$ un groupe topologique,
qui opère différentiablt sur $(X,S)$, avec
une action continue sur $S$. On

pose

(4) $$ \tilde{G} = \left\{ (u, \tilde{u}) \mid u \in G, \tilde{u} \text{ un relèvemt} \atop \text{continu de } u_S \text{ à } \tilde{S} \right\} $$

on trouve sur $\tilde{G}$ une topologie
[unclear]
canonique, qui soit telle que
l'hom. $\tilde{G} \to G$ fasse de
$\tilde{G}$, un
revêtement de $G$, de groupe $T_0^I$ :

(5) $$ 1 \to T_0^I \to \tilde{G} \to G \to 1 . $$

[MARGIN: NB $G$ opère sur $T_0^I$ via le caractère d'orientation]

Proposition Supposons que $X$ soit q. cpt,

\newpage

%% ===== Page 227 =====

761

et considérons l'hom. canonique
$$ G \longrightarrow G/G^\circ \xrightarrow{\sim} \mathcal{T}(X,S) \quad \text{(groupe de Teichmüller)} $$
on a alors un hom.
canonique d'extensions

[MARGIN: Pour la démonstration de la dite proposition voir [unclear] (p. 755 ff, p. 764 ff)]

(7)
\begin{tikzcd}
1 \ar[r] & T_0^I \ar[r] \ar[d] & \tilde{G} \ar[r] \ar[d] & G \ar[r] \ar[d] & 1 \\
1 \ar[r] & T_0^I \ar[r] & Sp\mathcal{T}(X,S) \ar[r] & \mathcal{T}(X,S) \ar[r] & 1 \ ,
\end{tikzcd}

$Sp\mathcal{T}(X,S)$ est le groupe de Teichmüller
spécial de $(X,S)$.

NB On notera que si pour tout $i \in I$ on
fixe une origine $e_i$ sur $S_i$, il
résulte de la propriété universelle de
$\tilde{S}$ que la donnée d'un relèvement
$\tilde{u}$ de $u \in G$ à $\tilde{S}$ revient à la donnée,
pour tout $i \in I$, d'un relèvement de
$u(e_i)$ en un point $\tilde{a}_i$ de $\tilde{S}_{\sigma(i)}$ (i.e., de $\tilde{S}_{u(i)}$).

\newpage

%% ===== Page 228 =====

2°) Le groupe $\operatorname{Gl}(2, \mathbb{R})^{\pm 1} \simeq$ [crossed out: $\operatorname{Gl}(2, \mathbb{Z})$]. On pose

(8) $\operatorname{Gl}(2, \mathbb{R})^{\pm 1} = \{ u \in \operatorname{Gl}(2, \mathbb{R}) \mid \det u \in \{\pm 1\} \}$

Ce groupe opère différentiablement sur
$\mathbb{C} \simeq \mathbb{R}^2$ en laissant fixe l'origine,
on peut donc appliquer la construction
des $\tilde{\mathbb{S}}_0$, $\tilde{T}_0$. $\mathbb{S}_0$ s'identifie au
cercle unité

(9) $\mathcal{U} \overset{\text{déf}}{=} \{ z \in \mathbb{C} \mid |z|=1 \} \longrightarrow \mathbb{S}_0$
$z \longmapsto \mathbb{R}^+ z$

et pour $u \in \operatorname{Gl}(2, \mathbb{R})^{\pm 1}$, son action
sur $\mathbb{S}_0$ s'identifie à l'action sur $\mathcal{U}$
donnée par

(10) $z \longmapsto \frac{u(z)}{|u(z)|}$
[unclear: $\in \mathcal{U}$]

Un relèvement de cette action en une
action sur $\tilde{\mathbb{S}}_0$ revient à : un relèvement
de $u(x_0=1)$ en un élément de $\tilde{\mathbb{S}}_0$.
Nous construirons [crossed out text] en choisissant
un pt. base $e_0$ de $\tilde{\mathbb{S}}_0$, au-dessus de $1$.

\newpage

%% ===== Page 229 =====

762

point $1$ de $\mathcal{U}$, de sorte que le rev. 
universel de $\mathbb{S}_0$ pointé au dessus de 
$e_0$ s'identifie au revêtement exponentiel

$$ (11) \qquad \mathbb{R} \longrightarrow \mathcal{U} \quad (\simeq \mathbb{S}_0) $$
$$ \qquad t \longmapsto \exp 2 i \pi t \quad (\mapsto \mathbb{R}^+ \exp 2 i \pi t) $$

Ainsi, un relèvement de $u_{\mathbb{S}_0}$ revient au 
choix d'un $t \in \mathbb{R}$ tel que

$$ (12) \qquad \exp 2 i \pi t = \frac{u(1)}{|u(1)|} $$

i.e. au choix d'un argument $t$ de $u(1)$.

On applique la construction de $\widetilde{G}$ du 1) au cas $G = \operatorname{Gl}(2, \mathbb{R})^{\pm 1}$, 
qui donne une structure d'extension

$$ (13) \qquad 1 \longrightarrow \mathbb{Z} \longrightarrow \operatorname{Gl}(2, \mathbb{R})^{\pm 1 \sim} \longrightarrow \operatorname{Gl}(2, \mathbb{R})^{\pm 1} \longrightarrow 1 $$

qui induit au dessus de la composante 
neutre $\operatorname{Sl}(2, \mathbb{R})$ de $\operatorname{Gl}(2, \mathbb{R})^{\pm 1}$ une 
structure d'extension

$$ (14) \qquad 1 \longrightarrow \mathbb{Z} \longrightarrow \operatorname{Sl}(2, \mathbb{R})^\sim \longrightarrow \operatorname{Sl}(2, \mathbb{R}) \longrightarrow 1 $$

Proposition Le groupe $\operatorname{Sl}(2, \mathbb{R})^\sim$ est connexe, 
et est [unclear: même] homéomorphe à $\mathbb{R}^3$.

\newpage

%% ===== Page 230 =====

Dém. Comme l'inclusion

(15) $$ \mathcal{U} \hookrightarrow \operatorname{Sl}(2, \mathbb{R}) $$

est une équivalence d'homotopie,
qui s'identifie topologiquement à l'inclusion

$$ \mathcal{U} \hookrightarrow \mathcal{U} \times \mathbb{R}^2 , $$

[unclear] que l'image inverse $\tilde{\mathcal{U}}$
de $\mathcal{U}$ dans $\widetilde{\operatorname{Sl}(2, \mathbb{R})}$ s'identifie au
rev. universel $\mathbb{R}$ de $\mathcal{U}$ (pointé
au dessus de $1$ par $0$ ...) Or si
$u \in \mathcal{U}$, le choix d'un $\tilde{u}$ au dessus
de $u$, i.e. d'un relèvement de $\frac{u(1)}{|u(1)|}$
dans $\mathbb{R}$ par la formule (12), s'identifie
(puisque $\frac{u(1)}{|u(1)|} = u(1) = u$) au choix
d'un $t \in \mathbb{R}$ tel que $\exp 2i\pi t = u$,

O.K.

Cor. 1 Le centre de $\widetilde{\operatorname{Gl}(2, \mathbb{R})}$ [unclear] est $\simeq \mathbb{Z}$ ($\exists 1$ générateur $w_0$ tel qu'on ait $w_0^2 = w'_0$, où $w'_0$ est le gén. de $\operatorname{Ker}(\widetilde{\operatorname{Sl}(2, \mathbb{R})} \to \operatorname{Sl}(2, \mathbb{R}))$.

Cor. 2 L'action d'un $g \in \operatorname{Gl}(2, \mathbb{R})^{-}$ sur $\widetilde{\operatorname{Sl}(2, \mathbb{R})}$ (ou sur le centre $w_0^\mathbb{Z}$ de $\widetilde{\operatorname{Sl}(2, \mathbb{R})}$) se lit : $w'_0 \mapsto {w'_0}^{-1}$ (resp. $w_0 \mapsto w_0^{-1}$) ...

[MARGIN: Expliquer un peu - le centre [unclear] de $\operatorname{Gl}(2, \mathbb{R})$ a 2 comp. conn. [unclear] un élément $w_0 \in \widetilde{\operatorname{Gl}(2, \mathbb{R})}$ au-dessus de $-1$ tel que $w_0^2 = w'_0$]

\newpage

%% ===== Page 231 =====

753

3) Le groupe $\operatorname{Gl}(2, \mathbb{Z})^\sim$.

Soit $\operatorname{Gl}(2, \mathbb{Z})^\sim$ l'image inverse de $\operatorname{Gl}(2, \mathbb{Z})$
dans $\operatorname{Gl}(2, \mathbb{R})^{\pm 1 \sim}$, de sorte qu'on a une
structure d'extension

(16) $$1 \longrightarrow \mathbb{Z} \longrightarrow \operatorname{Gl}(2, \mathbb{Z})^\sim \longrightarrow \operatorname{Gl}(2, \mathbb{Z}) \longrightarrow 1$$
[unclear: $n \mapsto w_0^n$]

Je voudrais décrire $\operatorname{Gl}(2, \mathbb{Z})^\sim$ par géné-
rateurs et relations. Pour ceci, je
note ceci :

Proposition. L'image inverse dans $\operatorname{Gl}(2, \mathbb{Z})^\sim$ des sous-groupes
cycliques $L_\rho \simeq \mathbb{Z}/6\mathbb{Z}$ et $L_\sigma \simeq \mathbb{Z}/4\mathbb{Z}$
engendrés par $\rho$ et $\sigma$ sont isomorphes
à $\mathbb{Z}$.

Ceci résulte du
Corollaire (à la prop. p. 762) : Pour tout sous-
groupe cyclique fini $Z \subset \operatorname{Sl}(2, \mathbb{R})$, l'image
inverse de $Z$ dans $\operatorname{Sl}(2, \mathbb{R})^\sim$ est
isomorphe à $\mathbb{Z}$.

En effet, p. "conjugaison" o.p.s. que
$Z \subset U$, i.e. $Z = \mu_n(\mathbb{C})$, mais alors
$\tilde{Z}$ s'identifie à l'ens des $t \in \mathbb{R}$ [unclear: où]
$(\exp 2 i \pi t)^n = 1$ i.e. $\exp 2 i \pi (nt) = 1$ i.e.

\newpage

%% ===== Page 232 =====

$$t \in \frac{1}{n} \mathbb{Z},$$

le groupe au dessus est $\simeq \mathbb{Z}$, cqfd.

[MARGIN: pour un choix de $\tilde{\rho}, \tilde{\sigma}$ plus symétrique (cf plus bas...)]
$\underline{Corollaire}$ Il existe $\tilde{\rho} \in \operatorname{Sl}(2,\mathbb{Z})^\sim$ et $\tilde{\sigma} \in \operatorname{Sl}(2,\mathbb{Z})^\sim$
au-dessus de $\rho$ et $\sigma$ respectivement, tels
qu'on ait
(17) $$\tilde{\rho}^3 = \tilde{\sigma}^2$$
L'ensemble de ces systèmes $(\tilde{\rho}, \tilde{\sigma})$ est un torseur
sous $\mathbb{Z}$, par l'opération
(18) $$(\tilde{\rho}, \tilde{\sigma}) \longmapsto (\tilde{\rho} \omega_0^{2n}, \tilde{\sigma} \omega_0^{3n}),$$
où $\omega_0$ est le générateur canonique de
$\ker( \operatorname{Sl}(2,\mathbb{Z})^\sim \to \operatorname{Sl}(2,\mathbb{Z}) )$.

[MARGIN: Il conviendrait mieux dire : l'image inverse du sous-groupe $Z$ du centre de $\operatorname{Sl}(2,\mathbb{Z})$ (eng. par $-1$) dans $\operatorname{Sl}(2,\mathbb{Z})^\sim$ est $\simeq \mathbb{Z}$ et on a alors la relation]
On peut dire aussi que l'image
inverse de l'élément central $-1$ de $\operatorname{Sl}(2,\mathbb{Z})$
[crossed out: est un groupe isomorphe à $\mathbb{Z}$, avec un unique]
avec un unique générateur $\omega_0$ tel que
(19) $$\omega_0^2 = z_0$$ (générateur can. de $\ker(\operatorname{Sl}(2,\mathbb{Z})^\sim \to \operatorname{Sl}(2,\mathbb{Z}))$)

et ceci étant, il est possible de
relever $\rho$ en $\tilde{\rho}$, $\sigma$ en $\tilde{\sigma}$ satisfaisant
(20) $$\tilde{\rho}^3 = \omega_0^{-1}, \tilde{\sigma}^2 = \omega_0$$
[MARGIN: Faire la vérification]
Et cette fois-ci, $\tilde{\rho}$ et $\tilde{\sigma}$ sont uniquement déter-

\newpage

%% ===== Page 233 =====

minés. Ceci résulte [crossed out: immédiatement] de la prop. précédente.

[MARGIN: NB $\tau = \tau_\infty = \begin{pmatrix} -1 & 0 \\ 0 & 1 \end{pmatrix}$ i.e. $\tau(x+iy) = -x+iy$ (symétrie p.r. axe des $y$)]

Considérons l'image inverse de $\{1, \tau\} \subset \operatorname{Gl}(2,\mathbb{Z})$ dans $\operatorname{Gl}(2,\mathbb{Z})^\sim$, extension de $\{1, \tau\} \simeq \mathbb{Z}/2\mathbb{Z}$ par $\mathbb{Z}$ - je dis que cette extension est scindée. En effet, l'action de $\tau$ sur $\mathcal{S}_0$ se relève en l'action de $\tilde{\tau}$ d'ordre 2 sur $\tilde{\mathcal{S}}_0 \simeq \mathbb{R}$, qui est donnée par

(21) $\tilde{\tau}t = \frac{1}{2} - t$

car on a pour $z \in \mathcal{H}$, $z = \exp 2i\pi t$

(22) $\tau z = -\bar{z} = -\exp(-2i\pi \bar{t}) = \exp(\pi i - 2i\pi \bar{t}) = \exp 2i\pi(\underbrace{\frac{1}{2}-t}_{\tilde{\tau}(t)})$

Il est peut-être plus commode de travailler par $\tau' = \tau'_0 = \begin{pmatrix} 0 & 1 \\ 1 & 0 \end{pmatrix}$, qui échange $e_0 = 1$ et $e_1 = i$, i.e. la symétrie par rapport à la diagonale principale. Cet automorphisme stabilise $\mathcal{H}$, et fixe le point $\exp \frac{\pi i}{4} = \frac{\sqrt{2}}{2}(1+i)$

il se relève en

(23) $\tilde{\tau}'(t) = \frac{1}{4} - t$.

Notons que [crossed out: toute] l'action de $\tilde{\tau}'$ (comme celle

\newpage

%% ===== Page 234 =====

de $\tilde{\tau}$) sur $\omega_0$ est la symétrie $w \mapsto w^{-1}$
(cf fin du $\S$ précédent). Je dis qu'on a :

(24) $\tilde{\tau}'(\tilde{\rho}) = \tilde{\rho}^{-1} \quad , \quad \tilde{\tau}'(\tilde{\sigma}) = \tilde{\sigma}^{-1}$

Ces formules s'écrivent aussi
$$\tilde{\tau}'(\tilde{\rho}^{-1}) = \tilde{\rho} \quad , \quad \tilde{\tau}'(\tilde{\sigma}^{-1}) = \tilde{\sigma}$$
elles sont valables dans $\operatorname{Gl}(2,\mathbb{Z})^\sim$, donc
pour voir qu'elles sont valables dans $\operatorname{Sl}(2,\mathbb{Z})^\sim$,
il suffit de vérifier que les premières
satisfont les relations caractéristiques
(20), ce qui s'écrit
$$\tilde{\tau}'( \underbrace{\tilde{\rho}^{-3}}_{\omega_0} ) = \omega_0^{-1} , \tilde{\tau}'( \underbrace{\tilde{\sigma}^{-2}}_{\omega_0^{-1}} ) = \omega_0$$
i.e. $\tilde{\tau}'(\omega_0) = \omega_0^{-1}$, OK.

Théorème Le groupe $\operatorname{Sl}(2,\mathbb{Z})^\sim$ admet
la présentation (à deux gén. $\tilde{\rho},\tilde{\sigma}$)

(25) $[ \tilde{\rho}^3 \cdot \tilde{\sigma}^{-2} = 1 ]$

Dém. Soit $\Pi$ le groupe défini
par cette présentation, et considérons
l'hom.

(26) $\Pi \to \operatorname{Sl}(2,\mathbb{Z})/\pm 1 = \{ \rho, \sigma | \rho^3 = \sigma^2 = 1 \}$
$(\tilde{\rho},\tilde{\sigma}) \mapsto (\rho,\sigma) \mod \pm 1)$

\newpage

%% ===== Page 235 =====

765

le groupe $\operatorname{Sl}(2, \mathbb{Z})/\pm 1$ s'identifie au
quotient de $\Pi$ par le sous-groupe distingué
engendré par l'élément $\omega_0 = \tilde{\rho}^3 = \tilde{\sigma}^2$.
Mais cet élément est au centre de $\Pi$, ce groupe
engendré par $\omega_0$ est [unclear: déjà] distingué.
donc on a une suite exacte
(27) $$ \mathbb{Z} \xrightarrow{\omega_0^n} \Pi \longrightarrow \operatorname{Sl}(2, \mathbb{Z})/\pm 1 \longrightarrow 1 $$
qui s'envoie dans la suite exacte
$$ 1 \longrightarrow \mathbb{Z} \longrightarrow \widetilde{\operatorname{Sl}(2, \mathbb{Z})} \longrightarrow \operatorname{Sl}(2, \mathbb{Z})/\pm 1 \longrightarrow 1 $$
avec $n \mapsto \tilde{\omega}_0^n$
en induisant des isomorphismes identiques sur les facteurs
extrêmes - donc par le lemme des 5
$\Pi \longrightarrow \widetilde{\operatorname{Sl}(2, \mathbb{Z})}$ est iso, cqfd.

Corollaire! Le groupe $\widetilde{\operatorname{Gl}(2, \mathbb{Z})}$ admet
la présentation (à trois gén.: $\tilde{\rho}, \tilde{\sigma}, \tilde{\tau}$)
(28) $[ \tilde{\rho}^3 = \tilde{\sigma}^2, \tilde{\tau}(\tilde{\rho}) = \tilde{\rho}^{-1}, \tilde{\tau}(\tilde{\sigma}) = \tilde{\sigma}^{-1} ]$

Je vais essayer de présenter le groupe
$\widetilde{\operatorname{Gl}(2, \mathbb{Z})}$ par trois générateurs involutifs,
$\tilde{v}_0, \tilde{\tau}, \tilde{\sigma}_0$ au-dessus des éléments
$v_0, \tau_0, \sigma_0$ de $\overline{\operatorname{Gl}(2, \mathbb{Z})}$, en m'inspirant
des formules de $\operatorname{Sl}(2, \mathbb{Z})$ :

\newpage

%% ===== Page 236 =====

$$
(29) \quad \begin{cases}
\Gamma_\infty = \tau_\infty = \begin{pmatrix} -1 & 0 \\ 0 & 1 \end{pmatrix} \\
\Gamma_0 = \tau'_\infty = \begin{pmatrix} 0 & 1 \\ 1 & 0 \end{pmatrix} \\
\Gamma_1 = -\tau'_0 = \begin{pmatrix} -1 & 1 \\ 0 & 1 \end{pmatrix}
\end{cases}
\qquad
(30) \quad \begin{cases}
\Gamma_0(\rho) = \rho^{-1}, \Gamma_0(\sigma) = \sigma^{-1} \\
\Gamma_1 = \rho^{-1} \Gamma_0 = \Gamma_0 \rho \\
\Gamma_\infty = \sigma \Gamma_0 = \Gamma_0 \sigma^{-1}
\end{cases}
$$

Notons que $(\tilde{\rho}, \tilde{\tau})$ engendre un groupe
isomorphe à $D_\infty$, que je note $D_{\tilde{\rho}}$,
et que $(\tilde{\sigma}, \tilde{\tau})$ engendre un groupe $D_{\tilde{\sigma}}$,
isomorphe à $D_\infty$. Enfin $(\omega_0, \tilde{\tau})$ enge-
ndre un groupe $D_{\omega_0}$, isomorphe encore
à $D_\infty$. On a un diagramme

$$
(31) \quad 
\begin{tikzcd}
& D_{\omega_0} \ar[dl, "i"'] \ar[dr, "j"] & \\
D_{\tilde{\rho}} & & D_{\tilde{\sigma}}
\end{tikzcd}
\qquad
\begin{array}{l}
\begin{cases} i(\omega_0) = \tilde{\rho}^3 \\ i(\tilde{\tau}) = \tilde{\tau} \end{cases} \\
\begin{cases} j(\omega_0) = \tilde{\sigma}^2 \\ j(\tilde{\tau}) = \tilde{\tau} \end{cases}
\end{array}
$$

et le cor. 1 s'écrit sous la forme

$$
(32) \quad \operatorname{Gl}(2, \mathbb{Z})^\sim \simeq D_{\tilde{\rho}} *_{D_{\omega_0}} D_{\tilde{\sigma}}
$$

(somme amalgamée de groupes).

Or on peut prendre dans $D_{\tilde{\rho}}$, $D_{\tilde{\sigma}}$ les
générateurs involutifs

$$
(33) \quad \begin{cases}
\tilde{\Gamma}_0 = \tilde{\tau} = \tilde{\tau}'_\infty \\
\tilde{\Gamma}_\rho = .......
\end{cases}
\text{et} \quad
\begin{matrix}
\tilde{\Gamma}_1 \overset{df}{=} \tilde{\rho}^{-1} \tilde{\Gamma}_0 = \tilde{\Gamma}_0 \tilde{\rho} \text{ dans } D_{\tilde{\rho}} \\
\text{et } \tilde{\Gamma}_\infty \overset{df}{=} \tilde{\sigma} \tilde{\Gamma}_0 = \tilde{\Gamma}_0 \tilde{\sigma}^{-1}
\end{matrix}
$$

\newpage

%% ===== Page 237 =====

755

et on trouve alors

Corollaire 2 $Cll(?, \mathbb{Z})^\sim$ admet une présentation
(à trois générateurs involutifs $\tilde{r}_0, \tilde{r}_1, \tilde{r}_\infty$) :

(34) $$ \tilde{r}_0^2 = \tilde{r}_1^2 = \tilde{r}_\infty^2 = 1 \quad , \quad (\underbrace{\tilde{r}_0 \tilde{r}_1}_{\tilde{\rho}})^3 = (\underbrace{\tilde{r}_\infty \tilde{r}_0}_{\tilde{\sigma}})^2 = 1 $$

\newpage

%% ===== Page 238 =====

4°) L'isomorphisme $\operatorname{Gl}(2, \mathbb{Z})^\sim \simeq$ [unclear: $\tilde{\Gamma}_{1,1}$]

Posons

(35) $X_1 = \mathbb{R}^2 / \mathbb{Z}^2 = (\mathbb{R}/\mathbb{Z})^2 \simeq \mathcal{U} \times \mathcal{U}$

(surface-type de genre $1$) , notons additive-
ment sa loi de composition, et soit

(36) $a_1 = 0 \in X_1$, $X_1 \setminus \{a_1\} = X_1 \setminus \{0\} = U_{1,1}$ ,

surface-type de type $(g,\nu) = (1,1)$.
Le groupe $\operatorname{Gl}(2, \mathbb{Z})$ s'identifie au groupe
des automorphismes de groupe topologi-
que $X_1$, et opère donc sur $X_1$ en laissant
fixe l'origine $0$, donc il opère
sur $U_{1,1}$ , d'où un hom.

(37) $\operatorname{Gl}(2, \mathbb{Z}) \xrightarrow{\sim} \Gamma_{1,1} = \Gamma(U_{1,1})$
groupe des classes d'isoto-
pie d'homéomorphismes
de $U_{1,1}$

dont on sait que c'est un isomorphis-
me. D'ailleurs, $G = \operatorname{Gl}(2, \mathbb{Z})$ opère effec-
tivement sur $X_1$, et l'action sur
l'espace tangent s'identifie à l'action
tautologique de $\operatorname{Gl}(2, \mathbb{Z}) \subset \operatorname{Gl}(2, \mathbb{R})$
sur $\mathbb{R}^2$. Donc le groupe noté $\tilde{G}$ dans

\newpage

%% ===== Page 239 =====

767

1°) $\tilde{G}$ est canoniquement isomorphe, en
tant qu'extension de $G$ par $T_0 \simeq \mathbb{Z}$ (NB.
$\mathbb{R}^2$ et son quotient $X_1$ sont canoniquement
orientés) au groupe $\operatorname{Gl}(2, \mathbb{Z})^\sim$ du 3°.
Ainsi on trouve [inserted: (cf. [unclear: th. 61])] le

Théorème. L'image inverse par l'isom.
(37) de l'extension $S \mathcal{T}_{1,1}$ de $\mathcal{T}_{1,1}$ par $\mathbb{Z}$
s'identifie canoniquement à $\operatorname{Gl}(2, \mathbb{Z})^\sim$,
dont il admet les présentations explicites
dans 3°

a) $\tilde{\rho}^3 \simeq \tilde{\sigma}^2$, $\tilde{\tau} \tilde{\rho} \tilde{\tau} = \tilde{\rho}^{-1}$, $\tilde{\tau} \tilde{\sigma} \tilde{\tau} = \tilde{\sigma}^{-1}$, $\tilde{\tau}^2 = 1$

b) $\tilde{\Gamma}_0^2 = \tilde{\Gamma}_1^2 = \tilde{\Gamma}_\infty^2 = 1$, $(\tilde{\Gamma}_0 \tilde{\Gamma}_1)^3 \simeq (\tilde{\Gamma}_\infty \tilde{\Gamma}_0)^2 \simeq 1$.

Interprétons maintenant $\mathcal{T}_{1,1}$ et $S \mathcal{T}_{1,1}$
en termes du groupe fondamental
de $\mathcal{U}_{1,1}$, que nous calculerons en
choisissant (judicieusement ?) un
pt base $a_{1,1}$ dans $\mathcal{U}_{1,1}$ :

(38) $\Pi_{1,1} = \pi_1(\mathcal{U}_{1,1}, a_{1,1})$,

(39) $\mathcal{T}_{1,1} \simeq Aut \, ext_{[unclear: loc]}(\Pi_{1,1})$

\newpage

%% ===== Page 240 =====

[MARGIN: générateur de l'inertie $l_0'$, au lieu de $l_0$ qui désigne déjà un élément bien determiné de $\operatorname{Gl}(2, \mathbb{Z})$ !]

Nous avons une présentation de $\Pi_{1,1}$ par

(40) $\Pi_{1,1} = \{ x, y, l_0' \mid [x,y] l_0' = 1 \}$

et désignant par $L$ le sous-groupe local de $\Pi_{1,1}$ engendré par $l_0'$, on aura

(41) $S\tau_{1,1} \simeq \operatorname{Norm}_{\operatorname{Aut}_{loc}(\Pi_{1,1})}(L)$.

Donc il faudrait expliciter l'action

$$S\tau_{1,1} \simeq Cl(2, \mathbb{Z})^\sim \to \operatorname{Aut}_{loc}(\Pi_{1,1}) ,$$

par l'action des générateurs $\tilde{\rho}, \tilde{\sigma}, \tilde{\tau}$
ou $\tilde{\Gamma}_0 = \tilde{\tau}, \tilde{\Gamma}_1, \tilde{\Gamma}_\infty$ - action qui doit
normaliser $L$, et de façon plus
précise, transformer $l_0'$ en $l_0'$ (pour $\tilde{\rho}, \tilde{\sigma}$)
ou en ${l_0'}^{-1}$ (pour $\tilde{\tau} = \tilde{\Gamma}_0, \tilde{\Gamma}_1, \tilde{\Gamma}_\infty$). Il faut
donc expliciter l'action des générateurs
sur $x$ et sur $y$. (NB $\Pi_{1,1}$ est le groupe
libre ayant comme générateurs $x$ et $y$)

Considérons le revêtement principal $U_{1,1}^\sim$
de groupe $\mathbb{Z}^2$ de $U_{1,1} = X_1 - \{0\}$, induit
par le revêtement principal $\mathbb{R}^2$ de $X_1 = \mathbb{R}^2/\mathbb{Z}^2$,
- l'image inverse de $U_{1,1}$ est $\mathbb{R}^2 \setminus \mathbb{Z}^2$.
Choisissons un pt base [unclear: $\bullet$]

\newpage

%% ===== Page 241 =====

[MARGIN: 268 crossed out]

$a \in \mathbb{R}^2 \setminus \mathbb{Z}^2$,
au dessus de $a_{1,1}$ , on a donc un iso. can.

(42) $$ \Pi_{1,1} \simeq \Pi_1(\mathbb{R}^2 \setminus \mathbb{Z}^2, \mathbb{Z}^2; a) $$

d'où une suite exacte

(43)
\begin{tikzcd}
1 \ar[r] & \Pi_1(\mathbb{R}^2 \setminus \mathbb{Z}^2) \ar[r] & \Pi_{1,1} \ar[r] \ar[d, equal] & \mathbb{Z}^2 \ar[r] \ar[d, "\simeq"] & 1 \\
& & \Pi_1(U_{1,1}=X_1 \setminus \{a_{1,1}\}) & \Pi_1(X_1) \ar[d, "\simeq"] \\
& & & (\Pi_{1,1})_{ab}
\end{tikzcd}

Les éléments du $\Pi_1$ mixte s'explicitent par le calcul de Jean Malgoire à l'aide de couples

(44) $$ (\gamma, l) \quad (\gamma \in \mathbb{Z}^2, l \text{ classe de chemin de } \mathbb{R}^2 \setminus \mathbb{Z}^2 \text{ de } a \text{ à } a+\gamma) $$

qu'on "compose comme les fonctions".

Nous prendrons, au dessus des générateurs $e_0, e_1$ de $\mathbb{Z}^2$, les éléments $x, y$ définis par les chemins rectilignes de $a$ à $a+e_0$, resp. $a+e_1$, à partir de l'origine

(45) $$ a = \tilde{a}_{1,1} = (+\frac{1}{2}, +\frac{1}{2}) $$

Alors l'élément $xyx^{-1}y^{-1}$ est figuré par le chemin ci-dessus, qui est visiblement opposé au lacet "rectiligne" direct de $a$ autour de l'origine $0$.

\newpage

%% ===== Page 242 =====

$5^\circ$) Digression sur les revêtements de germes d'espaces.

[MARGIN: A) Revêtements de germes. ici on s'en tient pour simplifier aux filtres induits [unclear] ($F \in \Phi$) formant syst fond d'un germe d'esp.]

Un "germe d'espace" est défini par un esp. topologique $S$ et un filtre $\Phi$ sur $S$, un morphisme de $(S,\Phi)$ ds $(S',\Phi')$ est un élément de

(46) $\varinjlim_{S_i \in \Phi} \operatorname{Hom}_{\Phi_i, \Phi'} (S_i, S')$ ,

où, pour $S_i \in \Phi$, $\Phi_i$ est le filtre induit par $\Phi$ sur $S_i$, et $\operatorname{Hom}_{\Phi_i, \Phi'}$ désigne l'ensemble des appl. continues de $S_i$ dans $S'$ compatibles avec $\Phi_i, \Phi'$, i.e. telles que $\forall S'_j \in \Phi'$, $f^{-1}(S'_j) \in \Phi_i$. La composition de morphismes de germes est évidente.

Un revêtement (étale) d'un germe $(S,\Phi)$ d'espace top. est un objet de la catégorie

(47) $\underline{Rev}(S,\Phi) = \varinjlim_{S_i \in \Phi} \underline{Rev}(S_i)$ ;

ces revêtements forment une catégorie, sans doute (au moins dans des "bons cas" ?) que nous pourrions interpréter comme la catégorie des revêtements d'un topos convenable (savoir le topos limite

\newpage

%% ===== Page 243 =====

[MARGIN: cela en fait par Olivier Leroy ?]
projective des (petits) topos étales des $S_i \in \Phi$ ).

Cette catégorie, "dans les bons cas", a tendance
à être multigaloisienne, voire galoisienne.

On dit qu'un germe d'espaces est
1-connexe, ou simplement connexe, si
le foncteur naturel

(48)
$$(Ens) \longrightarrow \underline{Rev}(S, \Phi)$$
$$E \longmapsto E_{(S, \Phi)} \text{ , rev. constant de fibre } E$$

est une équivalence de catégories. Ceci
signifie que

1°) ce foncteur est pl. fidèle (on dit
que $(S, \Phi)$ est $0$-connexe), ce qui signifie,
si $S$ est loc. connexe, que $\forall S_i \in \Phi$,
[MARGIN: dans ce cas on dit que le germe d'espaces est connexe] $\exists S_j \in \Phi$, $S_j \subset S_i$, tel que $\pi_0(S_j) \to \pi_0(S_i)$
soit une application constante.

2°) le foncteur est de plus ess. surjectif,
ce qui signifie que pour tout $S_i \in \Phi$,
$\exists S_j \in \Phi$, $S_j \subset S_i$, tel que pour toute composante
connexe $S_{j,\alpha}$ de $S_j$, l'hom. nat.
$\pi_1(S_{j,\alpha}) \to \pi_1(S_i)$ soit trivial.

Les revêtements de germes d'espaces
forment une catégorie fibrée sur celle

\newpage

%% ===== Page 244 =====

des germes d'espaces. La catégorie des espaces top. peut être considérée comme sous-catégorie pleine de celle des germes d'espaces topologiques, en associant à un espace $S$, le germe $(S, \Phi)$, où $\Phi$ est le filtre trivial sur $S$. Ceci étant, la cat. fibrée des revêtements de germes d'espaces induit celle des rev. d'espaces topologiques.

Soit $\mathcal{X}$ un germe d'espaces tel que $\underline{Rev} \mathcal{X}$ soit une catégorie multigaloisienne, soit $S$ un germe d'espace 1-connexe, $S \to \mathcal{X}$ un morphisme.

Alors -- 

(49) $$ \underline{Rev}(\mathcal{X}) \to \underline{Rev}(S) \simeq (Ens) $$

foncteur fibre sur $\underline{Rev}(\mathcal{X})$, ce qui permet de définir

(50) $$ \pi_1(\mathcal{X}, S) $$

comme le $\pi_1$ relatif à ce foncteur fibre. Si $\mathcal{X}$ est 0-connexe et $\underline{Rev}(\mathcal{X})$ est galoisienne, on aura donc, si $\pi \simeq \pi_1(\mathcal{X}, S)$ une équivalence de catégories

\newpage

%% ===== Page 245 =====

(51) 
$$ \underline{Rev}(\mathfrak{X}) \simeq \operatorname{Ens}(\Pi) . $$

Le groupe $\Pi_1(\mathfrak{X}, S)$ dépend fonctoriellement
des triples $(S, \mathfrak{X}, S \to \mathfrak{X})$ i.e. de
la flèche $S \to \mathfrak{X}$ dans la catégorie des
germes d'espaces, où $S$ la source est
$1$-connexe.

Soit $X$ un espace top., $Y$ une partie de
$X$. On appelle germe de $X$ autour de $Y$
le germe d'espaces $(X, \Phi)$, où $\Phi$ est le
filtre des voisinages de $Y$ dans $X$, germe
[MARGIN: On a un foncteur image inverse $Rev(X) \leftarrow Rev(\mathfrak{X})$ [unclear]]
résiduel de $X$ le long de $Y$ le germe

(52)
d'espaces $(X \setminus Y, \Phi_{X \setminus Y})$ induit sur $X \setminus Y$
(il est vide si $Y = X$). Nous
utiliserons cette notion surtout dans le cas
où $Y = \{x\}$, on parle ainsi de germe
et germe résiduel de $X$ en $x$. (NB le
germe de $X$ en $x$ est toujours $1$-connexe,
mais pas le germe résiduel). Les
germes d'espaces purs s'interprètent comme
germes résiduels d'un espace
$\bar{X} = S \cup \{\omega\}$ en un pt $\omega$ convenable...

\newpage

%% ===== Page 246 =====

Je m'intéresse au cas où $x \in X$ admet un
voisinage homéomorphe à $\mathbb{R}^2$ - i.e. $X$ est une
surface en $x$. Je désignerai par [crossed out: $D_{X,x}^*$]
[MARGIN: (53)] $D_{X,x}^*$ ou $D_x^*$
le germe résiduel de $X$ en $x$ (contenu
dans le germe $D_{X,x}$ de $X$ en $x$).
Je m'intéresse au choix d'un rev. universel
$\tilde{D}_x^* \to D_x^*$ - mais pour ça il me
faut faire une digression sur la notion
[unclear: bâtarde] de rev. universel d'un germe
d'espace.

Soit $(S, \Phi)$ un germe d'espace. Un
rev. de $(S, \Phi)$ provient d'un rev.
$S'_i$ d'un $S_i \in \Phi$, et définit un
germe d'espace $(S'_i, \Phi_{S'_i})$, où $\Phi_{S'_i}$ désigne
l'image inverse [crossed out: de $\Phi$ dans $S'_i$.] au dessus de $S = (S, \Phi)$, Ce
germe d'espace (à isomorphisme près
de germe) ne dépend pas des choix
de $S_i, S'_i$, mais seulement de
l'élément de $Rev(S)$. On trouve
ainsi un foncteur contravariant

\newpage

%% ===== Page 247 =====

771

$$ (54) \quad Rev(S) \to (Germesp)/S $$

qui est pleinement fidèle. Un germe d'espace $S'$ au dessus de $S$ est dit un revêtement étale de $S$, s'il est dans l'image essentielle. Ainsi la cat. $Rev(S)$ des revêtements étales du germe d'espace $S$ est équivalente à la sous-catégorie pleine de la cat. des germes d'espaces au-dessus du germe $S$, formée des rev. étales dudit germe $S$.

Supposons [unclear: pour fixer] que la cat. $Rev(S)$ est multigaloisienne. Les objets 1-connexes sont appelés des rev. universels de $S$ - [barré: si $Rev(S)$ est galoisienne]. Ils forment un groupoïde, connexe ssi $Rev(S)$ est galoisienne i.e. ssi le germe $S$ est connexe. Les rev. universels de $S$ peuvent s'interpréter comme des germes d'espaces. Soit $\tilde{S}$ un tel rev. universel, alors $\tilde{S}$ est bel et bien un germe 1-connexe, et $\tilde{S} \to S$ fournit la [unclear] $\Pi_1(S, \tilde{S})$. Cela posé, on a un iso. canon.
$$ \Pi_1(S, \tilde{S}) \simeq \operatorname{Aut}_{Rev(S)}(\tilde{S}) $$

[MARGIN: le groupoïde galoisien $\Pi_1(S, \tilde{S})$ n'est autre que le $\Pi_1$ de la cat. galoisienne "abstraite" $Rev(S)$, calculé en l'objet ind-1-connexe...]

\newpage

%% ===== Page 248 =====

[MARGIN: B) Cas des surfaces [unclear: ponctuées]]

Ceci dit, revenons à la situation de (53),
d'un espace topologique $X$, et d'un
$x \in X$ tel que $X$ soit une surface en $x$.
Alors la catégorie
$$ \underline{Rev} (D^*_x) $$
est une catégorie galoisienne, dont le
groupe fondamental s'identifie canoniquement
au module d'orientation $T_x$ de $X$ en $x$.

(56) $$ \pi_1(D^*_x) \overset{df}{=} \pi_1(\underline{Rev}(D^*_x)) \simeq T_x $$

Ceci provient du fait que le germe
$D^*_x$ [unclear: est formé d']
espaces $D^*_{x,i}$ [unclear: homéomorphes] à des disques
centrés en $x$, et les $\pi_1$ des dits
disques forment un syst. projectif
constant ...

Ainsi les rev. universels $\tilde{D}^*_x$ des
$D^*_x$ forment un groupoïde connexe,
de groupe fondamental $T_x$.
Choisissons en un tel $\tilde{D}^*_x$, et
considérons le morphisme composé

\newpage

%% ===== Page 249 =====

772

$$D_x^{*\sim} \to D_x^* \hookrightarrow X - \{x\}$$

- plus généralement, soit $S$ une partie de $X$, contenant $x$ comme pt isolé (p. ex. une partie discrète de $X$), et [unclear: posons]

(57) $$D_x^* \hookrightarrow U \overset{def}{=} X \setminus S$$

et considérons le composé

(58) $$D_x^{*\sim} \to D_x^* \hookrightarrow U$$ .

Si $\underline{Rev}(U)$ est multigaloisienne (p. ex. $U$ loc. connexe), on peut donc considérer le groupe fondamental correspondant de $U$

(59) $$\pi_1(U, D_x^{*\sim})$$ ,

je vais définir un hom. canonique

(60) $$K_x : T_x \longrightarrow \pi_1(U, D_x^{*\sim}) ,$$
$$ \simeq Aut \text{ du foncteur fibre } \underline{Rev}(U) \to \underline{Rev}(D_x^{*\sim}) \text{ sur } \underline{Rev}(U)$$

on peut le décrire p. ex. par

(61) $$\underset{\begin{subarray}{c} \simeq \\ T_x \end{subarray}}{\pi_1(D_x^*, D_x^{*\sim})} \longrightarrow \pi_1(U, D_x^{*\sim})$$

\newpage

%% ===== Page 250 =====

induit par l'inclusion
$$D_x^* \hookrightarrow \mathcal{U}$$
ou encore, pour en termes d'automorph. d'un
foncteur fibre, comme l'hom. opération de
$T_x$ sur $\Gamma(D_x^{*\sim}, \mathcal{U}' \times_\mathcal{U} D_x^{*\sim})$ provenant

\begin{tikzcd}
& \mathcal{U}'_{D_x^{*\sim}} \ar[dl] \ar[d] \\
D_x^{*\sim} \ar[dr] & \mathcal{U}' \ar[d] \\
& \mathcal{U}
\end{tikzcd}

de l'opération de $T_x$ sur
$D_x^{*\sim}$. [unclear: text crossed out]

Supposons qu'on ait déjà
[MARGIN: $\pi_1(\mathcal{U},\tilde{\mathcal{U}}) = \pi$]
choisi (par ailleurs) un [unclear: foncteur fibre sur] $\mathcal{U}$,
par la donnée d' un pt base, ..., ou plus
généralement, d'un rev. universel $\tilde{\mathcal{U}}$ de $\mathcal{U}$

Pour le comparer avec $\pi_1(\mathcal{U}, D_x^{*\sim})$,
- il suffira de donner une "classe de chemins"
de $D_x^{*\sim} / \mathcal{U}$ à $\tilde{\mathcal{U}} / \mathcal{U}$ , i.e. un
isomorphisme

(62) $$ \tilde{\mathcal{U}} \xrightarrow{\sim} [unclear: Rev](\mathcal{U}, D_x^{*\sim}) $$

[MARGIN: l'image dans $\mathcal{U}$ du [unclear] de (62) donne dans $\mathcal{U}$ un point $y$ : [unclear]]

[unclear: crossed out paragraph at the bottom]
NB si $D_x^{*\sim}$ défini
par un rev. universel $\widetilde{D_x^*}$
de $D_x^*$, et si $\tilde{\mathcal{U}}$
défini par un pt base $y \in \mathcal{U}$,
alors [unclear]

\newpage

%% ===== Page 251 =====

773

( l'indétermination est un autom.
int. de $\Pi_1$, ou de $\Pi_1(\mathcal{U}, D_x^{* \sim})$ ) qui
induit un isomorphisme

(63) $$ \pi = \Pi_1(\mathcal{U}, \tilde{\mathcal{U}}) \overset{\sim}{\longleftrightarrow} \Pi_1(\mathcal{U}, D_x^{* \sim}) $$

(qui est changé par un autom. intérieur,
quand on change (62) par un él. de $\pi$).
Composant (63) et (60), on trouve donc
un hom.

(64) $$ T_n \longrightarrow \pi = \Pi_1(\mathcal{U}, \tilde{\mathcal{U}}) , $$

qui dépend du choix de (62) — et
qui, indépendamment de ce choix, définit
un hom. ext. bien déterminé —
qui n'est autre que l'hom. can.

(65) 
\begin{tikzcd}
\Pi_1(D_x^*) \ar[r] \ar[d, "\simeq"'] & \Pi_1(\mathcal{U}) \ar[d, equal] \\
T_n & \pi
\end{tikzcd}

[MARGIN: C) Relation avec
germes d'arcs (approche
originale "par pts base")]
Je vais décrire une façon commode
de fixer un rev. universel $\tilde{D}_x^*$ de $D_x^*$,
par un germe d'arc de courbe [unclear: analyt.]

\newpage

%% ===== Page 252 =====

de $x$. De façon générale, soit $I = [0,1]$, et
$[I]$ le germe résiduel de $I$ en $0$. C'est
évidemment un germe 1-connexe.

[MARGIN: Plutôt que de prendre un [unclear] pour des germes [unclear] un voisinage de $0$ dans $[0,1]$, il serait plus [unclear] pour la suite de définir un [unclear] de germes de courbes dans $X$ issues de $x$, satisfaisant à [unclear]]

Soit $X$ un espace top. et $x \in X$, un
germe d'arc de courbe dans
$X$ issu de $x$, est un morphisme

(66) $$[I] \xrightarrow{\gamma} germe \ res_x (X) \overset{def}{=} D_x^*$$

du germe $[I]$ dans le germe résiduel $D_x^*$ de
$X$ en $x$. Un tel germe est donc décrit
par une application continue d'un
intervalle $[0, \varepsilon]$ ( $0 < \varepsilon \le 1$ )

(67) $$[0, \varepsilon] \xrightarrow{\varphi} X$$

telle que l'on ait

(68) $$\begin{cases} \varphi(0) = x \\ \varphi(t) \neq x \text{ si } t \in ]0, \varepsilon] \end{cases}$$

Dans le cas qui nous intéresse ($X$
une surface ou voisinage de $x$, donc
$Rev(D_x^*)$ galoisien abélien de groupe $T_n$),
on peut utiliser cet arc pour définir
un foncteur fibre sur $D_x^*$

(69) $$Rev(D_x^*) \to Rev([I]) \simeq (Ens)$$

\newpage

%% ===== Page 253 =====

774

d'où un rev. universel
(70) $$D_x^{*\sim} \overset{\text{déf}}{=} Rev_{univ}(D_x^*, [I], \varphi) \ .$$

Bien entendu, on aura une factorisation

\begin{tikzcd}
{[I]} \ar[r] \ar[rr, bend right, "\varphi"'] & D_x^{*\sim} \ar[r] & D_x^* \ (\hookrightarrow U)
\end{tikzcd}

donc un iso can

(71) $$Rev_{univ}(U, D_x^{*\sim}) \simeq Rev_{univ}(U, [I]) \ ,$$

où $I$ est considéré comme un germe
1-connexe sur $U$ par le morphisme composé
$$[I] \longrightarrow D_x^* \hookrightarrow U \ .$$

Donc . . .

(72) $$\pi_1(U, D_x^{*\sim}) = \pi_1(U, [I]) \overset{K_x}{\longleftarrow} T_x \ .$$

[MARGIN: D) Relation avec structure différentiable]
Lorsque $X$ est muni d'une structure
différentiable au voisi-nage de $x$, on dit
qu'un germe de courbe de $X$ sortant
de $x$ est différentiable, si l'application
continue
$$I \longrightarrow X$$
qui le définit est différentiable, et de plus

\newpage

%% ===== Page 254 =====

l'application tangente à l'origine : [unclear]

$$(73) \quad \mathbb{R} \longrightarrow Tang_x(X) \quad ,$$

est injective. Alors l'image de $\mathbb{R}^+$ est
une demi-droite homogène dans l'espace
tangent à $X$ en $x$, qu'on appelle la
demi-tangente du germe d'arc de courbe
envisagé en $x$.

Revenant pour un moment au cas
purement topologique, nous allons considérer
l'ensemble $Couronne(X,x)$ des germes d'arcs
dans $X$ issus de $x$, comme l'un des
objets d'un groupoïde (le groupoïde
fondamental local en $X$ en $x$), en prenant
comme morphismes de $\gamma$ sur $\gamma'$
les isomorphismes de $D_{x,\gamma}^{*\sim}$ sur $D_{x,\gamma'}^{*\sim}$

$$(74) \quad Isom(\gamma, \gamma') \overset{def}{=} Isom_{D_x^*}(D_{x,\gamma}^{*\sim}, D_{x,\gamma'}^{*\sim})$$

On trouve donc par construction une équivalence de
catégories (de groupoïdes connexes)

$$(75) \quad Couronne(X,x) \longrightarrow Revunies(D_x^*) \quad ,$$

\newpage

%% ===== Page 255 =====

Soit d'autre part, dans le cas où $X$ est diff.
en $x$,

(76) $$ \underline{Couronne}_{diff}(X,-) $$

la sous-catégorie pleine de $\underline{Couronne}(X,x)$
formée des germes d'arcs différentiables. Soit
$S_x$
l'espace des demi-droites de
l'espace tangent $Tang_x(X)$

(77) $$ S_x = \text{espace des demi-droites de } Tang_x(X) \simeq \mathbb{S}^1 $$

et considérons le groupoïde fondamental

(78) $$ \Pi_1(S_x) \overset{\simeq}{\longrightarrow} \underline{Revunis}(S_x) $$ .

Je vais définir une équivalence de
catégories

(79)
\begin{tikzcd}
\underline{Couronne}_{diff}(X,x) \ar[r, "\simeq"] \ar[d, "\simeq"'] & \Pi_1(S_x) \ar[d, "\simeq"] \\
\underline{Revunis}(D_x^{*\sim}) & \underline{Revunis}(S_x)
\end{tikzcd}

qui, sur les ensembles d'objets, soit
l'application associant à tout germe d'arc
différentiable en $x$ sa demi-tangente. Il
s'agit donc, si $\gamma, \gamma'$ sont deux tels
arcs, correspondants aux demi-tangentes $\Delta, \Delta'$
de définir une bijection canonique

\newpage

%% ===== Page 256 =====

(80) 
\begin{tikzcd}
Isom(\gamma, \gamma') \ar[r, "\sim"] \ar[d, equal, "\text{déf}"'] & Isom_{\Pi_1 \mathbb{S}_n}(\Delta, \Delta') \ar[d, equal, "\text{déf}"] \\
Isom(D^{*\sim}_{n,\gamma}, D^{*\sim}_{n,\gamma'}) & \text{classes de chemins dans } \mathbb{S}_n \text{ de } \Delta \text{ à } \Delta'
\end{tikzcd}

Au lieu de définir directement la bijection
(80), je vais définir plutôt un
foncteur (qui sera une équivalence de
catégories)

(81) $$ \Phi : \Pi_1(\mathbb{S}_n) \xrightarrow{(\approx)} Revouv(D^*_n) \ , $$

Soit

(82) $A_n$ ens des parties fermées (1-connexes) de $\mathbb{S}_n$
($i.e.$ des $\underline{arcs}$ de cercle de $\mathbb{S}_n$,
mais $\neq \emptyset$,)

c'est un ensemble ordonné, $\underline{donc}$
lequel définit une catégorie ordonnée $\underline{A}_n$
(les objets sont les $a \in A_n$, les flèches
sont les inclusions) -- et on sait que
le groupoïde enveloppant $Gr(\underline{A}_n)$ de cette catégorie
ordonnée est canoniquement équivalent à $\Pi_1(\mathbb{S}_n)$,

\newpage

%% ===== Page 257 =====

776

de façon précise, soit $\underline{Gr}_0(\underline{A}_n)$ la s.-catégorie pleine de $\underline{Gr}(\underline{A}_n)$ (dont l'ens. des objets est $\underline{A}_n$) dont les objets sont les [unclear: s.-ens.] ponctuels, & on a une équivalence canonique

(83) $$\underline{Gr}_0(\underline{A}_n) \overset{\sim}{\longrightarrow} \Pi_1(\$_n)$$

qui sur les objets se réduit à : 
$\{x\} \mapsto x$, et qui, pour tout arc $A$ d'extrémités $x,y$, d'où un isomorphisme dans $\underline{Gr}_0(\underline{A}_n)$

(84) $$ \left\{ \begin{matrix} x \overset{\sim}{\longrightarrow} A \overset{\sim}{\longleftarrow} y \\ \text{d'où } x \overset{\sim}{\longrightarrow} y \text{ dans } \underline{Gr}_0(\underline{A}_n) \end{matrix} \right. $$

associe : à cette flèche $x \overset{\sim}{\longrightarrow} y$ de $\underline{Gr}_0(\underline{A}_n)$ la flèche de $\Pi_1(\$_n)$, définie par l'arc $A$ considéré comme chemin de $x$ à $y$.

Ainsi, pour un groupoïde $\mathcal{G}$, le don d'un foncteur

(85) $$\Psi : \Pi_1(\$_n) \longrightarrow \mathcal{G}$$

équivaut à la donnée d'un foncteur

\newpage

%% ===== Page 258 =====

(86) $$ \Psi : \underline{A}_n \longrightarrow \mathcal{G} $$

et (85) est une équivalence de cat. ssi

$\mathcal{G}$ est un groupoïde connexe et le groupe $\pi_1$ est [MARGIN: cf 6,b) [unclear]]
$\simeq \mathbb{Z}$, et si, pour un couple de
pts $x, y$ distincts de $\mathbb{S}_n$, et pour
deux arcs $A$ et $B$ d'extrémités $x, y$,
déterminant l'automorphisme de
$Gr_{\bullet}(\underline{A}_n)$ [unclear: induit] dans $Gr(\underline{A}_n)$, composé de

(86 bis) $$ \{x\} \overset{\sim}{\longleftarrow} A \overset{\sim}{\longrightarrow} \{y\} \overset{\sim}{\longleftarrow} B \overset{\sim}{\longrightarrow} \{x\} $$

l'automorphisme correspondant de $\Psi(\{x\})$ dans $\mathcal{G}$
soit un *générateur* de $\operatorname{Aut}_{\mathcal{G}}(\Psi(\{x\}))$.

Nous allons donc définir un
foncteur

(87) $$ \underline{A}_n \longrightarrow \underline{Rev}_{sur} (D_n^*) $$

(satisfaisant la condition supplémentaire
énoncée). Pour définir un tel foncteur,
il suffit de définir un foncteur

(88) $$ \underline{A}_n \longrightarrow \text{catégorie des germes } 1\text{-connexes} $$
$$ \text{d'espaces au dessus des} $$
$$ \text{germes } D_n^* $$

\newpage

%% ===== Page 259 =====

777

en composant un foncteur avec le
foncteur

(89) $$ \text{cat. des germes 1-ann.} \atop \text{d'espaces au dessus} \atop \text{du germe } D^* $$ $\longrightarrow \underline{Rev}_{un} (D^*)$

En fait, (88) associe à $A \in \underline{A}_x$
filtre sur $X \setminus x$ tendant vers $x$ , [unclear]
[unclear] germes de sous-espaces de $X$ [unclear],
- [unclear] le germe résiduels
Un $A \in \underline{A}_x$ s'identifie évidemment
à une famille $C$ de $T_x(X)$, [unclear]
[unclear] cônes [crossed out text]
distincts de $\{0\}$ et de [unclear] entiers.
Soit $C_A$ l'un de ces cônes.
Considérons un isom. local

(90) $$ f : (T_x(X), 0) \xrightarrow{\sim} (X, x), $$ dont l'appl.
[crossed out text] l'identité. L'image de $C$
[unclear] $\in C A_x$, son image par un
tel isom. dépend du choix de l'iso. local
ainsi qu'. Mais considérons la base de filtre
des [unclear] $C \subset C A_x$ strictes
[unclear] correspondant à des [unclear] $A'$ qui soient des
voisinages de $A$, on aura tel que $C' \setminus \{0\}$
soit un voisinage de $C \setminus \{0\}$ , dans le

\newpage

%% ===== Page 260 =====

l'image [unclear: par] (90) [crossed out: de la] du filtre défini par le filtre
sur $T_{op,x}(X)$ formé des $V \cap (C \setminus \{0\}) : V$
parcourant les voisinages de $0$ dans $T_{op,x}(X)$,
est un filtre sur $X \setminus \{x\}$ qui ne dépend
plus du choix de $f$, et qui
tend vers $x$, donc définit un germe
$\tilde{A} \sim \tilde{C}$ d'espaces au-dessus de $D_x^*$. Ce germe
d'espace, [crossed out] germe d'espace
défini sur $T_{op,x}(X)$ muni du filtre
de base les $V \cap (C' \setminus \{0\})$ (où l'on peut
prendre des $V$ [unclear: ouverts formés de rayons y]
[unclear], des $V \cap (C' \setminus \{0\})$ [unclear] --
[unclear]. Pour $A \in \mathcal{C}$
variable i.e. pour $A_1 \subset A_2$ i.e. $C_1 \subset C_2$
[MARGIN: $\tilde{C}_1 \hookrightarrow \tilde{C}_2$]
la relation correspondante d'inclusion
sur les germes d'espaces sur $D_x^*$. On
trouve donc bien (88), et de (81), il se
vérifie que ce dernier a bien une
équivalence de catégories (ou groupoïdes, vu la vérification
par explicitation plus haut ((86 bis)).

\newpage

%% ===== Page 261 =====

778

Considérons maintenant, à nouveau, un
arc différentiable $\gamma$ dans $X$ issu de $x$,
et sa demi-tangente $\Delta$. Alors $\{\Delta\} \in A_x$,
et on peut considérer le germe d'espace
$\{\Delta\}^\sim$ au dessus de $D_x^*$. On a des
morphismes d'inclusion canoniques
(91) $$\gamma \hookrightarrow \{\Delta\}^\sim \hookrightarrow D_x^*$$
et on en déduit un isom. canonique
(92) $$Rev_{unis}(D_x^*, \gamma) \simeq Rev_{unis}(D_x^*, \{\Delta\}^\sim)$$
$$|| \text{déf} \quad \quad \quad \quad || \text{déf}$$
$$D_{x, \gamma}^{* \sim} \quad \quad \quad \quad \Theta(\Delta)$$

Si $\gamma'$ est un autre arc diff. en $x$, avec
pour [unclear: dir.] : la demi-tangente $\Delta'$, on a donc
$$D_{x, \gamma}^{* \sim} \simeq \Theta(\Delta), \quad D_{x, \gamma'}^{* \sim} \simeq \Theta(\Delta')$$

d'où
$$Isom(D_{x, \gamma}^{* \sim}, D_{x, \gamma'}^{* \sim}) \simeq Isom(\Theta(\Delta), \Theta(\Delta'))$$
$$|| \text{déf} \quad \quad \quad \quad \uparrow \Theta$$
$$Isom(\gamma, \gamma') \quad \quad Isom_{\Pi_1(S_x)}(\Delta, \Delta')$$

ce qui est la [unclear] bijection (80) cherchée.

Soit. En résumé, on a un diagramme
[unclear: can. comm.] d'équivalences de groupoïdes

\newpage

%% ===== Page 262 =====

(93)
\begin{tikzcd}
germ\, revdiff(X,x) \ar[r, "\simeq"', "\gamma \mapsto \Delta(\gamma)"] \ar[d, "\gamma \mapsto Revunis(D_x^*,\gamma) = D_{x,\gamma}^{*\sim}"'] & \Pi_1(S_x) \ar[d, "\simeq", "a \mapsto Revunis(S_x,-)"'] \\
Rev.unis(D_x^*) \ar[r, "\simeq"'] \ar[ur, dashed, "\Theta \simeq"'] & Revunis(S_x)
\end{tikzcd}

[MARGIN: Bien sûr, la difficulté en fait provenait de ce qu'on n'avait point d'interpr. géométrique explicite entre [unclear] l'espace $S_x$ ...]

Le seul foncteur qui n'ait pas une
explication est le foncteur
(94)
\begin{tikzcd}
Rev.unis.(D_x^*) \ar[r, "\simeq"] \ar[dr, "\Theta^{-1}"'] & Revunis(S_x) \\
& \Pi_1(S_x) \ar[u, "can"']
\end{tikzcd}

défini comme composé de $\Theta^{-1}$ avec
le foncteur can. $\Pi_1(S_x) \to Revunis(S_x)$.

Finalement, dans ces considérations
de géométrie topologico-différentielle, la
construction-clef est celle du foncteur
$\Theta$ (81). Nous avons pris quelques peines
à la définir de façon invariante. Dans
le cas où on dispose d'une iso locale
(95) $$f \colon (T_x(X),0) \to (X,x),$$ la
construction de $\Theta$ peut s'expliciter comme
une construction canonique dans un
espace vectoriel $V$. Toute demi-droite [unclear].

\newpage

%% ===== Page 263 =====

779

$\Delta$ dans $V$ définit un germe d'arc (diff.)
dans $V$ retiré de $0$, d'où un rev.
universel correspondant $D_{V,0;\Delta}^{* \sim}$ de $D_{V,0}^*$
donc un rev. universel $D_{x,\Delta}^{* \sim}$ de $D_x^*$.
Il faut voir comment ce $D_{x,\Delta}^{* \sim}$ dépend
"continûment" de $\Delta$, et ce qui
revient au même, le rev. universel
$V^{* \sim}_\Delta$ de $V^* = V \setminus \{0\}$ dépend continu-
ment de $\Delta$, ce qui est tout-à-fait
standard...

Soit maintenant $u$ un
[crossed out: automorphisme] difféomorphisme local de $X$, fixant $x$,
et différentiable en $x$. On a donc
des automorphismes induits
(95) $^u D_x^*$ autom. de $D_x^*$, $^u \$_{x}$ autom. de $\$_x$.

Considérons d'autre part un rev. universel
$\$_x^\sim$ de $\$_x$, et soit $D_x^{* \sim}$ un rev.
universel de $D_x^*$ associé, grâce à l'équiv.
(94). Je dis qu'on a alors un

\newpage

%% ===== Page 264 =====

bijection canonique

(95)

entre des relèvements de $^u\!D_x^*$ en un autom. $^u\!\tilde{D}_x^*$ du rev. univ. $\tilde{D}_x^*$
$\simeq$
avec des relèvements de $^u\!S_x$ en un autom. $^u\!\tilde{S}_x$ de $\tilde{S}_x$

bijection qui soit un isomorphisme de $T_n$-torseurs (où $T_n$ est le $\mathbb{Z}$-module d'orientation de $X$ en $n$, vu plus haut).

[MARGIN: la construction se fait dans un passage suivant, aboutissant à (122) p. 773]

[MARGIN: E) Amorce d'une corrélation entre rev. universels de $S_x$ et $D_x^*$]

J'ai envie de donner une description "géométrique" de (94) et (95).

[MARGIN: $D_x=$ germe "ouvert" ou "épointé" du filtre des voisinages de $x$ dans $X$]
1) Soit $(X,x)$ un germe de surface différentiable, d'où $V_{X,x} = V_x = V = Tang_x(X)$, et $S_x = V^* / \mathbb{R}^{*+}$ (où $V^* = V \setminus \{0\}$), ainsi que $D_{X,x}^* = D_x^*$. On désigne par $\tilde{S}_x$ (resp. $\tilde{D}_x^*$) des rev. univ. de $S_x$ (resp. $D_x^*$). Soit $\mathcal{A}_{X,x}$ ou $\mathcal{A}_x$ l'ens. des arcs de $S_x$, en corresp. biunivoque avec l'ens. $C_{X,x} = C_x$ des "cônes fermés" réguliers propres de $V$ ($\neq V$ et $\{0\}$). Pour $A \in \mathcal{A}_x$ (resp. $C \in C_x$), je désigne par $E_x(A)$ ou $E_x(C)$ l'ouvert correspondant de $D_x^*$. Si $\tilde{S}$ est un rev. univ. de $S$, $\tilde{D}_x^*$ un rev. universel de $D_x^*$, on considère les

[MARGIN: d'où un "arc" de $A$ (resp. de $C$) dans $\tilde{D}_x^*$]

\newpage

%% ===== Page 265 =====

(97) 
\begin{tikzcd}
\Gamma(A, \tilde{S}_x/S_x) \ar[d, equal] & , & \Gamma(E_p(A), \tilde{D}_x^*/D_x^*) \ar[d, equal] \\
\operatorname{Hom}_{S_x}(A, \tilde{S}_x) \ar[d, equal, "\text{déf}"'] & & \operatorname{Hom}_{D_x^*}(E_p(A), \tilde{D}_x^*) \ar[d, equal, "\text{déf}"'] \\
\tilde{S}_x(A) & & \tilde{D}_x^*(E_p(A))
\end{tikzcd}

qui sont l'un et l'autre des Torseurs sur le groupe $T_{X,x} = T_x$ le $\mathbb{Z}$-module des orientations de $X$ en $x$ (pour $\sim$). Considérons l'ens

(98) $$H(A) = Isom_{T_x}(\tilde{S}_x(A), \tilde{D}_x^*(E_p(A)))$$

qui est lui-même un $T_x$-torseur. Soit en effet une inclusion

$$A \subset A'$$

on en déduit, par restriction, des isomorph. de $T_x$-torseurs

(99) $$\tilde{S}_x(A') \simeq \tilde{S}_x(A), \quad \tilde{D}_x^*(E_p(A')) \simeq \tilde{D}_x^*(E_p(A))$$

d'où un isom. canonique

(100) $$H(A) \simeq H(A')$$

de sorte que 

(101) $$A \mapsto H(A), \text{ foncteur } A_x \xrightarrow{H} \underline{Tors}(T_x)$$

où $A_x$ est la catégorie ordonnée [unclear] associée à l'ens. ordonné $A_x$. On trouve donc par factorisation un foncteur de la catégorie groupoïde associé $Gr(A_x)$ de $A_x$ dans $\underline{Tors}(T_x)$,

(102) 
\begin{tikzcd}
A_x \ar[d] \ar[r, "H"] & \underline{Tors}(T_x) \\
Gr(A_x) \ar[ru, "Gr(H)"'] & 
\end{tikzcd}

\newpage

%% ===== Page 266 =====

[unclear: Je dis que]
On sait que $Gr(\underline{A}_x)$ est [unclear] un $T_x$-
groupoïde, de façon précise le foncteur

$$ (103) \quad \left\{ \begin{array}{rcl} \underline{A}_x & \longrightarrow & Rev_{univ}(\hat{S}_x) \overset{def}{=} \hat{\mathbb{S}}_x^\sim \\ A & \longmapsto & Rev_{univ}(\hat{S}, A) \end{array} \right. $$

qui se factorise en un foncteur de
groupoïdes

$$ (104) \quad Gr(\underline{A}_x) \overset{(\simeq)}{\longrightarrow} Rev_{univ}(\hat{S}_x) $$

[MARGIN: i.e. : pour tous $A, A' \in \operatorname{Ob}(\underline{A}_x)$ l'app. nat. $\operatorname{Hom}_{Gr(\underline{A}_x)}(A, A') \overset{\sim}{\longrightarrow} Isom_{Rev_{univ}(\hat{S}_x)}(\tilde{S}_{x,A}^\sim, \tilde{S}_{x,A'}^\sim)$]

donne une équivalence de cat. (104), d'où

$$ (105) \quad \operatorname{Hom}_{Gr(\underline{A}_x)}(A, A') \simeq \operatorname{Hom}_{[unclear]}(\tilde{S}_{x,A}^\sim, \tilde{S}_{x,A'}^\sim) $$

Ceci rappelé, je dis que le foncteur
$Gr(H)$ de (102) est essentiellement constant,
i.e. il définit sur les $\pi_1$ l'hom. nul $T_x \to T_\alpha$
- ce qui revient à dire, il y a entre
les $H(A)$ (en composant les isom (100)) un système transitif d'isom. -

On désignera donc par la "valeur"
commune de ces $H(A)$, qui ne dépend
que des choix de $\tilde{S}_x$ et $\tilde{D}_\alpha^\sim$, par
$H(\tilde{S}_x^\sim, \tilde{D}_\alpha^\sim)$. Cet ens. est muni de la
structure définie par des isomo. de $T_\alpha$-torseurs

$$ (106) \quad i_A: H(A) \overset{\sim}{\longrightarrow} H(\tilde{S}_x^\sim, \tilde{D}_\alpha^\sim) $$

satisfaisant à la condition de compatibilité, s'expri-

\newpage

%% ===== Page 267 =====

781

-ment par la commutation des diagrammes

(107)
\begin{tikzcd}
H(A) \ar[r, "i_A", "\sim"'] \ar[d, "(100)"'] & H(\tilde{S}_\alpha, \tilde{D}_\alpha^*) \\
H(A') \ar[ur, "\sim"', "i_{A'}"] &
\end{tikzcd}

Pour vérifier le fait que le foncteur $Gr(H)$
de (102) est [unclear: bien consistant], je considère les
deux foncteurs (dépendant du choix de
$\tilde{S}_\alpha, \tilde{D}_\alpha^*$)

(108)
$$ \begin{cases} 
A \longmapsto \tilde{S}_{\alpha,A} \quad , \quad A \longmapsto \tilde{D}_{\alpha,A}^* \\ 
\underline{A}_\alpha \longrightarrow \underline{Rev}_{univ}(S_\alpha) \quad \underline{A}_\alpha \longrightarrow \underline{Rev}_{univ}(D_\alpha^*) 
\end{cases} $$

qui se factorisent en des équivalences

(109)
$$ Gr(\underline{A}_\alpha) \xrightarrow{\sim} \underline{Rev}_{univ}(S_\alpha) \quad , \quad Gr(\underline{A}_\alpha) \xrightarrow{\sim} \underline{Rev}_{univ}(D_\alpha^*) $$
utilisant $\tilde{S}_\alpha, \tilde{D}_\alpha^*$.

Ceci dit, pour un $A \in Ob \, \underline{A}_\alpha = A_\alpha$ [unclear: donné],
des isom. fonctoriels en $A$

(110)
$$ \begin{cases} 
[unclear](A) \simeq \underline{Isom}_{\underline{Rev}_{univ}(S_\alpha)}(\underbrace{F(A)}_{\tilde{S}_{\alpha,A}}, \tilde{S}) \\ 
[unclear](E,A) \simeq \underline{Isom}_{\underline{Rev}_{univ}(D_\alpha^*)}(\underbrace{G(A)}_{\tilde{D}_{\alpha,A}^*}, \tilde{D}_\alpha^*) 
\end{cases} $$

et par la propriété univ. de $Gr(\underline{A}_\alpha)$,
ces isom. s'étendent sur la
catégorie $Gr(\underline{A}_\alpha)$. Dans la situation où
elle se réduit à un groupe profini $T (= T_\alpha)$
trois $T$-groupoïdes connexes $g, g', g''$

\newpage

%% ===== Page 268 =====

(savoir $\underline{Gr}(\underline{A}_n), \underline{Revuni}(\$_n), \underline{Revuni}(\mathbb{D}^*_n)$),
des équivalences de $T$-groupoïdes
(111) $$g \overset{F}{\simeq} g', \quad g \overset{G}{\simeq} g"$$
et des objets $X' \in \operatorname{Ob}\ g', X" \in \operatorname{Ob}\ g"$ (savoir $\tilde{\$}$ et $\tilde{\mathbb{D}}^*_n$), et où on considère les
foncteurs
(112) $$ \begin{cases} g \xrightarrow{f=f_{X'}} \underline{Tors}(T) \\ X \mapsto \operatorname{Hom}_{g'}(F(X), X') \end{cases} \quad , \quad \begin{cases} g \xrightarrow{g=g_{X"}} \underline{Tors}(T) \\ X \mapsto \operatorname{Hom}_{g"}(G(X), X") \end{cases} $$
et le foncteur correspondant
(113) $$ \begin{cases} g \xrightarrow{} \underline{Tors}(T) \\ X \mapsto \operatorname{Hom}_{\underline{Tors}(T)}(f(X), g(X)) \end{cases} $$
à prouver que ce dernier foncteur est
est constant. Or on a, si $F_1, F_2$ désignent
(114) $$ F_1 : g' \overset{\simeq}{\to} g \quad , \quad F_2 : g" \overset{\simeq}{\to} g $$
désignant des quasi-inverses de $F, G$, des
isom. fonctoriels en $X$
(115) $$ \begin{cases} f(X) \overset{\text{def}}{=} \operatorname{Hom}_{g'}(F(X), X') \simeq \operatorname{Hom}_g(X, F_1(X')) \\ g(X) \overset{\text{def}}{=} \operatorname{Hom}_{g"}(G(X), X") \simeq \operatorname{Hom}_g(X, F_2(X")) \end{cases} $$
d'où aussitôt un isom. de $T$-torseurs
(116) $$ \underline{Isom}(f(X), g(X)) \simeq \underline{Isom}_g(F_1(X'), F_2(X")) $$
d'où l'assertion.

\newpage

%% ===== Page 269 =====

7P2

Définition. — On dit qu'on a une corrélation entre un rev. univ. $\tilde{S}_x$ de $S_x$, et un rev. univ. $\tilde{D}_x^*$ de $D_x^*$, quand on s'est donné un élément de $H(\tilde{S}_x, \tilde{D}_x^*)$.

Donc, si $A \in A_a$, la donnée d'une corrélation revient à la donnée d'un élément de $H(A)$, i.e. d'un isomorphisme entre les $T_a$-torseurs $\tilde{S}_x^\sim(A)$ et $\tilde{D}_x^{*\sim}(E_p(A))$. Cette identification est compatible avec les flèches d'inclusion entre $A \subset A'$. Les corrélations entre $\tilde{S}_x$ et $\tilde{D}_x^*$ forment un torseur sous $T_a$. Pour $\tilde{S}_x$ fixé

Soit $Corr_x(X) = Cour_x(X)$ la catégorie des triplets $(\tilde{S}_x, \tilde{D}_x^*, h)$ de rev. univ. corrélés de $S_x$ et de $D_x^*$, on a des foncteurs $(\tilde{S}_x, \tilde{D}_x^*, h) \rightrightarrows {\tilde{S}_x \atop \tilde{D}_x^*}$

(117)
\begin{tikzcd}
& Corr_x(X) \ar[dl, "\simeq"'] \ar[dr, "\simeq"] & \\
Revuniv(S_x) & & Revuniv(D_x^*)
\end{tikzcd}

qui sont des équivalences de catégories.

\newpage

%% ===== Page 270 =====

Toutes ces constructions ont été faites de façon
intrinsèque, [unclear: et même] soigneusement pour passer
à un iso. de germes de surfaces diffi-
-rentiables $(X,x) \xrightarrow{\sim} (X',x')$. Supposons
donnés en plus des $(\hat{S}_x^\sim, \hat{D}_x^{*\sim})$ relatifs à $(X,x)$,
ainsi $(\hat{S}_{x'}^\sim, \hat{D}_{x'}^{*\sim})$ relatifs à $(X',x')$.
Cet iso. de germes étant donné, il induit
donc

(118) $$ \begin{cases} u_S : \hat{S}_x \xrightarrow{\sim} \hat{S}_{x'} & \text{iso. d'esp. top.} \\ u_{D^*} : \hat{D}_x^* \xrightarrow{\sim} \hat{D}_{x'}^* & \text{iso. de germes d'esp. top.} \\ u_T : T_x \xrightarrow{\sim} T_{x'} & \text{iso de } \mathbb{Z}\text{-modules} \\ & \text{libres de rg } 1 \end{cases} $$

Considérons alors les
deux ensembles

(119) $$ \begin{cases} F(u) = \text{ens des iso } \hat{S}_x^\sim \xrightarrow{\sim} \hat{S}_{x'}^\sim \text{ au-dessus de } u_S \\ G(u) = \text{ens des iso } \hat{D}_x^{*\sim} \xrightarrow{\sim} \hat{D}_{x'}^{*\sim} \text{ au-dessus de } u_{D^*} \end{cases} $$

ce sont l'un et l'autre des $T_x$-torseurs.
Ceci étant, je dis qu'on a un iso. can.
de $T_x$-torseurs

(120) $$ F(u) \xrightarrow{\sim} G(u) $$

i.e. "il revient au même" de relever $u_S$ en
un $\tilde{u}_S : \hat{S}_x^\sim \xrightarrow{\sim} \hat{S}_{x'}^\sim$, ou de relever $u_{D^*}$ en un
$\tilde{u}_{D^*} : \hat{D}_x^{*\sim} \xrightarrow{\sim} \hat{D}_{x'}^{*\sim}$. Pour le voir, on
note que par fonctorialité des structures les

\newpage

%% ===== Page 271 =====

78s

images inverses par $u_S$ et par $u_{D^*}$ de $\tilde{S}_{x'}^\sim$ et de $\tilde{D}_{x'}^{* \sim}$
soient $\tilde{S}_x^{\prime \sim}$, $\tilde{D}_x^{* \prime \sim}$, sont munis d'une covariation
que je désigne par $h_x^\prime$. Ceci posé, les theoremes
d'iso. de $T_x$-torseurs canoniques

$$ (121) \qquad F(u) \simeq Isom_{\tilde{S}_x}(\tilde{S}_x^\sim, \tilde{S}_x^{\prime \sim}) $$
$$ G(u) \simeq Isom_{\tilde{D}_x^*}(\tilde{D}_x^{*\sim}, \tilde{D}_x^{* \prime \sim}) $$

Or par le diagramme d'équivalence (117) on
a bien un iso de $T_x$-torseurs
$$ Isom((\tilde{S}_x^\sim, \tilde{D}_x^{*\sim}, h_x), (\tilde{S}_x^{\prime \sim}, \tilde{D}_x^{* \prime \sim}, h_x^\prime)) $$

(122)
\begin{tikzcd}
& Isom((\tilde{S}_x^\sim, \tilde{D}_x^{*\sim}, h_x), (\tilde{S}_x^{\prime \sim}, \tilde{D}_x^{* \prime \sim}, h_x^\prime)) \ar[dl] \ar[dr] & \\
F(u) \simeq Isom(\tilde{S}_x^\sim, \tilde{S}_x^{\prime \sim}) & & G(u) \simeq Isom(\tilde{D}_x^{* \sim}, \tilde{D}_x^{* \prime \sim})
\end{tikzcd}

d'où la bijection canonique cherchée (120) —
[unclear: des qui donne un justification (96).]
La bijection canonique (120) a la
propriété $\pm$ évidente suivante, qui la caractérise:

(123) $\begin{cases}
\text{Soient } u_{\tilde{S}} \text{ et } u_{\tilde{D}^*} \text{ des prolongements se correspondants} \\
\tilde{S}_x^\sim \xrightarrow{\sim} \tilde{S}_{x'}^\sim, \ \tilde{D}_x^{*\sim} \xrightarrow{\sim} \tilde{D}_{x'}^{* \sim}, \text{ associés par (120), soit } A \\
\in S_x, \ A' \in S_{x'} \text{ son image par } u_S \text{, soit } \alpha \text{ une} \\
\text{section de } \tilde{S}_x^\sim \text{ sur } A, \ \beta \text{ la section de } \tilde{D}_x^{*\sim} \text{ sur } Ep(A) \\
\text{associé par la covariation } h_x, \text{ soit } \alpha' \text{ la section} \\
\text{de } \tilde{S}_{x'}^\sim \text{ sur } A' \text{ déduite de } \alpha \text{ en appliquant } u_{\tilde{S}}, \\
\beta' \text{ la section de } \tilde{D}_{x'}^{* \sim} \text{ sur } Ep(A') \text{ déduite de } \beta \\
\text{en appliquant } u_{\tilde{D}^*}, \text{ alors } \alpha' \text{ et } \beta' \text{ sont associés} \\
\text{par la covariation } h'.
\end{cases}$

\newpage

%% ===== Page 272 =====

Ceci [unclear: montre] que $A, \alpha$ (fixés) [unclear: détermine] l'applica-
tion (120), car pour $u_{\tilde{S}}$ fixé, il détermine
$u_{\tilde{D}^*}$ ainsi : on a $\beta$ par $h_x$, $\alpha'$ produit de
$\alpha$ par $u_{\tilde{S}}$, $\beta'$ déduit de $\alpha'$ par $h_{x'}$, où
$\beta'$ doit être égal à la section déduite de
$\beta$ par un morphisme $u_{\tilde{D}^*} : \tilde{D}^*_x \to \tilde{D}^*_{x'}$, et ceci
détermine bien $u_{\tilde{D}^*}$ de façon unique...

NB On peut utiliser cette construction, pour
$A, \alpha$ fixés, pour définir l'application (120),
et prouver que c'est un $T_{\pi}$-morphisme
[unclear: Cela clair] et, [unclear: par suite], qu'elle ne
dépend pas du choix de $\alpha$. Mais a priori
il dépend peut-être du choix de $A$, et le
point qui demande vérification est que
cette bijection en fait ne dépend pas du
choix de $A$. On voit en effet que pour
$A \subset A'$, elle est la $\simeq$ [unclear: même] pour $A$ et $A'$ d'où
tout le résultat, car la relation d'équivalence engen-
drée sur les $A_n$ par la rel. d'inclusion est
la relation grossière ....

\newpage

%% ===== Page 273 =====

784

F) Application : interprétation différentielle des groupes de Teichmüller spéciaux.

Soit $X$ un espace topologique, muni d'une partie $I$ discrète finie, telle que $X$ soit une surface au voisinage de $I$. En plus, on se donne une structure différentiable ($C^\infty$) au voisinage de $I$. Soit $G$ un groupe topologique, opérant continûment sur $(X,I)$, et respectant la structure différentiable de $X$ au voisinage de $I$. Soit
(124) $$ \begin{cases} S_{X,I} = S_I = \coprod_{x \in I} S_x \\ D_{X,I}^* = D_I^* = \coprod_{x \in I} D_x^* \end{cases} $$
on suppose que l'opération de $G$ sur $S_I$ induite par l'action de $G$ soit également continue.

Soit, pour tout $x \in I$, $D_x^{*\sim}$ un rev. universel de $D_x^*$, et posons
(125) $$D_{X,I}^{*\sim} = D_I^{*\sim} = \coprod_{x \in I} D_x^{*\sim}$$ .

Oublions pour l'instant les données différentiables, et considérons le groupe
(126) $$\tilde{G} = \{ (u, \tilde{u}) / u \in G, \tilde{u} \text{ un relèvement de } u_{D_I^*} \text{ à } D_I^{*\sim} \}$$

On a un hom. $(u, \tilde{u}) \mapsto u$, $\tilde{G} \to G$, qui est

\newpage

%% ===== Page 274 =====

surj., et, posons
$$ (127) \quad\quad T_I = \prod_{x \in I} T_x \quad , $$
on a une suite exacte
$$ (128) \quad\quad 1 \to T_I \to \tilde{G} \to G \to 1 \quad . $$
Je voudrais introduire sur $\tilde{G}$ une
topologie telle que (128) devienne une
suite exacte de groupes topologiques.
Quitte à remplacer $G$ par le sous-groupe $G'$
d'indice fini noyau de l'hom.
canonique
$$ (129) \quad\quad G \to \mathfrak{S}_I \quad , $$
donc $\tilde{G}$ par l'image inverse $\tilde{G}'$ de $G'$,
on se ramène au cas où $G$ opère
trivialement sur $I$, Alors $\tilde{G}$ s'exprime
comme le produit fibré sur
$G$ de la famille, indexée par $I$,
des extensions
$$ (130) \quad\quad 1 \to T_x \to \tilde{G}_x \to G \to 1 $$
corresp. au cas $I_x \overset{\text{déf}}{=} \{x\}$. Donc
on se ramène au cas où $\text{card } I = 1$.

\newpage

%% ===== Page 275 =====

785

Soit $\Delta_{x,0}$
(131)
un voisinage de $x$ dans $X$
qui soit un disque ouvert. On prend alors
un rev. universel $\tilde{\Delta}_{x,0}^*$ de $\Delta_{x,0}^* = \Delta_{x,0} \setminus \{x\}$
qui induise $\tilde{D}_x^*$ sur $D_x^*$. Par hypothèse
de continuité, il existe un voisinage
$U$ de $1$ dans $G$, et un voisinage $\Delta_{x,1}$ de
$x$, tels que l'on ait
$$U \times \Delta_{x,1} \longrightarrow \Delta_{x,0}$$
(132) $$\cap \quad \quad \quad \quad \quad \cap$$
$$G \times X \longrightarrow X \quad ,$$
et on peut supposer que $\Delta_{x,1}$ soit un
disque ouvert. Ceci posé, pour $u \in U$
fixé, la donnée du relèvement $\tilde{u}$ de $u_{|D_x}$
à $\tilde{D}_x^*$ équivaut à celle d'un relèvement
de
(133) $$\tilde{\Delta}_{x,1} \stackrel{\tilde{u}_{0,1}}{\longrightarrow} \tilde{\Delta}_{x,0}$$
compatible avec
$$u | \Delta_{x,1}^* : \Delta_{x,1}^* \longrightarrow \Delta_{x,0}^* \quad .$$
Cela dit, soit $\tilde{U}$ l'image inverse dans
$\tilde{G}$ du voisinage $U$ de $1$ dans $G$, on a

\newpage

%% ===== Page 276 =====

introduire la [MARGIN: la moins fine rendant continues] top. de la [unclear]
[crossed out: de $\tilde{\mathcal{U}}$ et de $\mathcal{U}$]
les applications

(134)
$$
\begin{cases}
p : \tilde{\mathcal{U}} \longrightarrow \mathcal{U} \quad \text{induite par } \tilde{G} \to G \\
q : \tilde{\mathcal{U}} \longrightarrow \underline{Cont}(\Delta_{x,1}^{*\sim}, \Delta_{x,0}^{*\sim}) \quad (u,\tilde{u}) \longmapsto \tilde{u}_{x,0,1} \\
q' = q \circ \sigma, \quad - \quad \sigma(u,\tilde{u}) = (u,\tilde{u})^{-1} : \tilde{\mathcal{U}} \to \underline{Cont}(\Delta_{x,1}^{*\sim}, \Delta_{x,0}^{*\sim})
\end{cases}
$$

On aura pris soin de prendre $\tilde{\mathcal{U}}$ stable
par $u \mapsto u^{-1}$, ce qui implique qu'il
en est de même de $\mathcal{U}$.

On trouve donc, à l'aide de cette
topologie sur $\tilde{\mathcal{U}}$, un
filtre des voisinages de $1_{\tilde{G}}$ dans $\tilde{\mathcal{U}}$,
donc aussi le filtre image de celui-ci
dans $\tilde{G}$. Il faut vérifier que ces
derniers sont bien le filtre des voisinages
de $1$ pour une topologie de groupe
top., et que pour celle-ci, $\tilde{G} \to G$
est un rev. de $G$, [crossed out: j'admets]
pour que tout marche.

Quand $X$ est une surface compacte orientable, connexe,
et que si $T$ désigne [unclear: les modules] des

\newpage

%% ===== Page 277 =====

[MARGIN: 786]

orientations, alors tous les $T_x$ sont can.
isom. à $T$

(135) $$T_x \simeq T \quad \forall x \in \Gamma$$

et la suite exacte (128) prend la forme

(136) $$1 \longrightarrow T^\Gamma \longrightarrow \tilde{G} \longrightarrow G \longrightarrow 1$$

Le cas "universel" est celui où

(137) $$G = \mathcal{A} = Auttop(X,S)$$

(avec la top. de la conv. uniforme de
$u$ et de $u^{-1}$). Alors dans les cas

orientables (types de $(X,S)$ de genre $g,\nu$
et $2g-2+\nu > 0$ i.e. $2g+\nu \ge 3$

(138) i.e. $(g,\nu) \neq (0,0), (0,1), (0,2), (1,0)$ )

on sait que $\mathcal{A}^\circ$ est [unclear: paracompact]
(en fait = contractile) donc toute extension
de $\mathcal{A}^\circ$ par un groupe discret est
canoniquement scindée, et l'extension $\tilde{\mathcal{A}}$ de
$G$ provient donc d'une extension $\tilde{\mathcal{T}}$

(139) $$1 \longrightarrow T^\Gamma \longrightarrow \tilde{\mathcal{T}} \longrightarrow \mathcal{T} \longrightarrow 1$$

du groupe de Teichmüller

(140) $$\mathcal{T} = \mathcal{T}(X,S) = \mathcal{A}/\mathcal{A}^\circ = \pi_0(Auttop(X,S))$$

\newpage

%% ===== Page 278 =====

qui prend ainsi le nom de groupe de
Teichmüller spécial de $(X,S)$. Pour
$G$ quelconque et action de $G$ sur
$(X,S)$ orientée : la donnée d'un hom.
de groupes topologiques
$$G \to \mathcal{A}$$
s'intègre dans un diagramme de
groupes top.

(141)
\begin{tikzcd}
G \ar[r] \ar[d] & \mathcal{A} \ar[d] \\
\pi_0(G) \ar[r] & \tau = \pi_0(\mathcal{A})
\end{tikzcd}

et on trouve alors, par fonctorialité,
que le $\tilde{G}$ construit précédemment
est l'image inverse de l'ext. $\tilde{\tau}$ de $\tau$
par $\mathbb{Z}^I$, via $G \to \tau$.

Ces préliminaires étant posés,
revenons à la situation où $X$
admet une structure diff. $C^\infty$ au voisin-
age de $S$, et soit pour tout $x \in S$,

\newpage

%% ===== Page 279 =====

$\tilde{S}_x$ "rev" ouv. universel de $S_x$ associé au
rev. universel $\tilde{D}_x^{*\sim}$ de $D_x^*$. Posons

(142) $$ \tilde{S}_I = \coprod_{x \in I} \tilde{S}_x $$

et notons que l'on a, pour $u \in G$,
une correspondance biunivoque entre
l'ens. des relèvements $\tilde{u}_{\tilde{D}_I^*}$ de $u_{D_I^*}$
en un autom de $\tilde{D}_I^{*\sim}$, et l'ens des
relèvements $\tilde{u}_{\tilde{S}_I}$ de $u_S$ en autom. de $\tilde{S}_I$,
qui soit un isom. de $T_I$-torseurs.
Cette correspondance est compatible avec
les compositions, et on trouve ainsi
un isom. canonique

(143) $$ \tilde{G} \simeq \text{groupe des } \{ u, \tilde{u}_{\tilde{S}_I} | u \in G, \tilde{u}_{\tilde{S}_I} \text{ rel. de } u_S \text{ en un autom de } \tilde{S}_I \} $$

On vérifie que c'est là un isom.
top. en mettant sur le deuxième
membre la top. la moins fine rendant
continues les applications

(144) 
$p_1 : \tilde{G} \to G$
$q_1 : \tilde{G} \to \operatorname{Aut}\ top\ (\tilde{S}_I) \quad ( (u,\tilde{u}_{\tilde{S}}) \mapsto \tilde{u}_{\tilde{S}} )$

\newpage

%% ===== Page 280 =====

où $\operatorname{Aut}_{top} (\tilde{\mathfrak{S}}_I)$ est muni de la
convergence compacte.
Dans le cas où $X$ est une surface comp. et diff.
orientable, on trouve donc une interprétation
des groupes précédents, [unclear: purement abstraite, pour]
voire [unclear: effective], (comme images
inverses de sous groupes du Teichmüller
spécial $S\Gamma \to \Gamma$, par l'homom. canoni-
que $G \to \hat{\Gamma}$ dans (141).
